\documentclass[10pt]{article}
\usepackage[margin=0.6in]{geometry}
\usepackage{fontspec}
\setmainfont{Latin Modern Roman}
\setmonofont{Noto Sans Mono CJK SC}[Scale=0.82]
\usepackage{longtable,array,booktabs}
\usepackage[table]{xcolor}
\usepackage{titlesec,fancyhdr,enumitem}
\usepackage[hidelinks]{hyperref}
\definecolor{jhublue}{RGB}{0,45,114}
\definecolor{boxgray}{RGB}{244,246,250}
\titleformat{\section}{\large\bfseries\color{jhublue}}{\thesection}{0.6em}{}
\titleformat{\subsection}{\normalsize\bfseries}{\thesubsection}{0.6em}{}
\pagestyle{fancy}\fancyhf{}
\lhead{\small Solar Policy and Bird Diversity: line-by-line guide to the Stata code}
\rhead{\small Zhang et al.\ (2026) replication}
\cfoot{\small\thepage}
\setlength{\parindent}{0pt}\setlength{\parskip}{4pt}
\renewcommand{\arraystretch}{1.12}
\setlength{\LTpre}{2pt}\setlength{\LTpost}{6pt}
\newcommand{\codetable}[1]{%
{\footnotesize
\begin{longtable}{>{\raggedleft\arraybackslash}p{0.55cm}>{\raggedright\arraybackslash}p{8.3cm}>{\raggedright\arraybackslash}p{8.0cm}}
\toprule \textbf{Line} & \textbf{Code} & \textbf{What it does}\\ \midrule \endfirsthead
\toprule \textbf{Line} & \textbf{Code} & \textbf{What it does}\\ \midrule \endhead
\bottomrule \endfoot
#1
\end{longtable}}}
\begin{document}
\begin{center}
{\LARGE\bfseries Line-by-Line Guide to the Stata Code}\\[4pt]
{\large The authors' CODE.do and every do-file of our replication}\\[3pt]
{\small Replication of Zhang et al.\ (2026), ``China's solar expansion policy reduces bird diversity'', \textit{Science} 393, 831}
\end{center}
\tableofcontents
\newpage

\section{Part I: How to Read Stata Code}
\begin{itemize}[leftmargin=14pt,itemsep=1pt]
\item \textbf{One command per line.} The command name comes first (\texttt{reghdfe}, \texttt{gen}, \texttt{merge}), then variables, then an optional \texttt{if} condition, then \textbf{options after the comma}.
\item \textbf{\texttt{///}} at the end of a line means the command continues on the next line.
\item \textbf{Comments:} a line starting with \texttt{*} or \texttt{//}, text after \texttt{//}, and everything between \texttt{/*} and \texttt{*/} is ignored by Stata.
\item \textbf{Macros:} \texttt{\$name} is a \emph{global} macro (e.g.\ \texttt{\$control} = the 9 controls, \texttt{\$SAMPLE} = path of the analysis data). \texttt{\textasciigrave name\textquotesingle} is a \emph{local} macro (temporary, inside one do-file or loop).
\item \textbf{\texttt{preserve} / \texttt{restore}:} save a copy of the data, change it temporarily, then go back.
\item \textbf{\texttt{merge 1:1} / \texttt{m:1}:} join variables from another file; \texttt{m:1} copies one value (e.g.\ a county-year PSI) to many rows (all months of that year).
\item \textbf{\texttt{reghdfe y x controls, absorb(county year) cluster(county)}:} linear regression with county and year-month fixed effects removed by demeaning, and standard errors clustered by county. In this data the variable \texttt{year} is the \emph{year-month} (120 values); \texttt{YEAR} is the calendar year.
\item \textbf{\texttt{year\#urban}} means separate year-month effects for urban and rural units.
\item \textbf{\texttt{estadd}, \texttt{est store}, \texttt{esttab}:} attach notes to results, save results, and export them as a formatted table.
\item \textbf{\texttt{capture}} runs a command without stopping on error; \textbf{\texttt{\_rc}} is the error code (0 = success).
\item \textbf{\texttt{r()}, \texttt{e()}:} results left by the last command (e.g.\ \texttt{r(N)} count, \texttt{e(df\_r)} residual degrees of freedom, \texttt{\_b[PI]} coefficient, \texttt{\_se[PI]} standard error).
\end{itemize}
Chinese variable names in the authors' code: {\ttfamily 省} province, {\ttfamily 市} city, {\ttfamily 县} county, {\ttfamily 标准年月} year-month text, {\ttfamily 煤电厂} coal plant, {\ttfamily 风电数量} number of new wind projects, {\ttfamily 风电扩张} wind expansion, {\ttfamily 绿色金融} green finance, {\ttfamily 丰富度} richness, {\ttfamily 均匀度} evenness.

\paragraph{Not included.} The short demo files in \texttt{04\_Synthesis\_Report/show\_results/} and \texttt{DEMO\_show\_all\_results.do} only re-run or display results already produced by the files below. Long table-note text in the authors' \texttt{esttab} commands is abridged.
\section{Part II: The Authors' Code (CODE.do)}
The complete replication code published by the authors. Lines 1--245 build the analysis dataset from 46 source files; lines 246--975 produce every table and figure of the paper and supplement. Our copy 1\_1a\_authors\_code\_localpaths.do is identical except that the 6 \texttt{cd} / file-path lines point to our folders.

\codetable{1 & \texttt{*-{}-{}-{}-{}-{}-{}-{}-{}-{}-{}-{}-{}-{}-{}-{}-{}-{}-{}-{}-{}-{}-{}-{}-{}-{}-{}-{}-{}-{}-{}-{}-{}-{}-{}-{}-{}-{}-{}-{}-{}-{}-{}-{}-{}-{}-{}-{}-{}-{}-{}-{}-{}-{}-{}-{}-{}-{}\newline \hspace*{1em}-{}-{}-{}-{}-{}-{}-{}-{}-{}-{}-{}-{}-{}-{}-{}} & Divider line (comment).\\
2 & \texttt{* Data preparation} & Comment (not executed): explains the code that follows.\\
3 & \texttt{*-{}-{}-{}-{}-{}-{}-{}-{}-{}-{}-{}-{}-{}-{}-{}-{}-{}-{}-{}-{}-{}-{}-{}-{}-{}-{}-{}-{}-{}-{}-{}-{}-{}-{}-{}-{}-{}-{}-{}-{}-{}-{}-{}-{}-{}-{}-{}-{}-{}-{}-{}-{}-{}-{}-{}-{}-{}\newline \hspace*{1em}-{}-{}-{}-{}-{}-{}-{}-{}-{}-{}-{}-{}-{}-{}-{}} & Divider line (comment).\\
4 & \texttt{cd \textquotedbl{}/Users/zhangaixin/Library/CloudStorage/OneDrive-{}个人/BD/\newline \hspace*{1em}DATA\textquotedbl{}} & Changes the working directory to the authors' own folder (DATA); later file names are read from / saved there. This is one of the 6 path lines we changed.\\
6 & \texttt{*核心变量+控制变量} & Comment: "core variables + control variables".\\
7 & \texttt{use \textquotedbl{}Bird.dta\textquotedbl{},clear} & Loads the dataset Bird.dta into memory, replacing whatever was loaded.\\
8 & \texttt{merge 1:1 id year using \textquotedbl{}BTN.dta\textquotedbl{},keepus(BT BN)} & One-to-one merge on county code and year-month: attaches BT BN from BTN.dta.\\
9 & \texttt{drop if \_merge==2} & Drops rows that exist only in the merged-in file (no bird data), keeping the bird sample.\\
10 & \texttt{drop \_merge} & Deletes Stata's merge-status variable so the next merge can create a new one.\\
11 & \texttt{merge m:1 id YEAR using \textquotedbl{}PI.dta\textquotedbl{},keepus(PI)} & Many-to-one merge on county code and calendar year (each county-year value is copied to every month of that year): attaches PI from PI.dta.\\
12 & \texttt{drop if \_merge==2} & Drops rows that exist only in the merged-in file (no bird data), keeping the bird sample.\\
13 & \texttt{drop \_merge} & Deletes Stata's merge-status variable so the next merge can create a new one.\\
14 & \texttt{merge m:1 id YEAR using \textquotedbl{}Climate.dta\textquotedbl{},keepus(Avt Wind)} & Many-to-one merge on county code and calendar year (each county-year value is copied to every month of that year): attaches Avt Wind from Climate.dta.\\
15 & \texttt{drop if \_merge==2} & Drops rows that exist only in the merged-in file (no bird data), keeping the bird sample.\\
16 & \texttt{drop \_merge} & Deletes Stata's merge-status variable so the next merge can create a new one.\\
17 & \texttt{rename Avt Temp} & Renames the temperature variable Avt to Temp.\\
18 & \texttt{merge m:1 id YEAR using \textquotedbl{}PD.dta\textquotedbl{},keepus(PD)} & Many-to-one merge on county code and calendar year (each county-year value is copied to every month of that year): attaches PD from PD.dta.\\
19 & \texttt{drop if \_merge==2} & Drops rows that exist only in the merged-in file (no bird data), keeping the bird sample.\\
20 & \texttt{drop \_merge} & Deletes Stata's merge-status variable so the next merge can create a new one.\\
21 & \texttt{merge m:1 id YEAR using \textquotedbl{}air.dta\textquotedbl{},keepus(CO2 PM)} & Many-to-one merge on county code and calendar year (each county-year value is copied to every month of that year): attaches CO2 PM from air.dta.\\
22 & \texttt{drop if \_merge==2} & Drops rows that exist only in the merged-in file (no bird data), keeping the bird sample.\\
23 & \texttt{drop \_merge} & Deletes Stata's merge-status variable so the next merge can create a new one.\\
24 & \texttt{rename PM PM25} & Renames PM to PM25 (raw PM2.5 concentration).\\
25 & \texttt{merge m:1 id YEAR using \textquotedbl{}TD.dta\textquotedbl{},keepus(Water Green Farm\newline \hspace*{1em}Grass)} & Many-to-one merge on county code and calendar year (each county-year value is copied to every month of that year): attaches Water Green Farm Grass from TD.dta.\\
26 & \texttt{drop if \_merge==2} & Drops rows that exist only in the merged-in file (no bird data), keeping the bird sample.\\
27 & \texttt{drop \_merge} & Deletes Stata's merge-status variable so the next merge can create a new one.\\
30 & \texttt{*稳健性} & Comment: "robustness".\\
31 & \texttt{**替换PI} & Comment: "replace PI" (alternative PSI measures).\\
32 & \texttt{merge m:1 id YEAR using \textquotedbl{}PI1.dta\textquotedbl{},keepus(PI1)} & Many-to-one merge on county code and calendar year (each county-year value is copied to every month of that year): attaches PI1 from PI1.dta.\\
33 & \texttt{drop if \_merge==2} & Drops rows that exist only in the merged-in file (no bird data), keeping the bird sample.\\
34 & \texttt{drop \_merge} & Deletes Stata's merge-status variable so the next merge can create a new one.\\
35 & \texttt{merge m:1 id YEAR using \textquotedbl{}PI2.dta\textquotedbl{},keepus(PI2)} & Many-to-one merge on county code and calendar year (each county-year value is copied to every month of that year): attaches PI2 from PI2.dta.\\
36 & \texttt{drop if \_merge==2} & Drops rows that exist only in the merged-in file (no bird data), keeping the bird sample.\\
37 & \texttt{drop \_merge} & Deletes Stata's merge-status variable so the next merge can create a new one.\\
38 & \texttt{merge m:1 id YEAR using \textquotedbl{}PI3.dta\textquotedbl{},keepus(PI3)} & Many-to-one merge on county code and calendar year (each county-year value is copied to every month of that year): attaches PI3 from PI3.dta.\\
39 & \texttt{drop if \_merge==2} & Drops rows that exist only in the merged-in file (no bird data), keeping the bird sample.\\
40 & \texttt{drop \_merge} & Deletes Stata's merge-status variable so the next merge can create a new one.\\
41 & \texttt{merge m:1 id YEAR using \textquotedbl{}PI4.dta\textquotedbl{},keepus(PI4)} & Many-to-one merge on county code and calendar year (each county-year value is copied to every month of that year): attaches PI4 from PI4.dta.\\
42 & \texttt{drop if \_merge==2} & Drops rows that exist only in the merged-in file (no bird data), keeping the bird sample.\\
43 & \texttt{drop \_merge} & Deletes Stata's merge-status variable so the next merge can create a new one.\\
44 & \texttt{merge m:1 id YEAR using \textquotedbl{}PI5.dta\textquotedbl{},keepus(PI5)} & Many-to-one merge on county code and calendar year (each county-year value is copied to every month of that year): attaches PI5 from PI5.dta.\\
45 & \texttt{drop if \_merge==2} & Drops rows that exist only in the merged-in file (no bird data), keeping the bird sample.\\
46 & \texttt{drop \_merge} & Deletes Stata's merge-status variable so the next merge can create a new one.\\
47 & \texttt{merge m:1 id YEAR using \textquotedbl{}PI6.dta\textquotedbl{},keepus(PI6)} & Many-to-one merge on county code and calendar year (each county-year value is copied to every month of that year): attaches PI6 from PI6.dta.\\
48 & \texttt{drop if \_merge==2} & Drops rows that exist only in the merged-in file (no bird data), keeping the bird sample.\\
49 & \texttt{drop \_merge} & Deletes Stata's merge-status variable so the next merge can create a new one.\\
50 & \texttt{replace PI1 = 0 if missing(PI1)} & Sets missing values of PI1 to 0 (counties absent from that source file are treated as zero).\\
51 & \texttt{replace PI2 = 0 if missing(PI2)} & Sets missing values of PI2 to 0 (counties absent from that source file are treated as zero).\\
52 & \texttt{replace PI3 = 0 if missing(PI3)} & Sets missing values of PI3 to 0 (counties absent from that source file are treated as zero).\\
53 & \texttt{replace PI4 = 0 if missing(PI4)} & Sets missing values of PI4 to 0 (counties absent from that source file are treated as zero).\\
54 & \texttt{replace PI5 = 0 if missing(PI5)} & Sets missing values of PI5 to 0 (counties absent from that source file are treated as zero).\\
55 & \texttt{replace PI6 = 0 if missing(PI6)} & Sets missing values of PI6 to 0 (counties absent from that source file are treated as zero).\\
56 & \texttt{**王老师装机规模} & Comment: "Teacher Wang's installed capacity" (PV capacity data).\\
57 & \texttt{merge m:1 id YEAR using \textquotedbl{}IC.dta\textquotedbl{},keepus(ICtotal)} & Many-to-one merge on county code and calendar year (each county-year value is copied to every month of that year): attaches ICtotal from IC.dta.\\
58 & \texttt{drop if \_merge==2} & Drops rows that exist only in the merged-in file (no bird data), keeping the bird sample.\\
59 & \texttt{drop \_merge} & Deletes Stata's merge-status variable so the next merge can create a new one.\\
60 & \texttt{replace ICtotal=0 if missing(ICtotal)} & Sets missing values of ICtotal to 0 (counties absent from that source file are treated as zero).\\
61 & \texttt{**省会城市直辖市cities} & Comment: "provincial capitals and municipalities".\\
62 & \texttt{merge m:1 省 市 using \textquotedbl{}cities.dta\textquotedbl{},keepus(cities)} & Many-to-one merge on province and city name: attaches cities from cities.dta.\\
63 & \texttt{drop if \_merge==2} & Drops rows that exist only in the merged-in file (no bird data), keeping the bird sample.\\
64 & \texttt{drop \_merge} & Deletes Stata's merge-status variable so the next merge can create a new one.\\
65 & \texttt{replace cities = 0 if missing(cities)} & Sets missing values of cities to 0 (counties absent from that source file are treated as zero).\\
66 & \texttt{**环保督察政策CEI} & Comment: "central environmental inspection policy (CEI)".\\
67 & \texttt{merge m:1 省 YEAR using \textquotedbl{}enin.dta\textquotedbl{},keepus(enin)} & Many-to-one merge on province and year: attaches enin from enin.dta.\\
68 & \texttt{drop if \_merge==2} & Drops rows that exist only in the merged-in file (no bird data), keeping the bird sample.\\
69 & \texttt{drop \_merge} & Deletes Stata's merge-status variable so the next merge can create a new one.\\
70 & \texttt{rename enin CEI} & Renames the inspection variable to CEI.\\
71 & \texttt{**煤电厂} & Comment: "coal power plants".\\
72 & \texttt{merge m:1 id YEAR using \textquotedbl{}Coal.dta\textquotedbl{},keepus(FFP)} & Many-to-one merge on county code and calendar year (each county-year value is copied to every month of that year): attaches FFP from Coal.dta.\\
73 & \texttt{drop if \_merge==2} & Drops rows that exist only in the merged-in file (no bird data), keeping the bird sample.\\
74 & \texttt{drop \_merge} & Deletes Stata's merge-status variable so the next merge can create a new one.\\
75 & \texttt{gen 煤电厂=0} & Creates the coal-plant dummy (Chinese name "coal power plant") = 0 for everyone.\\
76 & \texttt{replace 煤电厂=1 if FFP !=.} & Sets the coal dummy to 1 where a coal-plant closure record exists (FFP not missing).\\
77 & \texttt{rename 煤电厂 Coal} & Renames it to Coal.\\
78 & \texttt{**风电扩张} & Comment: "wind power expansion".\\
79 & \texttt{merge m:1 id YEAR using \textquotedbl{}Windexp.dta\textquotedbl{},keepus(风电数量)} & Many-to-one merge on county code and calendar year (each county-year value is copied to every month of that year): attaches {\ttfamily 风电数量} from Windexp.dta.\\
80 & \texttt{drop if \_merge==2} & Drops rows that exist only in the merged-in file (no bird data), keeping the bird sample.\\
81 & \texttt{drop \_merge} & Deletes Stata's merge-status variable so the next merge can create a new one.\\
82 & \texttt{gen 风电扩张=0} & Creates the wind-expansion dummy = 0.\\
83 & \texttt{replace 风电扩张=1 if 风电数量 !=0} & Sets wind expansion = 1 if the number of new wind projects is not 0. BUG: in Stata a missing value counts as "not 0", so all 33,877 rows with missing wind data become 1.\\
84 & \texttt{rename 风电扩张 Windexp} & Renames it to Windexp.\\
85 & \texttt{**绿色金融} & Comment: "green finance".\\
86 & \texttt{gen 绿色金融 = 0} & Creates the green-finance pilot dummy = 0.\\
87 & \texttt{local group1 \textquotedbl{}湖州市 衢州市 广州市 哈密市 昌吉回族自治州 克拉玛依市\textquotedbl{}} & Lists the first group of green-finance pilot cities: Huzhou, Quzhou, Guangzhou, Hami, Changji, Karamay.\\
88 & \texttt{foreach city in \textasciigrave{}group1\textquotesingle{} \{} & Loops over those cities.\\
89 & \texttt{\hspace*{2.0em}replace 绿色金融 = 1 if strpos(市, \textquotedbl{}\textasciigrave{}city\textquotesingle{}\textquotedbl{}) > 0 \& (YEAR > 2017\newline \hspace*{1em}| (YEAR == 2017 \& month >= 7))} & Marks the city as a pilot zone from July 2017 onward (strpos finds the city name in the city variable).\\
90 & \texttt{\}} & Closes the loop / block.\\
91 & \texttt{*赣江新区} & Comment: "Ganjiang New Area" (green-finance pilot).\\
92 & \texttt{replace 绿色金融 = 1 if (strpos(市, \textquotedbl{}南昌市\textquotedbl{}) > 0 \& (strpos(县,\newline \hspace*{1em}\textquotedbl{}新建区\textquotedbl{}) > 0 | strpos(县, \textquotedbl{}青山湖区\textquotedbl{}) > 0)) ///} & Ganjiang New Area: Xinjian and Qingshanhu districts of Nanchang ...\\
93 & \texttt{\hspace*{4.0em}\& (YEAR > 2017 | (YEAR == 2017 \& month >= 7))} & Continuation of the condition: the date from which the pilot zone applies.\\
94 & \texttt{replace 绿色金融 = 1 if (strpos(市, \textquotedbl{}九江市\textquotedbl{}) > 0 \& (strpos(县,\newline \hspace*{1em}\textquotedbl{}共青城市\textquotedbl{}) > 0 | strpos(县, \textquotedbl{}永修县\textquotedbl{}) > 0)) ///} & ... and Gongqingcheng and Yongxiu in Jiujiang ...\\
95 & \texttt{\hspace*{4.0em}\& (YEAR > 2017 | (YEAR == 2017 \& month >= 7))} & Continuation of the condition: the date from which the pilot zone applies.\\
96 & \texttt{*贵安新区} & Comment: "Gui'an New Area" (green-finance pilot).\\
97 & \texttt{replace 绿色金融 = 1 if (strpos(市, \textquotedbl{}贵阳市\textquotedbl{}) > 0) ///} & Gui'an New Area: all of Guiyang ...\\
98 & \texttt{\hspace*{4.0em}\& (YEAR > 2017 | (YEAR == 2017 \& month >= 7))} & Continuation of the condition: the date from which the pilot zone applies.\\
99 & \texttt{replace 绿色金融 = 1 if (strpos(市, \textquotedbl{}安顺市\textquotedbl{}) > 0 \& strpos(县,\newline \hspace*{1em}\textquotedbl{}平坝区\textquotedbl{}) > 0) ///} & ... and Pingba district of Anshun ...\\
100 & \texttt{\hspace*{4.0em}\& (YEAR > 2017 | (YEAR == 2017 \& month >= 7))} & Continuation of the condition: the date from which the pilot zone applies.\\
101 & \texttt{* 兰州新区} & Comment: "Lanzhou New Area" (green-finance pilot).\\
102 & \texttt{replace 绿色金融 = 1 if (strpos(市, \textquotedbl{}兰州市\textquotedbl{}) > 0 \& (strpos(县,\newline \hspace*{1em}\textquotedbl{}皋兰县\textquotedbl{}) > 0 | strpos(县, \textquotedbl{}永登县\textquotedbl{}) > 0)) ///} & Lanzhou New Area: Gaolan and Yongdeng counties ...\\
103 & \texttt{\hspace*{4.0em}\& (YEAR > 2019 | (YEAR == 2019 \& month >= 12))} & Continuation of the condition: the date from which the pilot zone applies.\\
104 & \texttt{* 重庆市} & Comment: "Chongqing municipality" (green-finance pilot).\\
105 & \texttt{replace 绿色金融 = 1 if strpos(市, \textquotedbl{}重庆市\textquotedbl{}) > 0 ///} & Chongqing municipality ...\\
106 & \texttt{\hspace*{4.0em}\& (YEAR > 2022 | (YEAR == 2022 \& month >= 9))} & Continuation of the condition: the date from which the pilot zone applies.\\
107 & \texttt{rename 绿色金融 Greenfin} & Renames the dummy to Greenfin.\\
108 & \texttt{**IV} & Comment: instrumental variable inputs.\\
109 & \texttt{***sun算术平均1984-{}2013} & Comment: "sun = arithmetic mean 1984-2013".\\
110 & \texttt{merge m:1 id YEAR using \textquotedbl{}sunshine.dta\textquotedbl{},keepus(sun)} & Many-to-one merge on county code and calendar year (each county-year value is copied to every month of that year): attaches sun from sunshine.dta.\\
111 & \texttt{drop if \_merge==2} & Drops rows that exist only in the merged-in file (no bird data), keeping the bird sample.\\
112 & \texttt{drop \_merge} & Deletes Stata's merge-status variable so the next merge can create a new one.\\
113 & \texttt{***地级市CPU} & Comment: "prefecture-city CPU".\\
114 & \texttt{merge m:1 市 YEAR using \textquotedbl{}CCPU.dta\textquotedbl{},keepus(CCPU)} & Many-to-one merge on city name and year: attaches CCPU from CCPU.dta.\\
115 & \texttt{drop if \_merge==2} & Drops rows that exist only in the merged-in file (no bird data), keeping the bird sample.\\
116 & \texttt{drop \_merge} & Deletes Stata's merge-status variable so the next merge can create a new one.\\
117 & \texttt{**高PI县相邻县} & Comment: "counties adjacent to high-PI counties".\\
118 & \texttt{merge m:1 id using\newline \hspace*{1em}\textquotedbl{}highPI.dta\textquotedbl{},keepus(is\_neighbor\_of\_highPI)} & Many-to-one merge on county code (one value per county): attaches is\_neighbor\_of\_highPI from highPI.dta.\\
119 & \texttt{drop if \_merge==2} & Drops rows that exist only in the merged-in file (no bird data), keeping the bird sample.\\
120 & \texttt{drop \_merge} & Deletes Stata's merge-status variable so the next merge can create a new one.\\
121 & \texttt{replace is\_neighbor\_of\_highPI=0 if\newline \hspace*{1em}missing(is\_neighbor\_of\_highPI)} & Sets missing values of is\_neighbor\_of\_highPI to 0 (counties absent from that source file are treated as zero).\\
122 & \texttt{**剔除低于最低观鸟记录阈值的观测值} & Comment: "drop observations below the minimum birdwatching-record threshold".\\
123 & \texttt{merge 1:1 id year using \textquotedbl{}BirdDrop3.dta\textquotedbl{},keepus(ShannonBD3)} & One-to-one merge on county code and year-month: attaches ShannonBD3 from BirdDrop3.dta.\\
124 & \texttt{drop if \_merge==2} & Drops rows that exist only in the merged-in file (no bird data), keeping the bird sample.\\
125 & \texttt{drop \_merge} & Deletes Stata's merge-status variable so the next merge can create a new one.\\
126 & \texttt{merge 1:1 id year using \textquotedbl{}BirdDrop5.dta\textquotedbl{},keepus(ShannonBD5)} & One-to-one merge on county code and year-month: attaches ShannonBD5 from BirdDrop5.dta.\\
127 & \texttt{drop if \_merge==2} & Drops rows that exist only in the merged-in file (no bird data), keeping the bird sample.\\
128 & \texttt{drop \_merge} & Deletes Stata's merge-status variable so the next merge can create a new one.\\
129 & \texttt{merge 1:1 id year using\newline \hspace*{1em}\textquotedbl{}BirdDrop10.dta\textquotedbl{},keepus(ShannonBD10)} & One-to-one merge on county code and year-month: attaches ShannonBD10 from BirdDrop10.dta.\\
130 & \texttt{drop if \_merge==2} & Drops rows that exist only in the merged-in file (no bird data), keeping the bird sample.\\
131 & \texttt{drop \_merge} & Deletes Stata's merge-status variable so the next merge can create a new one.\\
132 & \texttt{merge 1:1 id year using\newline \hspace*{1em}\textquotedbl{}BirdDrop20.dta\textquotedbl{},keepus(ShannonBD20)} & One-to-one merge on county code and year-month: attaches ShannonBD20 from BirdDrop20.dta.\\
133 & \texttt{drop if \_merge==2} & Drops rows that exist only in the merged-in file (no bird data), keeping the bird sample.\\
134 & \texttt{drop \_merge} & Deletes Stata's merge-status variable so the next merge can create a new one.\\
137 & \texttt{*异质性} & Comment: "heterogeneity".\\
138 & \texttt{**贫困县poverty\_county、三北地区north\_region、低碳试点low\_carbon\_city} & Comment: "poverty counties, Three-North region, low-carbon pilot cities".\\
139 & \texttt{merge m:1 县 using \textquotedbl{}poverty.dta\textquotedbl{},keepus(poverty\_county)} & Many-to-one merge on county name: attaches poverty\_county from poverty.dta.\\
140 & \texttt{drop if \_merge==2} & Drops rows that exist only in the merged-in file (no bird data), keeping the bird sample.\\
141 & \texttt{drop \_merge} & Deletes Stata's merge-status variable so the next merge can create a new one.\\
142 & \texttt{replace poverty\_county = 0 if missing(poverty\_county)} & Sets missing values of poverty\_county to 0 (counties absent from that source file are treated as zero).\\
143 & \texttt{merge m:1 省 using \textquotedbl{}north\_region.dta\textquotedbl{},keepus(north\_region)} & Many-to-one merge on province name: attaches north\_region from north\_region.dta.\\
144 & \texttt{drop if \_merge==2} & Drops rows that exist only in the merged-in file (no bird data), keeping the bird sample.\\
145 & \texttt{drop \_merge} & Deletes Stata's merge-status variable so the next merge can create a new one.\\
146 & \texttt{replace north\_region = 0 if missing(north\_region)} & Sets missing values of north\_region to 0 (counties absent from that source file are treated as zero).\\
147 & \texttt{merge m:1 id using \textquotedbl{}Gebi.dta\textquotedbl{},keepus(desert\_gobi)} & Many-to-one merge on county code (one value per county): attaches desert\_gobi from Gebi.dta.\\
148 & \texttt{drop if \_merge==2} & Drops rows that exist only in the merged-in file (no bird data), keeping the bird sample.\\
149 & \texttt{drop \_merge} & Deletes Stata's merge-status variable so the next merge can create a new one.\\
150 & \texttt{**候鸟Migratory、留鸟Resident+特有物种Endemic、非特有物种non\_Endemic+国家保护\newline \hspace*{1em}Protected、非国家保护non\_Protected} & Comment: "migratory/resident, endemic/non-endemic, protected/non-protected birds".\\
151 & \texttt{merge 1:1 id year using \textquotedbl{}Migratory.dta\textquotedbl{}, keepus(Migratory)} & One-to-one merge on county code and year-month: attaches Migratory from Migratory.dta.\\
152 & \texttt{drop if \_merge==2} & Drops rows that exist only in the merged-in file (no bird data), keeping the bird sample.\\
153 & \texttt{drop \_merge} & Deletes Stata's merge-status variable so the next merge can create a new one.\\
154 & \texttt{merge 1:1 id year using \textquotedbl{}Resident.dta\textquotedbl{}, keepus(Resident)} & One-to-one merge on county code and year-month: attaches Resident from Resident.dta.\\
155 & \texttt{drop if \_merge==2} & Drops rows that exist only in the merged-in file (no bird data), keeping the bird sample.\\
156 & \texttt{drop \_merge} & Deletes Stata's merge-status variable so the next merge can create a new one.\\
157 & \texttt{merge 1:1 id year using \textquotedbl{}Endemic.dta\textquotedbl{}, keepus(Endemic)} & One-to-one merge on county code and year-month: attaches Endemic from Endemic.dta.\\
158 & \texttt{drop if \_merge==2} & Drops rows that exist only in the merged-in file (no bird data), keeping the bird sample.\\
159 & \texttt{drop \_merge} & Deletes Stata's merge-status variable so the next merge can create a new one.\\
160 & \texttt{merge 1:1 id year using \textquotedbl{}non\_Endemic.dta\textquotedbl{},\newline \hspace*{1em}keepus(non\_Endemic)} & One-to-one merge on county code and year-month: attaches non\_Endemic from non\_Endemic.dta.\\
161 & \texttt{drop if \_merge==2} & Drops rows that exist only in the merged-in file (no bird data), keeping the bird sample.\\
162 & \texttt{drop \_merge} & Deletes Stata's merge-status variable so the next merge can create a new one.\\
163 & \texttt{merge 1:1 id year using \textquotedbl{}Protected.dta\textquotedbl{}, keepus(Protected)} & One-to-one merge on county code and year-month: attaches Protected from Protected.dta.\\
164 & \texttt{drop if \_merge==2} & Drops rows that exist only in the merged-in file (no bird data), keeping the bird sample.\\
165 & \texttt{drop \_merge} & Deletes Stata's merge-status variable so the next merge can create a new one.\\
166 & \texttt{merge 1:1 id year using \textquotedbl{}non\_Protected.dta\textquotedbl{},\newline \hspace*{1em}keepus(non\_Protected)} & One-to-one merge on county code and year-month: attaches non\_Protected from non\_Protected.dta.\\
167 & \texttt{drop if \_merge==2} & Drops rows that exist only in the merged-in file (no bird data), keeping the bird sample.\\
168 & \texttt{drop \_merge} & Deletes Stata's merge-status variable so the next merge can create a new one.\\
169 & \texttt{**植物筑巢Nest\_veg、水上筑巢Land\_water+食肉杂食non\_herbivorous、食草Herbiv\newline \hspace*{1em}orous+小群落Small\_flock、大群落Large\_flock} & Comment: "vegetation/water nesters, carnivore-omnivore/herbivore, small/large flocks".\\
170 & \texttt{merge 1:1 id year using \textquotedbl{}Nest\_veg.dta\textquotedbl{}, keepus(Nest\_veg)} & One-to-one merge on county code and year-month: attaches Nest\_veg from Nest\_veg.dta.\\
171 & \texttt{drop if \_merge==2} & Drops rows that exist only in the merged-in file (no bird data), keeping the bird sample.\\
172 & \texttt{drop \_merge} & Deletes Stata's merge-status variable so the next merge can create a new one.\\
173 & \texttt{merge 1:1 id year using \textquotedbl{}Land\_water.dta\textquotedbl{},\newline \hspace*{1em}keepus(Land\_water)} & One-to-one merge on county code and year-month: attaches Land\_water from Land\_water.dta.\\
174 & \texttt{drop if \_merge==2} & Drops rows that exist only in the merged-in file (no bird data), keeping the bird sample.\\
175 & \texttt{drop \_merge} & Deletes Stata's merge-status variable so the next merge can create a new one.\\
176 & \texttt{merge 1:1 id year using \textquotedbl{}non\_Herbivorous.dta\textquotedbl{},\newline \hspace*{1em}keepus(non\_Herbivorous)} & One-to-one merge on county code and year-month: attaches non\_Herbivorous from non\_Herbivorous.dta.\\
177 & \texttt{drop if \_merge==2} & Drops rows that exist only in the merged-in file (no bird data), keeping the bird sample.\\
178 & \texttt{drop \_merge} & Deletes Stata's merge-status variable so the next merge can create a new one.\\
179 & \texttt{merge 1:1 id year using \textquotedbl{}Herbivorous.dta\textquotedbl{},\newline \hspace*{1em}keepus(Herbivorous)} & One-to-one merge on county code and year-month: attaches Herbivorous from Herbivorous.dta.\\
180 & \texttt{drop if \_merge==2} & Drops rows that exist only in the merged-in file (no bird data), keeping the bird sample.\\
181 & \texttt{drop \_merge} & Deletes Stata's merge-status variable so the next merge can create a new one.\\
182 & \texttt{merge 1:1 id year using \textquotedbl{}Small\_flock.dta\textquotedbl{},\newline \hspace*{1em}keepus(Small\_flock)} & One-to-one merge on county code and year-month: attaches Small\_flock from Small\_flock.dta.\\
183 & \texttt{drop if \_merge==2} & Drops rows that exist only in the merged-in file (no bird data), keeping the bird sample.\\
184 & \texttt{drop \_merge} & Deletes Stata's merge-status variable so the next merge can create a new one.\\
185 & \texttt{merge 1:1 id year using \textquotedbl{}Large\_flock.dta\textquotedbl{},\newline \hspace*{1em}keepus(Large\_flock)} & One-to-one merge on county code and year-month: attaches Large\_flock from Large\_flock.dta.\\
186 & \texttt{drop if \_merge==2} & Drops rows that exist only in the merged-in file (no bird data), keeping the bird sample.\\
187 & \texttt{drop \_merge} & Deletes Stata's merge-status variable so the next merge can create a new one.\\
190 & \texttt{*机制} & Comment: "mechanisms".\\
191 & \texttt{merge m:1 id YEAR using \textquotedbl{}Ndvi.dta\textquotedbl{},keepus(Ndvi)} & Many-to-one merge on county code and calendar year (each county-year value is copied to every month of that year): attaches Ndvi from Ndvi.dta.\\
192 & \texttt{drop if \_merge==2} & Drops rows that exist only in the merged-in file (no bird data), keeping the bird sample.\\
193 & \texttt{drop \_merge} & Deletes Stata's merge-status variable so the next merge can create a new one.\\
194 & \texttt{replace Ndvi = Ndvi*100} & Rescales NDVI from 0-1 to 0-100.\\
195 & \texttt{merge m:1 id YEAR using \textquotedbl{}NL.dta\textquotedbl{},keepus(Nightlight)} & Many-to-one merge on county code and calendar year (each county-year value is copied to every month of that year): attaches Nightlight from NL.dta.\\
196 & \texttt{drop if \_merge==2} & Drops rows that exist only in the merged-in file (no bird data), keeping the bird sample.\\
197 & \texttt{drop \_merge} & Deletes Stata's merge-status variable so the next merge can create a new one.\\
198 & \texttt{merge m:1 id YEAR using \textquotedbl{}LAI.dta\textquotedbl{},keepus(LAI)} & Many-to-one merge on county code and calendar year (each county-year value is copied to every month of that year): attaches LAI from LAI.dta.\\
199 & \texttt{drop if \_merge==2} & Drops rows that exist only in the merged-in file (no bird data), keeping the bird sample.\\
200 & \texttt{drop \_merge} & Deletes Stata's merge-status variable so the next merge can create a new one.\\
203 & \texttt{*经济后果} & Comment: "economic consequences".\\
204 & \texttt{**丰富度与均匀度} & Comment: "richness and evenness".\\
205 & \texttt{merge 1:1 id year using \textquotedbl{}Richeven.dta\textquotedbl{},keepus(丰富度)} & One-to-one merge on county code and year-month: attaches {\ttfamily 丰富度} from Richeven.dta.\\
206 & \texttt{drop if \_merge==2} & Drops rows that exist only in the merged-in file (no bird data), keeping the bird sample.\\
207 & \texttt{drop \_merge} & Deletes Stata's merge-status variable so the next merge can create a new one.\\
208 & \texttt{gen 均匀度 = ShannonBD/ln(丰富度)} & Creates evenness (Chinese "evenness") = Shannon / ln(richness): Pielou's evenness.\\
209 & \texttt{rename 丰富度 Richness} & Renames richness (Chinese "richness") to Richness.\\
210 & \texttt{rename 均匀度 Evenness} & Renames evenness to Evenness.\\
211 & \texttt{**农作物产量} & Comment: "crop yield".\\
212 & \texttt{merge m:1 id YEAR using \textquotedbl{}AP.dta\textquotedbl{},keepus(APtotal)} & Many-to-one merge on county code and calendar year (each county-year value is copied to every month of that year): attaches APtotal from AP.dta.\\
213 & \texttt{drop if \_merge==2} & Drops rows that exist only in the merged-in file (no bird data), keeping the bird sample.\\
214 & \texttt{drop \_merge} & Deletes Stata's merge-status variable so the next merge can create a new one.\\
217 & \texttt{*数据预处理} & Comment: "data pre-processing".\\
218 & \texttt{**生成人均观鸟时长} & Comment: "generate birdwatching time per person".\\
219 & \texttt{gen BTN=BT/BN} & Birdwatching time per birdwatcher = total time / number of birdwatchers.\\
220 & \texttt{**PI缺失值替换为0} & Comment: "replace missing PI with 0".\\
221 & \texttt{replace PI = 0 if missing(PI)} & KEY LINE: sets PSI to 0 whenever it is missing. 724 counties absent from the policy file (43.8\% of rows, mostly urban districts) get PSI = 0 and stay in every regression.\\
222 & \texttt{**剔除控制变量存在缺失的样本} & Comment: "drop observations with missing controls".\\
223 & \texttt{drop if missing(Temp, Wind, PD, BTN, CO2, Water, Green,\newline \hspace*{1em}Farm, Grass)} & Drops every row missing any control variable (48,900 bird rows become 46,528).\\
224 & \texttt{**王老师装机规模换算成面积} & Comment: "convert installed capacity to area".\\
225 & \texttt{gen Area = ICtotal*1000/0.15} & PV area: capacity in MW x 1000 = kW, divided by 0.15 kW per m2 (150 W/m2) = m2 ...\\
226 & \texttt{replace Area = Area / 1000000} & ... divided by 1,000,000 to convert m2 to km2.\\
227 & \texttt{label var Area \textquotedbl{}集中式光伏电站面积平方千米\textquotedbl{}} & Labels Area "centralized PV station area, km2".\\
228 & \texttt{**取对数} & Comment: "take logarithms".\\
229 & \texttt{gen Pop=ln(PD)} & Pop = log of population density.\\
230 & \texttt{gen Duration=ln(BTN+1)} & Duration = log(time per birdwatcher + 1): the observer-effort control (+1 avoids log of 0).\\
231 & \texttt{gen Carbon=ln(CO2)} & Carbon = log of CO2 emissions.\\
232 & \texttt{gen PM=ln(PM25)} & PM = log of PM2.5.\\
233 & \texttt{gen Output = ln(APtotal)} & Output = log of crop output.\\
234 & \texttt{**所有连续变量缩尾} & Comment: "winsorize all continuous variables".\\
235 & \texttt{local var \textquotedbl{}ShannonBD SimpsonBD PI Temp Wind Pop Duration\newline \hspace*{1em}Carbon Water Green Farm Grass PM ICtotal Area Ndvi\newline \hspace*{1em}Nightlight LAI PI1 PI2 PI3 PI4 PI5 PI6 Migratory Resident\newline \hspace*{1em}Endemic non\_Endemic Protected non\_Protected Nest\_veg\newline \hspace*{1em}Land\_water non\_Herbivorous Herbivorous Small\_flock\newline \hspace*{1em}Large\_flock Richness Evenness  ...} & Lists the 46 continuous variables to be winsorized on the next line (Shannon, PSI, controls, alternative PSIs, bird groups, mechanisms, instrument parts).\\
236 & \texttt{winsor2 \textasciigrave{}var\textquotesingle{}, cuts(1 99) replace} & Winsorizes every listed variable at the 1st and 99th percentiles: values beyond are replaced by those percentile values (replace = overwrite).\\
237 & \texttt{**生成IV} & Comment: "generate the IV".\\
238 & \texttt{gen lnsun = ln(sun)} & Log of long-run sunshine (the "share" of the shift-share instrument).\\
239 & \texttt{gen inv\_CCPU = 1 / CCPU} & Inverse of city climate policy uncertainty (the "shift").\\
240 & \texttt{gen Sun\_ccpu = lnsun / CCPU} & The instrument: ln(sunshine) / city CPU.\\
241 & \texttt{**生成county} & Comment: "generate county id".\\
242 & \texttt{egen county = group(id)} & county = 1, 2, ..., 2,501: a sequential number for each county code; used for fixed effects and clustering.\\
246 & \texttt{cd \textquotedbl{}/Users/zhangaixin/Library/CloudStorage/OneDrive-{}个人/BD/\newline \hspace*{1em}Documents\textquotedbl{}} & Changes the working directory to the authors' own folder (Documents); later file names are read from / saved there. This is one of the 6 path lines we changed.\\
247 & \texttt{*-{}-{}-{}-{}-{}-{}-{}-{}-{}-{}-{}-{}-{}-{}-{}-{}-{}-{}-{}-{}-{}-{}-{}-{}-{}-{}-{}-{}-{}-{}-{}-{}-{}-{}-{}-{}-{}-{}-{}-{}-{}-{}-{}-{}-{}-{}-{}-{}-{}-{}-{}-{}-{}-{}-{}-{}-{}\newline \hspace*{1em}-{}-{}-{}-{}-{}-{}-{}-{}-{}-{}-{}-{}-{}-{}-{}} & Divider line (comment).\\
248 & \texttt{* Descriptive statistics} & Comment (not executed): explains the code that follows.\\
249 & \texttt{*-{}-{}-{}-{}-{}-{}-{}-{}-{}-{}-{}-{}-{}-{}-{}-{}-{}-{}-{}-{}-{}-{}-{}-{}-{}-{}-{}-{}-{}-{}-{}-{}-{}-{}-{}-{}-{}-{}-{}-{}-{}-{}-{}-{}-{}-{}-{}-{}-{}-{}-{}-{}-{}-{}-{}-{}-{}\newline \hspace*{1em}-{}-{}-{}-{}-{}-{}-{}-{}-{}-{}-{}-{}-{}-{}-{}} & Divider line (comment).\\
250 & \texttt{preserve} & Saves a copy of the data in memory, so it can be restored after temporary changes.\\
251 & \texttt{local var \textquotedbl{}ShannonBD SimpsonBD PI Temp Wind Pop Duration\newline \hspace*{1em}Carbon Water Green Farm Grass CCPU lnsun Sun\_ccpu\textquotedbl{}} & Variables for Table S4.\\
252 & \texttt{estpost tabstat \textasciigrave{}var\textquotesingle{}, stats(count mean sd min max) c(s)} & Computes N, mean, SD, min, max for the listed variables and stores them for esttab (Table S4 Panel A).\\
253 & \texttt{eststo stats} & Stores them as "stats".\\
254 & \texttt{esttab  stats using sum\_stats.rtf, replace title(\textquotedbl{}Table\newline \hspace*{1em}S4. Descriptive statistics.\textquotedbl{}) ///} & Exports the stored results ( stats) into the Word table sum\_stats.rtf.\\
255 & \texttt{\hspace*{2.0em}cells(\textquotedbl{}count(fmt(\%9.0fc) label(N)) mean(fmt(\%9.4fc)\newline \hspace*{1em}label(Mean)) sd(fmt(\%9.4fc) label(S.D.)) min(fmt(\%9.4fc)\newline \hspace*{1em}label(Min)) max(fmt(\%9.4fc) label(Max))\textquotedbl{}) ///} & Continuation of the previous command: number formats for the bottom rows; which statistics to show and their formats.\\
256 & \texttt{\hspace*{2.0em}nomtitle nonumber noobs} & Continuation of the previous command: no model titles or numbers.\\
257 & \texttt{restore} & Restores the data saved by preserve (undoes changes made since).\\
258 & \texttt{/*} & Starts a block comment: everything until */ is ignored by Stata.\\
259 & \texttt{sum2docx ShannonBD SimpsonBD PI Temp Wind Pop Duration\newline \hspace*{1em}Carbon Water Green Farm Grass CCPU lnsun Sun\_ccpu using\newline \hspace*{1em}\textquotedbl{}summary.docx\textquotedbl{}, replace stats(N mean(\%9.4f) sd(\%9.4f)\newline \hspace*{1em}min(\%9.4f)  max(\%9.4f)) title(\textquotedbl{}Table S4：Descriptive\newline \hspace*{1em}statistics\textquotedbl{})} & Inside a block comment (not executed).\\
260 & \texttt{*/} & Ends the block comment.\\
264 & \texttt{*-{}-{}-{}-{}-{}-{}-{}-{}-{}-{}-{}-{}-{}-{}-{}-{}-{}-{}-{}-{}-{}-{}-{}-{}-{}-{}-{}-{}-{}-{}-{}-{}-{}-{}-{}-{}-{}-{}-{}-{}-{}-{}-{}-{}-{}-{}-{}-{}-{}-{}-{}-{}-{}-{}-{}-{}-{}\newline \hspace*{1em}-{}-{}-{}-{}-{}-{}-{}-{}-{}-{}-{}-{}-{}-{}-{}} & Divider line (comment).\\
265 & \texttt{* Pearson Correlation Matrix} & Comment (not executed): explains the code that follows.\\
266 & \texttt{*-{}-{}-{}-{}-{}-{}-{}-{}-{}-{}-{}-{}-{}-{}-{}-{}-{}-{}-{}-{}-{}-{}-{}-{}-{}-{}-{}-{}-{}-{}-{}-{}-{}-{}-{}-{}-{}-{}-{}-{}-{}-{}-{}-{}-{}-{}-{}-{}-{}-{}-{}-{}-{}-{}-{}-{}-{}\newline \hspace*{1em}-{}-{}-{}-{}-{}-{}-{}-{}-{}-{}-{}-{}-{}-{}-{}} & Divider line (comment).\\
267 & \texttt{preserve} & Saves a copy of the data in memory, so it can be restored after temporary changes.\\
268 & \texttt{estpost correlate ShannonBD PI Temp Wind Pop Duration\newline \hspace*{1em}Carbon Water Green Farm Grass CCPU lnsun, matrix listwise} & Pearson correlation matrix (Table S4 Panel B), using rows with all variables.\\
269 & \texttt{eststo corr} & Stores it.\\
270 & \texttt{esttab corr using sum\_stats.rtf, append title(\textquotedbl{}Table S4.\newline \hspace*{1em}Pearson Correlation Matrix.\textquotedbl{}) ///} & Exports the stored results (corr) into the Word table sum\_stats.rtf.\\
271 & \texttt{\hspace*{2.0em}star(* 0.1 ** 0.05 *** 0.01) b(\%6.2f) ///} & Continuation of the previous command: significance stars * 10\%, ** 5\%, *** 1\%.\\
272 & \texttt{\hspace*{2.0em}nomtitle nonumber noobs unstack compress nonotes not ///} & Continuation of the previous command: no default notes, constant omitted, compact layout; no model titles or numbers.\\
273 & \texttt{\hspace*{2.0em}note(\textquotedbl{}Notes: Panel A shows the descriptive sta ... \textquotedbl{})} & Continuation: the table note printed under the table (text abridged here).\\
274 & \texttt{restore} & Restores the data saved by preserve (undoes changes made since).\\
278 & \texttt{*-{}-{}-{}-{}-{}-{}-{}-{}-{}-{}-{}-{}-{}-{}-{}-{}-{}-{}-{}-{}-{}-{}-{}-{}-{}-{}-{}-{}-{}-{}-{}-{}-{}-{}-{}-{}-{}-{}-{}-{}-{}-{}-{}-{}-{}-{}-{}-{}-{}-{}-{}-{}-{}-{}-{}-{}-{}\newline \hspace*{1em}-{}-{}-{}-{}-{}-{}-{}-{}-{}-{}-{}-{}-{}-{}-{}} & Divider line (comment).\\
279 & \texttt{* Baseline Regression} & Comment (not executed): explains the code that follows.\\
280 & \texttt{*-{}-{}-{}-{}-{}-{}-{}-{}-{}-{}-{}-{}-{}-{}-{}-{}-{}-{}-{}-{}-{}-{}-{}-{}-{}-{}-{}-{}-{}-{}-{}-{}-{}-{}-{}-{}-{}-{}-{}-{}-{}-{}-{}-{}-{}-{}-{}-{}-{}-{}-{}-{}-{}-{}-{}-{}-{}\newline \hspace*{1em}-{}-{}-{}-{}-{}-{}-{}-{}-{}-{}-{}-{}-{}-{}-{}} & Divider line (comment).\\
281 & \texttt{preserve} & Saves a copy of the data in memory, so it can be restored after temporary changes.\\
282 & \texttt{reghdfe ShannonBD PI, absorb(county year) cluster(county)} & Regresses Shannon bird diversity on PSI (solar policy stringency); with county and year-month fixed effects; standard errors clustered by county. Table 1 col 1: -0.0157 (0.0037).\\
283 & \texttt{estadd local year \textquotedbl{}YES\textquotedbl{}} & Adds a "Year-month FE = YES" marker to the stored results (printed at the bottom of the table).\\
284 & \texttt{estadd local county \textquotedbl{}YES\textquotedbl{}} & Adds a "County FE = YES" marker to the stored results.\\
285 & \texttt{est store s1} & Saves the last estimates under the name s1 so they can be put in a table.\\
286 & \texttt{reghdfe ShannonBD PI Temp Wind Pop Duration Carbon,\newline \hspace*{1em}absorb(county year) cluster(county)} & Regresses Shannon bird diversity on PSI (solar policy stringency); controlling for Temp Wind Pop Duration Carbon; with county and year-month fixed effects; standard errors clustered by county. Table 1 col 2: -0.0127 (0.0036).\\
287 & \texttt{estadd local year \textquotedbl{}YES\textquotedbl{}} & Adds a "Year-month FE = YES" marker to the stored results (printed at the bottom of the table).\\
288 & \texttt{estadd local county \textquotedbl{}YES\textquotedbl{}} & Adds a "County FE = YES" marker to the stored results.\\
289 & \texttt{est store s2} & Saves the last estimates under the name s2 so they can be put in a table.\\
290 & \texttt{reghdfe ShannonBD PI Temp Wind Pop Duration Carbon Water\newline \hspace*{1em}Green Farm Grass, absorb(county year) cluster(county)} & Regresses Shannon bird diversity on PSI (solar policy stringency); controlling for Temp Wind Pop Duration Carbon Water Green Farm Grass; with county and year-month fixed effects; standard errors clustered by county. Table 1 col 3 = MAIN RESULT: -0.0125 (0.0037), N = 46,371, 2,344 counties.\\
291 & \texttt{estadd local year \textquotedbl{}YES\textquotedbl{}} & Adds a "Year-month FE = YES" marker to the stored results (printed at the bottom of the table).\\
292 & \texttt{estadd local county \textquotedbl{}YES\textquotedbl{}} & Adds a "County FE = YES" marker to the stored results.\\
293 & \texttt{est store s3} & Saves the last estimates under the name s3 so they can be put in a table.\\
294 & \texttt{esttab s1 s2 s3 using baseline.rtf, replace title(\textquotedbl{}Table\newline \hspace*{1em}1. Solar policy intensity significantly reduces bird\newline \hspace*{1em}diversity.\textquotedbl{}) ///} & Exports the stored results (s1 s2 s3) into the Word table baseline.rtf.\\
295 & \texttt{\hspace*{2.0em}stats(year county N r2\_a, labels(\textquotedbl{}Year-{}month FE\textquotedbl{} \textquotedbl{}County\newline \hspace*{1em}FE\textquotedbl{} \textquotedbl{}Observations\textquotedbl{} \textquotedbl{}R-{}squared\textquotedbl{}) fmt(0 0 0 4)) ///} & Continuation of the previous command: summary rows at the bottom (e.g. FE markers, observations, R-squared).\\
296 & \texttt{\hspace*{2.0em}star(* 0.10 ** 0.05 *** 0.01) nonotes nocons b(4) se\newline \hspace*{1em}compress nogaps ///} & Continuation of the previous command: significance stars * 10\%, ** 5\%, *** 1\%; 4-decimal coefficients with standard errors in parentheses; no default notes, constant omitted, compact layout.\\
297 & \texttt{\hspace*{2.0em}note(\textquotedbl{}Notes: This table presents the associati ... \textquotedbl{})} & Continuation: the table note printed under the table (text abridged here).\\
298 & \texttt{restore} & Restores the data saved by preserve (undoes changes made since).\\
302 & \texttt{*-{}-{}-{}-{}-{}-{}-{}-{}-{}-{}-{}-{}-{}-{}-{}-{}-{}-{}-{}-{}-{}-{}-{}-{}-{}-{}-{}-{}-{}-{}-{}-{}-{}-{}-{}-{}-{}-{}-{}-{}-{}-{}-{}-{}-{}-{}-{}-{}-{}-{}-{}-{}-{}-{}-{}-{}-{}\newline \hspace*{1em}-{}-{}-{}-{}-{}-{}-{}-{}-{}-{}-{}-{}-{}-{}-{}} & Divider line (comment).\\
303 & \texttt{* Endogeneity Treatment} & Comment (not executed): explains the code that follows.\\
304 & \texttt{*-{}-{}-{}-{}-{}-{}-{}-{}-{}-{}-{}-{}-{}-{}-{}-{}-{}-{}-{}-{}-{}-{}-{}-{}-{}-{}-{}-{}-{}-{}-{}-{}-{}-{}-{}-{}-{}-{}-{}-{}-{}-{}-{}-{}-{}-{}-{}-{}-{}-{}-{}-{}-{}-{}-{}-{}-{}\newline \hspace*{1em}-{}-{}-{}-{}-{}-{}-{}-{}-{}-{}-{}-{}-{}-{}-{}} & Divider line (comment).\\
305 & \texttt{*Instrumental variables regression} & Comment (not executed): explains the code that follows.\\
306 & \texttt{preserve} & Saves a copy of the data in memory, so it can be restored after temporary changes.\\
307 & \texttt{global control \textquotedbl{}Temp Wind Pop Duration Carbon Water Green\newline \hspace*{1em}Farm Grass\textquotedbl{}} & \$control = the authors' 9 control variables (Table 1, column 3).\\
308 & \texttt{reghdfe PI Sun\_ccpu \$control , absorb(county year)\newline \hspace*{1em}cluster(county) keepsingleton //Sun\_ccpu即Ln Sunshine /\newline \hspace*{1em}city CPU} & Regresses PSI (solar policy stringency) on the instrument ln(sunshine)/CPU; controlling for the 9 controls; with county and year-month fixed effects; standard errors clustered by county; keeping single-observation counties. Table S7 first stage: 0.0252 (0.0046).\\
309 & \texttt{estadd local year \textquotedbl{}YES\textquotedbl{}} & Adds a "Year-month FE = YES" marker to the stored results (printed at the bottom of the table).\\
310 & \texttt{estadd local county \textquotedbl{}YES\textquotedbl{}} & Adds a "County FE = YES" marker to the stored results.\\
311 & \texttt{est store s1} & Saves the last estimates under the name s1 so they can be put in a table.\\
312 & \texttt{ivreghdfe ShannonBD (PI = Sun\_ccpu) \$control ,\newline \hspace*{1em}absorb(county year) cluster(county) first keepsingleton} & Two-stage least squares (IV): Shannon bird diversity on PSI, with PSI instrumented by Sun\_ccpu; controls: the 9 controls; with county and year-month fixed effects; standard errors clustered by county; keeping single-observation counties; also reporting the first stage. Table S7 second stage: -0.1264 (0.0347); Kleibergen-Paap F = 30.18.\\
313 & \texttt{estadd local year \textquotedbl{}YES\textquotedbl{}} & Adds a "Year-month FE = YES" marker to the stored results (printed at the bottom of the table).\\
314 & \texttt{estadd local county \textquotedbl{}YES\textquotedbl{}} & Adds a "County FE = YES" marker to the stored results.\\
315 & \texttt{est store s2} & Saves the last estimates under the name s2 so they can be put in a table.\\
317 & \texttt{esttab s1 s2 using IV.rtf, replace title(\textquotedbl{}Table S7.\newline \hspace*{1em}Instrumental variable estimates.\textquotedbl{}) ///} & Exports the stored results (s1 s2) into the Word table IV.rtf.\\
318 & \texttt{\hspace*{2.0em}stats(year county N, labels(\textquotedbl{}Year-{}month FE\textquotedbl{} \textquotedbl{}County FE\textquotedbl{}\newline \hspace*{1em}\textquotedbl{}Observations\textquotedbl{}) fmt(0 0 0)) ///} & Continuation of the previous command: summary rows at the bottom (e.g. FE markers, observations, R-squared).\\
319 & \texttt{\hspace*{2.0em}star(* 0.10 ** 0.05 *** 0.01) nonotes nocons b(4) se\newline \hspace*{1em}compress nogaps ///} & Continuation of the previous command: significance stars * 10\%, ** 5\%, *** 1\%; 4-decimal coefficients with standard errors in parentheses; no default notes, constant omitted, compact layout.\\
320 & \texttt{\hspace*{2.0em}order(Sun\_ccpu PI \$control) ///} & Continuation of the previous command: row order.\\
321 & \texttt{\hspace*{2.0em}note(\textquotedbl{}Notes: Column (1) shows the result of th ... \textquotedbl{})} & Continuation: the table note printed under the table (text abridged here).\\
322 & \texttt{restore} & Restores the data saved by preserve (undoes changes made since).\\
325 & \texttt{*CPU Exogeneity Test (Point 2)} & Comment (not executed): explains the code that follows.\\
326 & \texttt{preserve} & Saves a copy of the data in memory, so it can be restored after temporary changes.\\
327 & \texttt{drop if missing(CCPU)} & Keeps rows with the CPU index.\\
328 & \texttt{collapse (mean) CCPU ShannonBD \$control, by(市 YEAR)} & Averages to city-year level.\\
329 & \texttt{egen city = group(市)} & Numeric city id.\\
330 & \texttt{*-{}-{}-{} 2a. Cross-{}sectional test -{}-{}-{}} & Comment (not executed): explains the code that follows.\\
331 & \texttt{tempfile citydata} & Temporary file.\\
332 & \texttt{save \textasciigrave{}citydata\textquotesingle{}} & Saves the city-year data.\\
333 & \texttt{collapse (mean) CCPU \$control ShannonBD, by(city)} & Averages to one row per city (cross-section).\\
334 & \texttt{reg CCPU Temp Wind Pop Duration Carbon Water Green Farm\newline \hspace*{1em}Grass, robust} & Table S8 col 1: do city characteristics predict average CPU?\\
335 & \texttt{est store cpu\_cross} & Stores it.\\
336 & \texttt{*-{}-{}-{} 2b. Panel test with city and year FE -{}-{}-{}} & Comment (not executed): explains the code that follows.\\
337 & \texttt{use \textasciigrave{}citydata\textquotesingle{}, clear} & Back to city-year data.\\
338 & \texttt{xtset city YEAR} & Declares the city-year panel.\\
339 & \texttt{reghdfe CCPU Temp Wind Pop Duration Carbon Water Green\newline \hspace*{1em}Farm Grass, absorb(city YEAR) cluster(city)} & Table S8 col 2: do changes in city characteristics predict changes in CPU (city and year FE)?\\
340 & \texttt{test Temp Wind Pop Duration Carbon Water Green Farm Grass} & Joint test that none of them predicts CPU.\\
341 & \texttt{estadd scalar joint\_F = r(F)} & Stores a number (here the joint F-test result) with the estimates for the table.\\
342 & \texttt{estadd scalar joint\_p = r(p)} & Stores a number (here the joint F-test result) with the estimates for the table.\\
343 & \texttt{estadd local year \textquotedbl{}YES\textquotedbl{}} & Adds a "Year-month FE = YES" marker to the stored results (printed at the bottom of the table).\\
344 & \texttt{estadd local city \textquotedbl{}YES\textquotedbl{}} & Adds a "City FE = YES" marker.\\
345 & \texttt{est store cpu\_panel} & Stores it.\\
347 & \texttt{esttab cpu\_cross cpu\_panel using CCPU\_exogeneity.rtf, ///} & Exports the stored results (cpu\_cross cpu\_panel) into the Word table CCPU\_exogeneity.rtf.\\
348 & \texttt{\hspace*{2.0em}replace title(\textquotedbl{}Table S8. Exogeneity test of the shift\newline \hspace*{1em}component (city CPU).\textquotedbl{}) ///} & Continuation of the previous command: table title.\\
349 & \texttt{\hspace*{2.0em}mtitles(\textquotedbl{}Cross-{}sectional OLS\textquotedbl{} \textquotedbl{}Panel with city and year\newline \hspace*{1em}FE\textquotedbl{}) ///} & Continuation of the previous command: column titles.\\
350 & \texttt{\hspace*{2.0em}stats(city year joint\_F joint\_p N r2 r2\_within, ///} & Continuation of the previous command: summary rows at the bottom (e.g. FE markers, observations, R-squared).\\
351 & \texttt{\hspace*{4.0em}labels(\textquotedbl{}City FE\textquotedbl{} \textquotedbl{}Year FE\textquotedbl{} \textquotedbl{}Joint F-{}statistic\textquotedbl{} \textquotedbl{}Joint F\newline \hspace*{1em}p-{}value\textquotedbl{} \textquotedbl{}Observations\textquotedbl{} \textquotedbl{}R-{}squared\textquotedbl{} \textquotedbl{}Within R-{}squared\textquotedbl{})\newline \hspace*{1em}///} & Continuation of the previous command: labels for the bottom rows.\\
352 & \texttt{\hspace*{4.0em}fmt(0 0 2 3 0 4 4)) ///} & Continuation of the previous command: number formats for the bottom rows.\\
353 & \texttt{\hspace*{2.0em}star(* 0.10 ** 0.05 *** 0.01) nonotes nocons b(4) se(4)\newline \hspace*{1em}compress nogaps ///} & Continuation of the previous command: significance stars * 10\%, ** 5\%, *** 1\%; 4-decimal coefficients with standard errors in parentheses; no default notes, constant omitted, compact layout.\\
354 & \texttt{\hspace*{2.0em}note(\textquotedbl{}Notes: This table tests whether the shif ... \textquotedbl{})} & Continuation: the table note printed under the table (text abridged here).\\
355 & \texttt{restore} & Restores the data saved by preserve (undoes changes made since).\\
358 & \texttt{*Shift-{}Share IV Diagnostics — GPSS (Point 4)} & Comment (not executed): explains the code that follows.\\
359 & \texttt{*-{}-{}-{} Overidentification test -{}-{}-{}} & Comment (not executed): explains the code that follows.\\
360 & \texttt{preserve} & Saves a copy of the data in memory, so it can be restored after temporary changes.\\
361 & \texttt{drop if missing(lnsun) | missing(CCPU)} & Keeps rows with both instrument parts.\\
362 & \texttt{xtile ccpu\_tercile = CCPU, nq(3)} & Splits CPU into three equal-sized groups (terciles).\\
363 & \texttt{gen iv\_ccpu\_t1 = Sun\_ccpu * (ccpu\_tercile == 1)} & Instrument active only in the low-CPU tercile ...\\
364 & \texttt{gen iv\_ccpu\_t2 = Sun\_ccpu * (ccpu\_tercile == 2)} & ... middle tercile ...\\
365 & \texttt{gen iv\_ccpu\_t3 = Sun\_ccpu * (ccpu\_tercile == 3)} & ... high tercile: three instruments for an over-identification test.\\
367 & \texttt{ivreghdfe ShannonBD (PI = iv\_ccpu\_t1 iv\_ccpu\_t2\newline \hspace*{1em}iv\_ccpu\_t3) \$control, absorb(county year) cluster(county)\newline \hspace*{1em}keepsingleton} & IV with the three tercile instruments; reports Hansen J (0.6356). Note: J tests whether the three versions of the same instrument agree, not whether the instrument is valid. Hansen J = 0.6356 (not rejected).\\
368 & \texttt{estadd local year \textquotedbl{}YES\textquotedbl{}} & Adds a "Year-month FE = YES" marker to the stored results (printed at the bottom of the table).\\
369 & \texttt{estadd local county \textquotedbl{}YES\textquotedbl{}} & Adds a "County FE = YES" marker to the stored results.\\
370 & \texttt{estadd local sample \textquotedbl{}Over-{}identification\textquotedbl{}} & Adds a text marker to the stored results for the table footer.\\
371 & \texttt{est store overid} & Stores it.\\
373 & \texttt{di \textquotedbl{}==============================================\textquotedbl{}} & Prints a divider.\\
374 & \texttt{di \textquotedbl{}Overidentification Test (Hansen J)\textquotedbl{}} & Prints a heading.\\
375 & \texttt{di \textquotedbl{}Hansen J statistic: \textquotedbl{} e(j)} & Prints J.\\
376 & \texttt{di \textquotedbl{}Hansen J p-{}value:   \textquotedbl{} e(jp)} & Prints its p-value.\\
377 & \texttt{di \textquotedbl{}==============================================\textquotedbl{}} & Prints a divider.\\
378 & \texttt{restore} & Restores the data saved by preserve (undoes changes made since).\\
379 & \texttt{*-{}-{}-{} Sensitivity analysis -{}-{}-{}} & Comment (not executed): explains the code that follows.\\
380 & \texttt{preserve} & Saves a copy of the data in memory, so it can be restored after temporary changes.\\
381 & \texttt{bysort county: egen county\_sun = mean(lnsun)} & County average ln(sunshine).\\
382 & \texttt{quietly sum county\_sun, detail} & Percentiles of county sunshine.\\
383 & \texttt{scalar p5  = r(p5)} & 5th percentile.\\
384 & \texttt{scalar p10 = r(p10)} & 10th percentile.\\
385 & \texttt{scalar p90 = r(p90)} & 90th percentile.\\
386 & \texttt{scalar p95 = r(p95)} & 95th percentile.\\
388 & \texttt{gen extreme\_sun = (county\_sun < p5 | county\_sun > p95)} & 1 for the sunniest and least sunny 5\% of counties.\\
389 & \texttt{gen high\_sun    = (county\_sun > p90)} & 1 for the top 10\%.\\
390 & \texttt{gen low\_sun     = (county\_sun < p10)} & 1 for the bottom 10\%.\\
392 & \texttt{ivreghdfe ShannonBD (PI = Sun\_ccpu) \$control,\newline \hspace*{1em}absorb(county year) cluster(county) keepsingleton} & Two-stage least squares (IV): Shannon bird diversity on PSI, with PSI instrumented by Sun\_ccpu; controls: the 9 controls; with county and year-month fixed effects; standard errors clustered by county; keeping single-observation counties.\\
393 & \texttt{estadd local year \textquotedbl{}YES\textquotedbl{}} & Adds a "Year-month FE = YES" marker to the stored results (printed at the bottom of the table).\\
394 & \texttt{estadd local county \textquotedbl{}YES\textquotedbl{}} & Adds a "County FE = YES" marker to the stored results.\\
395 & \texttt{estadd local controls \textquotedbl{}YES\textquotedbl{}} & Adds a text marker to the stored results for the table footer.\\
396 & \texttt{estadd local sample \textquotedbl{}Full sample\textquotedbl{}} & Adds a text marker to the stored results for the table footer.\\
397 & \texttt{est store sens\_full} & Saves the last estimates under the name sens\_full so they can be put in a table.\\
399 & \texttt{ivreghdfe ShannonBD (PI = Sun\_ccpu) \$control if\newline \hspace*{1em}extreme\_sun == 0, ///} & Two-stage least squares (IV): Shannon bird diversity on PSI, with PSI instrumented by Sun\_ccpu; controls: the 9 controls; excluding counties with extreme (top/bottom 5\%) sunshine.\\
400 & \texttt{\hspace*{2.0em}absorb(county year) cluster(county)} & Continuation of the previous command: estimation options (absorb(county year) cluster(county)).\\
401 & \texttt{estadd local year \textquotedbl{}YES\textquotedbl{}} & Adds a "Year-month FE = YES" marker to the stored results (printed at the bottom of the table).\\
402 & \texttt{estadd local county \textquotedbl{}YES\textquotedbl{}} & Adds a "County FE = YES" marker to the stored results.\\
403 & \texttt{estadd local controls \textquotedbl{}YES\textquotedbl{}} & Adds a text marker to the stored results for the table footer.\\
404 & \texttt{estadd local sample \textquotedbl{}Excl. extreme sunshine\textquotedbl{}} & Adds a text marker to the stored results for the table footer.\\
405 & \texttt{est store sens\_drop} & Saves the last estimates under the name sens\_drop so they can be put in a table.\\
407 & \texttt{ivreghdfe ShannonBD (PI = Sun\_ccpu) \$control if high\_sun\newline \hspace*{1em}== 0, ///} & Two-stage least squares (IV): Shannon bird diversity on PSI, with PSI instrumented by Sun\_ccpu; controls: the 9 controls; excluding the top 10\% sunshine counties.\\
408 & \texttt{\hspace*{2.0em}absorb(county year) cluster(county)} & Continuation of the previous command: estimation options (absorb(county year) cluster(county)).\\
409 & \texttt{estadd local year \textquotedbl{}YES\textquotedbl{}} & Adds a "Year-month FE = YES" marker to the stored results (printed at the bottom of the table).\\
410 & \texttt{estadd local county \textquotedbl{}YES\textquotedbl{}} & Adds a "County FE = YES" marker to the stored results.\\
411 & \texttt{estadd local controls \textquotedbl{}YES\textquotedbl{}} & Adds a text marker to the stored results for the table footer.\\
412 & \texttt{estadd local sample \textquotedbl{}Excl. top 10\% sunshine\textquotedbl{}} & Adds a text marker to the stored results for the table footer.\\
413 & \texttt{est store sens\_drop10} & Saves the last estimates under the name sens\_drop10 so they can be put in a table.\\
415 & \texttt{ivreghdfe ShannonBD (PI = Sun\_ccpu) \$control if low\_sun ==\newline \hspace*{1em}0, ///} & Two-stage least squares (IV): Shannon bird diversity on PSI, with PSI instrumented by Sun\_ccpu; controls: the 9 controls; excluding the bottom 10\% sunshine counties.\\
416 & \texttt{\hspace*{2.0em}absorb(county year) cluster(county)} & Continuation of the previous command: estimation options (absorb(county year) cluster(county)).\\
417 & \texttt{estadd local year \textquotedbl{}YES\textquotedbl{}} & Adds a "Year-month FE = YES" marker to the stored results (printed at the bottom of the table).\\
418 & \texttt{estadd local county \textquotedbl{}YES\textquotedbl{}} & Adds a "County FE = YES" marker to the stored results.\\
419 & \texttt{estadd local controls \textquotedbl{}YES\textquotedbl{}} & Adds a text marker to the stored results for the table footer.\\
420 & \texttt{estadd local sample \textquotedbl{}Excl. bottom 10\% sunshine\textquotedbl{}} & Adds a text marker to the stored results for the table footer.\\
421 & \texttt{est store sens\_drop10low} & Saves the last estimates under the name sens\_drop10low so they can be put in a table.\\
422 & \texttt{restore} & Restores the data saved by preserve (undoes changes made since).\\
423 & \texttt{*-{}-{}-{} Share exogeneity test (balance test) -{}-{}-{}} & Comment (not executed): explains the code that follows.\\
424 & \texttt{preserve} & Saves a copy of the data in memory, so it can be restored after temporary changes.\\
425 & \texttt{collapse (mean) lnsun Temp Wind Pop Duration Carbon Water\newline \hspace*{1em}Green Farm Grass ShannonBD PI, by(county)} & One row per county (cross-section).\\
427 & \texttt{reg lnsun Temp Wind Pop Duration Carbon Water Green Farm\newline \hspace*{1em}Grass, robust} & Table S10 col 1: is sunshine related to county characteristics (share balance test)?\\
428 & \texttt{estadd local note \textquotedbl{}Cross-{}sectional balance test\textquotedbl{}} & Adds a note marker.\\
429 & \texttt{est store balance} & Stores it.\\
430 & \texttt{restore} & Restores the data saved by preserve (undoes changes made since).\\
431 & \texttt{*-{}-{}-{} Shift-{}based IV: controlling for share time-{}varying\newline \hspace*{1em}effects -{}-{}-{}} & Comment (not executed): explains the code that follows.\\
432 & \texttt{* Panel A: IV +  Ln(Sunshine) × Year-{}month FE} & Comment (not executed): explains the code that follows.\\
433 & \texttt{preserve} & Saves a copy of the data in memory, so it can be restored after temporary changes.\\
434 & \texttt{reghdfe PI Sun\_ccpu \$control, absorb(county year\newline \hspace*{1em}i.year\#c.lnsun) cluster(county) keepsingleton} & First stage controlling for a separate ln(sunshine) slope in every year-month (removes any time-varying direct effect of sunshine).\\
435 & \texttt{estadd local yearfe \textquotedbl{}YES\textquotedbl{}} & Adds a "Year-month FE = YES" marker to the stored results (printed at the bottom of the table).\\
436 & \texttt{estadd local countyfe \textquotedbl{}YES\textquotedbl{}} & Adds a "County FE = YES" marker to the stored results.\\
437 & \texttt{estadd local sun\_year \textquotedbl{}YES\textquotedbl{}} & Adds a text marker to the stored results for the table footer.\\
438 & \texttt{est store shift\_s1a} & Saves the last estimates under the name shift\_s1a so they can be put in a table.\\
440 & \texttt{ivreghdfe ShannonBD (PI = Sun\_ccpu) \$control,\newline \hspace*{1em}absorb(county year i.year\#c.lnsun) cluster(county) first\newline \hspace*{1em}keepsingleton} & IV with the same extra control.\\
441 & \texttt{estadd local yearfe \textquotedbl{}YES\textquotedbl{}} & Adds a "Year-month FE = YES" marker to the stored results (printed at the bottom of the table).\\
442 & \texttt{estadd local countyfe \textquotedbl{}YES\textquotedbl{}} & Adds a "County FE = YES" marker to the stored results.\\
443 & \texttt{estadd local sun\_year \textquotedbl{}YES\textquotedbl{}} & Adds a text marker to the stored results for the table footer.\\
444 & \texttt{est store shift\_s2a} & Saves the last estimates under the name shift\_s2a so they can be put in a table.\\
445 & \texttt{* Panel B: only shift IV = 1/CPU} & Comment (not executed): explains the code that follows.\\
446 & \texttt{reghdfe PI inv\_CCPU \$control, absorb(county year)\newline \hspace*{1em}cluster(county)} & First stage using only 1/CPU (the shift) as instrument.\\
447 & \texttt{estadd local yearfe \textquotedbl{}YES\textquotedbl{}} & Adds a "Year-month FE = YES" marker to the stored results (printed at the bottom of the table).\\
448 & \texttt{estadd local countyfe \textquotedbl{}YES\textquotedbl{}} & Adds a "County FE = YES" marker to the stored results.\\
449 & \texttt{estadd local sun\_year \textquotedbl{}NO\textquotedbl{}} & Adds a text marker to the stored results for the table footer.\\
450 & \texttt{est store shift\_s1b} & Saves the last estimates under the name shift\_s1b so they can be put in a table.\\
452 & \texttt{ivreghdfe ShannonBD (PI = inv\_CCPU) \$control,\newline \hspace*{1em}absorb(county year) cluster(county) first} & IV using only 1/CPU.\\
453 & \texttt{estadd local yearfe \textquotedbl{}YES\textquotedbl{}} & Adds a "Year-month FE = YES" marker to the stored results (printed at the bottom of the table).\\
454 & \texttt{estadd local countyfe \textquotedbl{}YES\textquotedbl{}} & Adds a "County FE = YES" marker to the stored results.\\
455 & \texttt{estadd local sun\_year \textquotedbl{}NO\textquotedbl{}} & Adds a text marker to the stored results for the table footer.\\
456 & \texttt{est store shift\_s2b} & Saves the last estimates under the name shift\_s2b so they can be put in a table.\\
457 & \texttt{restore} & Restores the data saved by preserve (undoes changes made since).\\
458 & \texttt{* Output Tables：Overidentification + Sensitivity} & Comment (not executed): explains the code that follows.\\
459 & \texttt{esttab overid sens\_full sens\_drop sens\_drop10\newline \hspace*{1em}sens\_drop10low using SSIV\_table1.rtf, ///} & Exports the stored results (overid sens\_full sens\_drop sens\_drop10 sens\_drop10low) into the Word table SSIV\_table1.rtf.\\
460 & \texttt{\hspace*{2.0em}replace title(\textquotedbl{}Table S9. Shift-{}share IV diagnostics:\newline \hspace*{1em}Overidentification and sensitivity analysis.\textquotedbl{}) ///} & Continuation of the previous command: table title.\\
461 & \texttt{\hspace*{2.0em}mtitles(\textquotedbl{}Over-{}identification\textquotedbl{} \textquotedbl{}Full sample\textquotedbl{} \textquotedbl{}Excl. extreme\newline \hspace*{1em}5\%\textquotedbl{} \textquotedbl{}Excl. top 10\%\textquotedbl{} \textquotedbl{}Excl. bottom 10\%\textquotedbl{}) ///} & Continuation of the previous command: column titles.\\
462 & \texttt{\hspace*{2.0em}stats(year county sample N j jp, ///} & Continuation of the previous command: summary rows at the bottom (e.g. FE markers, observations, R-squared).\\
463 & \texttt{\hspace*{4.0em}labels(\textquotedbl{}Year-{}month FE\textquotedbl{} \textquotedbl{}County FE\textquotedbl{} \textquotedbl{}Sample restriction\textquotedbl{}\newline \hspace*{1em}\textquotedbl{}Observations\textquotedbl{} \textquotedbl{}Hansen J statistic\textquotedbl{} \textquotedbl{}Hansen J p-{}value\textquotedbl{})\newline \hspace*{1em}///} & Continuation of the previous command: labels for the bottom rows.\\
464 & \texttt{\hspace*{4.0em}fmt(0 0 0 0 4 4)) ///} & Continuation of the previous command: number formats for the bottom rows.\\
465 & \texttt{\hspace*{2.0em}star(* 0.10 ** 0.05 *** 0.01) nonotes nocons b(4) se\newline \hspace*{1em}compress nogaps ///} & Continuation of the previous command: significance stars * 10\%, ** 5\%, *** 1\%; 4-decimal coefficients with standard errors in parentheses; no default notes, constant omitted, compact layout.\\
466 & \texttt{\hspace*{2.0em}order(PI Temp Wind Pop Duration Carbon Water Green Farm\newline \hspace*{1em}Grass) ///} & Continuation of the previous command: row order.\\
467 & \texttt{\hspace*{2.0em}note(\textquotedbl{}Notes: This table reports the shift-{}shar ... \textquotedbl{})} & Continuation: the table note printed under the table (text abridged here).\\
468 & \texttt{* Output Tables：Balance test + Shift-{}based IV\newline \hspace*{1em}identification tests} & Comment (not executed): explains the code that follows.\\
469 & \texttt{esttab balance shift\_s1a shift\_s2a shift\_s1b shift\_s2b\newline \hspace*{1em}using SSIV\_table2.rtf, ///} & Exports the stored results (balance shift\_s1a shift\_s2a shift\_s1b shift\_s2b) into the Word table SSIV\_table2.rtf.\\
470 & \texttt{\hspace*{2.0em}replace title(\textquotedbl{}Table S10. Share exogeneity and shift-{}based\newline \hspace*{1em}IV identification tests.\textquotedbl{}) ///} & Continuation of the previous command: table title.\\
471 & \texttt{\hspace*{2.0em}mtitles(\textquotedbl{}Share exogeneity\textquotedbl{} \textquotedbl{}First stage\textquotedbl{} \textquotedbl{}Second stage\textquotedbl{}\newline \hspace*{1em}\textquotedbl{}First stage\textquotedbl{} \textquotedbl{}Second stage\textquotedbl{}) ///} & Continuation of the previous command: column titles.\\
472 & \texttt{\hspace*{2.0em}mgroups(\textquotedbl{}Balance test\textquotedbl{} \textquotedbl{}Panel A: [Ln(Sunshine) / city CPU]\newline \hspace*{1em}AND [Ln(Sunshine) x Year-{}month FE]\textquotedbl{} \textquotedbl{}Panel B: 1 / city CPU\newline \hspace*{1em}only\textquotedbl{}, pattern(1 1 0 1 0)) ///} & Continuation of the previous command: column-group headings.\\
473 & \texttt{\hspace*{2.0em}stats(N r2 F yearfe countyfe sun\_year, ///} & Continuation of the previous command: summary rows at the bottom (e.g. FE markers, observations, R-squared).\\
474 & \texttt{\hspace*{2.0em}labels(\textquotedbl{}Observations\textquotedbl{} \textquotedbl{}R-{}squared\textquotedbl{} \textquotedbl{}F-{}statistic\textquotedbl{}\newline \hspace*{1em}\textquotedbl{}Year-{}month FE\textquotedbl{} \textquotedbl{}County FE\textquotedbl{} \textquotedbl{}Ln(Sunshine) x Year-{}month\newline \hspace*{1em}FE\textquotedbl{}) ///} & Continuation of the previous command: labels for the bottom rows.\\
475 & \texttt{\hspace*{2.0em}fmt(0 4 2 0 0 0)) ///} & Continuation of the previous command: number formats for the bottom rows.\\
476 & \texttt{\hspace*{2.0em}star(* 0.10 ** 0.05 *** 0.01) nonotes nocons b(4) se(4)\newline \hspace*{1em}compress nogaps ///} & Continuation of the previous command: significance stars * 10\%, ** 5\%, *** 1\%; 4-decimal coefficients with standard errors in parentheses; no default notes, constant omitted, compact layout.\\
477 & \texttt{\hspace*{2.0em}order(Sun\_ccpu inv\_CCPU PI \$control) ///} & Continuation of the previous command: row order.\\
478 & \texttt{\hspace*{2.0em}keep(Sun\_ccpu inv\_CCPU PI \$control) ///} & Continuation of the previous command: which coefficients to show.\\
479 & \texttt{\hspace*{2.0em}note(\textquotedbl{}Notes: This table reports the share exog ... \textquotedbl{})} & Continuation: the table note printed under the table (text abridged here).\\
482 & \texttt{*psm} & Comment (not executed): explains the code that follows.\\
483 & \texttt{preserve} & Saves a copy of the data in memory, so it can be restored after temporary changes.\\
484 & \texttt{global control \textquotedbl{}Temp Wind Pop Duration Carbon Water Green\newline \hspace*{1em}Farm Grass\textquotedbl{}} & \$control = the authors' 9 control variables (Table 1, column 3).\\
485 & \texttt{gen treated = (PI > 0)} & Creates the treatment dummy for matching: 1 if PSI > 0, else 0.\\
486 & \texttt{**1:1 邻近匹配} & Comment: "1:1 nearest-neighbour matching".\\
487 & \texttt{psmatch2 treated \$control i.year, outcome(ShannonBD)\newline \hspace*{1em}neighbor(1) caliper(0.05) common} & Propensity-score matching: estimates the probability of PSI > 0 from the controls and year-month dummies, then matches each treated row to its 1 nearest untreated neighbour(s) within a caliper of 0.05; common = keep only common support. Creates \_weight (matching weights).\\
488 & \texttt{reghdfe ShannonBD PI \$control if \_weight != .,\newline \hspace*{1em}absorb(county year) cluster(county)} & Regresses Shannon bird diversity on PSI (solar policy stringency); controlling for the 9 controls; with county and year-month fixed effects; standard errors clustered by county; only observations kept by the matching. Table S11 col 1: -0.0072. PROBLEM: \_weight is only used to drop unmatched rows; the matching weights are never applied. With [aw=\_weight]: -0.0051, not significant.\\
489 & \texttt{estadd local year \textquotedbl{}YES\textquotedbl{}} & Adds a "Year-month FE = YES" marker to the stored results (printed at the bottom of the table).\\
490 & \texttt{estadd local county \textquotedbl{}YES\textquotedbl{}} & Adds a "County FE = YES" marker to the stored results.\\
491 & \texttt{est store s1} & Saves the last estimates under the name s1 so they can be put in a table.\\
492 & \texttt{**重新匹配，1:2 邻近匹配} & Comment: "re-match, 1:2 nearest-neighbour matching".\\
493 & \texttt{psmatch2 treated \$control i.year, outcome(ShannonBD)\newline \hspace*{1em}neighbor(2) caliper(0.05) common} & Propensity-score matching: estimates the probability of PSI > 0 from the controls and year-month dummies, then matches each treated row to its 2 nearest untreated neighbour(s) within a caliper of 0.05; common = keep only common support. Creates \_weight (matching weights).\\
494 & \texttt{reghdfe ShannonBD PI \$control if \_weight != .,\newline \hspace*{1em}absorb(county year) cluster(county)} & Regresses Shannon bird diversity on PSI (solar policy stringency); controlling for the 9 controls; with county and year-month fixed effects; standard errors clustered by county; only observations kept by the matching. Table S11 col 2: -0.0083. Same problem; weighted: -0.0052.\\
495 & \texttt{estadd local year \textquotedbl{}YES\textquotedbl{}} & Adds a "Year-month FE = YES" marker to the stored results (printed at the bottom of the table).\\
496 & \texttt{estadd local county \textquotedbl{}YES\textquotedbl{}} & Adds a "County FE = YES" marker to the stored results.\\
497 & \texttt{est store s2} & Saves the last estimates under the name s2 so they can be put in a table.\\
498 & \texttt{**核匹配} & Comment: "kernel matching".\\
499 & \texttt{psmatch2 treated \$control i.year, outcome(ShannonBD)\newline \hspace*{1em}kernel common} & Kernel propensity-score matching: estimates each row's probability of PSI > 0 from the controls and year-month dummies, then weights all untreated rows by closeness; common = keep only common support. Creates \_weight.\\
500 & \texttt{reghdfe ShannonBD PI \$control if \_weight != .,\newline \hspace*{1em}absorb(county year) cluster(county)} & Regresses Shannon bird diversity on PSI (solar policy stringency); controlling for the 9 controls; with county and year-month fixed effects; standard errors clustered by county; only observations kept by the matching. Table S11 col 3: -0.0125. Same problem; weighted: -0.0054.\\
501 & \texttt{estadd local year \textquotedbl{}YES\textquotedbl{}} & Adds a "Year-month FE = YES" marker to the stored results (printed at the bottom of the table).\\
502 & \texttt{estadd local county \textquotedbl{}YES\textquotedbl{}} & Adds a "County FE = YES" marker to the stored results.\\
503 & \texttt{est store s3} & Saves the last estimates under the name s3 so they can be put in a table.\\
504 & \texttt{esttab s1 s2 s3 using psm.rtf, replace title(\textquotedbl{}Table S11.\newline \hspace*{1em}Propensity score matching.\textquotedbl{}) ///} & Exports the stored results (s1 s2 s3) into the Word table psm.rtf.\\
505 & \texttt{\hspace*{2.0em}stats(year county N r2\_a, labels(\textquotedbl{}Year-{}month FE\textquotedbl{} \textquotedbl{}County\newline \hspace*{1em}FE\textquotedbl{} \textquotedbl{}Observations\textquotedbl{} \textquotedbl{}R-{}squared\textquotedbl{}) fmt(0 0 0 4)) ///} & Continuation of the previous command: summary rows at the bottom (e.g. FE markers, observations, R-squared).\\
506 & \texttt{\hspace*{2.0em}star(* 0.10 ** 0.05 *** 0.01) nonotes nocons b(4) se\newline \hspace*{1em}compress nogaps ///} & Continuation of the previous command: significance stars * 10\%, ** 5\%, *** 1\%; 4-decimal coefficients with standard errors in parentheses; no default notes, constant omitted, compact layout.\\
507 & \texttt{\hspace*{2.0em}note(\textquotedbl{}Notes: Column (1) presents the PSM resul ... \textquotedbl{})} & Continuation: the table note printed under the table (text abridged here).\\
508 & \texttt{restore} & Restores the data saved by preserve (undoes changes made since).\\
512 & \texttt{*-{}-{}-{}-{}-{}-{}-{}-{}-{}-{}-{}-{}-{}-{}-{}-{}-{}-{}-{}-{}-{}-{}-{}-{}-{}-{}-{}-{}-{}-{}-{}-{}-{}-{}-{}-{}-{}-{}-{}-{}-{}-{}-{}-{}-{}-{}-{}-{}-{}-{}-{}-{}-{}-{}-{}-{}-{}\newline \hspace*{1em}-{}-{}-{}-{}-{}-{}-{}-{}-{}-{}-{}-{}-{}-{}-{}} & Divider line (comment).\\
513 & \texttt{* Robustness Check} & Comment (not executed): explains the code that follows.\\
514 & \texttt{*-{}-{}-{}-{}-{}-{}-{}-{}-{}-{}-{}-{}-{}-{}-{}-{}-{}-{}-{}-{}-{}-{}-{}-{}-{}-{}-{}-{}-{}-{}-{}-{}-{}-{}-{}-{}-{}-{}-{}-{}-{}-{}-{}-{}-{}-{}-{}-{}-{}-{}-{}-{}-{}-{}-{}-{}-{}\newline \hspace*{1em}-{}-{}-{}-{}-{}-{}-{}-{}-{}-{}-{}-{}-{}-{}-{}} & Divider line (comment).\\
515 & \texttt{*Alternative variables} & Comment (not executed): explains the code that follows.\\
516 & \texttt{preserve} & Saves a copy of the data in memory, so it can be restored after temporary changes.\\
517 & \texttt{global control \textquotedbl{}Temp Wind Pop Duration Carbon Water Green\newline \hspace*{1em}Farm Grass\textquotedbl{}} & \$control = the authors' 9 control variables (Table 1, column 3).\\
518 & \texttt{reghdfe ShannonBD Area \$control , absorb(county year)\newline \hspace*{1em}cluster(county)} & Regresses Shannon bird diversity on PV station area; controlling for the 9 controls; with county and year-month fixed effects; standard errors clustered by county. Table S12: -0.0387.\\
519 & \texttt{estadd local year \textquotedbl{}YES\textquotedbl{}} & Adds a "Year-month FE = YES" marker to the stored results (printed at the bottom of the table).\\
520 & \texttt{estadd local county \textquotedbl{}YES\textquotedbl{}} & Adds a "County FE = YES" marker to the stored results.\\
521 & \texttt{est store s1} & Saves the last estimates under the name s1 so they can be put in a table.\\
522 & \texttt{reghdfe ShannonBD PI1 \$control , absorb(county year)\newline \hspace*{1em}cluster(county)} & Regresses Shannon bird diversity on alternative PSI 1-2-3 weighted; controlling for the 9 controls; with county and year-month fixed effects; standard errors clustered by county.\\
523 & \texttt{estadd local year \textquotedbl{}YES\textquotedbl{}} & Adds a "Year-month FE = YES" marker to the stored results (printed at the bottom of the table).\\
524 & \texttt{estadd local county \textquotedbl{}YES\textquotedbl{}} & Adds a "County FE = YES" marker to the stored results.\\
525 & \texttt{est store s2} & Saves the last estimates under the name s2 so they can be put in a table.\\
526 & \texttt{reghdfe ShannonBD PI2 \$control , absorb(county year)\newline \hspace*{1em}cluster(county)} & Regresses Shannon bird diversity on alternative PSI 1-2-4 weighted; controlling for the 9 controls; with county and year-month fixed effects; standard errors clustered by county.\\
527 & \texttt{estadd local year \textquotedbl{}YES\textquotedbl{}} & Adds a "Year-month FE = YES" marker to the stored results (printed at the bottom of the table).\\
528 & \texttt{estadd local county \textquotedbl{}YES\textquotedbl{}} & Adds a "County FE = YES" marker to the stored results.\\
529 & \texttt{est store s3} & Saves the last estimates under the name s3 so they can be put in a table.\\
530 & \texttt{reghdfe ShannonBD PI3 \$control , absorb(county year)\newline \hspace*{1em}cluster(county)} & Regresses Shannon bird diversity on alternative PSI 2-4-8 weighted; controlling for the 9 controls; with county and year-month fixed effects; standard errors clustered by county.\\
531 & \texttt{estadd local year \textquotedbl{}YES\textquotedbl{}} & Adds a "Year-month FE = YES" marker to the stored results (printed at the bottom of the table).\\
532 & \texttt{estadd local county \textquotedbl{}YES\textquotedbl{}} & Adds a "County FE = YES" marker to the stored results.\\
533 & \texttt{est store s4} & Saves the last estimates under the name s4 so they can be put in a table.\\
534 & \texttt{reghdfe ShannonBD PI4 \$control , absorb(county year)\newline \hspace*{1em}cluster(county)} & Regresses Shannon bird diversity on alternative PSI 1-2-3 unweighted; controlling for the 9 controls; with county and year-month fixed effects; standard errors clustered by county.\\
535 & \texttt{estadd local year \textquotedbl{}YES\textquotedbl{}} & Adds a "Year-month FE = YES" marker to the stored results (printed at the bottom of the table).\\
536 & \texttt{estadd local county \textquotedbl{}YES\textquotedbl{}} & Adds a "County FE = YES" marker to the stored results.\\
537 & \texttt{est store s5} & Saves the last estimates under the name s5 so they can be put in a table.\\
538 & \texttt{reghdfe ShannonBD PI5 \$control , absorb(county year)\newline \hspace*{1em}cluster(county)} & Regresses Shannon bird diversity on alternative PSI 1-2-4 unweighted; controlling for the 9 controls; with county and year-month fixed effects; standard errors clustered by county.\\
539 & \texttt{estadd local year \textquotedbl{}YES\textquotedbl{}} & Adds a "Year-month FE = YES" marker to the stored results (printed at the bottom of the table).\\
540 & \texttt{estadd local county \textquotedbl{}YES\textquotedbl{}} & Adds a "County FE = YES" marker to the stored results.\\
541 & \texttt{est store s6} & Saves the last estimates under the name s6 so they can be put in a table.\\
542 & \texttt{reghdfe ShannonBD PI6 \$control , absorb(county year)\newline \hspace*{1em}cluster(county)} & Regresses Shannon bird diversity on alternative PSI 2-4-8 unweighted; controlling for the 9 controls; with county and year-month fixed effects; standard errors clustered by county.\\
543 & \texttt{estadd local year \textquotedbl{}YES\textquotedbl{}} & Adds a "Year-month FE = YES" marker to the stored results (printed at the bottom of the table).\\
544 & \texttt{estadd local county \textquotedbl{}YES\textquotedbl{}} & Adds a "County FE = YES" marker to the stored results.\\
545 & \texttt{est store s7} & Saves the last estimates under the name s7 so they can be put in a table.\\
546 & \texttt{reghdfe SimpsonBD PI \$control , absorb(county year)\newline \hspace*{1em}cluster(county)} & Regresses Simpson bird diversity on PSI (solar policy stringency); controlling for the 9 controls; with county and year-month fixed effects; standard errors clustered by county. Table S12 col 8: -0.0014 (0.0008).\\
547 & \texttt{estadd local year \textquotedbl{}YES\textquotedbl{}} & Adds a "Year-month FE = YES" marker to the stored results (printed at the bottom of the table).\\
548 & \texttt{estadd local county \textquotedbl{}YES\textquotedbl{}} & Adds a "County FE = YES" marker to the stored results.\\
549 & \texttt{est store s8} & Saves the last estimates under the name s8 so they can be put in a table.\\
550 & \texttt{esttab s1 s2 s3 s4 s5 s6 s7 s8 using\newline \hspace*{1em}robust\_alternative\_var.rtf, replace title(\textquotedbl{}Table S12.\newline \hspace*{1em}Robustness checks using alternative variable measures.\textquotedbl{})\newline \hspace*{1em}///} & Exports the stored results (s1 s2 s3 s4 s5 s6 s7 s8) into the Word table robust\_alternative\_var.rtf.\\
551 & \texttt{\hspace*{2.0em}mtitles(\textquotedbl{}Central PV Areas\textquotedbl{} \textquotedbl{}Weighted Score 123\textquotedbl{} \textquotedbl{}Weighted\newline \hspace*{1em}Score 124\textquotedbl{} \textquotedbl{}Weighted Score 248\textquotedbl{} \textquotedbl{}Unweighted Score 123\textquotedbl{}\newline \hspace*{1em}\textquotedbl{}Unweighted Score 124\textquotedbl{} \textquotedbl{}Unweighted Score 248\textquotedbl{} \textquotedbl{}SimpsonBD\textquotedbl{})\newline \hspace*{1em}///} & Continuation of the previous command: column titles.\\
552 & \texttt{\hspace*{2.0em}stats(year county N r2\_a, labels(\textquotedbl{}Year-{}month FE\textquotedbl{} \textquotedbl{}County\newline \hspace*{1em}FE\textquotedbl{} \textquotedbl{}Observations\textquotedbl{} \textquotedbl{}R-{}squared\textquotedbl{}) fmt(0 0 0 4)) ///} & Continuation of the previous command: summary rows at the bottom (e.g. FE markers, observations, R-squared).\\
553 & \texttt{\hspace*{2.0em}star(* 0.10 ** 0.05 *** 0.01) nonotes nocons b(4) se\newline \hspace*{1em}compress nogaps ///} & Continuation of the previous command: significance stars * 10\%, ** 5\%, *** 1\%; 4-decimal coefficients with standard errors in parentheses; no default notes, constant omitted, compact layout.\\
554 & \texttt{\hspace*{2.0em}order(Area PI1 PI2 PI3 PI4 PI5 PI6 PI) ///} & Continuation of the previous command: row order.\\
555 & \texttt{\hspace*{2.0em}note(\textquotedbl{}Notes: This table presents the regressio ... \textquotedbl{})} & Continuation: the table note printed under the table (text abridged here).\\
556 & \texttt{restore} & Restores the data saved by preserve (undoes changes made since).\\
559 & \texttt{*Sample interval adjustment} & Comment (not executed): explains the code that follows.\\
560 & \texttt{preserve} & Saves a copy of the data in memory, so it can be restored after temporary changes.\\
561 & \texttt{xtset county year} & Declares the county x year-month panel.\\
562 & \texttt{global control \textquotedbl{}Temp Wind Pop Duration Carbon Water Green\newline \hspace*{1em}Farm Grass\textquotedbl{}} & \$control = the authors' 9 control variables (Table 1, column 3).\\
563 & \texttt{**5\%Shannon} & Comment: drop top/bottom 5\% by Shannon.\\
564 & \texttt{quietly egen mean\_diversity = mean(ShannonBD), by(county)} & County average Shannon.\\
565 & \texttt{centile mean\_diversity if !missing(mean\_diversity),\newline \hspace*{1em}centile(5 95)} & 5th and 95th percentiles of county-average Shannon.\\
566 & \texttt{local p5 = r(c\_1)} & Stores the 5th percentile.\\
567 & \texttt{local p95 = r(c\_2)} & Stores the 95th percentile.\\
568 & \texttt{gen to\_exclude = (mean\_diversity <= \textasciigrave{}p5\textquotesingle{} | mean\_diversity\newline \hspace*{1em}>= \textasciigrave{}p95\textquotesingle{}) if !missing(mean\_diversity)} & Flags counties in the top or bottom 5\% of average Shannon.\\
569 & \texttt{reghdfe ShannonBD PI \$control if  to\_exclude != 1,\newline \hspace*{1em}absorb(county year) cluster(county)} & Regresses Shannon bird diversity on PSI (solar policy stringency); controlling for the 9 controls; with county and year-month fixed effects; standard errors clustered by county; excluding counties in the top/bottom 5\% of mean Shannon. Table S13 col 1: -0.0152.\\
570 & \texttt{estadd local year \textquotedbl{}YES\textquotedbl{}} & Adds a "Year-month FE = YES" marker to the stored results (printed at the bottom of the table).\\
571 & \texttt{estadd local county \textquotedbl{}YES\textquotedbl{}} & Adds a "County FE = YES" marker to the stored results.\\
572 & \texttt{est store s1} & Saves the last estimates under the name s1 so they can be put in a table.\\
573 & \texttt{**5\%PI} & Comment: drop top 5\% by PSI.\\
574 & \texttt{quietly egen mean\_pi = mean(PI), by(county)} & County average PSI.\\
575 & \texttt{centile mean\_pi if !missing(mean\_pi), centile(95)} & 95th percentile of county-average PSI.\\
576 & \texttt{local p95 = r(c\_1)} & Stores it.\\
577 & \texttt{gen exclude\_high\_pi = (mean\_pi >= \textasciigrave{}p95\textquotesingle{}) if\newline \hspace*{1em}!missing(mean\_pi)} & Flags the top 5\% PSI counties.\\
578 & \texttt{reghdfe ShannonBD PI \$control if  exclude\_high\_pi!= 1,\newline \hspace*{1em}absorb(county year) cluster(county)} & Regresses Shannon bird diversity on PSI (solar policy stringency); controlling for the 9 controls; with county and year-month fixed effects; standard errors clustered by county; excluding the top 5\% of counties by mean PSI. Table S13 col 2: -0.0123.\\
579 & \texttt{estadd local year \textquotedbl{}YES\textquotedbl{}} & Adds a "Year-month FE = YES" marker to the stored results (printed at the bottom of the table).\\
580 & \texttt{estadd local county \textquotedbl{}YES\textquotedbl{}} & Adds a "County FE = YES" marker to the stored results.\\
581 & \texttt{est store s2} & Saves the last estimates under the name s2 so they can be put in a table.\\
582 & \texttt{**Provincial capitals and municipalities} & Comment (not executed): explains the code that follows.\\
583 & \texttt{reghdfe ShannonBD PI \$control if cities==0, absorb(county\newline \hspace*{1em}year) cluster(county)} & Regresses Shannon bird diversity on PSI (solar policy stringency); controlling for the 9 controls; with county and year-month fixed effects; standard errors clustered by county; excluding provincial capitals and the 4 municipalities. Table S13 col 3: -0.0099.\\
584 & \texttt{estadd local year \textquotedbl{}YES\textquotedbl{}} & Adds a "Year-month FE = YES" marker to the stored results (printed at the bottom of the table).\\
585 & \texttt{estadd local county \textquotedbl{}YES\textquotedbl{}} & Adds a "County FE = YES" marker to the stored results.\\
586 & \texttt{est store s3} & Saves the last estimates under the name s3 so they can be put in a table.\\
587 & \texttt{esttab s1 s2 s3 using robust\_alternative\_model.rtf,\newline \hspace*{1em}replace title(\textquotedbl{}Table S13. Robustness checks using\newline \hspace*{1em}alternative sampling criteria.\textquotedbl{}) ///} & Exports the stored results (s1 s2 s3) into the Word table robust\_alternative\_model.rtf.\\
588 & \texttt{\hspace*{2.0em}stats(year county N r2\_a, labels(\textquotedbl{}Year-{}month FE\textquotedbl{} \textquotedbl{}County\newline \hspace*{1em}FE\textquotedbl{} \textquotedbl{}Observations\textquotedbl{} \textquotedbl{}R-{}squared\textquotedbl{}) fmt(0 0 0 4)) ///} & Continuation of the previous command: summary rows at the bottom (e.g. FE markers, observations, R-squared).\\
589 & \texttt{\hspace*{2.0em}star(* 0.10 ** 0.05 *** 0.01) nonotes nocons b(4) se\newline \hspace*{1em}compress nogaps ///} & Continuation of the previous command: significance stars * 10\%, ** 5\%, *** 1\%; 4-decimal coefficients with standard errors in parentheses; no default notes, constant omitted, compact layout.\\
590 & \texttt{\hspace*{2.0em}note(\textquotedbl{}Notes: This table shows the results of r ... \textquotedbl{})} & Continuation: the table note printed under the table (text abridged here).\\
591 & \texttt{restore} & Restores the data saved by preserve (undoes changes made since).\\
594 & \texttt{*剔除低于最低观鸟记录阈值的观测值} & Comment: "exclude observations below minimum birdwatching-record thresholds" (Table S14).\\
595 & \texttt{preserve} & Saves a copy of the data in memory, so it can be restored after temporary changes.\\
596 & \texttt{global control \textquotedbl{}Temp Wind Pop Duration Carbon Water Green\newline \hspace*{1em}Farm Grass\textquotedbl{}} & \$control = the authors' 9 control variables (Table 1, column 3).\\
598 & \texttt{reghdfe ShannonBD20 PI \$control, absorb(county year)\newline \hspace*{1em}cluster(county)} & Regresses Shannon (record threshold 20) on PSI (solar policy stringency); controlling for the 9 controls; with county and year-month fixed effects; standard errors clustered by county. Table S14 col 1: -0.0079. Note: ShannonBD20 exists for only 17.7\% of rows, so this is not "dropping the bottom 20\%".\\
599 & \texttt{estadd local year \textquotedbl{}YES\textquotedbl{}} & Adds a "Year-month FE = YES" marker to the stored results (printed at the bottom of the table).\\
600 & \texttt{estadd local county \textquotedbl{}YES\textquotedbl{}} & Adds a "County FE = YES" marker to the stored results.\\
601 & \texttt{est store s1} & Saves the last estimates under the name s1 so they can be put in a table.\\
603 & \texttt{reghdfe ShannonBD10 PI \$control, absorb(county year)\newline \hspace*{1em}cluster(county)} & Regresses Shannon (threshold 10) on PSI (solar policy stringency); controlling for the 9 controls; with county and year-month fixed effects; standard errors clustered by county.\\
604 & \texttt{estadd local year \textquotedbl{}YES\textquotedbl{}} & Adds a "Year-month FE = YES" marker to the stored results (printed at the bottom of the table).\\
605 & \texttt{estadd local county \textquotedbl{}YES\textquotedbl{}} & Adds a "County FE = YES" marker to the stored results.\\
606 & \texttt{est store s2} & Saves the last estimates under the name s2 so they can be put in a table.\\
608 & \texttt{reghdfe ShannonBD5 PI \$control, absorb(county year)\newline \hspace*{1em}cluster(county)} & Regresses Shannon (threshold 5) on PSI (solar policy stringency); controlling for the 9 controls; with county and year-month fixed effects; standard errors clustered by county.\\
609 & \texttt{estadd local year \textquotedbl{}YES\textquotedbl{}} & Adds a "Year-month FE = YES" marker to the stored results (printed at the bottom of the table).\\
610 & \texttt{estadd local county \textquotedbl{}YES\textquotedbl{}} & Adds a "County FE = YES" marker to the stored results.\\
611 & \texttt{est store s3} & Saves the last estimates under the name s3 so they can be put in a table.\\
613 & \texttt{reghdfe ShannonBD3 PI \$control, absorb(county year)\newline \hspace*{1em}cluster(county)} & Regresses Shannon (threshold 3) on PSI (solar policy stringency); controlling for the 9 controls; with county and year-month fixed effects; standard errors clustered by county.\\
614 & \texttt{estadd local year \textquotedbl{}YES\textquotedbl{}} & Adds a "Year-month FE = YES" marker to the stored results (printed at the bottom of the table).\\
615 & \texttt{estadd local county \textquotedbl{}YES\textquotedbl{}} & Adds a "County FE = YES" marker to the stored results.\\
616 & \texttt{est store s4} & Saves the last estimates under the name s4 so they can be put in a table.\\
618 & \texttt{esttab s1 s2 s3 s4 using robust\_other2.rtf, replace\newline \hspace*{1em}title(\textquotedbl{}Table S14. Excluding county-{}year observations below\newline \hspace*{1em}minimum birdwatching record thresholds.\textquotedbl{}) ///} & Exports the stored results (s1 s2 s3 s4) into the Word table robust\_other2.rtf.\\
619 & \texttt{\hspace*{2.0em}mtitles(\textquotedbl{}Trim bottom 20\%\textquotedbl{} \textquotedbl{}Trim bottom 10\%\textquotedbl{} \textquotedbl{}Trim bottom\newline \hspace*{1em}5\%\textquotedbl{} \textquotedbl{}Trim bottom 3\%\textquotedbl{}) ///} & Continuation of the previous command: column titles.\\
620 & \texttt{\hspace*{2.0em}stats(year county N r2\_a, labels(\textquotedbl{}Year-{}month FE\textquotedbl{} \textquotedbl{}County\newline \hspace*{1em}FE\textquotedbl{} \textquotedbl{}Observations\textquotedbl{} \textquotedbl{}R-{}squared\textquotedbl{}) fmt(0 0 0 4)) ///} & Continuation of the previous command: summary rows at the bottom (e.g. FE markers, observations, R-squared).\\
621 & \texttt{\hspace*{2.0em}star(* 0.10 ** 0.05 *** 0.01) nonotes nocons b(4) se\newline \hspace*{1em}compress nogaps ///} & Continuation of the previous command: significance stars * 10\%, ** 5\%, *** 1\%; 4-decimal coefficients with standard errors in parentheses; no default notes, constant omitted, compact layout.\\
622 & \texttt{\hspace*{2.0em}note(\textquotedbl{}Notes:  Columns (1) to (4) report result ... \textquotedbl{})} & Continuation: the table note printed under the table (text abridged here).\\
623 & \texttt{restore} & Restores the data saved by preserve (undoes changes made since).\\
626 & \texttt{*Eliminate other policy interference} & Comment (not executed): explains the code that follows.\\
627 & \texttt{**CEPI policy} & Comment: CEI policy.\\
628 & \texttt{preserve} & Saves a copy of the data in memory, so it can be restored after temporary changes.\\
629 & \texttt{global control \textquotedbl{}Temp Wind Pop Duration Carbon Water Green\newline \hspace*{1em}Farm Grass\textquotedbl{}} & \$control = the authors' 9 control variables (Table 1, column 3).\\
630 & \texttt{reghdfe ShannonBD PI \$control CEI, absorb(county year)\newline \hspace*{1em}cluster(county)} & Regresses Shannon bird diversity on PSI (solar policy stringency); controlling for the 9 controls CEI; with county and year-month fixed effects; standard errors clustered by county.\\
631 & \texttt{estadd local year \textquotedbl{}YES\textquotedbl{}} & Adds a "Year-month FE = YES" marker to the stored results (printed at the bottom of the table).\\
632 & \texttt{estadd local county \textquotedbl{}YES\textquotedbl{}} & Adds a "County FE = YES" marker to the stored results.\\
633 & \texttt{est store s1} & Saves the last estimates under the name s1 so they can be put in a table.\\
634 & \texttt{**煤电厂} & Comment: "coal power plants".\\
635 & \texttt{reghdfe ShannonBD PI \$control Coal, absorb(county year)\newline \hspace*{1em}cluster(county)} & Regresses Shannon bird diversity on PSI (solar policy stringency); controlling for the 9 controls Coal; with county and year-month fixed effects; standard errors clustered by county.\\
636 & \texttt{estadd local year \textquotedbl{}YES\textquotedbl{}} & Adds a "Year-month FE = YES" marker to the stored results (printed at the bottom of the table).\\
637 & \texttt{estadd local county \textquotedbl{}YES\textquotedbl{}} & Adds a "County FE = YES" marker to the stored results.\\
638 & \texttt{est store s2} & Saves the last estimates under the name s2 so they can be put in a table.\\
639 & \texttt{**风电扩张} & Comment: "wind power expansion".\\
640 & \texttt{reghdfe ShannonBD PI \$control Windexp, absorb(county year)\newline \hspace*{1em}cluster(county)} & Regresses Shannon bird diversity on PSI (solar policy stringency); controlling for the 9 controls Windexp; with county and year-month fixed effects; standard errors clustered by county. Note: Windexp is wrongly 1 for all rows with missing wind data.\\
641 & \texttt{estadd local year \textquotedbl{}YES\textquotedbl{}} & Adds a "Year-month FE = YES" marker to the stored results (printed at the bottom of the table).\\
642 & \texttt{estadd local county \textquotedbl{}YES\textquotedbl{}} & Adds a "County FE = YES" marker to the stored results.\\
643 & \texttt{est store s3} & Saves the last estimates under the name s3 so they can be put in a table.\\
644 & \texttt{**绿色金融} & Comment: "green finance".\\
645 & \texttt{reghdfe ShannonBD PI \$control Greenfin, absorb(county\newline \hspace*{1em}year) cluster(county)} & Regresses Shannon bird diversity on PSI (solar policy stringency); controlling for the 9 controls Greenfin; with county and year-month fixed effects; standard errors clustered by county.\\
646 & \texttt{estadd local year \textquotedbl{}YES\textquotedbl{}} & Adds a "Year-month FE = YES" marker to the stored results (printed at the bottom of the table).\\
647 & \texttt{estadd local county \textquotedbl{}YES\textquotedbl{}} & Adds a "County FE = YES" marker to the stored results.\\
648 & \texttt{est store s4} & Saves the last estimates under the name s4 so they can be put in a table.\\
649 & \texttt{**全部放进来} & Comment: "put all of them in".\\
650 & \texttt{reghdfe ShannonBD PI \$control CEI Coal Windexp Greenfin,\newline \hspace*{1em}absorb(county year) cluster(county)} & Regresses Shannon bird diversity on PSI (solar policy stringency); controlling for the 9 controls CEI Coal Windexp Greenfin; with county and year-month fixed effects; standard errors clustered by county.\\
651 & \texttt{estadd local year \textquotedbl{}YES\textquotedbl{}} & Adds a "Year-month FE = YES" marker to the stored results (printed at the bottom of the table).\\
652 & \texttt{estadd local county \textquotedbl{}YES\textquotedbl{}} & Adds a "County FE = YES" marker to the stored results.\\
653 & \texttt{est store s5} & Saves the last estimates under the name s5 so they can be put in a table.\\
654 & \texttt{esttab s1 s2 s3 s4 s5 using robust\_eliminate\_policy.rtf,\newline \hspace*{1em}replace title(\textquotedbl{}Table S15. Ruling out confounding effects\newline \hspace*{1em}from related policies.\textquotedbl{}) ///} & Exports the stored results (s1 s2 s3 s4 s5) into the Word table robust\_eliminate\_policy.rtf.\\
655 & \texttt{\hspace*{2.0em}stats(year county N r2\_a, labels(\textquotedbl{}Year-{}month FE\textquotedbl{} \textquotedbl{}County\newline \hspace*{1em}FE\textquotedbl{} \textquotedbl{}Observations\textquotedbl{} \textquotedbl{}R-{}squared\textquotedbl{}) fmt(0 0 0 4)) ///} & Continuation of the previous command: summary rows at the bottom (e.g. FE markers, observations, R-squared).\\
656 & \texttt{\hspace*{2.0em}star(* 0.10 ** 0.05 *** 0.01) nonotes nocons b(4) se\newline \hspace*{1em}compress nogaps ///} & Continuation of the previous command: significance stars * 10\%, ** 5\%, *** 1\%; 4-decimal coefficients with standard errors in parentheses; no default notes, constant omitted, compact layout.\\
657 & \texttt{\hspace*{2.0em}order(PI CEI Coal Windexp Greenfin \$control) ///} & Continuation of the previous command: row order.\\
658 & \texttt{\hspace*{2.0em}note(\textquotedbl{}Notes: This table presents the regressio ... \textquotedbl{})} & Continuation: the table note printed under the table (text abridged here).\\
659 & \texttt{restore} & Restores the data saved by preserve (undoes changes made since).\\
662 & \texttt{*Other robustness check} & Comment (not executed): explains the code that follows.\\
663 & \texttt{preserve} & Saves a copy of the data in memory, so it can be restored after temporary changes.\\
664 & \texttt{global control \textquotedbl{}Temp Wind Pop Duration Carbon Water Green\newline \hspace*{1em}Farm Grass\textquotedbl{}} & \$control = the authors' 9 control variables (Table 1, column 3).\\
665 & \texttt{**dditional control variables} & Comment: "additional control variables" (PM2.5).\\
666 & \texttt{reghdfe ShannonBD PI \$control PM, absorb(county year)\newline \hspace*{1em}cluster(county)} & Regresses Shannon bird diversity on PSI (solar policy stringency); controlling for the 9 controls PM; with county and year-month fixed effects; standard errors clustered by county.\\
667 & \texttt{estadd local year \textquotedbl{}YES\textquotedbl{}} & Adds a "Year-month FE = YES" marker to the stored results (printed at the bottom of the table).\\
668 & \texttt{estadd local county \textquotedbl{}YES\textquotedbl{}} & Adds a "County FE = YES" marker to the stored results.\\
669 & \texttt{est store s1} & Saves the last estimates under the name s1 so they can be put in a table.\\
670 & \texttt{**剔除高PI相邻县域} & Comment: "exclude counties adjacent to high-PI counties".\\
671 & \texttt{reghdfe ShannonBD PI \$control if is\_neighbor\_of\_highPI==0,\newline \hspace*{1em}absorb(county year) cluster(county)} & Regresses Shannon bird diversity on PSI (solar policy stringency); controlling for the 9 controls; with county and year-month fixed effects; standard errors clustered by county; excluding neighbours of high-PSI counties.\\
672 & \texttt{estadd local year \textquotedbl{}YES\textquotedbl{}} & Adds a "Year-month FE = YES" marker to the stored results (printed at the bottom of the table).\\
673 & \texttt{estadd local county \textquotedbl{}YES\textquotedbl{}} & Adds a "County FE = YES" marker to the stored results.\\
674 & \texttt{est store s2} & Saves the last estimates under the name s2 so they can be put in a table.\\
675 & \texttt{/*} & Starts a block comment: everything until */ is ignored by Stata.\\
676 & \texttt{bysort YEAR: egen p90\_PI = pctile(PI), p(90)} & Inside a block comment (not executed).\\
677 & \texttt{gen high\_pi = (PI > p90\_PI)} & Inside a block comment (not executed).\\
678 & \texttt{*/} & Ends the block comment.\\
679 & \texttt{**Seasonal robustness check} & Comment (not executed): explains the code that follows.\\
680 & \texttt{keep if month>=5 \& month<=7} & Keeps only May-July (breeding season).\\
681 & \texttt{reghdfe ShannonBD PI \$control, absorb(county year)\newline \hspace*{1em}cluster(county)} & Regresses Shannon bird diversity on PSI (solar policy stringency); controlling for the 9 controls; with county and year-month fixed effects; standard errors clustered by county. Table S16 col 3 (breeding season): -0.0148.\\
682 & \texttt{estadd local year \textquotedbl{}YES\textquotedbl{}} & Adds a "Year-month FE = YES" marker to the stored results (printed at the bottom of the table).\\
683 & \texttt{estadd local county \textquotedbl{}YES\textquotedbl{}} & Adds a "County FE = YES" marker to the stored results.\\
684 & \texttt{est store s3} & Saves the last estimates under the name s3 so they can be put in a table.\\
685 & \texttt{esttab s1 s2 s3 using robust\_other.rtf, replace\newline \hspace*{1em}title(\textquotedbl{}Table S16. Other robustness checks.\textquotedbl{}) ///} & Exports the stored results (s1 s2 s3) into the Word table robust\_other.rtf.\\
686 & \texttt{\hspace*{2.0em}stats(year county N r2\_a, labels(\textquotedbl{}Year-{}month FE\textquotedbl{} \textquotedbl{}County\newline \hspace*{1em}FE\textquotedbl{} \textquotedbl{}Observations\textquotedbl{} \textquotedbl{}R-{}squared\textquotedbl{}) fmt(0 0 0 4)) ///} & Continuation of the previous command: summary rows at the bottom (e.g. FE markers, observations, R-squared).\\
687 & \texttt{\hspace*{2.0em}star(* 0.10 ** 0.05 *** 0.01) nonotes nocons b(4) se\newline \hspace*{1em}compress nogaps ///} & Continuation of the previous command: significance stars * 10\%, ** 5\%, *** 1\%; 4-decimal coefficients with standard errors in parentheses; no default notes, constant omitted, compact layout.\\
688 & \texttt{\hspace*{2.0em}note(\textquotedbl{}Notes: Column (1) includes PM2.5 as an a ... \textquotedbl{})} & Continuation: the table note printed under the table (text abridged here).\\
689 & \texttt{restore} & Restores the data saved by preserve (undoes changes made since).\\
693 & \texttt{*-{}-{}-{}-{}-{}-{}-{}-{}-{}-{}-{}-{}-{}-{}-{}-{}-{}-{}-{}-{}-{}-{}-{}-{}-{}-{}-{}-{}-{}-{}-{}-{}-{}-{}-{}-{}-{}-{}-{}-{}-{}-{}-{}-{}-{}-{}-{}-{}-{}-{}-{}-{}-{}-{}-{}-{}-{}\newline \hspace*{1em}-{}-{}-{}-{}-{}-{}-{}-{}-{}-{}-{}-{}-{}-{}-{}} & Divider line (comment).\\
694 & \texttt{* Heterogeneity Analysis1} & Comment (not executed): explains the code that follows.\\
695 & \texttt{*-{}-{}-{}-{}-{}-{}-{}-{}-{}-{}-{}-{}-{}-{}-{}-{}-{}-{}-{}-{}-{}-{}-{}-{}-{}-{}-{}-{}-{}-{}-{}-{}-{}-{}-{}-{}-{}-{}-{}-{}-{}-{}-{}-{}-{}-{}-{}-{}-{}-{}-{}-{}-{}-{}-{}-{}-{}\newline \hspace*{1em}-{}-{}-{}-{}-{}-{}-{}-{}-{}-{}-{}-{}-{}-{}-{}} & Divider line (comment).\\
696 & \texttt{preserve} & Saves a copy of the data in memory, so it can be restored after temporary changes.\\
697 & \texttt{global control \textquotedbl{}Temp Wind Pop Duration Carbon Water Green\newline \hspace*{1em}Farm Grass\textquotedbl{}} & \$control = the authors' 9 control variables (Table 1, column 3).\\
698 & \texttt{//bdiff, group(poverty\_county) model(reghdfe ShannonBD PI\newline \hspace*{1em}\$control, absorb(county year) vce(cluster county))\newline \hspace*{1em}reps(1000) first detail seed(1234)} & Comment (not executed): explains the code that follows.\\
699 & \texttt{reghdfe ShannonBD PI \$control if poverty\_county==0,\newline \hspace*{1em}absorb(county year) cluster(county)} & Regresses Shannon bird diversity on PSI (solar policy stringency); controlling for the 9 controls; with county and year-month fixed effects; standard errors clustered by county; only non-poverty counties.\\
700 & \texttt{estadd local year \textquotedbl{}YES\textquotedbl{}} & Adds a "Year-month FE = YES" marker to the stored results (printed at the bottom of the table).\\
701 & \texttt{estadd local county \textquotedbl{}YES\textquotedbl{}} & Adds a "County FE = YES" marker to the stored results.\\
702 & \texttt{est store s1} & Saves the last estimates under the name s1 so they can be put in a table.\\
703 & \texttt{reghdfe ShannonBD PI \$control if poverty\_county==1,\newline \hspace*{1em}absorb(county year) cluster(county)} & Regresses Shannon bird diversity on PSI (solar policy stringency); controlling for the 9 controls; with county and year-month fixed effects; standard errors clustered by county; only poverty counties.\\
704 & \texttt{estadd local year \textquotedbl{}YES\textquotedbl{}} & Adds a "Year-month FE = YES" marker to the stored results (printed at the bottom of the table).\\
705 & \texttt{estadd local county \textquotedbl{}YES\textquotedbl{}} & Adds a "County FE = YES" marker to the stored results.\\
706 & \texttt{est store s2} & Saves the last estimates under the name s2 so they can be put in a table.\\
707 & \texttt{//bdiff, group(north\_region) model(reghdfe ShannonBD PI\newline \hspace*{1em}\$control, absorb(county year) vce(cluster county))\newline \hspace*{1em}reps(1000) first detail seed(1234)} & Comment (not executed): explains the code that follows.\\
708 & \texttt{reghdfe ShannonBD PI \$control if north\_region==0,\newline \hspace*{1em}absorb(county year) cluster(county)} & Regresses Shannon bird diversity on PSI (solar policy stringency); controlling for the 9 controls; with county and year-month fixed effects; standard errors clustered by county; only counties outside the Three-North region.\\
709 & \texttt{estadd local year \textquotedbl{}YES\textquotedbl{}} & Adds a "Year-month FE = YES" marker to the stored results (printed at the bottom of the table).\\
710 & \texttt{estadd local county \textquotedbl{}YES\textquotedbl{}} & Adds a "County FE = YES" marker to the stored results.\\
711 & \texttt{est store s3} & Saves the last estimates under the name s3 so they can be put in a table.\\
712 & \texttt{reghdfe ShannonBD PI \$control if north\_region==1,\newline \hspace*{1em}absorb(county year) cluster(county)} & Regresses Shannon bird diversity on PSI (solar policy stringency); controlling for the 9 controls; with county and year-month fixed effects; standard errors clustered by county; only Three-North counties.\\
713 & \texttt{estadd local year \textquotedbl{}YES\textquotedbl{}} & Adds a "Year-month FE = YES" marker to the stored results (printed at the bottom of the table).\\
714 & \texttt{estadd local county \textquotedbl{}YES\textquotedbl{}} & Adds a "County FE = YES" marker to the stored results.\\
715 & \texttt{est store s4} & Saves the last estimates under the name s4 so they can be put in a table.\\
716 & \texttt{//bdiff, group(desert\_gobi) model(reghdfe ShannonBD PI\newline \hspace*{1em}\$control, absorb(county year) vce(cluster county))\newline \hspace*{1em}reps(1000) first detail seed(1234)} & Comment (not executed): explains the code that follows.\\
717 & \texttt{reghdfe ShannonBD PI \$control if desert\_gobi==0,\newline \hspace*{1em}absorb(county year) cluster(county)} & Regresses Shannon bird diversity on PSI (solar policy stringency); controlling for the 9 controls; with county and year-month fixed effects; standard errors clustered by county; only non-desert counties.\\
718 & \texttt{estadd local year \textquotedbl{}YES\textquotedbl{}} & Adds a "Year-month FE = YES" marker to the stored results (printed at the bottom of the table).\\
719 & \texttt{estadd local county \textquotedbl{}YES\textquotedbl{}} & Adds a "County FE = YES" marker to the stored results.\\
720 & \texttt{est store s5} & Saves the last estimates under the name s5 so they can be put in a table.\\
721 & \texttt{reghdfe ShannonBD PI \$control if desert\_gobi==1,\newline \hspace*{1em}absorb(county year) cluster(county)} & Regresses Shannon bird diversity on PSI (solar policy stringency); controlling for the 9 controls; with county and year-month fixed effects; standard errors clustered by county; only sandy/Gobi desert counties.\\
722 & \texttt{estadd local year \textquotedbl{}YES\textquotedbl{}} & Adds a "Year-month FE = YES" marker to the stored results (printed at the bottom of the table).\\
723 & \texttt{estadd local county \textquotedbl{}YES\textquotedbl{}} & Adds a "County FE = YES" marker to the stored results.\\
724 & \texttt{est store s6} & Saves the last estimates under the name s6 so they can be put in a table.\\
725 & \texttt{esttab s1 s2 s3 s4 s5 s6 using heterogeneity1.rtf, replace\newline \hspace*{1em}title(\textquotedbl{}Table S17. Cross-{}sectional heterogeneity.\textquotedbl{}) ///} & Exports the stored results (s1 s2 s3 s4 s5 s6) into the Word table heterogeneity1.rtf.\\
726 & \texttt{\hspace*{2.0em}mtitles(\textquotedbl{}Non-{}poverty counties\textquotedbl{} \textquotedbl{}Poverty counties\textquotedbl{}\newline \hspace*{1em}\textquotedbl{}Non-{}Three-{}North regions\textquotedbl{} \textquotedbl{}Three-{}North regions\textquotedbl{} \textquotedbl{}Non-{}sandy\newline \hspace*{1em}and gravel desert\textquotedbl{} \textquotedbl{}Sandy and gravel desert\textquotedbl{}) ///} & Continuation of the previous command: column titles.\\
727 & \texttt{\hspace*{2.0em}stats(year county N r2\_a, labels(\textquotedbl{}Year-{}month FE\textquotedbl{} \textquotedbl{}County\newline \hspace*{1em}FE\textquotedbl{} \textquotedbl{}Observations\textquotedbl{} \textquotedbl{}R-{}squared\textquotedbl{}) fmt(0 0 0 4)) ///} & Continuation of the previous command: summary rows at the bottom (e.g. FE markers, observations, R-squared).\\
728 & \texttt{\hspace*{2.0em}star(* 0.10 ** 0.05 *** 0.01) nonotes nocons b(4) se\newline \hspace*{1em}compress nogaps ///} & Continuation of the previous command: significance stars * 10\%, ** 5\%, *** 1\%; 4-decimal coefficients with standard errors in parentheses; no default notes, constant omitted, compact layout.\\
729 & \texttt{\hspace*{2.0em}note(\textquotedbl{}Notes: Specification 1 divides the sampl ... \textquotedbl{})} & Continuation: the table note printed under the table (text abridged here).\\
730 & \texttt{restore} & Restores the data saved by preserve (undoes changes made since).\\
734 & \texttt{*-{}-{}-{}-{}-{}-{}-{}-{}-{}-{}-{}-{}-{}-{}-{}-{}-{}-{}-{}-{}-{}-{}-{}-{}-{}-{}-{}-{}-{}-{}-{}-{}-{}-{}-{}-{}-{}-{}-{}-{}-{}-{}-{}-{}-{}-{}-{}-{}-{}-{}-{}-{}-{}-{}-{}-{}-{}\newline \hspace*{1em}-{}-{}-{}-{}-{}-{}-{}-{}-{}-{}-{}-{}-{}-{}-{}} & Divider line (comment).\\
735 & \texttt{* Heterogeneity Analysis2} & Comment (not executed): explains the code that follows.\\
736 & \texttt{*-{}-{}-{}-{}-{}-{}-{}-{}-{}-{}-{}-{}-{}-{}-{}-{}-{}-{}-{}-{}-{}-{}-{}-{}-{}-{}-{}-{}-{}-{}-{}-{}-{}-{}-{}-{}-{}-{}-{}-{}-{}-{}-{}-{}-{}-{}-{}-{}-{}-{}-{}-{}-{}-{}-{}-{}-{}\newline \hspace*{1em}-{}-{}-{}-{}-{}-{}-{}-{}-{}-{}-{}-{}-{}-{}-{}} & Divider line (comment).\\
737 & \texttt{preserve} & Saves a copy of the data in memory, so it can be restored after temporary changes.\\
738 & \texttt{global control \textquotedbl{}Temp Wind Pop Duration Carbon Water Green\newline \hspace*{1em}Farm Grass\textquotedbl{}} & \$control = the authors' 9 control variables (Table 1, column 3).\\
739 & \texttt{reghdfe Endemic PI \$control , absorb(county year)\newline \hspace*{1em}cluster(county)} & Regresses Shannon of endemic species on PSI (solar policy stringency); controlling for the 9 controls; with county and year-month fixed effects; standard errors clustered by county.\\
740 & \texttt{estadd local year \textquotedbl{}YES\textquotedbl{}} & Adds a "Year-month FE = YES" marker to the stored results (printed at the bottom of the table).\\
741 & \texttt{estadd local county \textquotedbl{}YES\textquotedbl{}} & Adds a "County FE = YES" marker to the stored results.\\
742 & \texttt{est store s1} & Saves the last estimates under the name s1 so they can be put in a table.\\
743 & \texttt{reghdfe non\_Endemic PI \$control , absorb(county year)\newline \hspace*{1em}cluster(county)} & Regresses Shannon of non-endemic species on PSI (solar policy stringency); controlling for the 9 controls; with county and year-month fixed effects; standard errors clustered by county.\\
744 & \texttt{estadd local year \textquotedbl{}YES\textquotedbl{}} & Adds a "Year-month FE = YES" marker to the stored results (printed at the bottom of the table).\\
745 & \texttt{estadd local county \textquotedbl{}YES\textquotedbl{}} & Adds a "County FE = YES" marker to the stored results.\\
746 & \texttt{est store s2} & Saves the last estimates under the name s2 so they can be put in a table.\\
747 & \texttt{reghdfe Migratory PI \$control , absorb(county year)\newline \hspace*{1em}cluster(county)} & Regresses Shannon of migratory species on PSI (solar policy stringency); controlling for the 9 controls; with county and year-month fixed effects; standard errors clustered by county.\\
748 & \texttt{estadd local year \textquotedbl{}YES\textquotedbl{}} & Adds a "Year-month FE = YES" marker to the stored results (printed at the bottom of the table).\\
749 & \texttt{estadd local county \textquotedbl{}YES\textquotedbl{}} & Adds a "County FE = YES" marker to the stored results.\\
750 & \texttt{est store s3} & Saves the last estimates under the name s3 so they can be put in a table.\\
751 & \texttt{reghdfe Resident PI \$control , absorb(county year)\newline \hspace*{1em}cluster(county)} & Regresses Shannon of resident species on PSI (solar policy stringency); controlling for the 9 controls; with county and year-month fixed effects; standard errors clustered by county.\\
752 & \texttt{estadd local year \textquotedbl{}YES\textquotedbl{}} & Adds a "Year-month FE = YES" marker to the stored results (printed at the bottom of the table).\\
753 & \texttt{estadd local county \textquotedbl{}YES\textquotedbl{}} & Adds a "County FE = YES" marker to the stored results.\\
754 & \texttt{est store s4} & Saves the last estimates under the name s4 so they can be put in a table.\\
755 & \texttt{reghdfe Protected PI \$control , absorb(county year)\newline \hspace*{1em}cluster(county)} & Regresses Shannon of protected species on PSI (solar policy stringency); controlling for the 9 controls; with county and year-month fixed effects; standard errors clustered by county.\\
756 & \texttt{estadd local year \textquotedbl{}YES\textquotedbl{}} & Adds a "Year-month FE = YES" marker to the stored results (printed at the bottom of the table).\\
757 & \texttt{estadd local county \textquotedbl{}YES\textquotedbl{}} & Adds a "County FE = YES" marker to the stored results.\\
758 & \texttt{est store s5} & Saves the last estimates under the name s5 so they can be put in a table.\\
759 & \texttt{reghdfe non\_Protected PI \$control , absorb(county year)\newline \hspace*{1em}cluster(county)} & Regresses Shannon of non-protected species on PSI (solar policy stringency); controlling for the 9 controls; with county and year-month fixed effects; standard errors clustered by county.\\
760 & \texttt{estadd local year \textquotedbl{}YES\textquotedbl{}} & Adds a "Year-month FE = YES" marker to the stored results (printed at the bottom of the table).\\
761 & \texttt{estadd local county \textquotedbl{}YES\textquotedbl{}} & Adds a "County FE = YES" marker to the stored results.\\
762 & \texttt{est store s6} & Saves the last estimates under the name s6 so they can be put in a table.\\
763 & \texttt{esttab s1 s2 s3 s4 s5 s6 using heterogeneity2.rtf, replace\newline \hspace*{1em}title(\textquotedbl{}Table S18. Bird species heterogeneity.\textquotedbl{}) ///} & Exports the stored results (s1 s2 s3 s4 s5 s6) into the Word table heterogeneity2.rtf.\\
764 & \texttt{\hspace*{2.0em}mtitles(\textquotedbl{}Endemic\textquotedbl{} \textquotedbl{}Non-{}endemic\textquotedbl{} \textquotedbl{}Migratory\textquotedbl{} \textquotedbl{}Resident\textquotedbl{}\newline \hspace*{1em}\textquotedbl{}Protected\textquotedbl{} \textquotedbl{}Non-{}protected\textquotedbl{}) ///} & Continuation of the previous command: column titles.\\
765 & \texttt{\hspace*{2.0em}stats(year county N r2\_a, labels(\textquotedbl{}Year-{}month FE\textquotedbl{} \textquotedbl{}County\newline \hspace*{1em}FE\textquotedbl{} \textquotedbl{}Observations\textquotedbl{} \textquotedbl{}R-{}squared\textquotedbl{}) fmt(0 0 0 4)) ///} & Continuation of the previous command: summary rows at the bottom (e.g. FE markers, observations, R-squared).\\
766 & \texttt{\hspace*{2.0em}star(* 0.10 ** 0.05 *** 0.01) nonotes nocons b(4) se\newline \hspace*{1em}compress nogaps ///} & Continuation of the previous command: significance stars * 10\%, ** 5\%, *** 1\%; 4-decimal coefficients with standard errors in parentheses; no default notes, constant omitted, compact layout.\\
767 & \texttt{\hspace*{2.0em}note(\textquotedbl{}Notes: Panel A presents the results clas ... \textquotedbl{})} & Continuation: the table note printed under the table (text abridged here).\\
768 & \texttt{restore} & Restores the data saved by preserve (undoes changes made since).\\
770 & \texttt{preserve} & Saves a copy of the data in memory, so it can be restored after temporary changes.\\
771 & \texttt{global control \textquotedbl{}Temp Wind Pop Duration Carbon Water Green\newline \hspace*{1em}Farm Grass\textquotedbl{}} & \$control = the authors' 9 control variables (Table 1, column 3).\\
772 & \texttt{reghdfe Land\_water PI \$control , absorb(county year)\newline \hspace*{1em}cluster(county)} & Regresses Shannon of ground/water nesters on PSI (solar policy stringency); controlling for the 9 controls; with county and year-month fixed effects; standard errors clustered by county.\\
773 & \texttt{estadd local year \textquotedbl{}YES\textquotedbl{}} & Adds a "Year-month FE = YES" marker to the stored results (printed at the bottom of the table).\\
774 & \texttt{estadd local county \textquotedbl{}YES\textquotedbl{}} & Adds a "County FE = YES" marker to the stored results.\\
775 & \texttt{est store s1} & Saves the last estimates under the name s1 so they can be put in a table.\\
776 & \texttt{reghdfe Nest\_veg PI \$control , absorb(county year)\newline \hspace*{1em}cluster(county)} & Regresses Shannon of vegetation nesters on PSI (solar policy stringency); controlling for the 9 controls; with county and year-month fixed effects; standard errors clustered by county.\\
777 & \texttt{estadd local year \textquotedbl{}YES\textquotedbl{}} & Adds a "Year-month FE = YES" marker to the stored results (printed at the bottom of the table).\\
778 & \texttt{estadd local county \textquotedbl{}YES\textquotedbl{}} & Adds a "County FE = YES" marker to the stored results.\\
779 & \texttt{est store s2} & Saves the last estimates under the name s2 so they can be put in a table.\\
780 & \texttt{reghdfe non\_Herbivorous PI \$control , absorb(county year)\newline \hspace*{1em}cluster(county)} & Regresses Shannon of carnivores/omnivores on PSI (solar policy stringency); controlling for the 9 controls; with county and year-month fixed effects; standard errors clustered by county.\\
781 & \texttt{estadd local year \textquotedbl{}YES\textquotedbl{}} & Adds a "Year-month FE = YES" marker to the stored results (printed at the bottom of the table).\\
782 & \texttt{estadd local county \textquotedbl{}YES\textquotedbl{}} & Adds a "County FE = YES" marker to the stored results.\\
783 & \texttt{est store s3} & Saves the last estimates under the name s3 so they can be put in a table.\\
784 & \texttt{reghdfe Herbivorous PI \$control , absorb(county year)\newline \hspace*{1em}cluster(county)} & Regresses Shannon of herbivores on PSI (solar policy stringency); controlling for the 9 controls; with county and year-month fixed effects; standard errors clustered by county.\\
785 & \texttt{estadd local year \textquotedbl{}YES\textquotedbl{}} & Adds a "Year-month FE = YES" marker to the stored results (printed at the bottom of the table).\\
786 & \texttt{estadd local county \textquotedbl{}YES\textquotedbl{}} & Adds a "County FE = YES" marker to the stored results.\\
787 & \texttt{est store s4} & Saves the last estimates under the name s4 so they can be put in a table.\\
788 & \texttt{reghdfe Small\_flock PI \$control , absorb(county year)\newline \hspace*{1em}cluster(county)} & Regresses Shannon of small-flock species on PSI (solar policy stringency); controlling for the 9 controls; with county and year-month fixed effects; standard errors clustered by county.\\
789 & \texttt{estadd local year \textquotedbl{}YES\textquotedbl{}} & Adds a "Year-month FE = YES" marker to the stored results (printed at the bottom of the table).\\
790 & \texttt{estadd local county \textquotedbl{}YES\textquotedbl{}} & Adds a "County FE = YES" marker to the stored results.\\
791 & \texttt{est store s5} & Saves the last estimates under the name s5 so they can be put in a table.\\
792 & \texttt{reghdfe Large\_flock PI \$control , absorb(county year)\newline \hspace*{1em}cluster(county)} & Regresses Shannon of large-flock species on PSI (solar policy stringency); controlling for the 9 controls; with county and year-month fixed effects; standard errors clustered by county.\\
793 & \texttt{estadd local year \textquotedbl{}YES\textquotedbl{}} & Adds a "Year-month FE = YES" marker to the stored results (printed at the bottom of the table).\\
794 & \texttt{estadd local county \textquotedbl{}YES\textquotedbl{}} & Adds a "County FE = YES" marker to the stored results.\\
795 & \texttt{est store s6} & Saves the last estimates under the name s6 so they can be put in a table.\\
796 & \texttt{esttab s1 s2 s3 s4 s5 s6 using heterogeneity3.rtf, replace\newline \hspace*{1em}title(\textquotedbl{}Table S18. Bird species heterogeneity.\textquotedbl{}) ///} & Exports the stored results (s1 s2 s3 s4 s5 s6) into the Word table heterogeneity3.rtf.\\
797 & \texttt{\hspace*{2.0em}mtitles(\textquotedbl{}Land-{}water\textquotedbl{} \textquotedbl{}Vegetation\textquotedbl{} \textquotedbl{}Carnivorous and\newline \hspace*{1em}omnivorous\textquotedbl{} \textquotedbl{}Herbivorous\textquotedbl{} \textquotedbl{}Small-{}flock\textquotedbl{} \textquotedbl{}Large-{}flock\textquotedbl{}) ///} & Continuation of the previous command: column titles.\\
798 & \texttt{\hspace*{2.0em}stats(year county N r2\_a, labels(\textquotedbl{}Year-{}month FE\textquotedbl{} \textquotedbl{}County\newline \hspace*{1em}FE\textquotedbl{} \textquotedbl{}Observations\textquotedbl{} \textquotedbl{}R-{}squared\textquotedbl{}) fmt(0 0 0 4)) ///} & Continuation of the previous command: summary rows at the bottom (e.g. FE markers, observations, R-squared).\\
799 & \texttt{\hspace*{2.0em}star(* 0.10 ** 0.05 *** 0.01) nonotes nocons b(4) se\newline \hspace*{1em}compress nogaps} & Continuation of the previous command: significance stars * 10\%, ** 5\%, *** 1\%; 4-decimal coefficients with standard errors in parentheses; no default notes, constant omitted, compact layout.\\
800 & \texttt{restore} & Restores the data saved by preserve (undoes changes made since).\\
804 & \texttt{*-{}-{}-{}-{}-{}-{}-{}-{}-{}-{}-{}-{}-{}-{}-{}-{}-{}-{}-{}-{}-{}-{}-{}-{}-{}-{}-{}-{}-{}-{}-{}-{}-{}-{}-{}-{}-{}-{}-{}-{}-{}-{}-{}-{}-{}-{}-{}-{}-{}-{}-{}-{}-{}-{}-{}-{}-{}\newline \hspace*{1em}-{}-{}-{}-{}-{}-{}-{}-{}-{}-{}-{}-{}-{}-{}-{}} & Divider line (comment).\\
805 & \texttt{* Plausible channels} & Comment (not executed): explains the code that follows.\\
806 & \texttt{*-{}-{}-{}-{}-{}-{}-{}-{}-{}-{}-{}-{}-{}-{}-{}-{}-{}-{}-{}-{}-{}-{}-{}-{}-{}-{}-{}-{}-{}-{}-{}-{}-{}-{}-{}-{}-{}-{}-{}-{}-{}-{}-{}-{}-{}-{}-{}-{}-{}-{}-{}-{}-{}-{}-{}-{}-{}\newline \hspace*{1em}-{}-{}-{}-{}-{}-{}-{}-{}-{}-{}-{}-{}-{}-{}-{}} & Divider line (comment).\\
807 & \texttt{preserve} & Saves a copy of the data in memory, so it can be restored after temporary changes.\\
808 & \texttt{global control \textquotedbl{}Temp Wind Pop Duration Carbon Water Green\newline \hspace*{1em}Farm Grass\textquotedbl{}} & \$control = the authors' 9 control variables (Table 1, column 3).\\
809 & \texttt{reghdfe Ndvi PI \$control , absorb(county year)\newline \hspace*{1em}cluster(county)} & Regresses NDVI (greenness) on PSI (solar policy stringency); controlling for the 9 controls; with county and year-month fixed effects; standard errors clustered by county. Table 2: NDVI -0.0609 (0.0138).\\
810 & \texttt{estadd local year \textquotedbl{}YES\textquotedbl{}} & Adds a "Year-month FE = YES" marker to the stored results (printed at the bottom of the table).\\
811 & \texttt{estadd local county \textquotedbl{}YES\textquotedbl{}} & Adds a "County FE = YES" marker to the stored results.\\
812 & \texttt{est store s1} & Saves the last estimates under the name s1 so they can be put in a table.\\
813 & \texttt{reghdfe Nightlight PI \$control , absorb(county year)\newline \hspace*{1em}cluster(county)} & Regresses night-time light on PSI (solar policy stringency); controlling for the 9 controls; with county and year-month fixed effects; standard errors clustered by county. Table 2: Nightlight -0.2220 (0.0256).\\
814 & \texttt{estadd local year \textquotedbl{}YES\textquotedbl{}} & Adds a "Year-month FE = YES" marker to the stored results (printed at the bottom of the table).\\
815 & \texttt{estadd local county \textquotedbl{}YES\textquotedbl{}} & Adds a "County FE = YES" marker to the stored results.\\
816 & \texttt{est store s2} & Saves the last estimates under the name s2 so they can be put in a table.\\
817 & \texttt{reghdfe LAI PI \$control , absorb(county year)\newline \hspace*{1em}cluster(county)} & Regresses leaf area index on PSI (solar policy stringency); controlling for the 9 controls; with county and year-month fixed effects; standard errors clustered by county. Table 2: LAI +0.0039 (0.0007).\\
818 & \texttt{estadd local year \textquotedbl{}YES\textquotedbl{}} & Adds a "Year-month FE = YES" marker to the stored results (printed at the bottom of the table).\\
819 & \texttt{estadd local county \textquotedbl{}YES\textquotedbl{}} & Adds a "County FE = YES" marker to the stored results.\\
820 & \texttt{est store s3} & Saves the last estimates under the name s3 so they can be put in a table.\\
821 & \texttt{esttab s1 s2 s3 using mechenism.rtf, replace title(\textquotedbl{}Table\newline \hspace*{1em}2. Statistical evidence for the inferior greening\newline \hspace*{1em}mechanism.\textquotedbl{}) ///} & Exports the stored results (s1 s2 s3) into the Word table mechenism.rtf.\\
822 & \texttt{\hspace*{2.0em}mtitles(\textquotedbl{}NDVI\textquotedbl{} \textquotedbl{}Nightlight\textquotedbl{} \textquotedbl{}Leaf area\textquotedbl{}) ///} & Continuation of the previous command: column titles.\\
823 & \texttt{\hspace*{2.0em}stats(year county N r2\_a, labels(\textquotedbl{}Year-{}month FE\textquotedbl{} \textquotedbl{}County\newline \hspace*{1em}FE\textquotedbl{} \textquotedbl{}Observations\textquotedbl{} \textquotedbl{}R-{}squared\textquotedbl{}) fmt(0 0 0 4)) ///} & Continuation of the previous command: summary rows at the bottom (e.g. FE markers, observations, R-squared).\\
824 & \texttt{\hspace*{2.0em}star(* 0.10 ** 0.05 *** 0.01) nonotes nocons b(4) se\newline \hspace*{1em}compress nogaps ///} & Continuation of the previous command: significance stars * 10\%, ** 5\%, *** 1\%; 4-decimal coefficients with standard errors in parentheses; no default notes, constant omitted, compact layout.\\
825 & \texttt{\hspace*{2.0em}note(\textquotedbl{}Notes: This table presents the results o ... \textquotedbl{})} & Continuation: the table note printed under the table (text abridged here).\\
826 & \texttt{restore} & Restores the data saved by preserve (undoes changes made since).\\
830 & \texttt{*-{}-{}-{}-{}-{}-{}-{}-{}-{}-{}-{}-{}-{}-{}-{}-{}-{}-{}-{}-{}-{}-{}-{}-{}-{}-{}-{}-{}-{}-{}-{}-{}-{}-{}-{}-{}-{}-{}-{}-{}-{}-{}-{}-{}-{}-{}-{}-{}-{}-{}-{}-{}-{}-{}-{}-{}-{}\newline \hspace*{1em}-{}-{}-{}-{}-{}-{}-{}-{}-{}-{}-{}-{}-{}-{}-{}} & Divider line (comment).\\
831 & \texttt{* Further discussion} & Comment (not executed): explains the code that follows.\\
832 & \texttt{*-{}-{}-{}-{}-{}-{}-{}-{}-{}-{}-{}-{}-{}-{}-{}-{}-{}-{}-{}-{}-{}-{}-{}-{}-{}-{}-{}-{}-{}-{}-{}-{}-{}-{}-{}-{}-{}-{}-{}-{}-{}-{}-{}-{}-{}-{}-{}-{}-{}-{}-{}-{}-{}-{}-{}-{}-{}\newline \hspace*{1em}-{}-{}-{}-{}-{}-{}-{}-{}-{}-{}-{}-{}-{}-{}-{}} & Divider line (comment).\\
833 & \texttt{*Socioeconomic consequences} & Comment (not executed): explains the code that follows.\\
834 & \texttt{preserve} & Saves a copy of the data in memory, so it can be restored after temporary changes.\\
835 & \texttt{keep id YEAR PI} & Keeps county, year and PSI.\\
836 & \texttt{duplicates drop id YEAR, force} & One row per county-year.\\
837 & \texttt{xtset id YEAR} & Declares the county-year panel.\\
838 & \texttt{gen PI\_LAG1 = L1.PI} & Last year's PSI.\\
839 & \texttt{tempfile lag1\_data} & Temporary file.\\
840 & \texttt{save \textasciigrave{}lag1\_data\textquotesingle{}} & Saves it.\\
841 & \texttt{restore} & Restores the data saved by preserve (undoes changes made since).\\
842 & \texttt{merge m:1 id YEAR using \textasciigrave{}lag1\_data\textquotesingle{},keepus(PI\_LAG1)} & Attaches last year's PSI to every month.\\
843 & \texttt{drop \_merge} & Deletes Stata's merge-status variable so the next merge can create a new one.\\
845 & \texttt{preserve} & Saves a copy of the data in memory, so it can be restored after temporary changes.\\
846 & \texttt{keep id YEAR Output} & Keeps county, year and crop output.\\
847 & \texttt{duplicates drop id YEAR, force} & One row per county-year.\\
848 & \texttt{xtset id YEAR} & Declares the county-year panel.\\
849 & \texttt{gen Output\_lead = F.Output} & Next year's crop output.\\
850 & \texttt{tempfile lead\_data} & Temporary file.\\
851 & \texttt{save \textasciigrave{}lead\_data\textquotesingle{}} & Saves it.\\
852 & \texttt{restore} & Restores the data saved by preserve (undoes changes made since).\\
853 & \texttt{merge m:1 id YEAR using \textasciigrave{}lead\_data\textquotesingle{},keepus(Output\_lead)} & Attaches next year's output.\\
854 & \texttt{drop \_merge} & Deletes Stata's merge-status variable so the next merge can create a new one.\\
856 & \texttt{preserve} & Saves a copy of the data in memory, so it can be restored after temporary changes.\\
857 & \texttt{global control \textquotedbl{}Temp Wind Pop Duration Carbon Water Green\newline \hspace*{1em}Farm Grass\textquotedbl{}} & \$control = the authors' 9 control variables (Table 1, column 3).\\
858 & \texttt{reghdfe ShannonBD PI PI\_LAG1 \$control , absorb(county\newline \hspace*{1em}year) cluster(county) //PI\_LAG1即Leg one PI} & Table S19 col 1: Shannon on current and last year's PSI (flow vs lagged policy). Table S19 col 1: PSI -0.0105; lagged PSI not significant.\\
859 & \texttt{estadd local year \textquotedbl{}YES\textquotedbl{}} & Adds a "Year-month FE = YES" marker to the stored results (printed at the bottom of the table).\\
860 & \texttt{estadd local county \textquotedbl{}YES\textquotedbl{}} & Adds a "County FE = YES" marker to the stored results.\\
861 & \texttt{est store s1} & Saves the last estimates under the name s1 so they can be put in a table.\\
862 & \texttt{reghdfe Richness PI \$control , absorb(county year)\newline \hspace*{1em}cluster(county)} & Regresses species richness on PSI (solar policy stringency); controlling for the 9 controls; with county and year-month fixed effects; standard errors clustered by county. Table S19: richness -0.975.\\
863 & \texttt{estadd local year \textquotedbl{}YES\textquotedbl{}} & Adds a "Year-month FE = YES" marker to the stored results (printed at the bottom of the table).\\
864 & \texttt{estadd local county \textquotedbl{}YES\textquotedbl{}} & Adds a "County FE = YES" marker to the stored results.\\
865 & \texttt{est store s2} & Saves the last estimates under the name s2 so they can be put in a table.\\
866 & \texttt{reghdfe Evenness PI \$control , absorb(county year)\newline \hspace*{1em}cluster(county)} & Regresses species evenness on PSI (solar policy stringency); controlling for the 9 controls; with county and year-month fixed effects; standard errors clustered by county. Table S19: evenness +0.0018.\\
867 & \texttt{estadd local year \textquotedbl{}YES\textquotedbl{}} & Adds a "Year-month FE = YES" marker to the stored results (printed at the bottom of the table).\\
868 & \texttt{estadd local county \textquotedbl{}YES\textquotedbl{}} & Adds a "County FE = YES" marker to the stored results.\\
869 & \texttt{est store s3} & Saves the last estimates under the name s3 so they can be put in a table.\\
870 & \texttt{reghdfe Output\_lead PI ShannonBD \$control , absorb(county\newline \hspace*{1em}year) cluster(county)} & Table S19 col 4: next-year crop output on PSI and Shannon. Table S19: crop output -0.0288.\\
871 & \texttt{estadd local year \textquotedbl{}YES\textquotedbl{}} & Adds a "Year-month FE = YES" marker to the stored results (printed at the bottom of the table).\\
872 & \texttt{estadd local county \textquotedbl{}YES\textquotedbl{}} & Adds a "County FE = YES" marker to the stored results.\\
873 & \texttt{est store s4} & Saves the last estimates under the name s4 so they can be put in a table.\\
874 & \texttt{esttab s1 s2 s3 s4 using outcomes.rtf, replace\newline \hspace*{1em}title(\textquotedbl{}Table S19. Further discussion.\textquotedbl{}) ///} & Exports the stored results (s1 s2 s3 s4) into the Word table outcomes.rtf.\\
875 & \texttt{\hspace*{2.0em}mtitles(\textquotedbl{}ShannonBD\textquotedbl{} \textquotedbl{}Species richness\textquotedbl{} \textquotedbl{}Species evenness\textquotedbl{}\newline \hspace*{1em}\textquotedbl{}Agricultural crop yield\textquotedbl{}) ///} & Continuation of the previous command: column titles.\\
876 & \texttt{\hspace*{2.0em}stats(year county N r2\_a, labels(\textquotedbl{}Year-{}month FE\textquotedbl{} \textquotedbl{}County\newline \hspace*{1em}FE\textquotedbl{} \textquotedbl{}Observations\textquotedbl{} \textquotedbl{}R-{}squared\textquotedbl{}) fmt(0 0 0 4)) ///} & Continuation of the previous command: summary rows at the bottom (e.g. FE markers, observations, R-squared).\\
877 & \texttt{\hspace*{2.0em}star(* 0.10 ** 0.05 *** 0.01) nonotes nocons b(4) se\newline \hspace*{1em}compress nogaps ///} & Continuation of the previous command: significance stars * 10\%, ** 5\%, *** 1\%; 4-decimal coefficients with standard errors in parentheses; no default notes, constant omitted, compact layout.\\
878 & \texttt{\hspace*{2.0em}order(PI PI\_LAG1 ShannonBD \$control) ///} & Continuation of the previous command: row order.\\
879 & \texttt{\hspace*{2.0em}note(\textquotedbl{}Notes: Column (1) includes Leg one PSI ( ... \textquotedbl{})} & Continuation: the table note printed under the table (text abridged here).\\
880 & \texttt{restore} & Restores the data saved by preserve (undoes changes made since).\\
884 & \texttt{*-{}-{}-{}-{}-{}-{}-{}-{}-{}-{}-{}-{}-{}-{}-{}-{}-{}-{}-{}-{}-{}-{}-{}-{}-{}-{}-{}-{}-{}-{}-{}-{}-{}-{}-{}-{}-{}-{}-{}-{}-{}-{}-{}-{}-{}-{}-{}-{}-{}-{}-{}-{}-{}-{}-{}-{}-{}\newline \hspace*{1em}-{}-{}-{}-{}-{}-{}-{}-{}-{}-{}-{}-{}-{}-{}-{}} & Divider line (comment).\\
885 & \texttt{* CPU within-{}city temporal variation} & Comment (not executed): explains the code that follows.\\
886 & \texttt{*-{}-{}-{}-{}-{}-{}-{}-{}-{}-{}-{}-{}-{}-{}-{}-{}-{}-{}-{}-{}-{}-{}-{}-{}-{}-{}-{}-{}-{}-{}-{}-{}-{}-{}-{}-{}-{}-{}-{}-{}-{}-{}-{}-{}-{}-{}-{}-{}-{}-{}-{}-{}-{}-{}-{}-{}-{}\newline \hspace*{1em}-{}-{}-{}-{}-{}-{}-{}-{}-{}-{}-{}-{}-{}-{}-{}} & Divider line (comment).\\
887 & \texttt{preserve} & Saves a copy of the data in memory, so it can be restored after temporary changes.\\
888 & \texttt{drop if missing(CCPU)} & Keeps rows with the CPU index.\\
889 & \texttt{collapse (mean) CCPU, by(市 YEAR)} & City-year average CPU.\\
890 & \texttt{egen city = group(市)} & Numeric city id.\\
891 & \texttt{xtset city YEAR} & Declares the city-year panel.\\
892 & \texttt{xtsum CCPU} & Between-city and within-city variation of CPU.\\
894 & \texttt{matrix decomp = J(8,1,.)} & Empty 8x1 matrix for the variance decomposition.\\
895 & \texttt{matrix decomp[1,1] = r(n)} & Number of cities.\\
896 & \texttt{matrix decomp[2,1] = r(Tbar)} & Average years per city.\\
897 & \texttt{matrix decomp[3,1] = r(mean)} & Overall mean.\\
898 & \texttt{matrix decomp[4,1] = r(sd)} & Overall SD.\\
899 & \texttt{matrix decomp[5,1] = r(sd\_b)} & Between-city SD.\\
900 & \texttt{matrix decomp[6,1] = r(sd\_w)} & Within-city SD.\\
901 & \texttt{matrix decomp[7,1] = r(sd\_b)\^{}2 / r(sd)\^{}2 * 100} & Between-city share of variance (\%).\\
902 & \texttt{matrix decomp[8,1] = r(sd\_w)\^{}2 / r(sd)\^{}2 * 100} & Within-city share of variance (\%).\\
903 & \texttt{matrix rownames decomp = \textquotedbl{}Number of cities\textquotedbl{} \textquotedbl{}Avg periods\newline \hspace*{1em}per city\textquotedbl{} \textquotedbl{}Overall mean\textquotedbl{} \textquotedbl{}Overall SD\textquotedbl{} \textquotedbl{}Between-{}city SD\textquotedbl{}\newline \hspace*{1em}\textquotedbl{}Within-{}city SD\textquotedbl{} \textquotedbl{}Between-{}city variance share (\%)\textquotedbl{}\newline \hspace*{1em}\textquotedbl{}Within-{}city variance share (\%)\textquotedbl{}} & Names the 8 rows of the decomposition table.\\
904 & \texttt{matrix colnames decomp = \textquotedbl{}Value\textquotedbl{}} & Column name.\\
905 & \texttt{* Within/Between variance decomposition} & Comment (not executed): explains the code that follows.\\
906 & \texttt{esttab matrix(decomp, fmt(\%9.4f)) using\newline \hspace*{1em}CPU\_within\_variation.rtf, replace ///} & Exports the stored results (matrix(decomp, fmt(\%9.4f))) into the Word table CPU\_within\_variation.rtf.\\
907 & \texttt{\hspace*{2.0em}title(\textquotedbl{}Table R1. Within-{}city temporal variation of the CPU\newline \hspace*{1em}index: Variance decomposition.\textquotedbl{}) ///} & Continuation of the previous command: table title.\\
908 & \texttt{\hspace*{2.0em}nomtitle nonumber ///} & Continuation of the previous command: no model titles or numbers.\\
909 & \texttt{\hspace*{2.0em}note(\textquotedbl{}Notes: This table reports the variance d ... \textquotedbl{})} & Continuation: the table note printed under the table (text abridged here).\\
910 & \texttt{* Year-{}by-{}year summary} & Comment (not executed): explains the code that follows.\\
911 & \texttt{estpost tabstat CCPU, by(YEAR) stats(count mean sd min\newline \hspace*{1em}max) c(s)} & CPU summary by year (Table R2).\\
912 & \texttt{eststo cpu\_year} & Stores it.\\
913 & \texttt{esttab cpu\_year using CPU\_annual\_summary.rtf, replace ///} & Exports the stored results (cpu\_year) into the Word table CPU\_annual\_summary.rtf.\\
914 & \texttt{\hspace*{2.0em}title(\textquotedbl{}Table R2. Within-{}city temporal variation of the CPU\newline \hspace*{1em}index: Annual summary statistics.\textquotedbl{}) ///} & Continuation of the previous command: table title.\\
915 & \texttt{\hspace*{2.0em}cells(\textquotedbl{}count(fmt(\%9.0fc) label(N)) mean(fmt(\%9.4fc)\newline \hspace*{1em}label(Mean)) sd(fmt(\%9.4fc) label(S.D.)) min(fmt(\%9.4fc)\newline \hspace*{1em}label(Min)) max(fmt(\%9.4fc) label(Max))\textquotedbl{}) ///} & Continuation of the previous command: number formats for the bottom rows; which statistics to show and their formats.\\
916 & \texttt{\hspace*{2.0em}nomtitle nonumber noobs ///} & Continuation of the previous command: no model titles or numbers.\\
917 & \texttt{\hspace*{2.0em}note(\textquotedbl{}Notes: This table reports the annual sum ... \textquotedbl{})} & Continuation: the table note printed under the table (text abridged here).\\
918 & \texttt{restore} & Restores the data saved by preserve (undoes changes made since).\\
922 & \texttt{*-{}-{}-{}-{}-{}-{}-{}-{}-{}-{}-{}-{}-{}-{}-{}-{}-{}-{}-{}-{}-{}-{}-{}-{}-{}-{}-{}-{}-{}-{}-{}-{}-{}-{}-{}-{}-{}-{}-{}-{}-{}-{}-{}-{}-{}-{}-{}-{}-{}-{}-{}-{}-{}-{}-{}-{}-{}\newline \hspace*{1em}-{}-{}-{}-{}-{}-{}-{}-{}-{}-{}-{}-{}-{}-{}-{}} & Divider line (comment).\\
923 & \texttt{* PV-{}Area/IV} & Comment (not executed): explains the code that follows.\\
924 & \texttt{*-{}-{}-{}-{}-{}-{}-{}-{}-{}-{}-{}-{}-{}-{}-{}-{}-{}-{}-{}-{}-{}-{}-{}-{}-{}-{}-{}-{}-{}-{}-{}-{}-{}-{}-{}-{}-{}-{}-{}-{}-{}-{}-{}-{}-{}-{}-{}-{}-{}-{}-{}-{}-{}-{}-{}-{}-{}\newline \hspace*{1em}-{}-{}-{}-{}-{}-{}-{}-{}-{}-{}-{}-{}-{}-{}-{}} & Divider line (comment).\\
925 & \texttt{preserve} & Saves a copy of the data in memory, so it can be restored after temporary changes.\\
926 & \texttt{reghdfe ShannonBD PI Temp Wind Pop Duration Carbon Water\newline \hspace*{1em}Green Farm Grass, absorb(county year) cluster(county)} & Regresses Shannon bird diversity on PSI (solar policy stringency); controlling for Temp Wind Pop Duration Carbon Water Green Farm Grass; with county and year-month fixed effects; standard errors clustered by county. (Here it is re-run only to mark the regression sample for the county-year validation file.)\\
927 & \texttt{keep if e(sample)==1} & Keeps the regression sample.\\
928 & \texttt{foreach v of varlist ShannonBD PI Area Temp Wind Pop\newline \hspace*{1em}Duration Carbon Water Green Farm Grass Sun\_ccpu\{} & For each variable ...\\
929 & \texttt{\hspace*{2.0em}bysort id YEAR: egen temp\_mean = mean(\textasciigrave{}v\textquotesingle{})} & ... county-year average ...\\
930 & \texttt{\hspace*{2.0em}replace \textasciigrave{}v\textquotesingle{} = temp\_mean} & ... replaces the monthly value ...\\
931 & \texttt{\hspace*{2.0em}drop temp\_mean} & ... helper removed.\\
932 & \texttt{\}} & Closes the loop / block.\\
933 & \texttt{duplicates drop id YEAR, force} & One row per county-year.\\
934 & \texttt{keep 省 市 县 id county YEAR ShannonBD PI Area Temp Wind Pop\newline \hspace*{1em}Duration Carbon Water Green Farm Grass Sun\_ccpu} & Keeps these variables.\\
935 & \texttt{save \textquotedbl{}/Users/zhangaixin/Library/CloudStorage/OneDrive-{}个人/B\newline \hspace*{1em}D/DATA/baseline.dta\textquotedbl{},replace} & Saves a county-year file (authors' path).\\
936 & \texttt{restore} & Restores the data saved by preserve (undoes changes made since).\\
938 & \texttt{use \textquotedbl{}/Users/zhangaixin/Library/CloudStorage/OneDrive-{}个人/BD\newline \hspace*{1em}/DATA/baseline.dta\textquotedbl{},clear} & Loads it.\\
939 & \texttt{**PI与Area} & Comment: "PI and Area".\\
940 & \texttt{preserve} & Saves a copy of the data in memory, so it can be restored after temporary changes.\\
941 & \texttt{binscatter Area PI, controls(Temp Wind Pop Duration Carbon\newline \hspace*{1em}Water Green Farm Grass) nquantiles(50) line(lfit)\newline \hspace*{1em}xtitle(\textquotedbl{}PV policy stringency\textquotedbl{}) ytitle(\textquotedbl{}PV area (km²)\textquotedbl{})} & Binned scatter plot of PV area against PSI (50 bins, controls partialled out, fitted line): Fig. S2.\\
942 & \texttt{graph export \textquotedbl{}FIGS1.pdf\textquotedbl{}, replace} & Saves Fig. S2 (PSI vs PV area).\\
943 & \texttt{reg Area PI Temp Wind Pop Duration Carbon Water Green Farm\newline \hspace*{1em}Grass} & PV area on PSI and controls, no fixed effects (the paper's validation slope 0.0222). Validation slope 0.0222.\\
944 & \texttt{restore} & Restores the data saved by preserve (undoes changes made since).\\
945 & \texttt{***滞后滞前阶数} & Comment: "lag and lead orders".\\
946 & \texttt{preserve} & Saves a copy of the data in memory, so it can be restored after temporary changes.\\
947 & \texttt{keep id YEAR PI Area} & Keeps county, year, PSI, area.\\
948 & \texttt{duplicates drop id YEAR, force} & One row per county-year.\\
949 & \texttt{xtset id YEAR} & Declares the county-year panel.\\
950 & \texttt{gen Area\_l1 = L1.Area} & Area one year earlier.\\
951 & \texttt{gen Area\_l2 = L2.Area} & Area two years earlier.\\
952 & \texttt{gen Area\_l3 = L3.Area} & Area three years earlier.\\
953 & \texttt{gen Area\_f1 = F1.Area} & Area one year later.\\
954 & \texttt{gen Area\_f2 = F2.Area} & Area two years later.\\
955 & \texttt{gen Area\_f3 = F3.Area} & Area three years later.\\
956 & \texttt{tempfile leadlag\_data} & Temporary file.\\
957 & \texttt{save \textasciigrave{}leadlag\_data\textquotesingle{}} & Saves it.\\
958 & \texttt{restore} & Restores the data saved by preserve (undoes changes made since).\\
959 & \texttt{merge m:1 id YEAR using \textasciigrave{}leadlag\_data\textquotesingle{},keepus(Area\_l1\newline \hspace*{1em}Area\_l2 Area\_l3 Area\_f1 Area\_f2 Area\_f3)} & Attaches leads and lags of area.\\
960 & \texttt{drop \_merge} & Deletes Stata's merge-status variable so the next merge can create a new one.\\
962 & \texttt{pwcorr PI Area\_l1,sig} & Correlation of PSI with area one year earlier (with p-value).\\
963 & \texttt{pwcorr PI Area\_l2,sig} & ... two years earlier.\\
964 & \texttt{pwcorr PI Area\_l3,sig} & ... three years earlier.\\
965 & \texttt{pwcorr PI Area,sig} & ... same year (0.041).\\
966 & \texttt{pwcorr PI Area\_f1,sig} & ... one year later.\\
967 & \texttt{pwcorr PI Area\_f2,sig} & ... two years later (largest, 0.053: the paper's "2-year lag").\\
968 & \texttt{pwcorr PI Area\_f3,sig} & ... three years later.\\
970 & \texttt{**PI与IV} & Comment: "PI and the IV".\\
971 & \texttt{preserve} & Saves a copy of the data in memory, so it can be restored after temporary changes.\\
972 & \texttt{binscatter PI Sun\_ccpu, controls(Temp Wind Pop Duration\newline \hspace*{1em}Carbon Water Green Farm Grass) nquantiles(50) line(lfit)\newline \hspace*{1em}xtitle(\textquotedbl{}Ln Sunshine duration / City-{}level climate policy\newline \hspace*{1em}uncertainty\textquotedbl{}) ytitle(\textquotedbl{}PV policy stringency\textquotedbl{})} & Binned scatter of PSI against the instrument (50 bins, controls partialled out): Fig. S3.\\
973 & \texttt{graph export FIGS2.pdf, replace} & Saves Fig. S3 (PSI vs instrument).\\
974 & \texttt{reg PI Sun\_ccpu Temp Wind Pop Duration Carbon Water Green\newline \hspace*{1em}Farm Grass} & PSI on the instrument without fixed effects (slope 0.0401 shown in Fig. S3).\\
975 & \texttt{restore} & Restores the data saved by preserve (undoes changes made since).\\}
\section{Part III: Our Code}
\subsection{run\_all.do}
Runs the whole project (Tasks 1--3) with one command.

\codetable{1 & \texttt{*=========================================================\newline \hspace*{1em}=================} & Divider line (comment).\\
2 & \texttt{* run\_all.do  —  reproduces the whole project from the\newline \hspace*{1em}original data} & Comment (not executed): explains the code that follows.\\
3 & \texttt{*   Task 1  Replication of Zhang et al. (2026)        \textasciitilde{}6\newline \hspace*{1em}min} & Comment (not executed): explains the code that follows.\\
4 & \texttt{*   Task 2  Cinelli-{}Hazlett sensitivity analysis      \textasciitilde{}1\newline \hspace*{1em}min} & Comment (not executed): explains the code that follows.\\
5 & \texttt{*   Task 3  Staggered difference-{}in-{}differences       \textasciitilde{}3\newline \hspace*{1em}min} & Comment (not executed): explains the code that follows.\\
6 & \texttt{* Usage (Stata command window):} & Comment (not executed): explains the code that follows.\\
7 & \texttt{*   do \textquotedbl{}/Users/sahiljangra/Desktop/Johns\newline \hspace*{1em}Hopkins/RA/Solar\_Birds\_Replication/run\_all.do\textquotedbl{}} & Comment (not executed): explains the code that follows.\\
8 & \texttt{*=========================================================\newline \hspace*{1em}=================} & Divider line (comment).\\
9 & \texttt{clear all} & Removes data (and with "all", also programs and stored results) from memory.\\
10 & \texttt{global ROOT \textquotedbl{}/Users/sahiljangra/Desktop/Johns\newline \hspace*{1em}Hopkins/RA/Solar\_Birds\_Replication\textquotedbl{}} & Defines \$ROOT = the project folder on the Mac. All other paths are built from it.\\
11 & \texttt{do \textquotedbl{}\$ROOT/01\_Task1\_Replication/code/1\_0\_run\_task1.do\textquotedbl{}} & Runs another do-file: \$ROOT/01\_Task1\_Replication/code/1\_0\_run\_task1.do\\
12 & \texttt{do \textquotedbl{}\$ROOT/02\_Task2\_Sensitivity/code/2\_0\_run\_task2.do\textquotedbl{}} & Runs another do-file: \$ROOT/02\_Task2\_Sensitivity/code/2\_0\_run\_task2.do\\
13 & \texttt{do \textquotedbl{}\$ROOT/03\_Task3\_Staggered\_DiD/code/3\_0\_run\_task3.do\textquotedbl{}} & Runs another do-file: \$ROOT/03\_Task3\_Staggered\_DiD/code/3\_0\_run\_task3.do\\
14 & \texttt{di as res \_n \textquotedbl{}ALL TASKS FINISHED\textquotedbl{}} & Prints a message when all three tasks have run.\\}
\subsection{00\_config.do}
Settings used by every step: folder paths, the control-variable lists, and automatic installation of the Stata packages.

\codetable{1 & \texttt{*=========================================================\newline \hspace*{1em}=================} & Divider line (comment).\\
2 & \texttt{* 00\_config.do} & Comment (not executed): explains the code that follows.\\
3 & \texttt{* Project : Replication of Zhang et al. (2026), \textquotedbl{}China\textquotesingle{}s\newline \hspace*{1em}solar expansion} & Comment (not executed): explains the code that follows.\\
4 & \texttt{*           policy reduces bird diversity\textquotedbl{}, Science 393,\newline \hspace*{1em}831.} & Comment (not executed): explains the code that follows.\\
5 & \texttt{* Purpose : One place for paths, globals and package\newline \hspace*{1em}installation.} & Comment (not executed): explains the code that follows.\\
6 & \texttt{*           Every task do-{}file starts with:  do\newline \hspace*{1em}\textquotedbl{}<ROOT>/00\_config.do\textquotedbl{}} & Comment (not executed): explains the code that follows.\\
7 & \texttt{* Author  : Sahil (Research Assistant, Johns Hopkins\newline \hspace*{1em}University)} & Comment (not executed): explains the code that follows.\\
8 & \texttt{*=========================================================\newline \hspace*{1em}=================} & Divider line (comment).\\
9 & \texttt{version 16} & Tells Stata to interpret the code with version-16 syntax, so results do not change with newer Stata versions.\\
10 & \texttt{set more off} & Stops Stata from pausing with --more-- when the screen fills.\\
11 & \texttt{set linesize 120} & Sets the width of the output to 120 characters.\\
13 & \texttt{*-{}-{}-{}-{}-{}-{}-{}-{}-{}-{}-{}-{}-{}-{}-{}-{}-{}-{}-{}-{}-{}-{}-{}-{}-{}-{}-{}-{}-{}-{}-{}-{}-{}-{}-{}-{}-{}-{}-{}-{}-{}-{}-{}-{}-{}-{}-{}-{}-{}-{}-{}-{}-{}-{}-{}-{}-{}\newline \hspace*{1em}-{}-{}-{}-{}-{}-{}-{} paths} & Divider line (comment).\\
14 & \texttt{global ROOT \textquotedbl{}/Users/sahiljangra/Desktop/Johns\newline \hspace*{1em}Hopkins/RA/Solar\_Birds\_Replication\textquotedbl{}} & Defines \$ROOT = the project folder on the Mac. All other paths are built from it.\\
15 & \texttt{global ORIG \textquotedbl{}\$ROOT/00\_Original\_Package\textquotedbl{}} & \$ORIG = folder with the untouched authors' package.\\
16 & \texttt{global DATA \textquotedbl{}\$ORIG/DATA\textquotedbl{}} & \$DATA = the authors' 46 .dta files.\\
17 & \texttt{global T1   \textquotedbl{}\$ROOT/01\_Task1\_Replication\textquotedbl{}} & \$T1 = Task 1 folder.\\
18 & \texttt{global T2   \textquotedbl{}\$ROOT/02\_Task2\_Sensitivity\textquotedbl{}} & \$T2 = Task 2 folder.\\
19 & \texttt{global T3   \textquotedbl{}\$ROOT/03\_Task3\_Staggered\_DiD\textquotedbl{}} & \$T3 = Task 3 folder.\\
21 & \texttt{* analysis sample built in Task 1 (used by all tasks)} & Comment (not executed): explains the code that follows.\\
22 & \texttt{global SAMPLE \textquotedbl{}\$T1/data/analysis\_sample.dta\textquotedbl{}} & \$SAMPLE = the analysis dataset built in step 1\_2; every later step loads it.\\
24 & \texttt{*-{}-{}-{}-{}-{}-{}-{}-{}-{}-{}-{}-{}-{}-{}-{}-{}-{}-{}-{}-{}-{}-{}-{}-{}-{}-{}-{}-{}-{}-{}-{}-{}-{}-{}-{}-{}-{}-{}-{}-{}-{}-{}-{}-{}-{}-{}-{}-{}-{}-{}-{}-{}-{}-{}-{}-{}-{}\newline \hspace*{1em}-{}-{}-{}-{}-{}-{}-{} model} & Divider line (comment).\\
25 & \texttt{* the authors\textquotesingle{} nine control variables (Table 1, column 3)} & Comment (not executed): explains the code that follows.\\
26 & \texttt{global control     \textquotedbl{}Temp Wind Pop Duration Carbon Water\newline \hspace*{1em}Green Farm Grass\textquotedbl{}} & \$control = the authors' 9 control variables (Table 1, column 3).\\
27 & \texttt{global control\_noD \textquotedbl{}Temp Wind Pop Carbon Water Green Farm\newline \hspace*{1em}Grass\textquotedbl{}} & \$control\_noD = the same controls without Duration; used when Duration (effort) is the outcome or is being tested as a bad control.\\
29 & \texttt{*-{}-{}-{}-{}-{}-{}-{}-{}-{}-{}-{}-{}-{}-{}-{}-{}-{}-{}-{}-{}-{}-{}-{}-{}-{}-{}-{}-{}-{}-{}-{}-{}-{}-{}-{}-{}-{}-{}-{}-{}-{}-{}-{}-{}-{}-{}-{}-{}-{}-{}-{}-{}-{}-{}-{}-{}-{}\newline \hspace*{1em}-{}-{}-{}-{}-{}-{}-{} packages} & Divider line (comment).\\
30 & \texttt{* installed only if missing} & Comment (not executed): explains the code that follows.\\
31 & \texttt{foreach p in require ftools reghdfe ivreg2 ranktest\newline \hspace*{1em}ivreghdfe estout winsor2 ///} & Starts a loop over every user-written package the project needs (continues on the next lines).\\
32 & \texttt{\hspace*{4.0em}psmatch2 binscatter sensemakr drdid csdid did\_imputation\newline \hspace*{1em}///} & Continuation of the package list.\\
33 & \texttt{\hspace*{4.0em}event\_plot avar eventstudyinteract \{} & End of the package list; the loop body follows.\\
34 & \texttt{\hspace*{2.0em}capture which \textasciigrave{}p\textquotesingle{}} & Checks whether package p is installed (capture = do not stop if it is not).\\
35 & \texttt{\hspace*{2.0em}if \_rc \{} & If that fails ...\\
36 & \texttt{\hspace*{4.0em}di as txt \textquotedbl{}installing \textasciigrave{}p\textquotesingle{} ...\textquotedbl{}} & ... print a message ...\\
37 & \texttt{\hspace*{4.0em}capture noisily ssc install \textasciigrave{}p\textquotesingle{}, replace} & ... and install it from the SSC archive.\\
38 & \texttt{\hspace*{2.0em}\}} & Closes the loop / block.\\
39 & \texttt{\}} & Closes the loop / block.\\
40 & \texttt{set scheme s2color} & Sets the default graph colour scheme.\\}
\subsection{1\_0\_run\_task1.do}
Runs the six Task 1 steps in order, each with its own log file; one failing step does not stop the others.

\codetable{1 & \texttt{*=========================================================\newline \hspace*{1em}=================} & Divider line (comment).\\
2 & \texttt{* 1\_0\_run\_task1.do  —  runs every Task 1 step in order,\newline \hspace*{1em}one log per step} & Comment (not executed): explains the code that follows.\\
3 & \texttt{*=========================================================\newline \hspace*{1em}=================} & Divider line (comment).\\
4 & \texttt{do \textquotedbl{}/Users/sahiljangra/Desktop/Johns\newline \hspace*{1em}Hopkins/RA/Solar\_Birds\_Replication/00\_config.do\textquotedbl{}} & Runs another do-file: /Users/sahiljangra/Desktop/Johns Hopkins/RA/Solar\_Birds\_Replication/00\_config.do\\
6 & \texttt{foreach step in 1\_1\_reproduce\_authors\_results\newline \hspace*{1em}1\_2\_build\_analysis\_sample ///} & Starts a loop: foreach step in 1\_1\_reproduce\_authors\_results 1\_2\_build\_analysis\_sample ///.\\
7 & \texttt{\hspace*{4.0em}1\_3\_data\_structure 1\_4\_reproduction\_check\newline \hspace*{1em}///} & Continuation: more Task 1 step names.\\
8 & \texttt{\hspace*{4.0em}1\_5\_twfe\_diagnostics 1\_6\_iv\_exclusion\_tests \{} & Continuation: last Task 1 step names; the loop body follows.\\
9 & \texttt{\hspace*{2.0em}capture log close \_all} & Closes any open log (capture = no error if none is open).\\
10 & \texttt{\hspace*{2.0em}log using \textquotedbl{}\$T1/logs/\textasciigrave{}step\textquotesingle{}.log\textquotedbl{}, replace text} & Opens a text log file that records all output of this step.\\
11 & \texttt{\hspace*{2.0em}di as res \_n \textquotedbl{}================ \textasciigrave{}step\textquotesingle{} ================\textquotedbl{}} & Prints the name of the step being run.\\
12 & \texttt{\hspace*{2.0em}clear} & Removes data (and with "all", also programs and stored results) from memory.\\
13 & \texttt{\hspace*{2.0em}estimates clear} & Clears stored estimates.\\
14 & \texttt{\hspace*{2.0em}do \textquotedbl{}\$ROOT/00\_config.do\textquotedbl{}} & Re-runs the settings file so every step starts with the same paths and macros.\\
15 & \texttt{\hspace*{2.0em}capture noisily do \textquotedbl{}\$T1/code/\textasciigrave{}step\textquotesingle{}.do\textquotedbl{}} & Runs the step's do-file; capture = an error does not stop the whole run, noisily = output is still shown.\\
16 & \texttt{\hspace*{2.0em}if \_rc di as err \textquotedbl{}STEP \textasciigrave{}step\textquotesingle{} FAILED with rc = \textquotedbl{} \_rc} & If the step failed (\_rc is the error code), prints a FAILED message.\\
17 & \texttt{\hspace*{2.0em}log close} & Closes the log file.\\
18 & \texttt{\}} & Closes the loop / block.\\}
\subsection{1\_1\_reproduce\_authors\_results.do}
Runs the authors' full code (only paths changed) to regenerate all their tables and figures.

\codetable{1 & \texttt{*=========================================================\newline \hspace*{1em}=================} & Divider line (comment).\\
2 & \texttt{* 1\_1\_reproduce\_authors\_results.do} & Comment (not executed): explains the code that follows.\\
3 & \texttt{* PURPOSE : Run the authors\textquotesingle{} complete replication code\newline \hspace*{1em}(paths changed only)} & Comment (not executed): explains the code that follows.\\
4 & \texttt{*           and produce every table and figure in the\newline \hspace*{1em}paper.} & Comment (not executed): explains the code that follows.\\
5 & \texttt{* INPUT   : 00\_Original\_Package/DATA (all .dta files)} & Comment (not executed): explains the code that follows.\\
6 & \texttt{* CODE    : 1\_1a\_authors\_code\_localpaths.do (authors\textquotesingle{}\newline \hspace*{1em}CODE.do; [PATH] lines)} & Comment (not executed): explains the code that follows.\\
7 & \texttt{* OUTPUT  : tables/authors\_tables (.rtf)  (Tables 1-{}2, S4,\newline \hspace*{1em}S7-{}S19, R1-{}R2)} & Comment (not executed): explains the code that follows.\\
8 & \texttt{*           figures/authors\_figures (.pdf) (Figs S2, S3)} & Comment (not executed): explains the code that follows.\\
9 & \texttt{* WHY     : Task 1, item 3 — check the published numbers\newline \hspace*{1em}can be reproduced} & Comment (not executed): explains the code that follows.\\
10 & \texttt{*           from the authors\textquotesingle{} own code and data before\newline \hspace*{1em}analysing them.} & Comment (not executed): explains the code that follows.\\
11 & \texttt{*=========================================================\newline \hspace*{1em}=================} & Divider line (comment).\\
12 & \texttt{do \textquotedbl{}\$T1/code/1\_1a\_authors\_code\_localpaths.do\textquotedbl{}} & Runs the authors' complete CODE.do (copy with only the 6 path lines changed), producing all their tables and figures.\\}
\subsection{1\_2\_build\_analysis\_sample.do}
Builds the analysis dataset exactly as the authors do (Part A, their code verbatim) and adds our own flags (Part B): whether the county is in the policy file, urban/rural, cumulative PSI, and ids.

\codetable{1 & \texttt{*=========================================================\newline \hspace*{1em}=================} & Divider line (comment).\\
2 & \texttt{* 1\_2\_build\_analysis\_sample.do} & Comment (not executed): explains the code that follows.\\
3 & \texttt{* PURPOSE : Build the authors\textquotesingle{} analysis sample (46,528\newline \hspace*{1em}county-{}months) using} & Comment (not executed): explains the code that follows.\\
4 & \texttt{*           their own data-{}preparation code VERBATIM (Part\newline \hspace*{1em}A), then add the} & Comment (not executed): explains the code that follows.\\
5 & \texttt{*           flags our later checks need (Part B). Saves\newline \hspace*{1em}data/analysis\_sample.dta} & Comment (not executed): explains the code that follows.\\
6 & \texttt{* WHY     : Every later step (Tasks 1-{}3) must start from\newline \hspace*{1em}exactly the sample} & Comment (not executed): explains the code that follows.\\
7 & \texttt{*           the authors analysed.} & Comment (not executed): explains the code that follows.\\
8 & \texttt{*=========================================================\newline \hspace*{1em}=================} & Divider line (comment).\\
10 & \texttt{*-{}-{}-{}-{}-{}-{}-{}-{}-{}-{}-{}-{}-{}-{}-{}-{}-{}-{}-{}-{}-{}-{}-{}-{}-{}-{}-{}-{}-{}-{}-{}-{}-{}-{}-{}-{}-{}-{}-{}-{}-{}-{}-{}-{}-{}-{}-{}-{}-{}-{}-{}-{}-{}-{}-{}-{}-{}\newline \hspace*{1em}-{}-{}-{}-{}-{}-{}-{}-{}-{}-{}-{}-{}-{}-{}-{}-{}} & Divider line (comment).\\
11 & \texttt{* PART A — authors\textquotesingle{} data preparation (CODE.do lines 1-{}245,\newline \hspace*{1em}unchanged)} & Comment (not executed): explains the code that follows.\\
12 & \texttt{*-{}-{}-{}-{}-{}-{}-{}-{}-{}-{}-{}-{}-{}-{}-{}-{}-{}-{}-{}-{}-{}-{}-{}-{}-{}-{}-{}-{}-{}-{}-{}-{}-{}-{}-{}-{}-{}-{}-{}-{}-{}-{}-{}-{}-{}-{}-{}-{}-{}-{}-{}-{}-{}-{}-{}-{}-{}\newline \hspace*{1em}-{}-{}-{}-{}-{}-{}-{}-{}-{}-{}-{}-{}-{}-{}-{}-{}} & Divider line (comment).\\
13 & \texttt{*-{}-{}-{}-{}-{}-{}-{}-{}-{}-{}-{}-{}-{}-{}-{}-{}-{}-{}-{}-{}-{}-{}-{}-{}-{}-{}-{}-{}-{}-{}-{}-{}-{}-{}-{}-{}-{}-{}-{}-{}-{}-{}-{}-{}-{}-{}-{}-{}-{}-{}-{}-{}-{}-{}-{}-{}-{}\newline \hspace*{1em}-{}-{}-{}-{}-{}-{}-{}-{}-{}-{}-{}-{}-{}-{}-{}} & Divider line (comment).\\
14 & \texttt{* Data preparation} & Comment (not executed): explains the code that follows.\\
15 & \texttt{*-{}-{}-{}-{}-{}-{}-{}-{}-{}-{}-{}-{}-{}-{}-{}-{}-{}-{}-{}-{}-{}-{}-{}-{}-{}-{}-{}-{}-{}-{}-{}-{}-{}-{}-{}-{}-{}-{}-{}-{}-{}-{}-{}-{}-{}-{}-{}-{}-{}-{}-{}-{}-{}-{}-{}-{}-{}\newline \hspace*{1em}-{}-{}-{}-{}-{}-{}-{}-{}-{}-{}-{}-{}-{}-{}-{}} & Divider line (comment).\\
16 & \texttt{cd \textquotedbl{}\$DATA\textquotedbl{}   // [PATH] originally the authors\textquotesingle{} local DATA\newline \hspace*{1em}folder} & Changes the working directory to the authors' own folder (DATA); later file names are read from / saved there. This is one of the 6 path lines we changed.\\
259 & \texttt{*-{}-{}-{}-{}-{}-{}-{}-{}-{}-{}-{}-{}-{}-{}-{}-{}-{}-{}-{}-{}-{}-{}-{}-{}-{}-{}-{}-{}-{}-{}-{}-{}-{}-{}-{}-{}-{}-{}-{}-{}-{}-{}-{}-{}-{}-{}-{}-{}-{}-{}-{}-{}-{}-{}-{}-{}-{}\newline \hspace*{1em}-{}-{}-{}-{}-{}-{}-{}-{}-{}-{}-{}-{}-{}-{}-{}-{}} & Divider line (comment).\\
260 & \texttt{* PART B — flags added for our analysis (not in the\newline \hspace*{1em}authors\textquotesingle{} code)} & Comment (not executed): explains the code that follows.\\
261 & \texttt{*-{}-{}-{}-{}-{}-{}-{}-{}-{}-{}-{}-{}-{}-{}-{}-{}-{}-{}-{}-{}-{}-{}-{}-{}-{}-{}-{}-{}-{}-{}-{}-{}-{}-{}-{}-{}-{}-{}-{}-{}-{}-{}-{}-{}-{}-{}-{}-{}-{}-{}-{}-{}-{}-{}-{}-{}-{}\newline \hspace*{1em}-{}-{}-{}-{}-{}-{}-{}-{}-{}-{}-{}-{}-{}-{}-{}-{}} & Divider line (comment).\\
262 & \texttt{* (1) Does the county appear in the PSI source file at\newline \hspace*{1em}all?} & Comment (not executed): explains the code that follows.\\
263 & \texttt{*     The authors set PSI = 0 whenever it is missing (line\newline \hspace*{1em}\textquotedbl{}replace PI = 0\textquotedbl{}).} & Comment (not executed): explains the code that follows.\\
264 & \texttt{preserve} & Saves a copy of the data in memory, so it can be restored after temporary changes.\\
265 & \texttt{\hspace*{2.0em}use \textquotedbl{}\$DATA/PI.dta\textquotedbl{}, clear} & Loads the authors' policy file PI.dta.\\
266 & \texttt{\hspace*{2.0em}keep id} & Keeps only the county code.\\
267 & \texttt{\hspace*{2.0em}duplicates drop} & Removes duplicate rows, leaving one row per county code.\\
268 & \texttt{\hspace*{2.0em}gen byte in\_PI\_file = 1} & Marks every county that appears in PI.dta with 1.\\
269 & \texttt{\hspace*{2.0em}tempfile pif} & Creates a temporary file name for this list.\\
270 & \texttt{\hspace*{2.0em}save \textasciigrave{}pif\textquotesingle{}} & Saves the list of counties present in PI.dta.\\
271 & \texttt{restore} & Restores the data saved by preserve (undoes changes made since).\\
272 & \texttt{merge m:1 id using \textasciigrave{}pif\textquotesingle{}, keep(master match) nogen} & Attaches the in\_PI\_file flag to the analysis sample by county code (keeps all sample rows; nogen = no \_merge variable).\\
273 & \texttt{replace in\_PI\_file = 0 if missing(in\_PI\_file)} & Counties not found in PI.dta get in\_PI\_file = 0 (their PSI was forced to 0).\\
274 & \texttt{label var in\_PI\_file \textquotedbl{}County present in PI.dta (1) or PSI\newline \hspace*{1em}imputed as 0 (0)\textquotedbl{}} & Attaches a descriptive label to a variable.\\
276 & \texttt{* (2) Urban district vs rural county (Chinese admin code:\newline \hspace*{1em}last two digits} & Comment (not executed): explains the code that follows.\\
277 & \texttt{*     01-{}20 = urban district \textquotesingle{}shixiaqu\textquotesingle{}; 21+ = county /\newline \hspace*{1em}county-{}level city)} & Comment (not executed): explains the code that follows.\\
278 & \texttt{gen byte urban = mod(id,100) <= 20} & urban = 1 if the last two digits of the county code are 01-20 (Chinese urban district, shixiaqu); mod(id,100) gives those last two digits.\\
279 & \texttt{label var urban \textquotedbl{}Urban district (admin code ending 01-{}20)\textquotedbl{}} & Attaches a descriptive label to a variable.\\
280 & \texttt{* name check: urban districts are named ...区 (district)} & Comment (not executed): explains the code that follows.\\
281 & \texttt{gen byte name\_is\_district = ustrregexm(县, \textquotedbl{}区\$\textquotedbl{})} & 1 if the county name ends with the character qu (district): a second check of the urban rule.\\
282 & \texttt{label var name\_is\_district \textquotedbl{}County name ends in 区\newline \hspace*{1em}(district)\textquotedbl{}} & Attaches a descriptive label to a variable.\\
284 & \texttt{* (3) Alternative PSI definitions kept in PI.dta but\newline \hspace*{1em}unused by the authors} & Comment (not executed): explains the code that follows.\\
285 & \texttt{merge m:1 id YEAR using \textquotedbl{}\$DATA/PI.dta\textquotedbl{}, keepusing(PI\_stock\newline \hspace*{1em}PI\_total) keep(master match) nogen} & Brings in the cumulative PSI versions (PI\_stock, PI\_total) that exist in PI.dta but the authors never used.\\
286 & \texttt{gen byte any\_PSI = (PI > 0)} & any\_PSI = 1 if the county has any solar policy that year.\\
287 & \texttt{label var any\_PSI \textquotedbl{}PSI > 0 in this year\textquotedbl{}} & Attaches a descriptive label to a variable.\\
289 & \texttt{* (4) Identifiers} & Comment (not executed): explains the code that follows.\\
290 & \texttt{egen prov = group(省)} & Numeric province id (1-31) from the province name.\\
291 & \texttt{egen city = group(市)} & Numeric city id.\\
292 & \texttt{gen month\_of\_year = month} & Copy of the calendar month (1-12), used to remove seasonality in Task 3.\\
293 & \texttt{label var prov \textquotedbl{}Province id\textquotedbl{}} & Attaches a descriptive label to a variable.\\
294 & \texttt{label var city \textquotedbl{}Prefecture city id\textquotedbl{}} & Attaches a descriptive label to a variable.\\
296 & \texttt{compress} & Stores variables in the smallest storage type to save memory.\\
297 & \texttt{save \textquotedbl{}\$SAMPLE\textquotedbl{}, replace} & Saves the analysis dataset used by every later step.\\
298 & \texttt{di as res \textquotedbl{}Analysis sample saved: \textquotedbl{} \_N \textquotedbl{} county-{}months\textquotedbl{}} & Prints the number of rows saved (46,528).\\}
\subsection{1\_3\_data\_structure.do}
Computes every fact about the data used in the reports (counts, PSI distribution, missing-PSI pattern, urban check, file counts).

\codetable{1 & \texttt{*=========================================================\newline \hspace*{1em}=================} & Divider line (comment).\\
2 & \texttt{* 1\_3\_data\_structure.do} & Comment (not executed): explains the code that follows.\\
3 & \texttt{* PURPOSE : Document the structure of the data and how\newline \hspace*{1em}variables behave.} & Comment (not executed): explains the code that follows.\\
4 & \texttt{*           Every fact quoted in the Task 1 report\newline \hspace*{1em}(Section 2) is computed here.} & Comment (not executed): explains the code that follows.\\
5 & \texttt{* OUTPUT  : tables/1\_3\_data\_structure\_facts.csv  (one row\newline \hspace*{1em}per fact)} & Comment (not executed): explains the code that follows.\\
6 & \texttt{* WHY     : Task 1, items 1-{}2 — understand the panel and\newline \hspace*{1em}the treatment} & Comment (not executed): explains the code that follows.\\
7 & \texttt{*           before interpreting any regression.} & Comment (not executed): explains the code that follows.\\
8 & \texttt{*=========================================================\newline \hspace*{1em}=================} & Divider line (comment).\\
9 & \texttt{tempname F} & Creates a temporary name (handle) for a results file or matrix.\\
10 & \texttt{postfile \textasciigrave{}F\textquotesingle{} str70 fact double value using\newline \hspace*{1em}\textquotedbl{}\$T1/tables/1\_3\_facts\_tmp.dta\textquotedbl{}, replace} & Opens a results file with these columns; each post adds a row.\\
12 & \texttt{*-{}-{}-{}-{}-{}-{}-{}-{}-{}-{}-{}-{}-{}-{}-{}-{}-{}-{}-{}-{}-{}-{}-{}-{}-{}-{}-{}-{}-{}-{}-{}-{}-{}-{}-{}-{}-{}-{}-{}-{}-{}-{}-{}-{}-{}-{}-{}-{} A. The\newline \hspace*{1em}raw bird file} & Divider line (comment).\\
13 & \texttt{use \textquotedbl{}\$DATA/Bird.dta\textquotedbl{}, clear} & Loads the raw bird file (48,900 county-months).\\
14 & \texttt{post \textasciigrave{}F\textquotesingle{} (\textquotedbl{}Bird.dta: county-{}month rows\textquotedbl{}) (\_N)} & Writes one row (label and numbers) to the results file opened with postfile.\\
15 & \texttt{egen tagc = tag(id)} & Marks the first row of each county (to count counties, not rows).\\
16 & \texttt{count if tagc} & Counts rows that satisfy the condition (stored in r(N)).\\
17 & \texttt{post \textasciigrave{}F\textquotesingle{} (\textquotedbl{}Bird.dta: counties\textquotedbl{}) (r(N))} & Writes one row (label and numbers) to the results file opened with postfile.\\
18 & \texttt{di \_n \textquotedbl{}Records (county-{}months) per year in Bird.dta:\textquotedbl{}} & Displays a message or result on screen.\\
19 & \texttt{tab YEAR} & Tabulates frequencies.\\
20 & \texttt{forvalues y = 2014/2023 \{} & Starts a loop: forvalues y = 2014/2023.\\
21 & \texttt{\hspace*{2.0em}count if YEAR == \textasciigrave{}y\textquotesingle{}} & Counts rows that satisfy the condition (stored in r(N)).\\
22 & \texttt{\hspace*{2.0em}post \textasciigrave{}F\textquotesingle{} (\textquotedbl{}Bird.dta: county-{}months in \textasciigrave{}y\textquotesingle{}\textquotedbl{}) (r(N))} & Writes one row (label and numbers) to the results file opened with postfile.\\
23 & \texttt{\}} & Closes the loop / block.\\
25 & \texttt{*-{}-{}-{}-{}-{}-{}-{}-{}-{}-{}-{}-{}-{}-{}-{}-{}-{}-{}-{}-{}-{}-{}-{}-{}-{}-{}-{}-{}-{}-{}-{}-{}-{}-{}-{}-{}-{}-{}-{}-{}-{}-{}-{}-{}-{}-{}-{}-{} B. The\newline \hspace*{1em}analysis sample} & Divider line (comment).\\
26 & \texttt{use \textquotedbl{}\$SAMPLE\textquotedbl{}, clear} & Loads the analysis dataset built in step 1\_2.\\
27 & \texttt{post \textasciigrave{}F\textquotesingle{} (\textquotedbl{}Analysis sample: county-{}months\textquotedbl{}) (\_N)} & Writes one row (label and numbers) to the results file opened with postfile.\\
28 & \texttt{egen tagc = tag(county)} & Marks one row per county.\\
29 & \texttt{count if tagc} & Counts rows that satisfy the condition (stored in r(N)).\\
30 & \texttt{post \textasciigrave{}F\textquotesingle{} (\textquotedbl{}Analysis sample: counties\textquotedbl{}) (r(N))} & Writes one row (label and numbers) to the results file opened with postfile.\\
31 & \texttt{bysort county: gen nobs = \_N} & Number of months each county is observed.\\
32 & \texttt{summarize nobs if tagc, detail} & Computes summary statistics (mean, SD, min, max, percentiles) and stores them in r().\\
33 & \texttt{post \textasciigrave{}F\textquotesingle{} (\textquotedbl{}Months observed per county: median\textquotedbl{}) (r(p50))} & Writes one row (label and numbers) to the results file opened with postfile.\\
34 & \texttt{post \textasciigrave{}F\textquotesingle{} (\textquotedbl{}Months observed per county: min\textquotedbl{}) (r(min))} & Writes one row (label and numbers) to the results file opened with postfile.\\
35 & \texttt{post \textasciigrave{}F\textquotesingle{} (\textquotedbl{}Months observed per county: max\textquotedbl{}) (r(max))} & Writes one row (label and numbers) to the results file opened with postfile.\\
36 & \texttt{count if tagc \& nobs == 1} & Counts rows that satisfy the condition (stored in r(N)).\\
37 & \texttt{post \textasciigrave{}F\textquotesingle{} (\textquotedbl{}Singleton counties (1 month only)\textquotedbl{}) (r(N))} & Writes one row (label and numbers) to the results file opened with postfile.\\
39 & \texttt{*-{}-{}-{}-{}-{}-{}-{}-{}-{}-{}-{}-{}-{}-{}-{}-{}-{}-{}-{}-{}-{}-{}-{}-{}-{}-{}-{}-{}-{}-{}-{}-{}-{}-{}-{}-{}-{}-{}-{}-{}-{}-{}-{}-{}-{}-{}-{}-{} C. The\newline \hspace*{1em}treatment PSI} & Divider line (comment).\\
40 & \texttt{summarize PI} & Computes summary statistics (mean, SD, min, max, percentiles) and stores them in r().\\
41 & \texttt{post \textasciigrave{}F\textquotesingle{} (\textquotedbl{}PSI mean (winsorized)\textquotedbl{}) (r(mean))} & Writes one row (label and numbers) to the results file opened with postfile.\\
42 & \texttt{post \textasciigrave{}F\textquotesingle{} (\textquotedbl{}PSI SD\textquotedbl{}) (r(sd))} & Writes one row (label and numbers) to the results file opened with postfile.\\
43 & \texttt{post \textasciigrave{}F\textquotesingle{} (\textquotedbl{}PSI max\textquotedbl{}) (r(max))} & Writes one row (label and numbers) to the results file opened with postfile.\\
44 & \texttt{count if PI == 0} & Counts rows that satisfy the condition (stored in r(N)).\\
45 & \texttt{post \textasciigrave{}F\textquotesingle{} (\textquotedbl{}Share of county-{}months with PSI = 0\textquotedbl{}) (r(N)/\_N)} & Writes one row (label and numbers) to the results file opened with postfile.\\
46 & \texttt{di \_n \textquotedbl{}PSI by year (median is zero every year):\textquotedbl{}} & Displays a message or result on screen.\\
47 & \texttt{tabstat PI, by(YEAR) stats(mean p50 p75 max) format(\%6.2f)} & Tabulates frequencies.\\
48 & \texttt{forvalues y = 2014/2023 \{} & Starts a loop: forvalues y = 2014/2023.\\
49 & \texttt{\hspace*{2.0em}quietly summarize PI if YEAR==\textasciigrave{}y\textquotesingle{}, detail} & Computes summary statistics (mean, SD, min, max, percentiles) and stores them in r().\\
50 & \texttt{\hspace*{2.0em}post \textasciigrave{}F\textquotesingle{} (\textquotedbl{}PSI median in \textasciigrave{}y\textquotesingle{}\textquotedbl{}) (r(p50))} & Writes one row (label and numbers) to the results file opened with postfile.\\
51 & \texttt{\}} & Closes the loop / block.\\
52 & \texttt{bysort county: egen sdPI = sd(PI)} & Standard deviation of PSI within each county (0 = PSI never changes).\\
53 & \texttt{count if tagc \& (sdPI == 0 | missing(sdPI))} & Counts rows that satisfy the condition (stored in r(N)).\\
54 & \texttt{post \textasciigrave{}F\textquotesingle{} (\textquotedbl{}Counties whose PSI never changes in sample\textquotedbl{})\newline \hspace*{1em}(r(N))} & Writes one row (label and numbers) to the results file opened with postfile.\\
55 & \texttt{corr PI Area} & Correlation between PSI and PV station area (result 0.072).\\
56 & \texttt{post \textasciigrave{}F\textquotesingle{} (\textquotedbl{}Correlation PSI vs PV area (county-{}month)\textquotedbl{})\newline \hspace*{1em}(r(rho))} & Writes one row (label and numbers) to the results file opened with postfile.\\
58 & \texttt{* PSI is a FLOW: it falls from one year to the next (raw\newline \hspace*{1em}PI.dta)} & Comment (not executed): explains the code that follows.\\
59 & \texttt{preserve} & Saves a copy of the data in memory, so it can be restored after temporary changes.\\
60 & \texttt{\hspace*{2.0em}use \textquotedbl{}\$DATA/PI.dta\textquotedbl{}, clear} & Loads the authors' policy file PI.dta.\\
61 & \texttt{\hspace*{2.0em}xtset id YEAR} & Declares the county-year panel.\\
62 & \texttt{\hspace*{2.0em}gen dPI = D.PI} & Year-to-year change in PSI (D. = first difference).\\
63 & \texttt{\hspace*{2.0em}count if dPI < 0} & Counts rows that satisfy the condition (stored in r(N)).\\
64 & \texttt{\hspace*{2.0em}local ndown = r(N)} & Stores the number of county-years where PSI fell.\\
65 & \texttt{\hspace*{2.0em}count if !missing(dPI)} & Counts rows that satisfy the condition (stored in r(N)).\\
66 & \texttt{\hspace*{2.0em}post \textasciigrave{}F\textquotesingle{} (\textquotedbl{}PI.dta: share of county-{}years where PSI falls\textquotedbl{})\newline \hspace*{1em}(\textasciigrave{}ndown\textquotesingle{}/r(N))} & Writes one row (label and numbers) to the results file opened with postfile.\\
67 & \texttt{\hspace*{2.0em}egen tagp = tag(id)} & Marks one row per county in PI.dta.\\
68 & \texttt{\hspace*{2.0em}count if tagp} & Counts rows that satisfy the condition (stored in r(N)).\\
69 & \texttt{\hspace*{2.0em}post \textasciigrave{}F\textquotesingle{} (\textquotedbl{}PI.dta: counties covered\textquotedbl{}) (r(N))} & Writes one row (label and numbers) to the results file opened with postfile.\\
70 & \texttt{\hspace*{2.0em}corr PI PI\_stock PI\_total} & Correlations between the yearly PSI and its cumulative versions.\\
71 & \texttt{restore} & Restores the data saved by preserve (undoes changes made since).\\
73 & \texttt{*-{}-{}-{}-{}-{}-{}-{}-{}-{}-{}-{}-{}-{}-{}-{}-{}-{}-{}-{}-{}-{}-{}-{}-{}-{}-{}-{}-{}-{}-{}-{}-{}-{}-{}-{}-{}-{}-{}-{}-{}-{}-{}-{}-{}-{}-{}-{}-{} D.\newline \hspace*{1em}Missing PSI coded as zero} & Divider line (comment).\\
74 & \texttt{count if tagc \& in\_PI\_file == 0} & Counts rows that satisfy the condition (stored in r(N)).\\
75 & \texttt{post \textasciigrave{}F\textquotesingle{} (\textquotedbl{}Counties absent from PI.dta (PSI forced to 0)\textquotedbl{})\newline \hspace*{1em}(r(N))} & Writes one row (label and numbers) to the results file opened with postfile.\\
76 & \texttt{count if in\_PI\_file == 0} & Counts rows that satisfy the condition (stored in r(N)).\\
77 & \texttt{post \textasciigrave{}F\textquotesingle{} (\textquotedbl{}Share of county-{}months with PSI imputed 0\textquotedbl{})\newline \hspace*{1em}(r(N)/\_N)} & Writes one row (label and numbers) to the results file opened with postfile.\\
78 & \texttt{di \_n \textquotedbl{}Coverage of PSI data by unit type:\textquotedbl{}} & Displays a message or result on screen.\\
79 & \texttt{tab urban in\_PI\_file, row} & Tabulates frequencies.\\
80 & \texttt{count if urban==1} & Counts rows that satisfy the condition (stored in r(N)).\\
81 & \texttt{local nu = r(N)} & Stores the number of urban rows.\\
82 & \texttt{count if urban==1 \& in\_PI\_file==0} & Counts rows that satisfy the condition (stored in r(N)).\\
83 & \texttt{post \textasciigrave{}F\textquotesingle{} (\textquotedbl{}Urban districts: share of obs with PSI imputed\newline \hspace*{1em}0\textquotedbl{}) (r(N)/\textasciigrave{}nu\textquotesingle{})} & Writes one row (label and numbers) to the results file opened with postfile.\\
84 & \texttt{count if urban==0} & Counts rows that satisfy the condition (stored in r(N)).\\
85 & \texttt{local nr = r(N)} & Stores the number of rural rows.\\
86 & \texttt{count if urban==0 \& in\_PI\_file==0} & Counts rows that satisfy the condition (stored in r(N)).\\
87 & \texttt{post \textasciigrave{}F\textquotesingle{} (\textquotedbl{}Rural counties: share of obs with PSI imputed\newline \hspace*{1em}0\textquotedbl{}) (r(N)/\textasciigrave{}nr\textquotesingle{})} & Writes one row (label and numbers) to the results file opened with postfile.\\
88 & \texttt{count if urban==1} & Counts rows that satisfy the condition (stored in r(N)).\\
89 & \texttt{post \textasciigrave{}F\textquotesingle{} (\textquotedbl{}Urban-{}district observations\textquotedbl{}) (r(N))} & Writes one row (label and numbers) to the results file opened with postfile.\\
90 & \texttt{count if urban==0} & Counts rows that satisfy the condition (stored in r(N)).\\
91 & \texttt{post \textasciigrave{}F\textquotesingle{} (\textquotedbl{}Rural-{}county observations\textquotedbl{}) (r(N))} & Writes one row (label and numbers) to the results file opened with postfile.\\
93 & \texttt{* validate the urban code rule with county names (…区 =\newline \hspace*{1em}district)} & Comment (not executed): explains the code that follows.\\
94 & \texttt{count if tagc \& urban==1} & Counts rows that satisfy the condition (stored in r(N)).\\
95 & \texttt{local cu = r(N)} & Stores the number of urban-coded counties.\\
96 & \texttt{count if tagc \& urban==1 \& name\_is\_district==1} & Counts rows that satisfy the condition (stored in r(N)).\\
97 & \texttt{post \textasciigrave{}F\textquotesingle{} (\textquotedbl{}Urban-{}coded units named ...区: share\textquotedbl{})\newline \hspace*{1em}(r(N)/\textasciigrave{}cu\textquotesingle{})} & Writes one row (label and numbers) to the results file opened with postfile.\\
98 & \texttt{count if tagc \& urban==0} & Counts rows that satisfy the condition (stored in r(N)).\\
99 & \texttt{local cr = r(N)} & Stores the number of rural-coded counties.\\
100 & \texttt{count if tagc \& urban==0 \& name\_is\_district==1} & Counts rows that satisfy the condition (stored in r(N)).\\
101 & \texttt{post \textasciigrave{}F\textquotesingle{} (\textquotedbl{}Rural-{}coded units named ...区: share\textquotedbl{})\newline \hspace*{1em}(r(N)/\textasciigrave{}cr\textquotesingle{})} & Writes one row (label and numbers) to the results file opened with postfile.\\
102 & \texttt{count if tagc \& in\_PI\_file==0} & Counts rows that satisfy the condition (stored in r(N)).\\
103 & \texttt{local cm = r(N)} & Stores the number of missing-PSI counties.\\
104 & \texttt{count if tagc \& in\_PI\_file==0 \& name\_is\_district==1} & Counts rows that satisfy the condition (stored in r(N)).\\
105 & \texttt{post \textasciigrave{}F\textquotesingle{} (\textquotedbl{}Missing-{}PSI counties named ...区: share\textquotedbl{})\newline \hspace*{1em}(r(N)/\textasciigrave{}cm\textquotesingle{})} & Writes one row (label and numbers) to the results file opened with postfile.\\
107 & \texttt{* is it a merge failure? do missing counties appear in\newline \hspace*{1em}PI.dta under their NAME?} & Comment (not executed): explains the code that follows.\\
108 & \texttt{preserve} & Saves a copy of the data in memory, so it can be restored after temporary changes.\\
109 & \texttt{\hspace*{2.0em}use \textquotedbl{}\$DATA/PI.dta\textquotedbl{}, clear} & Loads the authors' policy file PI.dta.\\
110 & \texttt{\hspace*{2.0em}keep 县} & Keeps only the county name.\\
111 & \texttt{\hspace*{2.0em}duplicates drop} & Removes duplicate rows, leaving one row per county code.\\
112 & \texttt{\hspace*{2.0em}gen byte name\_in\_PI = 1} & Flags names present in PI.dta.\\
113 & \texttt{\hspace*{2.0em}tempfile pin} & Temporary file for the list of names.\\
114 & \texttt{\hspace*{2.0em}save \textasciigrave{}pin\textquotesingle{}} & Saves the list of county names in PI.dta.\\
115 & \texttt{restore} & Restores the data saved by preserve (undoes changes made since).\\
116 & \texttt{preserve} & Saves a copy of the data in memory, so it can be restored after temporary changes.\\
117 & \texttt{\hspace*{2.0em}keep if tagc \& in\_PI\_file==0} & Keeps one row per missing-PSI county.\\
118 & \texttt{\hspace*{2.0em}merge m:1 县 using \textasciigrave{}pin\textquotesingle{}, keep(master match) nogen} & Checks whether those counties appear in PI.dta under their name (result: 0, so it is not a code-merge failure).\\
119 & \texttt{\hspace*{2.0em}count if name\_in\_PI == 1} & Counts rows that satisfy the condition (stored in r(N)).\\
120 & \texttt{\hspace*{2.0em}post \textasciigrave{}F\textquotesingle{} (\textquotedbl{}Missing-{}PSI counties whose NAME is in PI.dta\textquotedbl{})\newline \hspace*{1em}(r(N))} & Writes one row (label and numbers) to the results file opened with postfile.\\
121 & \texttt{restore} & Restores the data saved by preserve (undoes changes made since).\\
123 & \texttt{di \_n \textquotedbl{}Missing-{}PSI share by province:\textquotedbl{}} & Displays a message or result on screen.\\
124 & \texttt{tab 省 in\_PI\_file if tagc, row nofreq} & Tabulates frequencies.\\
126 & \texttt{*-{}-{}-{}-{}-{}-{}-{}-{}-{}-{}-{}-{}-{}-{}-{}-{}-{}-{}-{}-{}-{}-{}-{}-{}-{}-{}-{}-{}-{}-{}-{}-{}-{}-{}-{}-{}-{}-{}-{}-{}-{}-{}-{}-{}-{}-{}-{}-{} E.\newline \hspace*{1em}Package and variable facts} & Divider line (comment).\\
127 & \texttt{local dtafiles : dir \textquotedbl{}\$DATA\textquotedbl{} files \textquotedbl{}*.dta\textquotedbl{}} & Lists all .dta files in the authors' DATA folder.\\
128 & \texttt{post \textasciigrave{}F\textquotesingle{} (\textquotedbl{}Number of .dta files in the authors\textquotesingle{} package\textquotedbl{})\newline \hspace*{1em}(\textasciigrave{}: word count \textasciigrave{}dtafiles\textquotesingle{}\textquotesingle{})} & Writes one row (label and numbers) to the results file opened with postfile.\\
129 & \texttt{tempname fh} & Handle for reading a text file.\\
130 & \texttt{file open \textasciigrave{}fh\textquotesingle{} using \textquotedbl{}\$ORIG/authors\_CODE.do\textquotedbl{}, read text} & Opens the authors' CODE.do as text ...\\
131 & \texttt{local nl = 0} & ... starts a line counter ...\\
132 & \texttt{file read \textasciigrave{}fh\textquotesingle{} line} & ... reads one line ...\\
133 & \texttt{while r(eof) == 0 \{} & ... and loops until the end of the file ...\\
134 & \texttt{\hspace*{2.0em}local ++nl} & ... adding 1 per line ...\\
135 & \texttt{\hspace*{2.0em}file read \textasciigrave{}fh\textquotesingle{} line} & ... reads one line ...\\
136 & \texttt{\}} & Closes the loop / block.\\
137 & \texttt{file close \textasciigrave{}fh\textquotesingle{}} & ... then closes the file (result: 975 lines).\\
138 & \texttt{post \textasciigrave{}F\textquotesingle{} (\textquotedbl{}Lines in authors\textquotesingle{} CODE.do\textquotedbl{}) (\textasciigrave{}nl\textquotesingle{})} & Writes one row (label and numbers) to the results file opened with postfile.\\
139 & \texttt{use \textquotedbl{}\$SAMPLE\textquotedbl{}, clear} & Loads the analysis dataset built in step 1\_2.\\
140 & \texttt{summarize ShannonBD} & Computes summary statistics (mean, SD, min, max, percentiles) and stores them in r().\\
141 & \texttt{post \textasciigrave{}F\textquotesingle{} (\textquotedbl{}ShannonBD mean\textquotedbl{}) (r(mean))} & Writes one row (label and numbers) to the results file opened with postfile.\\
142 & \texttt{post \textasciigrave{}F\textquotesingle{} (\textquotedbl{}ShannonBD SD\textquotedbl{}) (r(sd))} & Writes one row (label and numbers) to the results file opened with postfile.\\
143 & \texttt{summarize Nightlight} & Computes summary statistics (mean, SD, min, max, percentiles) and stores them in r().\\
144 & \texttt{post \textasciigrave{}F\textquotesingle{} (\textquotedbl{}Nightlight min (winsorized)\textquotedbl{}) (r(min))} & Writes one row (label and numbers) to the results file opened with postfile.\\
145 & \texttt{post \textasciigrave{}F\textquotesingle{} (\textquotedbl{}Nightlight max (winsorized)\textquotedbl{}) (r(max))} & Writes one row (label and numbers) to the results file opened with postfile.\\
146 & \texttt{use \textquotedbl{}\$DATA/NL.dta\textquotedbl{}, clear} & Loads the night-light file to check its year range.\\
147 & \texttt{summarize YEAR} & Computes summary statistics (mean, SD, min, max, percentiles) and stores them in r().\\
148 & \texttt{post \textasciigrave{}F\textquotesingle{} (\textquotedbl{}NL.dta first year\textquotedbl{}) (r(min))} & Writes one row (label and numbers) to the results file opened with postfile.\\
149 & \texttt{post \textasciigrave{}F\textquotesingle{} (\textquotedbl{}NL.dta last year\textquotedbl{}) (r(max))} & Writes one row (label and numbers) to the results file opened with postfile.\\
150 & \texttt{summarize Nightlight} & Computes summary statistics (mean, SD, min, max, percentiles) and stores them in r().\\
151 & \texttt{post \textasciigrave{}F\textquotesingle{} (\textquotedbl{}NL.dta Nightlight max (raw)\textquotedbl{}) (r(max))} & Writes one row (label and numbers) to the results file opened with postfile.\\
153 & \texttt{postclose \textasciigrave{}F\textquotesingle{}} & Closes the results file so it can be loaded.\\
154 & \texttt{use \textquotedbl{}\$T1/tables/1\_3\_facts\_tmp.dta\textquotedbl{}, clear} & Loads the collected facts.\\
155 & \texttt{format value \%14.4f} & Sets display formats (decimals) for these variables.\\
156 & \texttt{list, noobs sep(0) abbreviate(70)} & Lists the rows/variables on screen.\\
157 & \texttt{export delimited using\newline \hspace*{1em}\textquotedbl{}\$T1/tables/1\_3\_data\_structure\_facts.csv\textquotedbl{}, replace} & Exports the current data to a CSV file.\\
158 & \texttt{erase \textquotedbl{}\$T1/tables/1\_3\_facts\_tmp.dta\textquotedbl{}} & Deletes a temporary file.\\}
\subsection{1\_4\_reproduction\_check.do}
Re-estimates 12 published numbers and compares them with the paper to four decimals (result: 0 mismatches).

\codetable{1 & \texttt{*=========================================================\newline \hspace*{1em}=================} & Divider line (comment).\\
2 & \texttt{* 1\_4\_reproduction\_check.do} & Comment (not executed): explains the code that follows.\\
3 & \texttt{* PURPOSE : Side-{}by-{}side comparison of the paper\textquotesingle{}s\newline \hspace*{1em}published estimates with} & Comment (not executed): explains the code that follows.\\
4 & \texttt{*           our re-{}estimation from the analysis sample.} & Comment (not executed): explains the code that follows.\\
5 & \texttt{* OUTPUT  : tables/1\_4\_reproduction\_check.csv / .rtf} & Comment (not executed): explains the code that follows.\\
6 & \texttt{* WHY     : Task 1, item 3 — a single table proving the\newline \hspace*{1em}results reproduce.} & Comment (not executed): explains the code that follows.\\
7 & \texttt{*           Paper values typed from Zhang et al. (2026),\newline \hspace*{1em}Tables 1, 2, S7, S11.} & Comment (not executed): explains the code that follows.\\
8 & \texttt{*=========================================================\newline \hspace*{1em}=================} & Divider line (comment).\\
9 & \texttt{use \textquotedbl{}\$SAMPLE\textquotedbl{}, clear} & Loads the analysis dataset built in step 1\_2.\\
11 & \texttt{tempname R} & Creates a temporary name (handle) for a results file or matrix.\\
12 & \texttt{postfile \textasciigrave{}R\textquotesingle{} str45 result double(paper\_b paper\_se our\_b\newline \hspace*{1em}our\_se our\_N) ///} & Opens a results file with columns: result name, paper estimate, paper SE, our estimate, our SE, our N ...\\
13 & \texttt{\hspace*{2.0em}using \textquotedbl{}\$T1/tables/1\_4\_tmp.dta\textquotedbl{}, replace} & ... saved to a temporary file.\\
15 & \texttt{* Table 1 (three columns)} & Comment (not executed): explains the code that follows.\\
16 & \texttt{reghdfe ShannonBD PI, absorb(county year) cluster(county)} & Regresses Shannon bird diversity on PSI (solar policy stringency); with county and year-month fixed effects; standard errors clustered by county.\\
17 & \texttt{post \textasciigrave{}R\textquotesingle{} (\textquotedbl{}Table 1 col 1: PSI, no controls\textquotedbl{}) (-{}0.0157)\newline \hspace*{1em}(0.0037) (\_b[PI]) (\_se[PI]) (e(N))} & Writes one row (label and numbers) to the results file opened with postfile.\\
18 & \texttt{reghdfe ShannonBD PI Temp Wind Pop Duration Carbon,\newline \hspace*{1em}absorb(county year) cluster(county)} & Regresses Shannon bird diversity on PSI (solar policy stringency); controlling for Temp Wind Pop Duration Carbon; with county and year-month fixed effects; standard errors clustered by county.\\
19 & \texttt{post \textasciigrave{}R\textquotesingle{} (\textquotedbl{}Table 1 col 2: PSI, 5 controls\textquotedbl{}) (-{}0.0127)\newline \hspace*{1em}(0.0036) (\_b[PI]) (\_se[PI]) (e(N))} & Writes one row (label and numbers) to the results file opened with postfile.\\
20 & \texttt{reghdfe ShannonBD PI \$control, absorb(county year)\newline \hspace*{1em}cluster(county)} & Regresses Shannon bird diversity on PSI (solar policy stringency); controlling for the 9 controls; with county and year-month fixed effects; standard errors clustered by county.\\
21 & \texttt{post \textasciigrave{}R\textquotesingle{} (\textquotedbl{}Table 1 col 3: PSI, 9 controls (MAIN)\textquotedbl{})\newline \hspace*{1em}(-{}0.0125) (0.0037) (\_b[PI]) (\_se[PI]) (e(N))} & Writes one row (label and numbers) to the results file opened with postfile.\\
23 & \texttt{* Table S7 (IV)} & Comment (not executed): explains the code that follows.\\
24 & \texttt{reghdfe PI Sun\_ccpu \$control, absorb(county year)\newline \hspace*{1em}cluster(county) keepsingleton} & Regresses PSI (solar policy stringency) on the instrument ln(sunshine)/CPU; controlling for the 9 controls; with county and year-month fixed effects; standard errors clustered by county; keeping single-observation counties.\\
25 & \texttt{post \textasciigrave{}R\textquotesingle{} (\textquotedbl{}Table S7 first stage: Sun\_ccpu\textquotedbl{}) (0.0252)\newline \hspace*{1em}(0.0046) (\_b[Sun\_ccpu]) (\_se[Sun\_ccpu]) (e(N))} & Writes one row (label and numbers) to the results file opened with postfile.\\
26 & \texttt{ivreghdfe ShannonBD (PI = Sun\_ccpu) \$control,\newline \hspace*{1em}absorb(county year) cluster(county) keepsingleton} & Two-stage least squares (IV): Shannon bird diversity on PSI, with PSI instrumented by Sun\_ccpu; controls: the 9 controls; with county and year-month fixed effects; standard errors clustered by county; keeping single-observation counties.\\
27 & \texttt{post \textasciigrave{}R\textquotesingle{} (\textquotedbl{}Table S7 second stage: PSI\textquotedbl{}) (-{}0.1264) (0.0347)\newline \hspace*{1em}(\_b[PI]) (\_se[PI]) (e(N))} & Writes one row (label and numbers) to the results file opened with postfile.\\
28 & \texttt{di \textquotedbl{}Kleibergen-{}Paap F = \textquotedbl{} e(widstat) \textquotedbl{}   (paper: 30.181)\textquotedbl{}} & Displays a message or result on screen.\\
30 & \texttt{* Table S11 (PSM, authors\textquotesingle{} implementation)} & Comment (not executed): explains the code that follows.\\
31 & \texttt{gen treated = (PI > 0)} & Creates the treatment dummy for matching: 1 if PSI > 0, else 0.\\
32 & \texttt{psmatch2 treated \$control i.year, outcome(ShannonBD)\newline \hspace*{1em}neighbor(1) caliper(0.05) common} & Propensity-score matching: estimates the probability of PSI > 0 from the controls and year-month dummies, then matches each treated row to its 1 nearest untreated neighbour(s) within a caliper of 0.05; common = keep only common support. Creates \_weight (matching weights).\\
33 & \texttt{reghdfe ShannonBD PI \$control if \_weight != .,\newline \hspace*{1em}absorb(county year) cluster(county)} & Regresses Shannon bird diversity on PSI (solar policy stringency); controlling for the 9 controls; with county and year-month fixed effects; standard errors clustered by county; only observations kept by the matching.\\
34 & \texttt{post \textasciigrave{}R\textquotesingle{} (\textquotedbl{}Table S11 col 1: 1:1 NN matching\textquotedbl{}) (-{}0.0072)\newline \hspace*{1em}(0.0039) (\_b[PI]) (\_se[PI]) (e(N))} & Writes one row (label and numbers) to the results file opened with postfile.\\
35 & \texttt{psmatch2 treated \$control i.year, outcome(ShannonBD)\newline \hspace*{1em}neighbor(2) caliper(0.05) common} & Propensity-score matching: estimates the probability of PSI > 0 from the controls and year-month dummies, then matches each treated row to its 2 nearest untreated neighbour(s) within a caliper of 0.05; common = keep only common support. Creates \_weight (matching weights).\\
36 & \texttt{reghdfe ShannonBD PI \$control if \_weight != .,\newline \hspace*{1em}absorb(county year) cluster(county)} & Regresses Shannon bird diversity on PSI (solar policy stringency); controlling for the 9 controls; with county and year-month fixed effects; standard errors clustered by county; only observations kept by the matching.\\
37 & \texttt{post \textasciigrave{}R\textquotesingle{} (\textquotedbl{}Table S11 col 2: 1:2 NN matching\textquotedbl{}) (-{}0.0083)\newline \hspace*{1em}(0.0038) (\_b[PI]) (\_se[PI]) (e(N))} & Writes one row (label and numbers) to the results file opened with postfile.\\
38 & \texttt{psmatch2 treated \$control i.year, outcome(ShannonBD)\newline \hspace*{1em}kernel common} & Kernel propensity-score matching: estimates each row's probability of PSI > 0 from the controls and year-month dummies, then weights all untreated rows by closeness; common = keep only common support. Creates \_weight.\\
39 & \texttt{reghdfe ShannonBD PI \$control if \_weight != .,\newline \hspace*{1em}absorb(county year) cluster(county)} & Regresses Shannon bird diversity on PSI (solar policy stringency); controlling for the 9 controls; with county and year-month fixed effects; standard errors clustered by county; only observations kept by the matching.\\
40 & \texttt{post \textasciigrave{}R\textquotesingle{} (\textquotedbl{}Table S11 col 3: kernel matching\textquotedbl{}) (-{}0.0125)\newline \hspace*{1em}(0.0037) (\_b[PI]) (\_se[PI]) (e(N))} & Writes one row (label and numbers) to the results file opened with postfile.\\
41 & \texttt{drop treated \_*} & Deletes the matching variables created by psmatch2.\\
43 & \texttt{* Table S12 col 8 (Simpson index)} & Comment (not executed): explains the code that follows.\\
44 & \texttt{reghdfe SimpsonBD PI \$control, absorb(county year)\newline \hspace*{1em}cluster(county)} & Regresses Simpson bird diversity on PSI (solar policy stringency); controlling for the 9 controls; with county and year-month fixed effects; standard errors clustered by county.\\
45 & \texttt{post \textasciigrave{}R\textquotesingle{} (\textquotedbl{}Table S12 col 8: SimpsonBD\textquotedbl{}) (-{}0.0014) (0.0008)\newline \hspace*{1em}(\_b[PI]) (\_se[PI]) (e(N))} & Writes one row (label and numbers) to the results file opened with postfile.\\
46 & \texttt{di as res \textquotedbl{}SimpsonBD: p-{}value = \textquotedbl{} \%6.4f 2*ttail(e(df\_r),\newline \hspace*{1em}abs(\_b[PI]/\_se[PI])) \textquotedbl{}  (paper: p < 0.10)\textquotedbl{}} & Displays a message or result on screen.\\
48 & \texttt{* Table 2 (mechanisms)} & Comment (not executed): explains the code that follows.\\
49 & \texttt{reghdfe Ndvi PI \$control, absorb(county year)\newline \hspace*{1em}cluster(county)} & Regresses NDVI (greenness) on PSI (solar policy stringency); controlling for the 9 controls; with county and year-month fixed effects; standard errors clustered by county.\\
50 & \texttt{post \textasciigrave{}R\textquotesingle{} (\textquotedbl{}Table 2: NDVI\textquotedbl{}) (-{}0.0609) (0.0138) (\_b[PI])\newline \hspace*{1em}(\_se[PI]) (e(N))} & Writes one row (label and numbers) to the results file opened with postfile.\\
51 & \texttt{reghdfe Nightlight PI \$control, absorb(county year)\newline \hspace*{1em}cluster(county)} & Regresses night-time light on PSI (solar policy stringency); controlling for the 9 controls; with county and year-month fixed effects; standard errors clustered by county.\\
52 & \texttt{post \textasciigrave{}R\textquotesingle{} (\textquotedbl{}Table 2: Nightlight\textquotedbl{}) (-{}0.2220) (0.0256)\newline \hspace*{1em}(\_b[PI]) (\_se[PI]) (e(N))} & Writes one row (label and numbers) to the results file opened with postfile.\\
53 & \texttt{reghdfe LAI PI \$control, absorb(county year)\newline \hspace*{1em}cluster(county)} & Regresses leaf area index on PSI (solar policy stringency); controlling for the 9 controls; with county and year-month fixed effects; standard errors clustered by county.\\
54 & \texttt{post \textasciigrave{}R\textquotesingle{} (\textquotedbl{}Table 2: Leaf area index\textquotedbl{}) (0.0039) (0.0007)\newline \hspace*{1em}(\_b[PI]) (\_se[PI]) (e(N))} & Writes one row (label and numbers) to the results file opened with postfile.\\
56 & \texttt{postclose \textasciigrave{}R\textquotesingle{}} & Closes the results file so it can be loaded.\\
57 & \texttt{use \textquotedbl{}\$T1/tables/1\_4\_tmp.dta\textquotedbl{}, clear} & Loads the comparison table.\\
58 & \texttt{gen diff\_b  = round(our\_b,  0.0001) -{} paper\_b} & Difference between our estimate (rounded to 4 decimals) and the paper.\\
59 & \texttt{gen diff\_se = round(our\_se, 0.0001) -{} paper\_se} & Same for the standard error.\\
60 & \texttt{gen match = (abs(diff\_b) < 0.00005 \& abs(diff\_se) <\newline \hspace*{1em}0.00005)} & match = 1 if both coefficient and SE agree to 4 decimals.\\
61 & \texttt{format paper\_b paper\_se our\_b our\_se diff\_b diff\_se \%9.4f} & Sets display formats (decimals) for these variables.\\
62 & \texttt{list, noobs sep(0) abbreviate(40)} & Lists the rows/variables on screen.\\
63 & \texttt{count if match == 0} & Counts rows that satisfy the condition (stored in r(N)).\\
64 & \texttt{di as res \textquotedbl{}Results NOT matching the paper (to 4 decimals):\newline \hspace*{1em}\textquotedbl{} r(N)} & Displays a message or result on screen.\\
65 & \texttt{export delimited using\newline \hspace*{1em}\textquotedbl{}\$T1/tables/1\_4\_reproduction\_check.csv\textquotedbl{}, replace} & Exports the current data to a CSV file.\\
66 & \texttt{erase \textquotedbl{}\$T1/tables/1\_4\_tmp.dta\textquotedbl{}} & Deletes a temporary file.\\}
\subsection{1\_5\_twfe\_diagnostics.do}
Tests whether the TWFE comparison is fair (Panel A), alternative treatment definitions (B), observer effort as a bad control (C), and matching with the weights actually applied (D).

\codetable{1 & \texttt{*=========================================================\newline \hspace*{1em}=================} & Divider line (comment).\\
2 & \texttt{* 1\_5\_twfe\_diagnostics.do} & Comment (not executed): explains the code that follows.\\
3 & \texttt{* PURPOSE : Probe the identification of the main TWFE\newline \hspace*{1em}estimate.} & Comment (not executed): explains the code that follows.\\
4 & \texttt{*   Panel A  Who drives the effect? (missing-{}PSI counties,\newline \hspace*{1em}urban vs rural)} & Comment (not executed): explains the code that follows.\\
5 & \texttt{*   Panel B  Alternative treatment definitions (stock,\newline \hspace*{1em}binary)} & Comment (not executed): explains the code that follows.\\
6 & \texttt{*   Panel C  Observer effort (Duration) as a possible \textquotedbl{}bad\newline \hspace*{1em}control\textquotedbl{}} & Comment (not executed): explains the code that follows.\\
7 & \texttt{*   Panel D  Propensity-{}score matching with the matching\newline \hspace*{1em}weights APPLIED} & Comment (not executed): explains the code that follows.\\
8 & \texttt{* OUTPUT  : tables/1\_5\_twfe\_diagnostics.csv / .rtf} & Comment (not executed): explains the code that follows.\\
9 & \texttt{* WHY     : Task 1, item 5 — the TWFE design is only\newline \hspace*{1em}credible if the} & Comment (not executed): explains the code that follows.\\
10 & \texttt{*           comparison group is valid; these checks test\newline \hspace*{1em}that directly.} & Comment (not executed): explains the code that follows.\\
11 & \texttt{*=========================================================\newline \hspace*{1em}=================} & Divider line (comment).\\
12 & \texttt{use \textquotedbl{}\$SAMPLE\textquotedbl{}, clear} & Loads the analysis dataset built in step 1\_2.\\
14 & \texttt{*-{}-{}-{}-{}-{}-{}-{}-{}-{}-{}-{}-{}-{}-{}-{}-{}-{}-{}-{}-{}-{}-{}-{}-{}-{}-{}-{}-{}-{}-{}-{}-{}-{}-{}-{}-{}-{}-{}-{}-{}-{}-{}-{}-{}-{}-{}-{}-{} Panel A} & Divider line (comment).\\
15 & \texttt{eststo clear} & Clears previously stored estimates.\\
16 & \texttt{eststo A1: reghdfe ShannonBD PI \$control,\newline \hspace*{1em}absorb(county year)       cluster(county)} & Regresses Shannon bird diversity on PSI (solar policy stringency); controlling for the 9 controls; with county and year-month fixed effects; standard errors clustered by county. Result: -0.0125 (0.0037), p = 0.001.\\
17 & \texttt{eststo A2: reghdfe ShannonBD PI \$control if in\_PI\_file==1,\newline \hspace*{1em}absorb(county year)       cluster(county)} & Regresses Shannon bird diversity on PSI (solar policy stringency); controlling for the 9 controls; with county and year-month fixed effects; standard errors clustered by county; only counties that really appear in the policy file. Result: -0.0027 (0.0038), p = 0.49, N = 26,007.\\
18 & \texttt{eststo A3: reghdfe ShannonBD PI \$control if urban==0,\newline \hspace*{1em}absorb(county year)       cluster(county)} & Regresses Shannon bird diversity on PSI (solar policy stringency); controlling for the 9 controls; with county and year-month fixed effects; standard errors clustered by county; only rural counties. Result: -0.0056 (0.0042), p = 0.19.\\
19 & \texttt{eststo A4: reghdfe ShannonBD PI \$control if urban==1,\newline \hspace*{1em}absorb(county year)       cluster(county)} & Regresses Shannon bird diversity on PSI (solar policy stringency); controlling for the 9 controls; with county and year-month fixed effects; standard errors clustered by county; only urban districts. Result: +0.0056 (0.0079), p = 0.48.\\
20 & \texttt{eststo A5: reghdfe ShannonBD PI \$control,\newline \hspace*{1em}absorb(county year\#urban) cluster(county)} & Regresses Shannon bird diversity on PSI (solar policy stringency); controlling for the 9 controls; with county fixed effects plus separate year-month effects for urban and rural units; standard errors clustered by county. KEY RESULT: -0.0027 (0.0037), p = 0.47.\\
21 & \texttt{eststo A6: reghdfe ShannonBD PI \$control,\newline \hspace*{1em}absorb(county year\#prov)  cluster(county)} & Regresses Shannon bird diversity on PSI (solar policy stringency); controlling for the 9 controls; with county fixed effects plus separate year-month effects for each province; standard errors clustered by county. Result: -0.0146 (0.0047), p = 0.002.\\
22 & \texttt{esttab A1 A2 A3 A4 A5 A6, keep(PI) b(4) se(4) star(* 0.10\newline \hspace*{1em}** 0.05 *** 0.01) ///} & Prints a formatted results table of the stored estimates on screen.\\
23 & \texttt{\hspace*{2.0em}stats(N N\_clust, fmt(0 0) labels(\textquotedbl{}Obs\textquotedbl{} \textquotedbl{}Counties\textquotedbl{})) ///} & Continuation of the previous command: summary rows at the bottom (e.g. FE markers, observations, R-squared).\\
24 & \texttt{\hspace*{2.0em}mtitles(\textquotedbl{}Baseline\textquotedbl{} \textquotedbl{}PSI-{}data counties\textquotedbl{} \textquotedbl{}Rural only\textquotedbl{} \textquotedbl{}Urban\newline \hspace*{1em}only\textquotedbl{} \textquotedbl{}Urban x YM FE\textquotedbl{} \textquotedbl{}Province x YM FE\textquotedbl{}) ///} & Continuation of the previous command: column titles.\\
25 & \texttt{\hspace*{2.0em}title(\textquotedbl{}Panel A. Who drives the TWFE estimate?\textquotedbl{})} & Continuation of the previous command: table title.\\
26 & \texttt{esttab A1 A2 A3 A4 A5 A6 using\newline \hspace*{1em}\textquotedbl{}\$T1/tables/1\_5\_twfe\_diagnostics.rtf\textquotedbl{}, replace keep(PI)\newline \hspace*{1em}b(4) se(4) ///} & Exports the stored results (A1 A2 A3 A4 A5 A6) into the Word table "\$T1/tables/1\_5\_twfe\_diagnostics.rtf".\\
27 & \texttt{\hspace*{2.0em}star(* 0.10 ** 0.05 *** 0.01) stats(N N\_clust, fmt(0 0)\newline \hspace*{1em}labels(\textquotedbl{}Obs\textquotedbl{} \textquotedbl{}Counties\textquotedbl{})) ///} & Continuation of the previous command: summary rows at the bottom (e.g. FE markers, observations, R-squared); significance stars * 10\%, ** 5\%, *** 1\%.\\
28 & \texttt{\hspace*{2.0em}mtitles(\textquotedbl{}Baseline\textquotedbl{} \textquotedbl{}PSI-{}data counties\textquotedbl{} \textquotedbl{}Rural only\textquotedbl{} \textquotedbl{}Urban\newline \hspace*{1em}only\textquotedbl{} \textquotedbl{}Urban x YM FE\textquotedbl{} \textquotedbl{}Province x YM FE\textquotedbl{}) ///} & Continuation of the previous command: column titles.\\
29 & \texttt{\hspace*{2.0em}title(\textquotedbl{}Panel A. Who drives the TWFE estimate?\textquotedbl{})} & Continuation of the previous command: table title.\\
31 & \texttt{*-{}-{}-{}-{}-{}-{}-{}-{}-{}-{}-{}-{}-{}-{}-{}-{}-{}-{}-{}-{}-{}-{}-{}-{}-{}-{}-{}-{}-{}-{}-{}-{}-{}-{}-{}-{}-{}-{}-{}-{}-{}-{}-{}-{}-{}-{}-{}-{} Panel B} & Divider line (comment).\\
32 & \texttt{eststo clear} & Clears previously stored estimates.\\
33 & \texttt{eststo B1: reghdfe ShannonBD PI\_stock \$control if\newline \hspace*{1em}in\_PI\_file==1, absorb(county year) cluster(county)} & Regresses Shannon bird diversity on cumulative PSI stock; controlling for the 9 controls; with county and year-month fixed effects; standard errors clustered by county; only counties that really appear in the policy file.\\
34 & \texttt{eststo B2: reghdfe ShannonBD any\_PSI  \$control,\newline \hspace*{1em}absorb(county year) cluster(county)} & Regresses Shannon bird diversity on dummy PSI>0; controlling for the 9 controls; with county and year-month fixed effects; standard errors clustered by county.\\
35 & \texttt{eststo B3: reghdfe ShannonBD any\_PSI  \$control if\newline \hspace*{1em}in\_PI\_file==1, absorb(county year) cluster(county)} & Regresses Shannon bird diversity on dummy PSI>0; controlling for the 9 controls; with county and year-month fixed effects; standard errors clustered by county; only counties that really appear in the policy file.\\
36 & \texttt{esttab B1 B2 B3, keep(PI\_stock any\_PSI) b(4) se(4) star(*\newline \hspace*{1em}0.10 ** 0.05 *** 0.01) ///} & Prints a formatted results table of the stored estimates on screen.\\
37 & \texttt{\hspace*{2.0em}stats(N N\_clust, fmt(0 0)) mtitles(\textquotedbl{}PSI stock (PSI-{}data)\textquotedbl{}\newline \hspace*{1em}\textquotedbl{}Any PSI>0 (all)\textquotedbl{} \textquotedbl{}Any PSI>0 (PSI-{}data)\textquotedbl{}) ///} & Continuation of the previous command: column titles; summary rows at the bottom (e.g. FE markers, observations, R-squared).\\
38 & \texttt{\hspace*{2.0em}title(\textquotedbl{}Panel B. Alternative treatment definitions\textquotedbl{})} & Continuation of the previous command: table title.\\
39 & \texttt{esttab B1 B2 B3 using\newline \hspace*{1em}\textquotedbl{}\$T1/tables/1\_5\_twfe\_diagnostics.rtf\textquotedbl{}, append\newline \hspace*{1em}keep(PI\_stock any\_PSI) b(4) se(4) ///} & Exports the stored results (B1 B2 B3) into the Word table "\$T1/tables/1\_5\_twfe\_diagnostics.rtf".\\
40 & \texttt{\hspace*{2.0em}star(* 0.10 ** 0.05 *** 0.01) stats(N N\_clust, fmt(0 0))\newline \hspace*{1em}///} & Continuation of the previous command: summary rows at the bottom (e.g. FE markers, observations, R-squared); significance stars * 10\%, ** 5\%, *** 1\%.\\
41 & \texttt{\hspace*{2.0em}mtitles(\textquotedbl{}PSI stock (PSI-{}data)\textquotedbl{} \textquotedbl{}Any PSI>0 (all)\textquotedbl{} \textquotedbl{}Any\newline \hspace*{1em}PSI>0 (PSI-{}data)\textquotedbl{}) ///} & Continuation of the previous command: column titles.\\
42 & \texttt{\hspace*{2.0em}title(\textquotedbl{}Panel B. Alternative treatment definitions\textquotedbl{})} & Continuation of the previous command: table title.\\
44 & \texttt{*-{}-{}-{}-{}-{}-{}-{}-{}-{}-{}-{}-{}-{}-{}-{}-{}-{}-{}-{}-{}-{}-{}-{}-{}-{}-{}-{}-{}-{}-{}-{}-{}-{}-{}-{}-{}-{}-{}-{}-{}-{}-{}-{}-{}-{}-{}-{}-{} Panel C} & Divider line (comment).\\
45 & \texttt{eststo clear} & Clears previously stored estimates.\\
46 & \texttt{eststo C1: reghdfe ShannonBD PI \$control\_noD,\newline \hspace*{1em}absorb(county year) cluster(county)} & Regresses Shannon bird diversity on PSI (solar policy stringency); controlling for the 8 controls without Duration; with county and year-month fixed effects; standard errors clustered by county.\\
47 & \texttt{eststo C2: reghdfe Duration  PI \$control\_noD,\newline \hspace*{1em}absorb(county year) cluster(county)} & Regresses birdwatching effort (Duration) on PSI (solar policy stringency); controlling for the 8 controls without Duration; with county and year-month fixed effects; standard errors clustered by county.\\
48 & \texttt{esttab C1 C2, keep(PI) b(4) se(4) star(* 0.10 ** 0.05 ***\newline \hspace*{1em}0.01) stats(N N\_clust, fmt(0 0)) ///} & Prints a formatted results table of the stored estimates on screen.\\
49 & \texttt{\hspace*{2.0em}mtitles(\textquotedbl{}Shannon, no Duration\textquotedbl{} \textquotedbl{}Outcome = Duration\textquotedbl{})\newline \hspace*{1em}title(\textquotedbl{}Panel C. Observer effort\textquotedbl{})} & Continuation of the previous command: column titles.\\
50 & \texttt{esttab C1 C2 using \textquotedbl{}\$T1/tables/1\_5\_twfe\_diagnostics.rtf\textquotedbl{},\newline \hspace*{1em}append keep(PI) b(4) se(4) ///} & Exports the stored results (C1 C2) into the Word table "\$T1/tables/1\_5\_twfe\_diagnostics.rtf".\\
51 & \texttt{\hspace*{2.0em}star(* 0.10 ** 0.05 *** 0.01) stats(N N\_clust, fmt(0 0))\newline \hspace*{1em}///} & Continuation of the previous command: summary rows at the bottom (e.g. FE markers, observations, R-squared); significance stars * 10\%, ** 5\%, *** 1\%.\\
52 & \texttt{\hspace*{2.0em}mtitles(\textquotedbl{}Shannon, no Duration\textquotedbl{} \textquotedbl{}Outcome = Duration\textquotedbl{})\newline \hspace*{1em}title(\textquotedbl{}Panel C. Observer effort\textquotedbl{})} & Continuation of the previous command: column titles.\\
54 & \texttt{*-{}-{}-{}-{}-{}-{}-{}-{}-{}-{}-{}-{}-{}-{}-{}-{}-{}-{}-{}-{}-{}-{}-{}-{}-{}-{}-{}-{}-{}-{}-{}-{}-{}-{}-{}-{}-{}-{}-{}-{}-{}-{}-{}-{}-{}-{}-{}-{} Panel D} & Divider line (comment).\\
55 & \texttt{* The authors run psmatch2 and then regress on the matched\newline \hspace*{1em}sample WITHOUT} & Comment (not executed): explains the code that follows.\\
56 & \texttt{* the weights (if \_weight!=.). Here the weights are\newline \hspace*{1em}applied.} & Comment (not executed): explains the code that follows.\\
57 & \texttt{gen byte treated = (PI > 0)} & Treatment dummy for matching: 1 if PSI > 0.\\
58 & \texttt{eststo clear} & Clears previously stored estimates.\\
59 & \texttt{psmatch2 treated \$control i.year, outcome(ShannonBD)\newline \hspace*{1em}neighbor(1) caliper(0.05) common} & Propensity-score matching: estimates the probability of PSI > 0 from the controls and year-month dummies, then matches each treated row to its 1 nearest untreated neighbour(s) within a caliper of 0.05; common = keep only common support. Creates \_weight (matching weights).\\
60 & \texttt{eststo D1: reghdfe ShannonBD PI \$control [aw=\_weight] if\newline \hspace*{1em}\_weight!=., absorb(county year) cluster(county)} & Regresses Shannon bird diversity on PSI (solar policy stringency); controlling for the 9 controls; with county and year-month fixed effects; standard errors clustered by county; only observations kept by the matching; weighting each observation by \_weight (the matching weight).\\
61 & \texttt{drop \_*} & Deletes psmatch2's variables before the next matching.\\
62 & \texttt{psmatch2 treated \$control i.year, outcome(ShannonBD)\newline \hspace*{1em}neighbor(2) caliper(0.05) common} & Propensity-score matching: estimates the probability of PSI > 0 from the controls and year-month dummies, then matches each treated row to its 2 nearest untreated neighbour(s) within a caliper of 0.05; common = keep only common support. Creates \_weight (matching weights).\\
63 & \texttt{eststo D2: reghdfe ShannonBD PI \$control [aw=\_weight] if\newline \hspace*{1em}\_weight!=., absorb(county year) cluster(county)} & Regresses Shannon bird diversity on PSI (solar policy stringency); controlling for the 9 controls; with county and year-month fixed effects; standard errors clustered by county; only observations kept by the matching; weighting each observation by \_weight (the matching weight).\\
64 & \texttt{drop \_*} & Deletes psmatch2's variables before the next matching.\\
65 & \texttt{psmatch2 treated \$control i.year, outcome(ShannonBD)\newline \hspace*{1em}kernel common} & Kernel propensity-score matching: estimates each row's probability of PSI > 0 from the controls and year-month dummies, then weights all untreated rows by closeness; common = keep only common support. Creates \_weight.\\
66 & \texttt{di \_n \textquotedbl{}Kernel weights of matched controls (far from\newline \hspace*{1em}uniform):\textquotedbl{}} & Displays a message or result on screen.\\
67 & \texttt{summarize \_weight if treated==0 \& \_weight!=., detail} & Computes summary statistics (mean, SD, min, max, percentiles) and stores them in r().\\
68 & \texttt{eststo D3: reghdfe ShannonBD PI \$control [aw=\_weight] if\newline \hspace*{1em}\_weight!=., absorb(county year) cluster(county)} & Regresses Shannon bird diversity on PSI (solar policy stringency); controlling for the 9 controls; with county and year-month fixed effects; standard errors clustered by county; only observations kept by the matching; weighting each observation by \_weight (the matching weight).\\
69 & \texttt{drop \_*} & Deletes psmatch2's variables before the next matching.\\
70 & \texttt{esttab D1 D2 D3, keep(PI) b(4) se(4) star(* 0.10 ** 0.05\newline \hspace*{1em}*** 0.01) stats(N N\_clust, fmt(0 0)) ///} & Prints a formatted results table of the stored estimates on screen.\\
71 & \texttt{\hspace*{2.0em}mtitles(\textquotedbl{}1:1 NN, weighted\textquotedbl{} \textquotedbl{}1:2 NN, weighted\textquotedbl{} \textquotedbl{}Kernel,\newline \hspace*{1em}weighted\textquotedbl{}) ///} & Continuation of the previous command: column titles.\\
72 & \texttt{\hspace*{2.0em}title(\textquotedbl{}Panel D. Table S11 with matching weights applied\textquotedbl{})} & Continuation of the previous command: table title.\\
73 & \texttt{esttab D1 D2 D3 using\newline \hspace*{1em}\textquotedbl{}\$T1/tables/1\_5\_twfe\_diagnostics.rtf\textquotedbl{}, append keep(PI)\newline \hspace*{1em}b(4) se(4) ///} & Exports the stored results (D1 D2 D3) into the Word table "\$T1/tables/1\_5\_twfe\_diagnostics.rtf".\\
74 & \texttt{\hspace*{2.0em}star(* 0.10 ** 0.05 *** 0.01) stats(N N\_clust, fmt(0 0))\newline \hspace*{1em}///} & Continuation of the previous command: summary rows at the bottom (e.g. FE markers, observations, R-squared); significance stars * 10\%, ** 5\%, *** 1\%.\\
75 & \texttt{\hspace*{2.0em}mtitles(\textquotedbl{}1:1 NN, weighted\textquotedbl{} \textquotedbl{}1:2 NN, weighted\textquotedbl{} \textquotedbl{}Kernel,\newline \hspace*{1em}weighted\textquotedbl{}) ///} & Continuation of the previous command: column titles.\\
76 & \texttt{\hspace*{2.0em}title(\textquotedbl{}Panel D. Table S11 with matching weights applied\textquotedbl{})} & Continuation of the previous command: table title.\\}
\subsection{1\_6\_iv\_exclusion\_tests.do}
Tests the IV exclusion restriction: reduced form, zero-first-stage test, whether the instrument predicts other channels, and the IV under city clustering and other specifications.

\codetable{1 & \texttt{*=========================================================\newline \hspace*{1em}=================} & Divider line (comment).\\
2 & \texttt{* 1\_6\_iv\_exclusion\_tests.do} & Comment (not executed): explains the code that follows.\\
3 & \texttt{* PURPOSE : Empirical checks on the IV exclusion\newline \hspace*{1em}restriction.} & Comment (not executed): explains the code that follows.\\
4 & \texttt{*   Instrument Z = Sun\_ccpu = ln(mean sunshine 1984-{}2013)\newline \hspace*{1em}/ city CPU\_t} & Comment (not executed): explains the code that follows.\\
5 & \texttt{*   Panel A  Reduced form and ZERO-{}FIRST-{}STAGE test: in\newline \hspace*{1em}counties where PSI} & Comment (not executed): explains the code that follows.\\
6 & \texttt{*            cannot change, Z must have no effect on\newline \hspace*{1em}Shannon if exclusion holds} & Comment (not executed): explains the code that follows.\\
7 & \texttt{*            (van Kippersluis \& Rietveld 2018).} & Comment (not executed): explains the code that follows.\\
8 & \texttt{*   Panel B  Does Z predict other channels (effort, other\newline \hspace*{1em}policies, land)?} & Comment (not executed): explains the code that follows.\\
9 & \texttt{*   Panel C  IV under city clustering (level of the shift)\newline \hspace*{1em}and other specs.} & Comment (not executed): explains the code that follows.\\
10 & \texttt{* OUTPUT  : tables/1\_6\_iv\_exclusion\_tests.rtf} & Comment (not executed): explains the code that follows.\\
11 & \texttt{* WHY     : Task 1, item 6 — logic and possible violations\newline \hspace*{1em}of exclusion.} & Comment (not executed): explains the code that follows.\\
12 & \texttt{*=========================================================\newline \hspace*{1em}=================} & Divider line (comment).\\
13 & \texttt{use \textquotedbl{}\$SAMPLE\textquotedbl{}, clear} & Loads the analysis dataset built in step 1\_2.\\
14 & \texttt{bysort county: egen sdPI = sd(PI)} & Standard deviation of PSI within each county (0 = PSI never changes).\\
15 & \texttt{gen byte PI\_constant = (sdPI==0 | missing(sdPI))} & 1 if PSI never changes within the county.\\
17 & \texttt{*-{}-{}-{}-{}-{}-{}-{}-{}-{}-{}-{}-{}-{}-{}-{}-{}-{}-{}-{}-{}-{}-{}-{}-{}-{}-{}-{}-{}-{}-{}-{}-{}-{}-{}-{}-{}-{}-{}-{}-{}-{}-{}-{}-{}-{}-{}-{}-{} Panel A} & Divider line (comment).\\
18 & \texttt{eststo clear} & Clears previously stored estimates.\\
19 & \texttt{eststo A1: reghdfe ShannonBD Sun\_ccpu \$control,\newline \hspace*{1em}absorb(county year) cluster(county)} & Regresses Shannon bird diversity on the instrument ln(sunshine)/CPU; controlling for the 9 controls; with county and year-month fixed effects; standard errors clustered by county.\\
20 & \texttt{* first stage is exactly zero where PSI is imputed 0 (PSI\newline \hspace*{1em}does not vary):} & Comment (not executed): explains the code that follows.\\
21 & \texttt{reghdfe PI Sun\_ccpu \$control if in\_PI\_file==0,\newline \hspace*{1em}absorb(county year) cluster(county)} & Regresses PSI (solar policy stringency) on the instrument ln(sunshine)/CPU; controlling for the 9 controls; with county and year-month fixed effects; standard errors clustered by county; only counties missing from the policy file (PSI forced to 0). Result: reduced form -0.0032 (0.0007).\\
22 & \texttt{eststo A2: reghdfe ShannonBD Sun\_ccpu \$control if\newline \hspace*{1em}in\_PI\_file==0, absorb(county year) cluster(county)} & Regresses Shannon bird diversity on the instrument ln(sunshine)/CPU; controlling for the 9 controls; with county and year-month fixed effects; standard errors clustered by county; only counties missing from the policy file (PSI forced to 0).\\
23 & \texttt{eststo A3: reghdfe ShannonBD Sun\_ccpu \$control if\newline \hspace*{1em}PI\_constant==1,absorb(county year) cluster(county)} & Regresses Shannon bird diversity on the instrument ln(sunshine)/CPU; controlling for the 9 controls; with county and year-month fixed effects; standard errors clustered by county; only counties whose PSI never changes.\\
24 & \texttt{eststo A4: reghdfe ShannonBD Sun\_ccpu \$control if\newline \hspace*{1em}PI\_constant==0,absorb(county year) cluster(county)} & Regresses Shannon bird diversity on the instrument ln(sunshine)/CPU; controlling for the 9 controls; with county and year-month fixed effects; standard errors clustered by county; only counties whose PSI changes. KEY RESULT (zero-first-stage test): -0.0021 (0.0011), p = 0.048, N = 17,497.\\
25 & \texttt{esttab A1 A2 A3 A4, keep(Sun\_ccpu) b(4) se(4) star(* 0.10\newline \hspace*{1em}** 0.05 *** 0.01) ///} & Prints a formatted results table of the stored estimates on screen.\\
26 & \texttt{\hspace*{2.0em}stats(N N\_clust, fmt(0 0)) ///} & Continuation of the previous command: summary rows at the bottom (e.g. FE markers, observations, R-squared).\\
27 & \texttt{\hspace*{2.0em}mtitles(\textquotedbl{}Reduced form\textquotedbl{} \textquotedbl{}PSI imputed 0\textquotedbl{} \textquotedbl{}PSI constant\textquotedbl{} \textquotedbl{}PSI\newline \hspace*{1em}varies\textquotedbl{}) ///} & Continuation of the previous command: column titles.\\
28 & \texttt{\hspace*{2.0em}title(\textquotedbl{}Panel A. Reduced form and zero-{}first-{}stage tests\newline \hspace*{1em}(outcome: ShannonBD)\textquotedbl{})} & Continuation of the previous command: table title.\\
29 & \texttt{esttab A1 A2 A3 A4 using\newline \hspace*{1em}\textquotedbl{}\$T1/tables/1\_6\_iv\_exclusion\_tests.rtf\textquotedbl{}, replace\newline \hspace*{1em}keep(Sun\_ccpu) b(4) se(4) ///} & Exports the stored results (A1 A2 A3 A4) into the Word table "\$T1/tables/1\_6\_iv\_exclusion\_tests.rtf".\\
30 & \texttt{\hspace*{2.0em}star(* 0.10 ** 0.05 *** 0.01) stats(N N\_clust, fmt(0 0))\newline \hspace*{1em}///} & Continuation of the previous command: summary rows at the bottom (e.g. FE markers, observations, R-squared); significance stars * 10\%, ** 5\%, *** 1\%.\\
31 & \texttt{\hspace*{2.0em}mtitles(\textquotedbl{}Reduced form\textquotedbl{} \textquotedbl{}PSI imputed 0\textquotedbl{} \textquotedbl{}PSI constant\textquotedbl{} \textquotedbl{}PSI\newline \hspace*{1em}varies\textquotedbl{}) ///} & Continuation of the previous command: column titles.\\
32 & \texttt{\hspace*{2.0em}title(\textquotedbl{}Panel A. Reduced form and zero-{}first-{}stage tests\newline \hspace*{1em}(outcome: ShannonBD)\textquotedbl{})} & Continuation of the previous command: table title.\\
34 & \texttt{*-{}-{}-{}-{}-{}-{}-{}-{}-{}-{}-{}-{}-{}-{}-{}-{}-{}-{}-{}-{}-{}-{}-{}-{}-{}-{}-{}-{}-{}-{}-{}-{}-{}-{}-{}-{}-{}-{}-{}-{}-{}-{}-{}-{}-{}-{}-{}-{} Panel B} & Divider line (comment).\\
35 & \texttt{eststo clear} & Clears previously stored estimates.\\
36 & \texttt{eststo B1: reghdfe Duration Sun\_ccpu \$control\_noD,\newline \hspace*{1em}absorb(county year) cluster(county)} & Regresses birdwatching effort (Duration) on the instrument ln(sunshine)/CPU; controlling for the 8 controls without Duration; with county and year-month fixed effects; standard errors clustered by county.\\
37 & \texttt{foreach y in CEI Coal Windexp Greenfin Ndvi LAI Nightlight\newline \hspace*{1em}\{} & Loops over other possible channels from the instrument to birds.\\
38 & \texttt{\hspace*{2.0em}eststo B\_\textasciigrave{}y\textquotesingle{}: reghdfe \textasciigrave{}y\textquotesingle{} Sun\_ccpu \$control, absorb(county\newline \hspace*{1em}year) cluster(county)} & Regresses each other channel (CEI, Coal, Windexp, Greenfin, NDVI, LAI, Nightlight) on the instrument: does the instrument predict them?\\
39 & \texttt{\}} & Closes the loop / block.\\
40 & \texttt{esttab B1 B\_*, keep(Sun\_ccpu) b(4) se(4) star(* 0.10 **\newline \hspace*{1em}0.05 *** 0.01) ///} & Prints a formatted results table of the stored estimates on screen.\\
41 & \texttt{\hspace*{2.0em}stats(N N\_clust, fmt(0 0)) ///} & Continuation of the previous command: summary rows at the bottom (e.g. FE markers, observations, R-squared).\\
42 & \texttt{\hspace*{2.0em}mtitles(\textquotedbl{}Duration\textquotedbl{} \textquotedbl{}CEI\textquotedbl{} \textquotedbl{}Coal\textquotedbl{} \textquotedbl{}Windexp\textquotedbl{} \textquotedbl{}Greenfin\textquotedbl{}\newline \hspace*{1em}\textquotedbl{}NDVI\textquotedbl{} \textquotedbl{}LAI\textquotedbl{} \textquotedbl{}Nightlight\textquotedbl{}) ///} & Continuation of the previous command: column titles.\\
43 & \texttt{\hspace*{2.0em}title(\textquotedbl{}Panel B. Does the instrument predict other\newline \hspace*{1em}channels?\textquotedbl{})} & Continuation of the previous command: table title.\\
44 & \texttt{esttab B1 B\_* using\newline \hspace*{1em}\textquotedbl{}\$T1/tables/1\_6\_iv\_exclusion\_tests.rtf\textquotedbl{}, append\newline \hspace*{1em}keep(Sun\_ccpu) b(4) se(4) ///} & Exports the stored results (B1 B\_*) into the Word table "\$T1/tables/1\_6\_iv\_exclusion\_tests.rtf".\\
45 & \texttt{\hspace*{2.0em}star(* 0.10 ** 0.05 *** 0.01) stats(N N\_clust, fmt(0 0))\newline \hspace*{1em}///} & Continuation of the previous command: summary rows at the bottom (e.g. FE markers, observations, R-squared); significance stars * 10\%, ** 5\%, *** 1\%.\\
46 & \texttt{\hspace*{2.0em}mtitles(\textquotedbl{}Duration\textquotedbl{} \textquotedbl{}CEI\textquotedbl{} \textquotedbl{}Coal\textquotedbl{} \textquotedbl{}Windexp\textquotedbl{} \textquotedbl{}Greenfin\textquotedbl{}\newline \hspace*{1em}\textquotedbl{}NDVI\textquotedbl{} \textquotedbl{}LAI\textquotedbl{} \textquotedbl{}Nightlight\textquotedbl{}) ///} & Continuation of the previous command: column titles.\\
47 & \texttt{\hspace*{2.0em}title(\textquotedbl{}Panel B. Does the instrument predict other\newline \hspace*{1em}channels?\textquotedbl{})} & Continuation of the previous command: table title.\\
49 & \texttt{*-{}-{}-{}-{}-{}-{}-{}-{}-{}-{}-{}-{}-{}-{}-{}-{}-{}-{}-{}-{}-{}-{}-{}-{}-{}-{}-{}-{}-{}-{}-{}-{}-{}-{}-{}-{}-{}-{}-{}-{}-{}-{}-{}-{}-{}-{}-{}-{} Panel C} & Divider line (comment).\\
50 & \texttt{eststo clear} & Clears previously stored estimates.\\
51 & \texttt{eststo C1: ivreghdfe ShannonBD (PI = Sun\_ccpu) \$control,\newline \hspace*{1em}absorb(county year) cluster(county) keepsingleton} & Two-stage least squares (IV): Shannon bird diversity on PSI, with PSI instrumented by Sun\_ccpu; controls: the 9 controls; with county and year-month fixed effects; standard errors clustered by county; keeping single-observation counties.\\
52 & \texttt{eststo C2: ivreghdfe ShannonBD (PI = Sun\_ccpu) \$control,\newline \hspace*{1em}absorb(county year) cluster(city) keepsingleton} & Two-stage least squares (IV): Shannon bird diversity on PSI, with PSI instrumented by Sun\_ccpu; controls: the 9 controls; with county and year-month fixed effects; standard errors clustered by city; keeping single-observation counties.\\
53 & \texttt{eststo C3: ivreghdfe ShannonBD (PI = Sun\_ccpu) \$control,\newline \hspace*{1em}absorb(county year\#urban) cluster(county) keepsingleton} & Two-stage least squares (IV): Shannon bird diversity on PSI, with PSI instrumented by Sun\_ccpu; controls: the 9 controls; with county fixed effects plus separate year-month effects for urban and rural units; standard errors clustered by county; keeping single-observation counties.\\
54 & \texttt{eststo C4: ivreghdfe ShannonBD (PI = Sun\_ccpu) \$control if\newline \hspace*{1em}in\_PI\_file==1, absorb(county year) cluster(county)\newline \hspace*{1em}keepsingleton} & Two-stage least squares (IV): Shannon bird diversity on PSI, with PSI instrumented by Sun\_ccpu; controls: the 9 controls; with county and year-month fixed effects; standard errors clustered by county; keeping single-observation counties; only counties that really appear in the policy file.\\
55 & \texttt{eststo C5: ivreghdfe ShannonBD (PI = Sun\_ccpu)\newline \hspace*{1em}\$control\_noD, absorb(county year) cluster(county)\newline \hspace*{1em}keepsingleton} & Two-stage least squares (IV): Shannon bird diversity on PSI, with PSI instrumented by Sun\_ccpu; controls: the 8 controls without Duration; with county and year-month fixed effects; standard errors clustered by county; keeping single-observation counties.\\
56 & \texttt{esttab C1 C2 C3 C4 C5, keep(PI) b(4) se(4) star(* 0.10 **\newline \hspace*{1em}0.05 *** 0.01) ///} & Prints a formatted results table of the stored estimates on screen.\\
57 & \texttt{\hspace*{2.0em}stats(N N\_clust widstat, fmt(0 0 2) labels(\textquotedbl{}Obs\textquotedbl{}\newline \hspace*{1em}\textquotedbl{}Clusters\textquotedbl{} \textquotedbl{}Kleibergen-{}Paap F\textquotedbl{})) ///} & Continuation of the previous command: summary rows at the bottom (e.g. FE markers, observations, R-squared).\\
58 & \texttt{\hspace*{2.0em}mtitles(\textquotedbl{}Paper\textquotedbl{} \textquotedbl{}Cluster city\textquotedbl{} \textquotedbl{}Urban x YM FE\textquotedbl{} \textquotedbl{}PSI-{}data\newline \hspace*{1em}only\textquotedbl{} \textquotedbl{}No Duration\textquotedbl{}) ///} & Continuation of the previous command: column titles.\\
59 & \texttt{\hspace*{2.0em}title(\textquotedbl{}Panel C. IV estimate: inference and specification\textquotedbl{})} & Continuation of the previous command: table title.\\
60 & \texttt{esttab C1 C2 C3 C4 C5 using\newline \hspace*{1em}\textquotedbl{}\$T1/tables/1\_6\_iv\_exclusion\_tests.rtf\textquotedbl{}, append keep(PI)\newline \hspace*{1em}b(4) se(4) ///} & Exports the stored results (C1 C2 C3 C4 C5) into the Word table "\$T1/tables/1\_6\_iv\_exclusion\_tests.rtf".\\
61 & \texttt{\hspace*{2.0em}star(* 0.10 ** 0.05 *** 0.01) stats(N N\_clust widstat,\newline \hspace*{1em}fmt(0 0 2) labels(\textquotedbl{}Obs\textquotedbl{} \textquotedbl{}Clusters\textquotedbl{} \textquotedbl{}Kleibergen-{}Paap F\textquotedbl{}))\newline \hspace*{1em}///} & Continuation of the previous command: summary rows at the bottom (e.g. FE markers, observations, R-squared); significance stars * 10\%, ** 5\%, *** 1\%.\\
62 & \texttt{\hspace*{2.0em}mtitles(\textquotedbl{}Paper\textquotedbl{} \textquotedbl{}Cluster city\textquotedbl{} \textquotedbl{}Urban x YM FE\textquotedbl{} \textquotedbl{}PSI-{}data\newline \hspace*{1em}only\textquotedbl{} \textquotedbl{}No Duration\textquotedbl{}) ///} & Continuation of the previous command: column titles.\\
63 & \texttt{\hspace*{2.0em}title(\textquotedbl{}Panel C. IV estimate: inference and specification\textquotedbl{})} & Continuation of the previous command: table title.\\}
\subsection{2\_0\_run\_task2.do}
Runs Task 2 with a log file.

\codetable{1 & \texttt{*=========================================================\newline \hspace*{1em}=================} & Divider line (comment).\\
2 & \texttt{* 2\_0\_run\_task2.do  —  runs Task 2 (needs Task 1\textquotesingle{}s\newline \hspace*{1em}analysis\_sample.dta)} & Comment (not executed): explains the code that follows.\\
3 & \texttt{*=========================================================\newline \hspace*{1em}=================} & Divider line (comment).\\
4 & \texttt{do \textquotedbl{}/Users/sahiljangra/Desktop/Johns\newline \hspace*{1em}Hopkins/RA/Solar\_Birds\_Replication/00\_config.do\textquotedbl{}} & Runs another do-file: /Users/sahiljangra/Desktop/Johns Hopkins/RA/Solar\_Birds\_Replication/00\_config.do\\
5 & \texttt{capture log close \_all} & Closes any open log (capture = no error if none is open).\\
6 & \texttt{log using \textquotedbl{}\$T2/logs/2\_1\_sensitivity\_analysis.log\textquotedbl{}, replace\newline \hspace*{1em}text} & Opens a text log file that records all output of this step.\\
7 & \texttt{do \textquotedbl{}\$T2/code/2\_1\_sensitivity\_analysis.do\textquotedbl{}} & Runs another do-file: \$T2/code/2\_1\_sensitivity\_analysis.do\\
8 & \texttt{log close} & Closes the log file.\\}
\subsection{2\_1\_sensitivity\_analysis.do}
Cinelli--Hazlett sensitivity analysis of the TWFE estimate: robustness values, benchmark bounds, the clustered-SE variant, contour plots and a sensemakr cross-check.

\codetable{1 & \texttt{*=========================================================\newline \hspace*{1em}=================} & Divider line (comment).\\
2 & \texttt{* 2\_1\_sensitivity\_analysis.do} & Comment (not executed): explains the code that follows.\\
3 & \texttt{* PURPOSE : Cinelli \& Hazlett (2020) sensitivity of the\newline \hspace*{1em}TWFE estimate to an} & Comment (not executed): explains the code that follows.\\
4 & \texttt{*           unobserved confounder Z, for two\newline \hspace*{1em}specifications:} & Comment (not executed): explains the code that follows.\\
5 & \texttt{*             S1  paper\textquotesingle{}s baseline: county + year-{}month FE} & Comment (not executed): explains the code that follows.\\
6 & \texttt{*             S2  urban x year-{}month FE (Task 1\newline \hspace*{1em}diagnostic)} & Comment (not executed): explains the code that follows.\\
7 & \texttt{* STEPS   : A  sensitivity statistics: partial R2 of PSI\newline \hspace*{1em}with Shannon,} & Comment (not executed): explains the code that follows.\\
8 & \texttt{*              robustness values RV\_q=1 and\newline \hspace*{1em}RV\_q=1,alpha=0.05  (C-{}H eqs 18-{}20)} & Comment (not executed): explains the code that follows.\\
9 & \texttt{*           B  benchmark bounds: Z \textquotedbl{}k times as strong\textquotedbl{} as\newline \hspace*{1em}each of the nine} & Comment (not executed): explains the code that follows.\\
10 & \texttt{*              controls (k = 1,2,3), adjusted estimate \& t\newline \hspace*{1em}(C-{}H eqs 12-{}13, 21-{}22)} & Comment (not executed): explains the code that follows.\\
11 & \texttt{*           C  group benchmark: Z as strong as urban x\newline \hspace*{1em}year-{}month trends} & Comment (not executed): explains the code that follows.\\
12 & \texttt{*           D  clustered-{}inference variant (significance\newline \hspace*{1em}threshold scaled by} & Comment (not executed): explains the code that follows.\\
13 & \texttt{*              lambda = clustered SE / classical SE)} & Comment (not executed): explains the code that follows.\\
14 & \texttt{*           E  contour plots of the adjusted estimate and\newline \hspace*{1em}t-{}value} & Comment (not executed): explains the code that follows.\\
15 & \texttt{*           F  cross-{}check with the official -{}sensemakr-{}\newline \hspace*{1em}on FE-{}demeaned data} & Comment (not executed): explains the code that follows.\\
16 & \texttt{* NOTE    : The \textasciitilde{}2,460 fixed effects are absorbed\newline \hspace*{1em}(Frisch-{}Waugh-{}Lovell); all} & Comment (not executed): explains the code that follows.\\
17 & \texttt{*           statistics use reghdfe\textquotesingle{}s exact residual\newline \hspace*{1em}degrees of freedom.} & Comment (not executed): explains the code that follows.\\
18 & \texttt{* OUTPUT  : tables/2\_1\_robustness\_values.csv,\newline \hspace*{1em}2\_1\_benchmark\_bounds.csv} & Comment (not executed): explains the code that follows.\\
19 & \texttt{*           figures/2\_1\_contour\_estimate\_S*.pdf,\newline \hspace*{1em}2\_1\_contour\_tvalue\_S*.pdf} & Comment (not executed): explains the code that follows.\\
20 & \texttt{*=========================================================\newline \hspace*{1em}=================} & Divider line (comment).\\
22 & \texttt{*-{}-{}-{}-{}-{}-{}-{}-{}-{}-{}-{}-{}-{}-{}-{}-{}-{}-{}-{}-{}-{}-{}-{}-{}-{}-{}-{}-{}-{}-{}-{}-{}-{}-{}-{}-{}-{}-{}-{}-{}-{}-{}-{}-{}-{}-{}-{}-{}-{}-{}-{}-{}-{}-{}-{}-{}-{}\newline \hspace*{1em}-{}-{}-{}-{}-{}-{}-{}-{}-{}-{}-{}-{}-{}-{}-{}-{}-{}} & Divider line (comment).\\
23 & \texttt{* Program 1: robustness values from a t-{}statistic and df\newline \hspace*{1em}(C-{}H eqs 18-{}20)} & Comment (not executed): explains the code that follows.\\
24 & \texttt{*-{}-{}-{}-{}-{}-{}-{}-{}-{}-{}-{}-{}-{}-{}-{}-{}-{}-{}-{}-{}-{}-{}-{}-{}-{}-{}-{}-{}-{}-{}-{}-{}-{}-{}-{}-{}-{}-{}-{}-{}-{}-{}-{}-{}-{}-{}-{}-{}-{}-{}-{}-{}-{}-{}-{}-{}-{}\newline \hspace*{1em}-{}-{}-{}-{}-{}-{}-{}-{}-{}-{}-{}-{}-{}-{}-{}-{}-{}} & Divider line (comment).\\
25 & \texttt{capture program drop ch\_rv} & Removes an old version of the program if it exists.\\
26 & \texttt{program define ch\_rv, rclass} & Defines our program ch\_rv that computes the sensitivity statistics (rclass = returns results in r()).\\
27 & \texttt{\hspace*{2.0em}args t df q alpha lambda} & Program inputs: t-statistic, residual df, q (share of the estimate to remove; 1 = all), alpha (0.05), lambda (clustered/classical SE ratio).\\
28 & \texttt{\hspace*{2.0em}if \textquotedbl{}\textasciigrave{}lambda\textquotesingle{}\textquotedbl{} == \textquotedbl{}\textquotedbl{} local lambda 1} & If no lambda is given, use 1 (classical SEs).\\
29 & \texttt{\hspace*{2.0em}tempname f fq fcrit fqa rv rva} & Temporary scalar names.\\
30 & \texttt{\hspace*{2.0em}scalar \textasciigrave{}f\textquotesingle{}     = abs(\textasciigrave{}t\textquotesingle{})/sqrt(\textasciigrave{}df\textquotesingle{})\newline \hspace*{1em}// partial Cohen f of D with Y} & Partial Cohen's f of PSI with Shannon = |t| / sqrt(df).\\
31 & \texttt{\hspace*{2.0em}scalar \textasciigrave{}fq\textquotesingle{}    = \textasciigrave{}q\textquotesingle{}*\textasciigrave{}f\textquotesingle{}} & Scales f by q.\\
32 & \texttt{\hspace*{2.0em}scalar \textasciigrave{}rv\textquotesingle{}    = 0.5*(sqrt(\textasciigrave{}fq\textquotesingle{}\^{}4 + 4*\textasciigrave{}fq\textquotesingle{}\^{}2) -{} \textasciigrave{}fq\textquotesingle{}\^{}2)\newline \hspace*{1em}// eq 18} & Robustness value RV (Cinelli-Hazlett eq. 18): the equal partial R2 a confounder needs with PSI and Shannon to remove the estimate.\\
33 & \texttt{\hspace*{2.0em}* lambda = clustered SE / classical SE: inflates only the\newline \hspace*{1em}significance} & Comment (not executed): explains the code that follows.\\
34 & \texttt{\hspace*{2.0em}* threshold; the confounding-{}strength scale (partial R2)\newline \hspace*{1em}is unchanged} & Comment (not executed): explains the code that follows.\\
35 & \texttt{\hspace*{2.0em}scalar \textasciigrave{}fcrit\textquotesingle{} = \textasciigrave{}lambda\textquotesingle{}*abs(invttail(\textasciigrave{}df\textquotesingle{}-{}1,\newline \hspace*{1em}\textasciigrave{}alpha\textquotesingle{}/2))/sqrt(\textasciigrave{}df\textquotesingle{}-{}1)} & Critical f for significance at alpha; multiplied by lambda so clustered inference raises the threshold.\\
36 & \texttt{\hspace*{2.0em}scalar \textasciigrave{}fqa\textquotesingle{}   = \textasciigrave{}fq\textquotesingle{} -{} \textasciigrave{}fcrit\textquotesingle{}\newline \hspace*{1em}// eq 19} & f left above the significance threshold (eq. 19).\\
37 & \texttt{\hspace*{2.0em}if \textasciigrave{}fqa\textquotesingle{} <= 0 scalar \textasciigrave{}rva\textquotesingle{} = 0} & If already insignificant, RV-alpha = 0.\\
38 & \texttt{\hspace*{2.0em}else \{} & Otherwise ...\\
39 & \texttt{\hspace*{4.0em}scalar \textasciigrave{}rva\textquotesingle{} = 0.5*(sqrt(\textasciigrave{}fqa\textquotesingle{}\^{}4 + 4*\textasciigrave{}fqa\textquotesingle{}\^{}2) -{} \textasciigrave{}fqa\textquotesingle{}\^{}2)\newline \hspace*{1em}// eq 20} & RV-alpha (eq. 20): strength that makes the estimate insignificant.\\
40 & \texttt{\hspace*{4.0em}if \textasciigrave{}rva\textquotesingle{} <= 1 -{} 1/\textasciigrave{}fq\textquotesingle{}\^{}2 \& \textasciigrave{}fq\textquotesingle{} > 1 scalar \textasciigrave{}rva\textquotesingle{} = (\textasciigrave{}fq\textquotesingle{}\^{}2\newline \hspace*{1em}-{} \textasciigrave{}fcrit\textquotesingle{}\^{}2)/(1 + \textasciigrave{}fq\textquotesingle{}\^{}2)} & Special case of the formula when the effect is very strong (not triggered here).\\
41 & \texttt{\hspace*{2.0em}\}} & Closes the loop / block.\\
42 & \texttt{\hspace*{2.0em}return scalar r2yd = \textasciigrave{}t\textquotesingle{}\^{}2/(\textasciigrave{}t\textquotesingle{}\^{}2 + \textasciigrave{}df\textquotesingle{})\newline \hspace*{1em}// partial R2 of D with Y} & Returns the partial R2 of PSI with Shannon = t\^{}2/(t\^{}2+df).\\
43 & \texttt{\hspace*{2.0em}return scalar rv   = \textasciigrave{}rv\textquotesingle{}} & Returns RV.\\
44 & \texttt{\hspace*{2.0em}return scalar rva  = \textasciigrave{}rva\textquotesingle{}} & Returns RV-alpha.\\
45 & \texttt{end} & Ends the program definition.\\
47 & \texttt{*-{}-{}-{}-{}-{}-{}-{}-{}-{}-{}-{}-{}-{}-{}-{}-{}-{}-{}-{}-{}-{}-{}-{}-{}-{}-{}-{}-{}-{}-{}-{}-{}-{}-{}-{}-{}-{}-{}-{}-{}-{}-{}-{}-{}-{}-{}-{}-{}-{}-{}-{}-{}-{}-{}-{}-{}-{}\newline \hspace*{1em}-{}-{}-{}-{}-{}-{}-{}-{}-{}-{}-{}-{}-{}-{}-{}-{}-{}} & Divider line (comment).\\
48 & \texttt{* Program 2: bound + adjusted estimate for Z \textquotedbl{}k times\textquotedbl{} a\newline \hspace*{1em}benchmark} & Comment (not executed): explains the code that follows.\\
49 & \texttt{*            (C-{}H eqs 21-{}22; supplement section B; =\newline \hspace*{1em}sensemakr\textquotesingle{}s bounds)} & Comment (not executed): explains the code that follows.\\
50 & \texttt{*-{}-{}-{}-{}-{}-{}-{}-{}-{}-{}-{}-{}-{}-{}-{}-{}-{}-{}-{}-{}-{}-{}-{}-{}-{}-{}-{}-{}-{}-{}-{}-{}-{}-{}-{}-{}-{}-{}-{}-{}-{}-{}-{}-{}-{}-{}-{}-{}-{}-{}-{}-{}-{}-{}-{}-{}-{}\newline \hspace*{1em}-{}-{}-{}-{}-{}-{}-{}-{}-{}-{}-{}-{}-{}-{}-{}-{}-{}} & Divider line (comment).\\
51 & \texttt{capture program drop ch\_bound} & Removes an old version of the program if it exists.\\
52 & \texttt{program define ch\_bound, rclass} & Defines our program ch\_bound that computes benchmark bounds.\\
53 & \texttt{\hspace*{2.0em}args r2dxj r2yxj kd ky est se df} & Inputs: benchmark partial R2 with PSI and with Shannon, multipliers kd and ky, estimate, SE, df.\\
54 & \texttt{\hspace*{2.0em}tempname r2dz r2zxj r2yz bias adj adjse} & Temporary scalar names.\\
55 & \texttt{\hspace*{2.0em}scalar \textasciigrave{}r2dz\textquotesingle{}  = \textasciigrave{}kd\textquotesingle{}*\textasciigrave{}r2dxj\textquotesingle{}/(1 -{} \textasciigrave{}r2dxj\textquotesingle{})} & Partial R2 of the confounder with PSI if it is kd times as strong as the benchmark.\\
56 & \texttt{\hspace*{2.0em}scalar \textasciigrave{}r2zxj\textquotesingle{} = \textasciigrave{}kd\textquotesingle{}*\textasciigrave{}r2dxj\textquotesingle{}\^{}2/((1 -{} \textasciigrave{}kd\textquotesingle{}*\textasciigrave{}r2dxj\textquotesingle{})*(1 -{}\newline \hspace*{1em}\textasciigrave{}r2dxj\textquotesingle{}))} & Helper term from the Cinelli-Hazlett bound formula.\\
57 & \texttt{\hspace*{2.0em}scalar \textasciigrave{}r2yz\textquotesingle{}  = ((sqrt(\textasciigrave{}ky\textquotesingle{}) + sqrt(\textasciigrave{}r2zxj\textquotesingle{}))/sqrt(1 -{}\newline \hspace*{1em}\textasciigrave{}r2zxj\textquotesingle{}))\^{}2 * (\textasciigrave{}r2yxj\textquotesingle{}/(1 -{} \textasciigrave{}r2yxj\textquotesingle{}))} & Implied partial R2 of the confounder with Shannon.\\
58 & \texttt{\hspace*{2.0em}if \textasciigrave{}r2yz\textquotesingle{} > 1 scalar \textasciigrave{}r2yz\textquotesingle{} = 1} & Caps it at 1.\\
59 & \texttt{\hspace*{2.0em}* omitted-{}variable bias (eq 12) and adjusted SE (eq 13)} & Comment (not executed): explains the code that follows.\\
60 & \texttt{\hspace*{2.0em}scalar \textasciigrave{}bias\textquotesingle{}  = \textasciigrave{}se\textquotesingle{}*sqrt(\textasciigrave{}r2yz\textquotesingle{}*\textasciigrave{}r2dz\textquotesingle{}/(1 -{}\newline \hspace*{1em}\textasciigrave{}r2dz\textquotesingle{})*\textasciigrave{}df\textquotesingle{})} & Omitted-variable bias of such a confounder (eq. 12).\\
61 & \texttt{\hspace*{2.0em}scalar \textasciigrave{}adj\textquotesingle{}   = sign(\textasciigrave{}est\textquotesingle{})*(abs(\textasciigrave{}est\textquotesingle{}) -{} \textasciigrave{}bias\textquotesingle{})\newline \hspace*{1em}// Z pushes toward zero} & Adjusted estimate: the bias pushes the estimate toward zero.\\
62 & \texttt{\hspace*{2.0em}scalar \textasciigrave{}adjse\textquotesingle{} = \textasciigrave{}se\textquotesingle{}*sqrt((1 -{} \textasciigrave{}r2yz\textquotesingle{})/(1 -{}\newline \hspace*{1em}\textasciigrave{}r2dz\textquotesingle{})*\textasciigrave{}df\textquotesingle{}/(\textasciigrave{}df\textquotesingle{} -{} 1))} & Adjusted standard error (eq. 13).\\
63 & \texttt{\hspace*{2.0em}return scalar r2dz  = \textasciigrave{}r2dz\textquotesingle{}} & Returns the confounder's partial R2 with PSI.\\
64 & \texttt{\hspace*{2.0em}return scalar r2yz  = \textasciigrave{}r2yz\textquotesingle{}} & Returns its partial R2 with Shannon.\\
65 & \texttt{\hspace*{2.0em}return scalar adj   = \textasciigrave{}adj\textquotesingle{}} & Returns the adjusted estimate.\\
66 & \texttt{\hspace*{2.0em}return scalar adjse = \textasciigrave{}adjse\textquotesingle{}} & Returns the adjusted SE.\\
67 & \texttt{\hspace*{2.0em}return scalar adjt  = \textasciigrave{}adj\textquotesingle{}/\textasciigrave{}adjse\textquotesingle{}} & Returns the adjusted t-statistic.\\
68 & \texttt{end} & Ends the program definition.\\
70 & \texttt{*-{}-{}-{}-{}-{}-{}-{}-{}-{}-{}-{}-{}-{}-{}-{}-{}-{}-{}-{}-{}-{}-{}-{}-{}-{}-{}-{}-{}-{}-{}-{}-{}-{}-{}-{}-{}-{}-{}-{}-{}-{}-{}-{}-{}-{}-{}-{}-{}-{}-{}-{}-{}-{}-{}-{}-{}-{}\newline \hspace*{1em}-{}-{}-{}-{}-{}-{}-{}-{}-{}-{}-{}-{}-{}-{}-{}-{}-{}} & Divider line (comment).\\
71 & \texttt{use \textquotedbl{}\$SAMPLE\textquotedbl{}, clear} & Loads the analysis dataset built in step 1\_2.\\
72 & \texttt{local fe1 \textquotedbl{}county year\textquotedbl{}} & Fixed effects of specification S1 (paper): county + year-month.\\
73 & \texttt{local fe2 \textquotedbl{}county year\#urban\textquotedbl{}} & Fixed effects of S2: county + separate year-month effects for urban and rural.\\
74 & \texttt{local nm1 \textquotedbl{}S1 baseline (county + YM FE)\textquotedbl{}} & Label of S1.\\
75 & \texttt{local nm2 \textquotedbl{}S2 urban x YM FE\textquotedbl{}} & Label of S2.\\
77 & \texttt{tempname RV BD} & Handles for two results files.\\
78 & \texttt{postfile \textasciigrave{}RV\textquotesingle{} str32 spec double(est se\_classical\newline \hspace*{1em}t\_classical df R2\_YD RV\_q1 RV\_q1\_a05 ///} & Opens the robustness-value results file (columns listed) ...\\
79 & \texttt{\hspace*{2.0em}se\_clustered t\_clustered G lambda RV\_q1\_a05\_clustered)\newline \hspace*{1em}using \textquotedbl{}\$T2/tables/rv\_tmp.dta\textquotedbl{}, replace} & ... continued.\\
80 & \texttt{postfile \textasciigrave{}BD\textquotesingle{} str32 spec str32 benchmark double(k R2\_D\_X\newline \hspace*{1em}R2\_Y\_X R2\_D\_Z R2\_Y\_Z adj\_est adj\_se adj\_t adj\_t\_clustered)\newline \hspace*{1em}///} & Opens the benchmark-bounds results file ...\\
81 & \texttt{\hspace*{2.0em}using \textquotedbl{}\$T2/tables/bd\_tmp.dta\textquotedbl{}, replace} & ... continued.\\
83 & \texttt{forvalues s = 1/2 \{} & Loops over the two specifications S1 and S2.\\
84 & \texttt{\hspace*{2.0em}di as res \_n \textquotedbl{}======================= \textasciigrave{}nm\textasciigrave{}s\textquotesingle{}\textquotesingle{}\newline \hspace*{1em}=======================\textquotedbl{}} & Prints which specification is running.\\
86 & \texttt{\hspace*{2.0em}*-{}-{}-{}-{}-{}-{}-{}-{}-{}-{}-{}-{}-{}-{}-{}-{} clustered SE first (the paper\textquotesingle{}s\newline \hspace*{1em}inference)} & Divider line (comment).\\
87 & \texttt{\hspace*{2.0em}reghdfe ShannonBD PI \$control, absorb(\textasciigrave{}fe\textasciigrave{}s\textquotesingle{}\textquotesingle{})\newline \hspace*{1em}cluster(county)} & Regresses Shannon bird diversity on PSI (solar policy stringency); controlling for the 9 controls; with the fixed effects of specification s (S1: county + year-month; S2: county + urban-by-year-month); standard errors clustered by county.\\
88 & \texttt{\hspace*{2.0em}keep if e(sample)} & Keeps exactly the rows used by the regression (drops singletons), so all later regressions use the same sample.\\
89 & \texttt{\hspace*{2.0em}scalar secl\_\textasciigrave{}s\textquotesingle{} = \_se[PI]} & Stores the clustered SE of PSI.\\
90 & \texttt{\hspace*{2.0em}scalar tcl\_\textasciigrave{}s\textquotesingle{}  = \_b[PI]/\_se[PI]} & Stores the clustered t.\\
91 & \texttt{\hspace*{2.0em}scalar G\_\textasciigrave{}s\textquotesingle{}    = e(N\_clust)} & Stores the number of clusters (2,344).\\
93 & \texttt{\hspace*{2.0em}*-{}-{}-{}-{}-{}-{}-{}-{}-{}-{}-{}-{}-{}-{}-{}-{} A. sensitivity statistics (classical SE,\newline \hspace*{1em}exact df)} & Divider line (comment).\\
94 & \texttt{\hspace*{2.0em}reghdfe ShannonBD PI \$control, absorb(\textasciigrave{}fe\textasciigrave{}s\textquotesingle{}\textquotesingle{})\newline \hspace*{1em}vce(unadjusted)} & Regresses Shannon bird diversity on PSI (solar policy stringency); controlling for the 9 controls; with the fixed effects of specification s (S1: county + year-month; S2: county + urban-by-year-month); classical (unclustered) standard errors, as Cinelli-Hazlett formulas require.\\
95 & \texttt{\hspace*{2.0em}scalar est\_\textasciigrave{}s\textquotesingle{} = \_b[PI]} & Stores the PSI estimate.\\
96 & \texttt{\hspace*{2.0em}scalar se\_\textasciigrave{}s\textquotesingle{}  = \_se[PI]} & Stores the classical SE.\\
97 & \texttt{\hspace*{2.0em}scalar df\_\textasciigrave{}s\textquotesingle{}  = e(df\_r)} & Stores reghdfe's exact residual degrees of freedom (43,898), which already subtracts the absorbed fixed effects.\\
98 & \texttt{\hspace*{2.0em}scalar t\_\textasciigrave{}s\textquotesingle{}   = \_b[PI]/\_se[PI]} & Stores the classical t (-4.58).\\
99 & \texttt{\hspace*{2.0em}scalar lam\_\textasciigrave{}s\textquotesingle{} = secl\_\textasciigrave{}s\textquotesingle{}/se\_\textasciigrave{}s\textquotesingle{}} & lambda = clustered SE / classical SE (1.337 for S1).\\
100 & \texttt{\hspace*{2.0em}foreach b in \$control \{} & Loops over the 9 controls.\\
101 & \texttt{\hspace*{4.0em}scalar tY = \_b[\textasciigrave{}b\textquotesingle{}]/\_se[\textasciigrave{}b\textquotesingle{}]} & t-statistic of control b in the Shannon regression.\\
102 & \texttt{\hspace*{4.0em}scalar r2y\_\textasciigrave{}b\textquotesingle{} = tY\^{}2/(tY\^{}2 + df\_\textasciigrave{}s\textquotesingle{})        // R2 of Y\newline \hspace*{1em}with Xj | D, X-{}j, FE} & Partial R2 of control b with Shannon (given PSI, other controls and FE).\\
103 & \texttt{\hspace*{2.0em}\}} & Closes the loop / block.\\
104 & \texttt{\hspace*{2.0em}ch\_rv t\_\textasciigrave{}s\textquotesingle{} df\_\textasciigrave{}s\textquotesingle{} 1 0.05} & Runs ch\_rv with classical t and df: q = 1 (remove all), alpha = 0.05.\\
105 & \texttt{\hspace*{2.0em}scalar r2yd\_\textasciigrave{}s\textquotesingle{} = r(r2yd)} & Saves the partial R2 of PSI with Shannon.\\
106 & \texttt{\hspace*{2.0em}scalar rv\_\textasciigrave{}s\textquotesingle{}   = r(rv)} & Saves RV.\\
107 & \texttt{\hspace*{2.0em}scalar rva\_\textasciigrave{}s\textquotesingle{}  = r(rva)} & Saves RV-alpha.\\
108 & \texttt{\hspace*{2.0em}di \textquotedbl{}Partial R2 of PSI with Shannon : \textquotedbl{} \%9.6f r(r2yd)} & Prints the partial R2.\\
109 & \texttt{\hspace*{2.0em}di \textquotedbl{}Robustness value RV(q=1)       : \textquotedbl{} \%9.6f r(rv)} & Prints RV.\\
110 & \texttt{\hspace*{2.0em}di \textquotedbl{}Robustness value RV(q=1,a=.05) : \textquotedbl{} \%9.6f r(rva)} & Prints RV-alpha.\\
112 & \texttt{\hspace*{2.0em}reghdfe PI \$control, absorb(\textasciigrave{}fe\textasciigrave{}s\textquotesingle{}\textquotesingle{}) vce(unadjusted)} & Regresses PSI (solar policy stringency) on \$control; with the fixed effects of specification s (S1: county + year-month; S2: county + urban-by-year-month); classical (unclustered) standard errors, as Cinelli-Hazlett formulas require.\\
113 & \texttt{\hspace*{2.0em}scalar rssD\_\textasciigrave{}s\textquotesingle{} = e(rss)} & Stores the residual sum of squares of the PSI regression.\\
114 & \texttt{\hspace*{2.0em}foreach b in \$control \{} & Loops over the 9 controls.\\
115 & \texttt{\hspace*{4.0em}scalar tD = \_b[\textasciigrave{}b\textquotesingle{}]/\_se[\textasciigrave{}b\textquotesingle{}]} & t-statistic of control b in the PSI regression.\\
116 & \texttt{\hspace*{4.0em}scalar r2d\_\textasciigrave{}b\textquotesingle{} = tD\^{}2/(tD\^{}2 + e(df\_r))       // R2 of D\newline \hspace*{1em}with Xj | X-{}j, FE} & Partial R2 of control b with PSI.\\
117 & \texttt{\hspace*{2.0em}\}} & Closes the loop / block.\\
119 & \texttt{\hspace*{2.0em}*-{}-{}-{}-{}-{}-{}-{}-{}-{}-{}-{}-{}-{}-{}-{}-{} B. benchmark bounds} & Divider line (comment).\\
120 & \texttt{\hspace*{2.0em}foreach b in \$control \{} & Loops over the 9 controls.\\
121 & \texttt{\hspace*{4.0em}foreach k in 1 2 3 \{} & Loops over k = 1, 2, 3 times the benchmark strength.\\
122 & \texttt{\hspace*{4.0em}ch\_bound r2d\_\textasciigrave{}b\textquotesingle{} r2y\_\textasciigrave{}b\textquotesingle{} \textasciigrave{}k\textquotesingle{} \textasciigrave{}k\textquotesingle{} est\_\textasciigrave{}s\textquotesingle{} se\_\textasciigrave{}s\textquotesingle{} df\_\textasciigrave{}s\textquotesingle{}} & Computes the adjusted estimate if a confounder were k times as strong as control b.\\
123 & \texttt{\hspace*{4.0em}post \textasciigrave{}BD\textquotesingle{} (\textquotedbl{}\textasciigrave{}nm\textasciigrave{}s\textquotesingle{}\textquotesingle{}\textquotedbl{}) (\textquotedbl{}\textasciigrave{}b\textquotesingle{}\textquotedbl{}) (\textasciigrave{}k\textquotesingle{}) (r2d\_\textasciigrave{}b\textquotesingle{}) (r2y\_\textasciigrave{}b\textquotesingle{})\newline \hspace*{1em}(r(r2dz)) (r(r2yz)) ///} & Writes the benchmark result row ...\\
124 & \texttt{\hspace*{4.0em}(r(adj)) (r(adjse)) (r(adjt)) (r(adjt)/lam\_\textasciigrave{}s\textquotesingle{})} & ... including the adjusted t with clustered SEs (adjusted t / lambda).\\
125 & \texttt{\hspace*{4.0em}di \textquotedbl{}\textasciigrave{}b\textquotesingle{} k=\textasciigrave{}k\textquotesingle{}: R2\_DZ=\textquotedbl{} \%7.4f r(r2dz) \textquotedbl{} R2\_YZ=\textquotedbl{} \%7.4f\newline \hspace*{1em}r(r2yz) ///} & Prints the benchmark result ...\\
126 & \texttt{\hspace*{4.0em}\textquotedbl{} adj.est=\textquotedbl{} \%9.5f r(adj) \textquotedbl{} adj.t=\textquotedbl{} \%6.2f r(adjt)} & ... continued.\\
127 & \texttt{\hspace*{4.0em}\}} & Closes the loop / block.\\
128 & \texttt{\hspace*{2.0em}\}} & Closes the loop / block.\\
130 & \texttt{\hspace*{2.0em}*-{}-{}-{}-{}-{}-{}-{}-{}-{}-{}-{}-{}-{}-{}-{}-{} D. clustered variant: same partial R2s,\newline \hspace*{1em}significance} & Divider line (comment).\\
131 & \texttt{\hspace*{2.0em}*                   threshold inflated by lambda =\newline \hspace*{1em}clustered/classical SE} & Comment (not executed): explains the code that follows.\\
132 & \texttt{\hspace*{2.0em}ch\_rv t\_\textasciigrave{}s\textquotesingle{} df\_\textasciigrave{}s\textquotesingle{} 1 0.05 lam\_\textasciigrave{}s\textquotesingle{}} & Re-runs ch\_rv with lambda: RV-alpha under clustered inference (0.93\% for S1).\\
133 & \texttt{\hspace*{2.0em}scalar rvacl\_\textasciigrave{}s\textquotesingle{} = r(rva)} & Saves the clustered RV-alpha.\\
134 & \texttt{\hspace*{2.0em}di \textquotedbl{}Clustered t = \textquotedbl{} \%6.3f tcl\_\textasciigrave{}s\textquotesingle{} \textquotedbl{}  lambda = \textquotedbl{} \%5.3f\newline \hspace*{1em}lam\_\textasciigrave{}s\textquotesingle{} ///} & Prints the clustered results ...\\
135 & \texttt{\hspace*{3.5em}\textquotedbl{}  RV(q=1,a=.05, clustered) = \textquotedbl{} \%8.6f r(rva)} & ... continued.\\
137 & \texttt{\hspace*{2.0em}post \textasciigrave{}RV\textquotesingle{} (\textquotedbl{}\textasciigrave{}nm\textasciigrave{}s\textquotesingle{}\textquotesingle{}\textquotedbl{}) (est\_\textasciigrave{}s\textquotesingle{}) (se\_\textasciigrave{}s\textquotesingle{}) (t\_\textasciigrave{}s\textquotesingle{}) (df\_\textasciigrave{}s\textquotesingle{})\newline \hspace*{1em}(r2yd\_\textasciigrave{}s\textquotesingle{}) (rv\_\textasciigrave{}s\textquotesingle{}) (rva\_\textasciigrave{}s\textquotesingle{}) ///} & Writes the robustness-value row ...\\
138 & \texttt{\hspace*{4.0em}(secl\_\textasciigrave{}s\textquotesingle{}) (tcl\_\textasciigrave{}s\textquotesingle{}) (G\_\textasciigrave{}s\textquotesingle{}) (lam\_\textasciigrave{}s\textquotesingle{}) (rvacl\_\textasciigrave{}s\textquotesingle{})} & ... continued.\\
139 & \texttt{\hspace*{2.0em}use \textquotedbl{}\$SAMPLE\textquotedbl{}, clear} & Loads the analysis dataset built in step 1\_2.\\
140 & \texttt{\}} & Closes the loop / block.\\
142 & \texttt{*-{}-{}-{}-{}-{}-{}-{}-{}-{}-{}-{}-{}-{}-{}-{}-{} C. group benchmark: urban x year-{}month\newline \hspace*{1em}trends} & Divider line (comment).\\
143 & \texttt{* partial R2 that urban x YM FE add, beyond county + YM FE\newline \hspace*{1em}and controls,} & Comment (not executed): explains the code that follows.\\
144 & \texttt{*   for PSI (treatment model) and for Shannon given PSI\newline \hspace*{1em}(outcome model)} & Comment (not executed): explains the code that follows.\\
145 & \texttt{reghdfe PI \$control, absorb(county year) vce(unadjusted)} & Regresses PSI (solar policy stringency) on \$control; with county and year-month fixed effects; classical (unclustered) standard errors, as Cinelli-Hazlett formulas require.\\
146 & \texttt{keep if e(sample)} & Keeps exactly the rows used by the regression (drops singletons), so all later regressions use the same sample.\\
147 & \texttt{scalar rr = e(rss)} & Residual sum of squares of the model without urban x month effects.\\
148 & \texttt{reghdfe PI \$control, absorb(county year\#urban)\newline \hspace*{1em}vce(unadjusted)} & Regresses PSI (solar policy stringency) on \$control; with county fixed effects plus separate year-month effects for urban and rural units; classical (unclustered) standard errors, as Cinelli-Hazlett formulas require.\\
149 & \texttt{scalar r2d\_urb = (rr -{} e(rss))/rr} & Share of PSI residual variance explained by urban x month effects (12.3\%).\\
150 & \texttt{reghdfe ShannonBD PI \$control, absorb(county year)\newline \hspace*{1em}vce(unadjusted)} & Paper's model with classical SEs (for Cinelli-Hazlett).\\
151 & \texttt{scalar rr = e(rss)} & Residual sum of squares of the model without urban x month effects.\\
152 & \texttt{reghdfe ShannonBD PI \$control, absorb(county year\#urban)\newline \hspace*{1em}vce(unadjusted)} & Regresses Shannon bird diversity on PSI (solar policy stringency); controlling for the 9 controls; with county fixed effects plus separate year-month effects for urban and rural units; classical (unclustered) standard errors, as Cinelli-Hazlett formulas require.\\
153 & \texttt{scalar r2y\_urb = (rr -{} e(rss))/rr} & Share of Shannon residual variance explained by urban x month effects (0.59\%).\\
154 & \texttt{di as res \_n \textquotedbl{}Urban x YM trends: partial R2 with PSI = \textquotedbl{}\newline \hspace*{1em}\%7.4f r2d\_urb \textquotedbl{}   with Shannon | PSI = \textquotedbl{} \%7.4f r2y\_urb} & Prints the two partial R2 of the urban benchmark.\\
155 & \texttt{foreach k in 0.25 0.5 1 \{} & Loops over 0.25, 0.5 and 1 times the urban benchmark.\\
156 & \texttt{\hspace*{2.0em}ch\_bound r2d\_urb r2y\_urb \textasciigrave{}k\textquotesingle{} \textasciigrave{}k\textquotesingle{} est\_1 se\_1 df\_1} & Adjusted estimate if a confounder were k times as strong as urban trends.\\
157 & \texttt{\hspace*{2.0em}post \textasciigrave{}BD\textquotesingle{} (\textquotedbl{}\textasciigrave{}nm1\textquotesingle{}\textquotedbl{}) (\textquotedbl{}Urban x YM trends (group)\textquotedbl{}) (\textasciigrave{}k\textquotesingle{})\newline \hspace*{1em}(r2d\_urb) (r2y\_urb) (r(r2dz)) (r(r2yz)) ///} & Writes the urban-benchmark row ...\\
158 & \texttt{\hspace*{4.0em}(r(adj)) (r(adjse)) (r(adjt)) (r(adjt)/lam\_1)} & ... continued.\\
159 & \texttt{\hspace*{2.0em}di \textquotedbl{}urban x YM k=\textasciigrave{}k\textquotesingle{}: adj.est=\textquotedbl{} \%9.5f r(adj) \textquotedbl{}  adj.t=\textquotedbl{}\newline \hspace*{1em}\%6.2f r(adjt)} & Prints it.\\
160 & \texttt{\}} & Closes the loop / block.\\
162 & \texttt{postclose \textasciigrave{}RV\textquotesingle{}} & Closes the results file so it can be loaded.\\
163 & \texttt{postclose \textasciigrave{}BD\textquotesingle{}} & Closes the results file so it can be loaded.\\
165 & \texttt{use \textquotedbl{}\$T2/tables/rv\_tmp.dta\textquotedbl{}, clear} & Loads the dataset \$T2/tables/rv\_tmp.dta into memory, replacing whatever was loaded.\\
166 & \texttt{format est se\_classical \%9.5f} & Sets display formats (decimals) for these variables.\\
167 & \texttt{format R2\_YD RV\_q1 RV\_q1\_a05 RV\_q1\_a05\_clustered \%9.5f} & Sets display formats (decimals) for these variables.\\
168 & \texttt{list, noobs sep(0) abbreviate(20)} & Lists the rows/variables on screen.\\
169 & \texttt{export delimited using\newline \hspace*{1em}\textquotedbl{}\$T2/tables/2\_1\_robustness\_values.csv\textquotedbl{}, replace} & Exports the current data to a CSV file.\\
170 & \texttt{erase \textquotedbl{}\$T2/tables/rv\_tmp.dta\textquotedbl{}} & Deletes a temporary file.\\
172 & \texttt{use \textquotedbl{}\$T2/tables/bd\_tmp.dta\textquotedbl{}, clear} & Loads the dataset \$T2/tables/bd\_tmp.dta into memory, replacing whatever was loaded.\\
173 & \texttt{format R2\_* \%8.5f} & Sets display formats (decimals) for these variables.\\
174 & \texttt{format adj\_est adj\_se \%9.5f} & Sets display formats (decimals) for these variables.\\
175 & \texttt{format adj\_t adj\_t\_clustered \%6.2f} & Sets display formats (decimals) for these variables.\\
176 & \texttt{list, noobs sep(0) abbreviate(24)} & Lists the rows/variables on screen.\\
177 & \texttt{export delimited using\newline \hspace*{1em}\textquotedbl{}\$T2/tables/2\_1\_benchmark\_bounds.csv\textquotedbl{}, replace} & Exports the current data to a CSV file.\\
178 & \texttt{save \textquotedbl{}\$T2/tables/bounds\_for\_plot.dta\textquotedbl{}, replace} & Keeps the bounds for the contour plots.\\
179 & \texttt{erase \textquotedbl{}\$T2/tables/bd\_tmp.dta\textquotedbl{}} & Deletes a temporary file.\\
181 & \texttt{*-{}-{}-{}-{}-{}-{}-{}-{}-{}-{}-{}-{}-{}-{}-{}-{}-{}-{}-{}-{}-{}-{}-{}-{}-{}-{}-{}-{}-{}-{}-{}-{}-{}-{}-{}-{}-{}-{}-{}-{}-{}-{}-{}-{}-{}-{}-{}-{}-{}-{}-{}-{}-{}-{}-{}-{}-{}\newline \hspace*{1em}-{}-{}-{}-{}-{}-{}-{}-{}-{}-{}-{}-{}-{}-{}-{}-{}-{}} & Divider line (comment).\\
182 & \texttt{* E. Contour plots (adjusted estimate and adjusted t) for\newline \hspace*{1em}each spec} & Comment (not executed): explains the code that follows.\\
183 & \texttt{*-{}-{}-{}-{}-{}-{}-{}-{}-{}-{}-{}-{}-{}-{}-{}-{}-{}-{}-{}-{}-{}-{}-{}-{}-{}-{}-{}-{}-{}-{}-{}-{}-{}-{}-{}-{}-{}-{}-{}-{}-{}-{}-{}-{}-{}-{}-{}-{}-{}-{}-{}-{}-{}-{}-{}-{}-{}\newline \hspace*{1em}-{}-{}-{}-{}-{}-{}-{}-{}-{}-{}-{}-{}-{}-{}-{}-{}-{}} & Divider line (comment).\\
184 & \texttt{forvalues s = 1/2 \{} & Loops over the two specifications S1 and S2.\\
185 & \texttt{\hspace*{2.0em}clear} & Removes data (and with "all", also programs and stored results) from memory.\\
186 & \texttt{\hspace*{2.0em}local n = 61} & Grid size: 61 x 61 points.\\
187 & \texttt{\hspace*{2.0em}set obs \textasciigrave{}=\textasciigrave{}n\textquotesingle{}*\textasciigrave{}n\textquotesingle{}\textquotesingle{}} & Creates 3,721 empty rows for the grid.\\
188 & \texttt{\hspace*{2.0em}gen r2dz = 0.0025*floor((\_n-{}1)/\textasciigrave{}n\textquotesingle{})                 // 0\newline \hspace*{1em}.. 0.15} & x-coordinate: confounder partial R2 with PSI, 0 to 0.15.\\
189 & \texttt{\hspace*{2.0em}gen r2yz = 0.0025*mod(\_n-{}1, \textasciigrave{}n\textquotesingle{})} & y-coordinate: partial R2 with Shannon, 0 to 0.15.\\
190 & \texttt{\hspace*{2.0em}replace r2dz = 0.0001 if r2dz == 0} & Avoids exact 0.\\
191 & \texttt{\hspace*{2.0em}replace r2yz = 0.0001 if r2yz == 0} & Avoids exact 0.\\
192 & \texttt{\hspace*{2.0em}gen bias  = se\_\textasciigrave{}s\textquotesingle{}*sqrt(r2yz*r2dz/(1 -{} r2dz)*df\_\textasciigrave{}s\textquotesingle{})} & Bias at each grid point.\\
193 & \texttt{\hspace*{2.0em}gen adj   = sign(est\_\textasciigrave{}s\textquotesingle{})*(abs(est\_\textasciigrave{}s\textquotesingle{}) -{} bias)} & Adjusted estimate at each grid point.\\
194 & \texttt{\hspace*{2.0em}gen adjse = se\_\textasciigrave{}s\textquotesingle{}*sqrt((1 -{} r2yz)/(1 -{}\newline \hspace*{1em}r2dz)*df\_\textasciigrave{}s\textquotesingle{}/(df\_\textasciigrave{}s\textquotesingle{} -{} 1))} & Adjusted SE at each grid point.\\
195 & \texttt{\hspace*{2.0em}gen adjt  = adj/adjse} & Adjusted t at each grid point.\\
196 & \texttt{\hspace*{2.0em}* benchmark points (Duration k=1,2,3) for spec 1 only} & Comment (not executed): explains the code that follows.\\
197 & \texttt{\hspace*{2.0em}if \textasciigrave{}s\textquotesingle{} == 1 \{} & For S1 only, add the benchmark points:\\
198 & \texttt{\hspace*{4.0em}append using \textquotedbl{}\$T2/tables/bounds\_for\_plot.dta\textquotedbl{},\newline \hspace*{1em}keep(benchmark k R2\_D\_Z R2\_Y\_Z spec)} & Appends the benchmark coordinates.\\
199 & \texttt{\hspace*{4.0em}keep if missing(spec) | (strpos(spec,\textquotedbl{}S1\textquotedbl{}) \&\newline \hspace*{1em}(benchmark==\textquotedbl{}Duration\textquotedbl{} | strpos(benchmark,\textquotedbl{}Urban\textquotedbl{})))} & Keeps grid rows plus the Duration and urban benchmark points.\\
200 & \texttt{\hspace*{4.0em}gen pt\_x = R2\_D\_Z} & x of each benchmark point.\\
201 & \texttt{\hspace*{4.0em}gen pt\_y = R2\_Y\_Z} & y of each benchmark point.\\
202 & \texttt{\hspace*{4.0em}gen str20 pt\_lab = cond(benchmark==\textquotedbl{}Duration\textquotedbl{}, string(k) +\newline \hspace*{1em}\textquotedbl{}x Duration\textquotedbl{}, ///} & Label of each point, e.g. "2x Duration" ...\\
203 & \texttt{\hspace*{4.0em}cond(benchmark!=\textquotedbl{}\textquotedbl{}, string(k) + \textquotedbl{}x urban-{}YM\textquotedbl{}, \textquotedbl{}\textquotedbl{}))} & ... or "0.5x urban-YM".\\
204 & \texttt{\hspace*{4.0em}replace pt\_x = . if pt\_x > 0.15 | pt\_y > 0.15} & Hides points outside the plot range.\\
205 & \texttt{\hspace*{2.0em}\}} & Closes the loop / block.\\
206 & \texttt{\hspace*{2.0em}else \{} & Otherwise ...\\
207 & \texttt{\hspace*{4.0em}gen pt\_x = .} & S2: no benchmark points.\\
208 & \texttt{\hspace*{4.0em}gen pt\_y = .} & S2: no benchmark points.\\
209 & \texttt{\hspace*{4.0em}gen pt\_lab = \textquotedbl{}\textquotedbl{}} & S2: no labels.\\
210 & \texttt{\hspace*{2.0em}\}} & Closes the loop / block.\\
211 & \texttt{\hspace*{2.0em}local e = string(est\_\textasciigrave{}s\textquotesingle{}, \textquotedbl{}\%6.4f\textquotedbl{})} & Estimate as text for the subtitle.\\
212 & \texttt{\hspace*{2.0em}* Contours of the adjusted estimate are closed-{}form: an\newline \hspace*{1em}adjusted estimate} & Comment (not executed): explains the code that follows.\\
213 & \texttt{\hspace*{2.0em}* equal to share f of the original needs bias =\newline \hspace*{1em}(1-{}f)|est|, i.e.} & Comment (not executed): explains the code that follows.\\
214 & \texttt{\hspace*{2.0em}*   R2\_YZ = B\^{}2 (1 -{} R2\_DZ) / R2\_DZ,   B = (1-{}f)|est| /\newline \hspace*{1em}(se*sqrt(df))} & Comment (not executed): explains the code that follows.\\
215 & \texttt{\hspace*{2.0em}local K = se\_\textasciigrave{}s\textquotesingle{}*sqrt(df\_\textasciigrave{}s\textquotesingle{})} & Constant used in the closed-form contour curves.\\
216 & \texttt{\hspace*{2.0em}local plots \textquotedbl{}\textquotedbl{}} & Starts an empty list of curve layers.\\
217 & \texttt{\hspace*{2.0em}local txt \textquotedbl{}\textquotedbl{}} & Starts an empty list of curve labels.\\
218 & \texttt{\hspace*{2.0em}foreach f in 0.8 0.6 0.4 0.2 0 \{} & Loops over contour levels: 80\%, 60\%, 40\%, 20\% of the estimate left, and 0 (zero).\\
219 & \texttt{\hspace*{4.0em}local B  = (1-{}\textasciigrave{}f\textquotesingle{})*abs(est\_\textasciigrave{}s\textquotesingle{})/\textasciigrave{}K\textquotesingle{}} & Bias needed to remove (1-f) of the estimate, scaled.\\
220 & \texttt{\hspace*{4.0em}local x0 = max(0.0002, \textasciigrave{}B\textquotesingle{}\^{}2/(0.15 + \textasciigrave{}B\textquotesingle{}\^{}2))} & Start of the curve on the x-axis.\\
221 & \texttt{\hspace*{4.0em}local col = cond(\textasciigrave{}f\textquotesingle{}==0, \textquotedbl{}red\textquotedbl{}, \textquotedbl{}navy\textquotedbl{})} & Zero curve red, others navy.\\
222 & \texttt{\hspace*{4.0em}local pat = cond(\textasciigrave{}f\textquotesingle{}==0, \textquotedbl{}dash\textquotedbl{}, \textquotedbl{}solid\textquotedbl{})} & Zero curve dashed.\\
223 & \texttt{\hspace*{4.0em}local wid = cond(\textasciigrave{}f\textquotesingle{}==0, \textquotedbl{}thick\textquotedbl{}, \textquotedbl{}medium\textquotedbl{})} & Zero curve thick.\\
224 & \texttt{\hspace*{4.0em}local plots \textasciigrave{}plots\textquotesingle{} (function y = (\textasciigrave{}B\textquotesingle{})\^{}2*(1-{}x)/x,\newline \hspace*{1em}range(\textasciigrave{}x0\textquotesingle{} 0.15) lcolor(\textasciigrave{}col\textquotesingle{}) lpattern(\textasciigrave{}pat\textquotesingle{})\newline \hspace*{1em}lwidth(\textasciigrave{}wid\textquotesingle{}))} & Adds the curve R2\_YZ = B\^{}2 (1 - R2\_DZ)/R2\_DZ: all confounders that leave exactly share f of the estimate.\\
225 & \texttt{\hspace*{4.0em}local yl = (\textasciigrave{}B\textquotesingle{})\^{}2*(0.85)/0.15} & Height at which to place the curve label.\\
226 & \texttt{\hspace*{4.0em}if \textasciigrave{}yl\textquotesingle{} <= 0.15 local txt \textasciigrave{}txt\textquotesingle{} text(\textasciigrave{}yl\textquotesingle{} 0.151\newline \hspace*{1em}\textquotedbl{}\textasciigrave{}=round(100*\textasciigrave{}f\textquotesingle{})\textquotesingle{}\% left\textquotedbl{}, place(e) size(vsmall)\newline \hspace*{1em}color(\textasciigrave{}col\textquotesingle{}))} & Adds a label "xx\% left" at the right edge.\\
227 & \texttt{\hspace*{2.0em}\}} & Closes the loop / block.\\
228 & \texttt{\hspace*{2.0em}twoway \textasciigrave{}plots\textquotesingle{} ///} & Draws the contour curves ...\\
229 & \texttt{\hspace*{4.0em}(scatter pt\_y pt\_x, msymbol(D) mcolor(black)\newline \hspace*{1em}mlabel(pt\_lab) mlabcolor(black) mlabsize(vsmall)), ///} & ... plus the benchmark diamonds with labels ...\\
230 & \texttt{\hspace*{4.0em}\textasciigrave{}txt\textquotesingle{} xscale(range(0 0.17)) ///} & ... plus curve labels; x-axis range.\\
231 & \texttt{\hspace*{4.0em}xtitle(\textquotedbl{}Partial R\{superscript:2\} of confounder with PSI\textquotedbl{})\newline \hspace*{1em}///} & Continuation of the previous command: table title.\\
232 & \texttt{\hspace*{4.0em}ytitle(\textquotedbl{}Partial R\{superscript:2\} of confounder with\newline \hspace*{1em}Shannon\textquotedbl{}) ///} & Continuation of the previous command: table title.\\
233 & \texttt{\hspace*{4.0em}title(\textquotedbl{}Adjusted PSI estimate: \textasciigrave{}nm\textasciigrave{}s\textquotesingle{}\textquotesingle{}\textquotedbl{}, size(medium)) ///} & Continuation of the previous command: table title.\\
234 & \texttt{\hspace*{4.0em}subtitle(\textquotedbl{}Estimate = \textasciigrave{}e\textquotesingle{}. Curves: share of the estimate\newline \hspace*{1em}left after confounding; red = zero\textquotedbl{}, size(vsmall)) ///} & Continuation of the previous command: table title.\\
235 & \texttt{\hspace*{4.0em}xlabel(0(0.03)0.15) ylabel(0(0.03)0.15) legend(off)} & Continuation of the graph command: axis titles, reference lines, labels and legend.\\
236 & \texttt{\hspace*{2.0em}graph export \textquotedbl{}\$T2/figures/2\_1\_contour\_estimate\_S\textasciigrave{}s\textquotesingle{}.pdf\textquotedbl{},\newline \hspace*{1em}replace} & Saves the current graph as a PDF.\\
237 & \texttt{\hspace*{2.0em}* t-{}value: red dashed line where the adjusted t reaches\newline \hspace*{1em}-{}1.96} & Comment (not executed): explains the code that follows.\\
238 & \texttt{\hspace*{2.0em}twoway (contourline adjt r2yz r2dz if !missing(r2dz),\newline \hspace*{1em}ccuts(-{}1.96) ///} & Draws the t = -1.96 contour line (significance boundary) ...\\
239 & \texttt{\hspace*{4.0em}ccolors(red) clwidths(thick) clpatterns(dash)) ///} & ... red, thick, dashed.\\
240 & \texttt{\hspace*{4.0em}(scatter pt\_y pt\_x, msymbol(D) mcolor(black)\newline \hspace*{1em}mlabel(pt\_lab) mlabcolor(black) mlabsize(vsmall)), ///} & ... plus the benchmark diamonds with labels ...\\
241 & \texttt{\hspace*{4.0em}xtitle(\textquotedbl{}Partial R\{superscript:2\} of confounder with PSI\textquotedbl{})\newline \hspace*{1em}///} & Continuation of the previous command: table title.\\
242 & \texttt{\hspace*{4.0em}ytitle(\textquotedbl{}Partial R\{superscript:2\} of confounder with\newline \hspace*{1em}Shannon\textquotedbl{}) ///} & Continuation of the previous command: table title.\\
243 & \texttt{\hspace*{4.0em}title(\textquotedbl{}Statistical significance: \textasciigrave{}nm\textasciigrave{}s\textquotesingle{}\textquotesingle{}\textquotedbl{}, size(medium))\newline \hspace*{1em}///} & Continuation of the previous command: table title.\\
244 & \texttt{\hspace*{4.0em}subtitle(\textquotedbl{}Beyond the red dashed line the adjusted estimate\newline \hspace*{1em}is not significant at 5\% (classical SE)\textquotedbl{}, size(vsmall))\newline \hspace*{1em}///} & Continuation of the previous command: table title.\\
245 & \texttt{\hspace*{4.0em}xlabel(0(0.03)0.15) ylabel(0(0.03)0.15) legend(off)} & Continuation of the graph command: axis titles, reference lines, labels and legend.\\
246 & \texttt{\hspace*{2.0em}graph export \textquotedbl{}\$T2/figures/2\_1\_contour\_tvalue\_S\textasciigrave{}s\textquotesingle{}.pdf\textquotedbl{},\newline \hspace*{1em}replace} & Saves the current graph as a PDF.\\
247 & \texttt{\}} & Closes the loop / block.\\
248 & \texttt{erase \textquotedbl{}\$T2/tables/bounds\_for\_plot.dta\textquotedbl{}} & Deletes the temporary file.\\
250 & \texttt{*-{}-{}-{}-{}-{}-{}-{}-{}-{}-{}-{}-{}-{}-{}-{}-{}-{}-{}-{}-{}-{}-{}-{}-{}-{}-{}-{}-{}-{}-{}-{}-{}-{}-{}-{}-{}-{}-{}-{}-{}-{}-{}-{}-{}-{}-{}-{}-{}-{}-{}-{}-{}-{}-{}-{}-{}-{}\newline \hspace*{1em}-{}-{}-{}-{}-{}-{}-{}-{}-{}-{}-{}-{}-{}-{}-{}-{}-{}} & Divider line (comment).\\
251 & \texttt{* F. Cross-{}check with the official sensemakr (Stata) on\newline \hspace*{1em}FE-{}demeaned data.} & Comment (not executed): explains the code that follows.\\
252 & \texttt{*    Point estimate and SE match exactly by FWL;\newline \hspace*{1em}sensemakr\textquotesingle{}s df ignores the} & Comment (not executed): explains the code that follows.\\
253 & \texttt{*    absorbed FE (N-{}k-{}1), so its RVs differ slightly from\newline \hspace*{1em}Part A.} & Comment (not executed): explains the code that follows.\\
254 & \texttt{*-{}-{}-{}-{}-{}-{}-{}-{}-{}-{}-{}-{}-{}-{}-{}-{}-{}-{}-{}-{}-{}-{}-{}-{}-{}-{}-{}-{}-{}-{}-{}-{}-{}-{}-{}-{}-{}-{}-{}-{}-{}-{}-{}-{}-{}-{}-{}-{}-{}-{}-{}-{}-{}-{}-{}-{}-{}\newline \hspace*{1em}-{}-{}-{}-{}-{}-{}-{}-{}-{}-{}-{}-{}-{}-{}-{}-{}-{}} & Divider line (comment).\\
255 & \texttt{use \textquotedbl{}\$SAMPLE\textquotedbl{}, clear} & Loads the analysis dataset built in step 1\_2.\\
256 & \texttt{reghdfe ShannonBD PI \$control, absorb(county year)\newline \hspace*{1em}vce(unadjusted)} & Paper's model with classical SEs (for Cinelli-Hazlett).\\
257 & \texttt{keep if e(sample)} & Keeps exactly the rows used by the regression (drops singletons), so all later regressions use the same sample.\\
258 & \texttt{foreach v in ShannonBD PI \$control \{} & For Shannon, PSI and each control ...\\
259 & \texttt{\hspace*{2.0em}quietly reghdfe \textasciigrave{}v\textquotesingle{}, absorb(county year) resid(r\_\textasciigrave{}v\textquotesingle{})} & ... removes county and year-month means (FWL): r\_v is the demeaned variable.\\
260 & \texttt{\}} & Closes the loop / block.\\
261 & \texttt{capture noisily sensemakr r\_ShannonBD r\_PI r\_Temp r\_Wind\newline \hspace*{1em}r\_Pop r\_Duration r\_Carbon ///} & Runs the official sensemakr on the demeaned variables as a cross-check ...\\
262 & \texttt{\hspace*{2.0em}r\_Water r\_Green r\_Farm r\_Grass, treat(r\_PI)\newline \hspace*{1em}benchmark(r\_Duration) kd(1 2 3)} & ... treatment PSI, Duration as benchmark at 1, 2, 3 times.\\
263 & \texttt{if \_rc di as err \textquotedbl{}sensemakr cross-{}check could not run\newline \hspace*{1em}(rc=\textquotedbl{} \_rc \textquotedbl{}); Part A results stand.\textquotedbl{}} & If sensemakr fails, says so; main results do not depend on it.\\}
\subsection{3\_0\_run\_task3.do}
Runs the five Task 3 steps in order, each with its own log.

\codetable{1 & \texttt{*=========================================================\newline \hspace*{1em}=================} & Divider line (comment).\\
2 & \texttt{* 3\_0\_run\_task3.do  —  runs every Task 3 step in order,\newline \hspace*{1em}one log per step} & Comment (not executed): explains the code that follows.\\
3 & \texttt{*=========================================================\newline \hspace*{1em}=================} & Divider line (comment).\\
4 & \texttt{do \textquotedbl{}/Users/sahiljangra/Desktop/Johns\newline \hspace*{1em}Hopkins/RA/Solar\_Birds\_Replication/00\_config.do\textquotedbl{}} & Runs another do-file: /Users/sahiljangra/Desktop/Johns Hopkins/RA/Solar\_Birds\_Replication/00\_config.do\\
6 & \texttt{foreach step in 3\_1\_adoption\_timing 3\_2\_twfe\_event\_study\newline \hspace*{1em}3\_3\_callaway\_santanna ///} & Starts a loop: foreach step in 3\_1\_adoption\_timing 3\_2\_twfe\_event\_study 3\_3\_callaway\_santanna ///.\\
7 & \texttt{\hspace*{4.0em}3\_4\_imputation\_bjs 3\_5\_sun\_abraham \{} & Continuation: last Task 3 step names; the loop body follows.\\
8 & \texttt{\hspace*{2.0em}capture log close \_all} & Closes any open log (capture = no error if none is open).\\
9 & \texttt{\hspace*{2.0em}log using \textquotedbl{}\$T3/logs/\textasciigrave{}step\textquotesingle{}.log\textquotedbl{}, replace text} & Opens a text log file that records all output of this step.\\
10 & \texttt{\hspace*{2.0em}di as res \_n \textquotedbl{}================ \textasciigrave{}step\textquotesingle{} ================\textquotedbl{}} & Prints the name of the step being run.\\
11 & \texttt{\hspace*{2.0em}clear} & Removes data (and with "all", also programs and stored results) from memory.\\
12 & \texttt{\hspace*{2.0em}estimates clear} & Clears stored estimates.\\
13 & \texttt{\hspace*{2.0em}do \textquotedbl{}\$ROOT/00\_config.do\textquotedbl{}} & Re-runs the settings file so every step starts with the same paths and macros.\\
14 & \texttt{\hspace*{2.0em}capture noisily do \textquotedbl{}\$T3/code/\textasciigrave{}step\textquotesingle{}.do\textquotedbl{}} & Runs the step's do-file; capture = an error does not stop the whole run, noisily = output is still shown.\\
15 & \texttt{\hspace*{2.0em}if \_rc di as err \textquotedbl{}STEP \textasciigrave{}step\textquotesingle{} FAILED with rc = \textquotedbl{} \_rc} & If the step failed (\_rc is the error code), prints a FAILED message.\\
16 & \texttt{\hspace*{2.0em}log close} & Closes the log file.\\
17 & \texttt{\}} & Closes the loop / block.\\}
\subsection{3\_0b\_rerun\_cs\_bjs.do}
Helper that re-runs only the Callaway--Sant'Anna and BJS steps.

\codetable{1 & \texttt{* helper: rerun only the Callaway-{}Sant\textquotesingle{}Anna and imputation\newline \hspace*{1em}steps} & Comment (not executed): explains the code that follows.\\
2 & \texttt{do \textquotedbl{}/Users/sahiljangra/Desktop/Johns\newline \hspace*{1em}Hopkins/RA/Solar\_Birds\_Replication/00\_config.do\textquotedbl{}} & Runs another do-file: /Users/sahiljangra/Desktop/Johns Hopkins/RA/Solar\_Birds\_Replication/00\_config.do\\
3 & \texttt{foreach step in 3\_3\_callaway\_santanna 3\_4\_imputation\_bjs \{} & Starts a loop: foreach step in 3\_3\_callaway\_santanna 3\_4\_imputation\_bjs.\\
4 & \texttt{\hspace*{2.0em}capture log close \_all} & Closes any open log (capture = no error if none is open).\\
5 & \texttt{\hspace*{2.0em}log using \textquotedbl{}\$T3/logs/\textasciigrave{}step\textquotesingle{}.log\textquotedbl{}, replace text} & Opens a text log file that records all output of this step.\\
6 & \texttt{\hspace*{2.0em}clear} & Removes data (and with "all", also programs and stored results) from memory.\\
7 & \texttt{\hspace*{2.0em}estimates clear} & Clears stored estimates.\\
8 & \texttt{\hspace*{2.0em}do \textquotedbl{}\$ROOT/00\_config.do\textquotedbl{}} & Re-runs the settings file so every step starts with the same paths and macros.\\
9 & \texttt{\hspace*{2.0em}capture noisily do \textquotedbl{}\$T3/code/\textasciigrave{}step\textquotesingle{}.do\textquotedbl{}} & Runs the step's do-file; capture = an error does not stop the whole run, noisily = output is still shown.\\
10 & \texttt{\hspace*{2.0em}if \_rc di as err \textquotedbl{}STEP \textasciigrave{}step\textquotesingle{} FAILED with rc = \textquotedbl{} \_rc} & If the step failed (\_rc is the error code), prints a FAILED message.\\
11 & \texttt{\hspace*{2.0em}log close} & Closes the log file.\\
12 & \texttt{\}} & Closes the loop / block.\\}
\subsection{3\_1\_adoption\_timing.do}
Turns PSI into a staggered-adoption design: adoption year, cohorts for the main and paper-style samples, event time, deseasonalised Shannon, and a county-year panel.

\codetable{1 & \texttt{*=========================================================\newline \hspace*{1em}=================} & Divider line (comment).\\
2 & \texttt{* 3\_1\_adoption\_timing.do} & Comment (not executed): explains the code that follows.\\
3 & \texttt{* PURPOSE : Turn the continuous, annual PSI into a\newline \hspace*{1em}STAGGERED ADOPTION design.} & Comment (not executed): explains the code that follows.\\
4 & \texttt{*   Adoption year g = first year (2014-{}2023) in which a\newline \hspace*{1em}county\textquotesingle{}s PSI > 0,} & Comment (not executed): explains the code that follows.\\
5 & \texttt{*   taken from the full annual policy file PI.dta (not\newline \hspace*{1em}only the months in} & Comment (not executed): explains the code that follows.\\
6 & \texttt{*   which birds were recorded). Treatment is treated as\newline \hspace*{1em}absorbing: once a} & Comment (not executed): explains the code that follows.\\
7 & \texttt{*   county has issued solar policy it stays \textquotedbl{}treated\textquotedbl{}.} & Comment (not executed): explains the code that follows.\\
8 & \texttt{* SAMPLES :} & Comment (not executed): explains the code that follows.\\
9 & \texttt{*   main  : counties present in PI.dta (PSI actually\newline \hspace*{1em}measured).} & Comment (not executed): explains the code that follows.\\
10 & \texttt{*           never-{}treated = counties in PI.dta with PSI =\newline \hspace*{1em}0 in every year.} & Comment (not executed): explains the code that follows.\\
11 & \texttt{*   paper : adds the 724 counties absent from PI.dta as\newline \hspace*{1em}never-{}treated} & Comment (not executed): explains the code that follows.\\
12 & \texttt{*           (this is what the authors\textquotesingle{} \textquotedbl{}PSI = 0 if\newline \hspace*{1em}missing\textquotedbl{} implies).} & Comment (not executed): explains the code that follows.\\
13 & \texttt{*   Counties already treated in 2014 have no pre-{}period\newline \hspace*{1em}and are dropped.} & Comment (not executed): explains the code that follows.\\
14 & \texttt{* OUTPUT  : data/staggered\_monthly.dta,\newline \hspace*{1em}data/staggered\_countyyear.dta} & Comment (not executed): explains the code that follows.\\
15 & \texttt{*           tables/3\_1\_cohort\_sizes.csv} & Comment (not executed): explains the code that follows.\\
16 & \texttt{*=========================================================\newline \hspace*{1em}=================} & Divider line (comment).\\
17 & \texttt{capture mkdir \textquotedbl{}\$T3/data\textquotedbl{}} & Creates the Task 3 data folder if it does not exist.\\
19 & \texttt{*-{}-{}-{}-{}-{}-{}-{}-{}-{}-{}-{}-{}-{}-{}-{}-{}-{}-{}-{}-{}-{}-{}-{}-{}-{}-{}-{}-{}-{}-{}-{}-{}-{}-{}-{}-{}-{}-{}-{}-{}-{}-{}-{}-{}-{}-{}-{}-{} adoption\newline \hspace*{1em}year from PI.dta} & Divider line (comment).\\
20 & \texttt{use \textquotedbl{}\$DATA/PI.dta\textquotedbl{}, clear} & Loads the authors' policy file PI.dta.\\
21 & \texttt{keep id YEAR PI} & Keeps county, year and PSI.\\
22 & \texttt{bysort id: egen first\_pos = min(cond(PI > 0, YEAR, .))} & Adoption year = the earliest year in which PSI > 0 (missing if never).\\
23 & \texttt{bysort id (YEAR): gen byte tag = (\_n == 1)} & Marks the first year row of each county.\\
24 & \texttt{* how often does PSI return to 0 after adoption?\newline \hspace*{1em}(absorbing-{}treatment check)} & Comment (not executed): explains the code that follows.\\
25 & \texttt{count if YEAR > first\_pos \& !missing(first\_pos)} & Counts post-adoption county-years ...\\
26 & \texttt{local npost = r(N)} & ... and stores the number.\\
27 & \texttt{count if YEAR > first\_pos \& !missing(first\_pos) \& PI == 0} & Counts post-adoption county-years with PSI back at 0 ...\\
28 & \texttt{di as res \textquotedbl{}Share of post-{}adoption county-{}years with PSI =\newline \hspace*{1em}0: \textquotedbl{} \%5.3f r(N)/\textasciigrave{}npost\textquotesingle{}} & ... and prints the share (15.3\%): treatment is not truly absorbing.\\
29 & \texttt{keep if tag} & One row per county.\\
30 & \texttt{keep id first\_pos} & Keeps county and adoption year.\\
31 & \texttt{tempfile adopt} & Temporary file.\\
32 & \texttt{save \textasciigrave{}adopt\textquotesingle{}} & Saves adoption years.\\
34 & \texttt{*-{}-{}-{}-{}-{}-{}-{}-{}-{}-{}-{}-{}-{}-{}-{}-{}-{}-{}-{}-{}-{}-{}-{}-{}-{}-{}-{}-{}-{}-{}-{}-{}-{}-{}-{}-{}-{}-{}-{}-{}-{}-{}-{}-{}-{}-{}-{}-{} attach\newline \hspace*{1em}to analysis sample} & Divider line (comment).\\
35 & \texttt{use \textquotedbl{}\$SAMPLE\textquotedbl{}, clear} & Loads the analysis dataset built in step 1\_2.\\
36 & \texttt{merge m:1 id using \textasciigrave{}adopt\textquotesingle{}, keep(master match) nogen} & Attaches each county's adoption year to the analysis sample.\\
38 & \texttt{* g  : adoption year (csdid convention: 0 = never treated)} & Comment (not executed): explains the code that follows.\\
39 & \texttt{gen g\_main  = cond(in\_PI\_file==1, cond(missing(first\_pos),\newline \hspace*{1em}0, first\_pos), .)} & Cohort for the MAIN sample: adoption year, 0 = never treated; missing for counties absent from PI.dta (excluded).\\
40 & \texttt{gen g\_paper = cond(in\_PI\_file==1, g\_main, 0)} & Cohort for the PAPER-STYLE sample: missing-PSI counties are set to 0 (never treated).\\
41 & \texttt{label var g\_main  \textquotedbl{}Adoption year; 0 = never (PSI-{}data\newline \hspace*{1em}counties only)\textquotedbl{}} & Label.\\
42 & \texttt{label var g\_paper \textquotedbl{}Adoption year; 0 = never (missing-{}PSI\newline \hspace*{1em}counties = never)\textquotedbl{}} & Label.\\
44 & \texttt{* drop counties treated in the first year (no pre-{}period)} & Comment (not executed): explains the code that follows.\\
45 & \texttt{gen byte always\_treated = (first\_pos == 2014)} & 1 if already treated in 2014 (no pre-period).\\
46 & \texttt{egen tagc = tag(county)} & Marks one row per county.\\
47 & \texttt{count if tagc \& always\_treated} & Counts them (838 counties) ...\\
48 & \texttt{di as res \textquotedbl{}Counties already treated in 2014 (dropped): \textquotedbl{}\newline \hspace*{1em}r(N)} & ... and prints the number.\\
49 & \texttt{drop if always\_treated} & Drops them.\\
51 & \texttt{* relative time and indicators} & Comment (not executed): explains the code that follows.\\
52 & \texttt{gen rel\_year = YEAR -{} first\_pos if !missing(first\_pos)} & Event time: years since adoption (negative = before).\\
53 & \texttt{gen byte post = (YEAR >= first\_pos) if !missing(first\_pos)} & post = 1 from the adoption year on.\\
54 & \texttt{replace post = 0 if missing(first\_pos)} & Never-treated counties: post = 0.\\
55 & \texttt{gen byte never\_main  = (g\_main  == 0) if !missing(g\_main)} & Never-treated flag (main sample).\\
56 & \texttt{gen byte never\_paper = (g\_paper == 0)} & Never-treated flag (paper-style sample).\\
58 & \texttt{* cohort sizes (counties)} & Comment (not executed): explains the code that follows.\\
59 & \texttt{preserve} & Saves a copy of the data in memory, so it can be restored after temporary changes.\\
60 & \texttt{\hspace*{2.0em}keep if tagc} & Keeps only rows where tagc.\\
61 & \texttt{\hspace*{2.0em}gen sample\_main = !missing(g\_main)} & 1 if the county is in the main sample.\\
62 & \texttt{\hspace*{2.0em}collapse (count) counties\_paper = id (sum) counties\_main =\newline \hspace*{1em}sample\_main, by(g\_paper)} & Counts counties per cohort in each sample.\\
63 & \texttt{\hspace*{2.0em}rename g\_paper adoption\_year} & Renames for the output table.\\
64 & \texttt{\hspace*{2.0em}list, noobs sep(0)} & Lists the rows/variables on screen.\\
65 & \texttt{\hspace*{2.0em}export delimited using \textquotedbl{}\$T3/tables/3\_1\_cohort\_sizes.csv\textquotedbl{},\newline \hspace*{1em}replace} & Exports the current data to a CSV file.\\
66 & \texttt{restore} & Restores the data saved by preserve (undoes changes made since).\\
67 & \texttt{tab g\_main if tagc, missing} & Cohort sizes on screen.\\
69 & \texttt{* deseasonalised Shannon (removes month-{}of-{}year\newline \hspace*{1em}composition before} & Comment (not executed): explains the code that follows.\\
70 & \texttt{* aggregating months to county-{}years)} & Comment (not executed): explains the code that follows.\\
71 & \texttt{reg ShannonBD i.month\_of\_year} & Regresses Shannon on month-of-year dummies (seasonality).\\
72 & \texttt{predict double sh\_res, residuals} & Shannon with seasonal pattern removed (residuals).\\
73 & \texttt{summarize ShannonBD, meanonly} & Overall mean of Shannon.\\
74 & \texttt{gen double shannon\_sa = sh\_res + r(mean)} & Deseasonalised Shannon = residual + overall mean (keeps the original level).\\
75 & \texttt{drop sh\_res} & Deletes these variables.\\
76 & \texttt{label var shannon\_sa \textquotedbl{}Shannon, month-{}of-{}year effects\newline \hspace*{1em}removed\textquotedbl{}} & Label.\\
78 & \texttt{compress} & Stores variables in the smallest storage type to save memory.\\
79 & \texttt{save \textquotedbl{}\$T3/data/staggered\_monthly.dta\textquotedbl{}, replace} & Saves the monthly staggered dataset.\\
81 & \texttt{*-{}-{}-{}-{}-{}-{}-{}-{}-{}-{}-{}-{}-{}-{}-{}-{}-{}-{}-{}-{}-{}-{}-{}-{}-{}-{}-{}-{}-{}-{}-{}-{}-{}-{}-{}-{}-{}-{}-{}-{}-{}-{}-{}-{}-{}-{}-{}-{}\newline \hspace*{1em}county-{}year panel (for CS)} & Divider line (comment).\\
82 & \texttt{collapse (mean) shannon\_sa ShannonBD Duration (count)\newline \hspace*{1em}n\_months = ShannonBD ///} & Collapses months to county-years: means of Shannon and Duration, and the number of months ...\\
83 & \texttt{\hspace*{2.0em}(first) g\_main g\_paper first\_pos urban prov in\_PI\_file,\newline \hspace*{1em}by(county id YEAR)} & ... keeping cohort and flags, one row per county-year.\\
84 & \texttt{xtset county YEAR} & Declares the county-year panel.\\
85 & \texttt{di as res \textquotedbl{}County-{}years: \textquotedbl{} \_N} & Prints the number of county-years (7,984).\\
86 & \texttt{compress} & Stores variables in the smallest storage type to save memory.\\
87 & \texttt{save \textquotedbl{}\$T3/data/staggered\_countyyear.dta\textquotedbl{}, replace} & Saves the county-year dataset (for Callaway-Sant'Anna).\\}
\subsection{3\_2\_twfe\_event\_study.do}
Conventional TWFE benchmark: static (post dummy) and dynamic (leads and lags) event study.

\codetable{1 & \texttt{*=========================================================\newline \hspace*{1em}=================} & Divider line (comment).\\
2 & \texttt{* 3\_2\_twfe\_event\_study.do} & Comment (not executed): explains the code that follows.\\
3 & \texttt{* PURPOSE : Conventional (\textquotedbl{}naive\textquotedbl{}) TWFE estimates on the\newline \hspace*{1em}binary adoption} & Comment (not executed): explains the code that follows.\\
4 & \texttt{*           design, as the benchmark that the\newline \hspace*{1em}heterogeneity-{}robust} & Comment (not executed): explains the code that follows.\\
5 & \texttt{*           estimators (3\_3-{}3\_5) are compared against.} & Comment (not executed): explains the code that follows.\\
6 & \texttt{*   (a) static TWFE:   Shannon on post-{}adoption dummy} & Comment (not executed): explains the code that follows.\\
7 & \texttt{*   (b) dynamic TWFE:  event-{}time dummies, -{}5..+5,\newline \hspace*{1em}reference year -{}1,} & Comment (not executed): explains the code that follows.\\
8 & \texttt{*                      endpoints binned} & Comment (not executed): explains the code that follows.\\
9 & \texttt{* WHY     : Under staggered timing and heterogeneous\newline \hspace*{1em}effects, TWFE uses} & Comment (not executed): explains the code that follows.\\
10 & \texttt{*           already-{}treated counties as controls\newline \hspace*{1em}(\textquotedbl{}forbidden comparisons\textquotedbl{};} & Comment (not executed): explains the code that follows.\\
11 & \texttt{*           Goodman-{}Bacon 2021; Baker et al. 2025), so it\newline \hspace*{1em}can be biased.} & Comment (not executed): explains the code that follows.\\
12 & \texttt{* OUTPUT  : tables/3\_2\_twfe.csv,\newline \hspace*{1em}figures/3\_2\_event\_study\_twfe.pdf} & Comment (not executed): explains the code that follows.\\
13 & \texttt{*=========================================================\newline \hspace*{1em}=================} & Divider line (comment).\\
14 & \texttt{use \textquotedbl{}\$T3/data/staggered\_monthly.dta\textquotedbl{}, clear} & Loads the monthly staggered dataset.\\
15 & \texttt{keep if !missing(g\_main)                       // main\newline \hspace*{1em}sample} & Main sample only (drops the 724 missing-PSI counties).\\
17 & \texttt{*-{}-{}-{}-{}-{}-{}-{}-{}-{}-{}-{}-{}-{}-{}-{}-{}-{}-{}-{}-{}-{}-{}-{}-{}-{}-{}-{}-{}-{}-{}-{}-{}-{}-{}-{}-{}-{}-{}-{}-{}-{}-{}-{}-{}-{}-{}-{}-{} (a)\newline \hspace*{1em}static} & Divider line (comment).\\
18 & \texttt{reghdfe ShannonBD post \$control, absorb(county year)\newline \hspace*{1em}cluster(county)} & Static TWFE: Shannon on the post-adoption dummy, 9 controls, county + year-month FE (result +0.032).\\
19 & \texttt{scalar static\_b = \_b[post]} & Stores the static estimate.\\
20 & \texttt{scalar static\_se = \_se[post]} & Stores its SE.\\
21 & \texttt{scalar static\_N = e(N)} & Stores N.\\
23 & \texttt{*-{}-{}-{}-{}-{}-{}-{}-{}-{}-{}-{}-{}-{}-{}-{}-{}-{}-{}-{}-{}-{}-{}-{}-{}-{}-{}-{}-{}-{}-{}-{}-{}-{}-{}-{}-{}-{}-{}-{}-{}-{}-{}-{}-{}-{}-{}-{}-{} (b)\newline \hspace*{1em}dynamic} & Divider line (comment).\\
24 & \texttt{gen ev = rel\_year} & Event time copy.\\
25 & \texttt{replace ev = -{}5 if ev < -{}5 \& !missing(ev)} & Bins all years before -5 into -5.\\
26 & \texttt{replace ev =  5 if ev >  5 \& !missing(ev)} & Bins all years after +5 into +5.\\
27 & \texttt{forvalues k = 5(-{}1)2 \{} & For k = 5, 4, 3, 2 ...\\
28 & \texttt{\hspace*{2.0em}gen byte lead\textasciigrave{}k\textquotesingle{} = (ev == -{}\textasciigrave{}k\textquotesingle{})} & ... lead dummy: k years before adoption.\\
29 & \texttt{\}} & Closes the loop / block.\\
30 & \texttt{forvalues k = 0/5 \{} & For k = 0 ... 5 ...\\
31 & \texttt{\hspace*{2.0em}gen byte lag\textasciigrave{}k\textquotesingle{} = (ev == \textasciigrave{}k\textquotesingle{})} & ... lag dummy: k years after adoption.\\
32 & \texttt{\}} & Closes the loop / block.\\
33 & \texttt{reghdfe ShannonBD lead5 lead4 lead3 lead2 lag0-{}lag5\newline \hspace*{1em}\$control, absorb(county year) cluster(county)} & Dynamic TWFE event study: leads and lags (year -1 omitted as reference), controls, county + year-month FE. No urban trend.\\
34 & \texttt{test lead5 lead4 lead3 lead2} & Joint pre-trend test: are all leads zero?\\
35 & \texttt{scalar pre\_p = r(p)} & Stores the p-value (0.341).\\
36 & \texttt{di as res \textquotedbl{}Joint test of pre-{}trends (leads = 0): p = \textquotedbl{}\newline \hspace*{1em}\%5.3f pre\_p} & Prints it.\\
38 & \texttt{tempname P} & Creates a temporary name (handle) for a results file or matrix.\\
39 & \texttt{postfile \textasciigrave{}P\textquotesingle{} str12 estimator int event\_time double(b se)\newline \hspace*{1em}using \textquotedbl{}\$T3/tables/es\_twfe.dta\textquotedbl{}, replace} & Opens the results file for event-study coefficients.\\
40 & \texttt{forvalues k = 5(-{}1)2 \{} & For k = 5, 4, 3, 2 ...\\
41 & \texttt{\hspace*{2.0em}post \textasciigrave{}P\textquotesingle{} (\textquotedbl{}TWFE\textquotedbl{}) (-{}\textasciigrave{}k\textquotesingle{}) (\_b[lead\textasciigrave{}k\textquotesingle{}]) (\_se[lead\textasciigrave{}k\textquotesingle{}])} & Writes each lead coefficient.\\
42 & \texttt{\}} & Closes the loop / block.\\
43 & \texttt{post \textasciigrave{}P\textquotesingle{} (\textquotedbl{}TWFE\textquotedbl{}) (-{}1) (0) (0)} & Writes the reference year -1 as 0.\\
44 & \texttt{forvalues k = 0/5 \{} & For k = 0 ... 5 ...\\
45 & \texttt{\hspace*{2.0em}post \textasciigrave{}P\textquotesingle{} (\textquotedbl{}TWFE\textquotedbl{}) (\textasciigrave{}k\textquotesingle{}) (\_b[lag\textasciigrave{}k\textquotesingle{}]) (\_se[lag\textasciigrave{}k\textquotesingle{}])} & Writes each lag coefficient.\\
46 & \texttt{\}} & Closes the loop / block.\\
47 & \texttt{post \textasciigrave{}P\textquotesingle{} (\textquotedbl{}TWFE\_static\textquotedbl{}) (99) (static\_b) (static\_se)} & Writes the static estimate (event\_time 99 = overall).\\
48 & \texttt{postclose \textasciigrave{}P\textquotesingle{}} & Closes the results file so it can be loaded.\\
50 & \texttt{use \textquotedbl{}\$T3/tables/es\_twfe.dta\textquotedbl{}, clear} & Loads the coefficients.\\
51 & \texttt{gen lo = b -{} 1.96*se} & Lower 95\% confidence bound.\\
52 & \texttt{gen hi = b + 1.96*se} & Upper 95\% confidence bound.\\
53 & \texttt{export delimited using \textquotedbl{}\$T3/tables/3\_2\_twfe.csv\textquotedbl{}, replace} & Exports the current data to a CSV file.\\
54 & \texttt{twoway (rcap hi lo event\_time if estimator==\textquotedbl{}TWFE\textquotedbl{},\newline \hspace*{1em}lcolor(navy)) ///} & Plots confidence bars ...\\
55 & \texttt{\hspace*{3.5em}(scatter b event\_time if estimator==\textquotedbl{}TWFE\textquotedbl{}, mcolor(navy)),\newline \hspace*{1em}///} & ... and point estimates.\\
56 & \texttt{\hspace*{2.0em}yline(0, lpattern(dash)) xline(-{}0.5, lpattern(dot)) ///} & Continuation of the graph command: axis titles, reference lines, labels and legend.\\
57 & \texttt{\hspace*{2.0em}xtitle(\textquotedbl{}Years relative to first solar policy\textquotedbl{})\newline \hspace*{1em}ytitle(\textquotedbl{}Effect on Shannon diversity\textquotedbl{}) ///} & Continuation of the previous command: table title.\\
58 & \texttt{\hspace*{2.0em}title(\textquotedbl{}Dynamic TWFE event study (PSI-{}data counties)\textquotedbl{},\newline \hspace*{1em}size(medium)) ///} & Continuation of the previous command: table title.\\
59 & \texttt{\hspace*{2.0em}subtitle(\textquotedbl{}Reference year = -{}1; 95\% CI; county-{}clustered\newline \hspace*{1em}SE\textquotedbl{}, size(small)) ///} & Continuation of the previous command: table title.\\
60 & \texttt{\hspace*{2.0em}xlabel(-{}5(1)5) legend(off)} & Continuation of the graph command: axis titles, reference lines, labels and legend.\\
61 & \texttt{graph export \textquotedbl{}\$T3/figures/3\_2\_event\_study\_twfe.pdf\textquotedbl{},\newline \hspace*{1em}replace} & Saves the current graph as a PDF.\\}
\subsection{3\_3\_callaway\_santanna.do}
Callaway--Sant'Anna estimator with not-yet-treated controls for three specifications, with a helper program to aggregate results.

\codetable{1 & \texttt{*=========================================================\newline \hspace*{1em}=================} & Divider line (comment).\\
2 & \texttt{* 3\_3\_callaway\_santanna.do} & Comment (not executed): explains the code that follows.\\
3 & \texttt{* PURPOSE : Callaway \& Sant\textquotesingle{}Anna (2021) group-{}time\newline \hspace*{1em}ATT(g,t), aggregated to} & Comment (not executed): explains the code that follows.\\
4 & \texttt{*           an overall ATT and an event study, with\newline \hspace*{1em}NOT-{}YET-{}TREATED counties} & Comment (not executed): explains the code that follows.\\
5 & \texttt{*           as controls (never-{}treated included).} & Comment (not executed): explains the code that follows.\\
6 & \texttt{* DATA    : county-{}year panel (monthly Shannon averaged\newline \hspace*{1em}within year after} & Comment (not executed): explains the code that follows.\\
7 & \texttt{*           removing month-{}of-{}year effects). The panel is\newline \hspace*{1em}UNBALANCED: by default} & Comment (not executed): explains the code that follows.\\
8 & \texttt{*           csdid would balance it by dropping every\newline \hspace*{1em}county not observed in all} & Comment (not executed): explains the code that follows.\\
9 & \texttt{*           10 years (903 of 939 counties). bal(pair)\newline \hspace*{1em}instead balances each 2x2} & Comment (not executed): explains the code that follows.\\
10 & \texttt{*           comparison separately, so every county\newline \hspace*{1em}observed in both years of a} & Comment (not executed): explains the code that follows.\\
11 & \texttt{*           comparison is used.} & Comment (not executed): explains the code that follows.\\
12 & \texttt{* SPECS   : M1 main sample, unconditional parallel trends} & Comment (not executed): explains the code that follows.\\
13 & \texttt{*           M2 main sample, parallel trends conditional on\newline \hspace*{1em}urban} & Comment (not executed): explains the code that follows.\\
14 & \texttt{*           P1 paper-{}style sample (missing-{}PSI counties =\newline \hspace*{1em}never treated)} & Comment (not executed): explains the code that follows.\\
15 & \texttt{* OUTPUT  : tables/3\_3\_cs\_att.csv,\newline \hspace*{1em}tables/3\_3\_cs\_event.csv,} & Comment (not executed): explains the code that follows.\\
16 & \texttt{*           figures/3\_3\_event\_study\_cs\_*.pdf} & Comment (not executed): explains the code that follows.\\
17 & \texttt{*=========================================================\newline \hspace*{1em}=================} & Divider line (comment).\\
18 & \texttt{*-{}-{}-{}-{}-{}-{}-{}-{}-{}-{}-{}-{}-{}-{}-{}-{}-{}-{}-{}-{}-{}-{}-{}-{}-{}-{}-{}-{}-{}-{}-{}-{}-{}-{}-{}-{}-{}-{}-{}-{}-{}-{}-{}-{}-{}-{}-{}-{}-{}-{}-{}-{}-{}-{}-{}-{}-{}\newline \hspace*{1em}-{}-{}-{}-{}-{}-{}-{}-{}-{}-{}-{}-{}-{}-{}-{}-{}-{}} & Divider line (comment).\\
19 & \texttt{* Helper: aggregate csdid\textquotesingle{}s ATT(g,t) cells. Cells that\newline \hspace*{1em}cannot be estimated} & Comment (not executed): explains the code that follows.\\
20 & \texttt{* (tiny cohorts) must be excluded with -{}dropmissing-{}; the\newline \hspace*{1em}option syntax differs} & Comment (not executed): explains the code that follows.\\
21 & \texttt{* across csdid versions, so each variant is tried and the\newline \hspace*{1em}one used is logged.} & Comment (not executed): explains the code that follows.\\
22 & \texttt{* Returns row vectors \_B (estimates) and \_SE (standard\newline \hspace*{1em}errors).} & Comment (not executed): explains the code that follows.\\
23 & \texttt{*-{}-{}-{}-{}-{}-{}-{}-{}-{}-{}-{}-{}-{}-{}-{}-{}-{}-{}-{}-{}-{}-{}-{}-{}-{}-{}-{}-{}-{}-{}-{}-{}-{}-{}-{}-{}-{}-{}-{}-{}-{}-{}-{}-{}-{}-{}-{}-{}-{}-{}-{}-{}-{}-{}-{}-{}-{}\newline \hspace*{1em}-{}-{}-{}-{}-{}-{}-{}-{}-{}-{}-{}-{}-{}-{}-{}-{}-{}} & Divider line (comment).\\
24 & \texttt{capture program drop cs\_agg} & Removes an old version of the helper program.\\
25 & \texttt{program define cs\_agg} & Defines helper cs\_agg: aggregates Callaway-Sant'Anna ATT(g,t) results robustly across csdid versions.\\
26 & \texttt{\hspace*{2.0em}syntax anything(name=type) [, Window(string)]} & Inputs: aggregation type (simple, event, group) and an optional event window.\\
27 & \texttt{\hspace*{2.0em}local w \textquotedbl{}\textquotedbl{}} & Empty window option by default.\\
28 & \texttt{\hspace*{2.0em}if \textquotedbl{}\textasciigrave{}window\textquotesingle{}\textquotedbl{} != \textquotedbl{}\textquotedbl{} local w \textquotedbl{}window(\textasciigrave{}window\textquotesingle{})\textquotedbl{}} & Builds the window option if given.\\
29 & \texttt{\hspace*{2.0em}local used \textquotedbl{}\textquotedbl{}} & Records which option set worked.\\
30 & \texttt{\hspace*{2.0em}foreach opt in \textquotedbl{}post dropmissing\textquotedbl{} \textquotedbl{}dropmissing post\textquotedbl{}\newline \hspace*{1em}\textquotedbl{}dropmissing\textquotedbl{} \textquotedbl{}post\textquotedbl{} \{} & Tries several option orders because csdid versions differ.\\
31 & \texttt{\hspace*{4.0em}capture noisily estat \textasciigrave{}type\textquotesingle{}, \textasciigrave{}w\textquotesingle{} \textasciigrave{}opt\textquotesingle{}} & Runs the aggregation (estat simple/event/group); dropmissing skips cells that could not be estimated.\\
32 & \texttt{\hspace*{4.0em}if \_rc == 0 \{} & If it worked ...\\
33 & \texttt{\hspace*{4.0em}local used \textquotedbl{}\textasciigrave{}opt\textquotesingle{}\textquotedbl{}} & ... remember the options ...\\
34 & \texttt{\hspace*{4.0em}continue, break} & ... and stop trying.\\
35 & \texttt{\hspace*{4.0em}\}} & Closes the loop / block.\\
36 & \texttt{\hspace*{2.0em}\}} & Closes the loop / block.\\
37 & \texttt{\hspace*{2.0em}if \textquotedbl{}\textasciigrave{}used\textquotesingle{}\textquotedbl{} == \textquotedbl{}\textquotedbl{} \{} & If none worked ...\\
38 & \texttt{\hspace*{4.0em}di as err \textquotedbl{}estat \textasciigrave{}type\textquotesingle{} failed with every option set\textquotedbl{}} & ... print an error ...\\
39 & \texttt{\hspace*{4.0em}exit 498} & ... and stop with error 498.\\
40 & \texttt{\hspace*{2.0em}\}} & Closes the loop / block.\\
41 & \texttt{\hspace*{2.0em}di as res \textquotedbl{}estat \textasciigrave{}type\textquotesingle{} ran with options: \textasciigrave{}used\textquotesingle{}\textquotedbl{}} & Logs which options were used.\\
42 & \texttt{\hspace*{2.0em}if strpos(\textquotedbl{}\textasciigrave{}used\textquotesingle{}\textquotedbl{}, \textquotedbl{}post\textquotedbl{}) \{} & If results were posted to e() ...\\
43 & \texttt{\hspace*{4.0em}matrix \_B = e(b)} & ... take estimates from e(b) ...\\
44 & \texttt{\hspace*{4.0em}matrix \_V = e(V)} & ... and the variance matrix from e(V).\\
45 & \texttt{\hspace*{2.0em}\}} & Closes the loop / block.\\
46 & \texttt{\hspace*{2.0em}else \{} & Otherwise ...\\
47 & \texttt{\hspace*{4.0em}capture confirm matrix r(b)} & Otherwise check whether r(b) exists.\\
48 & \texttt{\hspace*{4.0em}if !\_rc \{} & If it does ...\\
49 & \texttt{\hspace*{4.0em}matrix \_B = r(b)} & ... take estimates from r(b) ...\\
50 & \texttt{\hspace*{4.0em}matrix \_V = r(V)} & ... and r(V).\\
51 & \texttt{\hspace*{4.0em}\}} & Closes the loop / block.\\
52 & \texttt{\hspace*{4.0em}else \{} & Otherwise ...\\
53 & \texttt{\hspace*{4.0em}matrix \_T = r(table)} & Otherwise use the r(table) matrix ...\\
54 & \texttt{\hspace*{4.0em}matrix \_B = \_T[1, 1...]} & ... row 1 = estimates ...\\
55 & \texttt{\hspace*{4.0em}matrix \_V = diag(\_T[2, 1...])} & ... row 2 = standard errors on a diagonal ...\\
56 & \texttt{\hspace*{4.0em}forvalues i = 1/\textasciigrave{}=colsof(\_V)\textquotesingle{} \{} & ... loop over them ...\\
57 & \texttt{\hspace*{4.0em}matrix \_V[\textasciigrave{}i\textquotesingle{},\textasciigrave{}i\textquotesingle{}] = \_V[\textasciigrave{}i\textquotesingle{},\textasciigrave{}i\textquotesingle{}]\^{}2} & ... squared to variances.\\
58 & \texttt{\hspace*{4.0em}\}} & Closes the loop / block.\\
59 & \texttt{\hspace*{4.0em}\}} & Closes the loop / block.\\
60 & \texttt{\hspace*{2.0em}\}} & Closes the loop / block.\\
61 & \texttt{\hspace*{2.0em}matrix \_SE = \_B} & Matrix for standard errors, same size as \_B.\\
62 & \texttt{\hspace*{2.0em}forvalues i = 1/\textasciigrave{}=colsof(\_B)\textquotesingle{} \{} & Loop over estimates ...\\
63 & \texttt{\hspace*{4.0em}matrix \_SE[1,\textasciigrave{}i\textquotesingle{}] = sqrt(\_V[\textasciigrave{}i\textquotesingle{},\textasciigrave{}i\textquotesingle{}])} & ... SE = square root of the variance.\\
64 & \texttt{\hspace*{2.0em}\}} & Closes the loop / block.\\
65 & \texttt{\hspace*{2.0em}matrix list \_B} & Shows the estimates.\\
66 & \texttt{\hspace*{2.0em}matrix list \_SE} & Shows the SEs.\\
67 & \texttt{end} & Ends the program definition.\\
69 & \texttt{tempname A} & Handle for the overall-ATT results file.\\
70 & \texttt{postfile \textasciigrave{}A\textquotesingle{} str40 spec double(att se N\_cy) using\newline \hspace*{1em}\textquotedbl{}\$T3/tables/cs\_att.dta\textquotedbl{}, replace} & Opens the overall-ATT results file.\\
71 & \texttt{tempname E} & Handle for the event-study results file.\\
72 & \texttt{postfile \textasciigrave{}E\textquotesingle{} str40 spec int event\_time double(b se) using\newline \hspace*{1em}\textquotedbl{}\$T3/tables/cs\_event.dta\textquotedbl{}, replace} & Opens the event-study results file.\\
74 & \texttt{local specs \textquotedbl{}M1 M2 P1\textquotedbl{}} & Three specifications: M1 main unconditional, M2 main conditional on urban, P1 paper-style.\\
75 & \texttt{local lab\_M1 \textquotedbl{}Main sample, unconditional\textquotedbl{}} & Label M1.\\
76 & \texttt{local lab\_M2 \textquotedbl{}Main sample, conditional on urban\textquotedbl{}} & Label M2.\\
77 & \texttt{local lab\_P1 \textquotedbl{}Paper-{}style sample (missing PSI = never)\textquotedbl{}} & Label P1.\\
79 & \texttt{foreach s of local specs \{} & Loops over the three specifications.\\
80 & \texttt{\hspace*{2.0em}use \textquotedbl{}\$T3/data/staggered\_countyyear.dta\textquotedbl{}, clear} & Loads the county-year panel.\\
81 & \texttt{\hspace*{2.0em}if \textquotedbl{}\textasciigrave{}s\textquotesingle{}\textquotedbl{} == \textquotedbl{}P1\textquotedbl{} local gv g\_paper} & P1 uses the paper-style cohort variable ...\\
82 & \texttt{\hspace*{2.0em}else \{} & Otherwise ...\\
83 & \texttt{\hspace*{4.0em}local gv g\_main} & ... otherwise the main cohort variable ...\\
84 & \texttt{\hspace*{4.0em}keep if !missing(g\_main)} & Main sample only (drops the 724 missing-PSI counties).\\
85 & \texttt{\hspace*{2.0em}\}} & Closes the loop / block.\\
86 & \texttt{\hspace*{2.0em}local x \textquotedbl{}\textquotedbl{}} & No covariates by default.\\
87 & \texttt{\hspace*{2.0em}if \textquotedbl{}\textasciigrave{}s\textquotesingle{}\textquotedbl{} == \textquotedbl{}M2\textquotedbl{} local x \textquotedbl{}urban\textquotedbl{}} & M2: parallel trends conditional on urban.\\
88 & \texttt{\hspace*{2.0em}di as res \_n \textquotedbl{}================ CS: \textasciigrave{}lab\_\textasciigrave{}s\textquotesingle{}\textquotesingle{}\newline \hspace*{1em}================\textquotedbl{}} & Prints the specification.\\
89 & \texttt{\hspace*{2.0em}* how unbalanced is the panel? (counties observed in all\newline \hspace*{1em}10 years)} & Comment (not executed): explains the code that follows.\\
90 & \texttt{\hspace*{2.0em}preserve} & Saves a copy of the data in memory, so it can be restored after temporary changes.\\
91 & \texttt{\hspace*{4.0em}bysort county: gen byte nyears = \_N} & Years observed per county.\\
92 & \texttt{\hspace*{4.0em}egen byte tagc = tag(county)} & One row per county.\\
93 & \texttt{\hspace*{4.0em}count if tagc} & Counts rows that satisfy the condition (stored in r(N)).\\
94 & \texttt{\hspace*{4.0em}local nc = r(N)} & Number of counties.\\
95 & \texttt{\hspace*{4.0em}count if tagc \& nyears < 10} & Counties not observed in all 10 years ...\\
96 & \texttt{\hspace*{4.0em}di as res \textquotedbl{}Counties: \textasciigrave{}nc\textquotesingle{}; not observed in all 10 years: \textquotedbl{}\newline \hspace*{1em}r(N)} & ... printed (903 of 939): why bal(pair) is needed.\\
97 & \texttt{\hspace*{2.0em}restore} & Restores the data saved by preserve (undoes changes made since).\\
99 & \texttt{\hspace*{2.0em}set seed 20261107} & Fixes the random seed so bootstrap SEs are reproducible.\\
100 & \texttt{\hspace*{2.0em}capture noisily csdid shannon\_sa \textasciigrave{}x\textquotesingle{}, ivar(county)\newline \hspace*{1em}time(YEAR) gvar(\textasciigrave{}gv\textquotesingle{}) notyet bal(pair) rseed(20261107)} & Callaway-Sant'Anna: outcome deseasonalised Shannon; panel id county, time YEAR, cohort gv; notyet = not-yet-treated counties as controls; bal(pair) = balance each 2x2 comparison separately (keeps unbalanced counties); rseed = bootstrap seed.\\
101 & \texttt{\hspace*{2.0em}if \_rc csdid shannon\_sa \textasciigrave{}x\textquotesingle{}, ivar(county) time(YEAR)\newline \hspace*{1em}gvar(\textasciigrave{}gv\textquotesingle{}) notyet bal(pair)} & Fallback without rseed for older csdid versions.\\
102 & \texttt{\hspace*{2.0em}* number of counties and county-{}years entering the\newline \hspace*{1em}estimation} & Comment (not executed): explains the code that follows.\\
103 & \texttt{\hspace*{2.0em}quietly count if e(sample)} & County-years used ...\\
104 & \texttt{\hspace*{2.0em}local Ncy = r(N)} & ... stored ...\\
105 & \texttt{\hspace*{2.0em}di as res \textquotedbl{}County-{}years used: \textasciigrave{}Ncy\textquotesingle{}\textquotedbl{}} & ... and printed (3,907 main).\\
106 & \texttt{\hspace*{2.0em}* pre-{}test of parallel trends (all pre-{}treatment ATT(g,t)\newline \hspace*{1em}= 0)} & Comment (not executed): explains the code that follows.\\
107 & \texttt{\hspace*{2.0em}capture noisily estat pretrend} & Tries the pre-trend test command (csdid prints the pre-test with the ATT table anyway; p < 0.001).\\
108 & \texttt{\hspace*{2.0em}estimates store cs\_full\_\textasciigrave{}s\textquotesingle{}} & Saves the full ATT(g,t) results.\\
110 & \texttt{\hspace*{2.0em}* overall ATT (simple aggregation)} & Comment (not executed): explains the code that follows.\\
111 & \texttt{\hspace*{2.0em}cs\_agg simple} & Overall ATT (simple average over all treated cells).\\
112 & \texttt{\hspace*{2.0em}post \textasciigrave{}A\textquotesingle{} (\textquotedbl{}\textasciigrave{}lab\_\textasciigrave{}s\textquotesingle{}\textquotesingle{}\textquotedbl{}) (\_B[1,1]) (\_SE[1,1]) (\textasciigrave{}Ncy\textquotesingle{})} & Writes the overall ATT row.\\
113 & \texttt{\hspace*{2.0em}di \textquotedbl{}Overall ATT = \textquotedbl{} \%8.4f \_B[1,1] \textquotedbl{}  (SE \textquotedbl{} \%6.4f \_SE[1,1]\newline \hspace*{1em}\textquotedbl{})\textquotedbl{}} & Prints it (+0.057 main).\\
115 & \texttt{\hspace*{2.0em}* event-{}study aggregation, window -{}5..+5} & Comment (not executed): explains the code that follows.\\
116 & \texttt{\hspace*{2.0em}estimates restore cs\_full\_\textasciigrave{}s\textquotesingle{}} & Re-loads the full results.\\
117 & \texttt{\hspace*{2.0em}cs\_agg event, window(-{}5 5)} & Event-study aggregation for years -5 to +5.\\
118 & \texttt{\hspace*{2.0em}local names : colnames \_B} & Names of the event-time estimates (Tm5 ... Tp5).\\
119 & \texttt{\hspace*{2.0em}local j = 0} & Counter.\\
120 & \texttt{\hspace*{2.0em}foreach nm of local names \{} & Loop over names ...\\
121 & \texttt{\hspace*{4.0em}local ++j} & ... counter + 1 ...\\
122 & \texttt{\hspace*{4.0em}if ustrregexm(\textquotedbl{}\textasciigrave{}nm\textquotesingle{}\textquotedbl{}, \textquotedbl{}\^{}[Tt]m([0-{}9]+)\$\textquotedbl{}) \{} & ... a "Tm\#" name is a pre-period (minus) ...\\
123 & \texttt{\hspace*{4.0em}post \textasciigrave{}E\textquotesingle{} (\textquotedbl{}\textasciigrave{}lab\_\textasciigrave{}s\textquotesingle{}\textquotesingle{}\textquotedbl{}) (-{}real(ustrregexs(1))) (\_B[1,\textasciigrave{}j\textquotesingle{}])\newline \hspace*{1em}(\_SE[1,\textasciigrave{}j\textquotesingle{}])} & ... write it with negative event time ...\\
124 & \texttt{\hspace*{4.0em}\}} & Closes the loop / block.\\
125 & \texttt{\hspace*{4.0em}else if ustrregexm(\textquotedbl{}\textasciigrave{}nm\textquotesingle{}\textquotedbl{}, \textquotedbl{}\^{}[Tt]p([0-{}9]+)\$\textquotedbl{}) \{} & ... a "Tp\#" name is a post-period (plus) ...\\
126 & \texttt{\hspace*{4.0em}post \textasciigrave{}E\textquotesingle{} (\textquotedbl{}\textasciigrave{}lab\_\textasciigrave{}s\textquotesingle{}\textquotesingle{}\textquotedbl{}) (real(ustrregexs(1))) (\_B[1,\textasciigrave{}j\textquotesingle{}])\newline \hspace*{1em}(\_SE[1,\textasciigrave{}j\textquotesingle{}])} & ... write it with positive event time ...\\
127 & \texttt{\hspace*{4.0em}\}} & Closes the loop / block.\\
128 & \texttt{\hspace*{4.0em}else di \textquotedbl{}\textasciigrave{}nm\textquotesingle{} = \textquotedbl{} \%8.4f \_B[1,\textasciigrave{}j\textquotesingle{}] \textquotedbl{}  (SE \textquotedbl{} \%6.4f\newline \hspace*{1em}\_SE[1,\textasciigrave{}j\textquotesingle{}] \textquotedbl{})\textquotedbl{}} & ... otherwise (pre/post averages) just print it.\\
129 & \texttt{\hspace*{2.0em}\}} & Closes the loop / block.\\
131 & \texttt{\hspace*{2.0em}* cohort-{}specific ATTs} & Comment (not executed): explains the code that follows.\\
132 & \texttt{\hspace*{2.0em}estimates restore cs\_full\_\textasciigrave{}s\textquotesingle{}} & Re-loads the full results.\\
133 & \texttt{\hspace*{2.0em}capture noisily cs\_agg group} & Cohort-specific ATTs (printed only).\\
134 & \texttt{\}} & Closes the loop / block.\\
135 & \texttt{postclose \textasciigrave{}A\textquotesingle{}} & Closes the ATT file.\\
136 & \texttt{postclose \textasciigrave{}E\textquotesingle{}} & Closes the event file.\\
138 & \texttt{use \textquotedbl{}\$T3/tables/cs\_att.dta\textquotedbl{}, clear} & Loads overall results.\\
139 & \texttt{list, noobs} & Lists the rows/variables on screen.\\
140 & \texttt{export delimited using \textquotedbl{}\$T3/tables/3\_3\_cs\_att.csv\textquotedbl{},\newline \hspace*{1em}replace} & Exports the current data to a CSV file.\\
142 & \texttt{use \textquotedbl{}\$T3/tables/cs\_event.dta\textquotedbl{}, clear} & Loads event-study results.\\
143 & \texttt{gen lo = b -{} 1.96*se} & Lower 95\% confidence bound.\\
144 & \texttt{gen hi = b + 1.96*se} & Upper 95\% confidence bound.\\
145 & \texttt{export delimited using \textquotedbl{}\$T3/tables/3\_3\_cs\_event.csv\textquotedbl{},\newline \hspace*{1em}replace} & Exports the current data to a CSV file.\\
146 & \texttt{local i = 0} & Counter for figure file names.\\
147 & \texttt{foreach sp in \textquotedbl{}Main sample, unconditional\textquotedbl{} \textquotedbl{}Main sample,\newline \hspace*{1em}conditional on urban\textquotedbl{} \textquotedbl{}Paper-{}style sample (missing PSI =\newline \hspace*{1em}never)\textquotedbl{} \{} & Draws one event-study graph per specification.\\
148 & \texttt{\hspace*{2.0em}local ++i} & Counter + 1.\\
149 & \texttt{\hspace*{2.0em}twoway (rcap hi lo event\_time if spec==\textquotedbl{}\textasciigrave{}sp\textquotesingle{}\textquotedbl{},\newline \hspace*{1em}lcolor(maroon)) ///} & Confidence bars ...\\
150 & \texttt{\hspace*{4.0em}(scatter b event\_time if spec==\textquotedbl{}\textasciigrave{}sp\textquotesingle{}\textquotedbl{}, mcolor(maroon)),\newline \hspace*{1em}///} & ... and point estimates.\\
151 & \texttt{\hspace*{4.0em}yline(0, lpattern(dash)) xline(-{}0.5, lpattern(dot)) ///} & Continuation of the graph command: axis titles, reference lines, labels and legend.\\
152 & \texttt{\hspace*{4.0em}xtitle(\textquotedbl{}Years relative to first solar policy\textquotedbl{}) ytitle(\textquotedbl{}ATT\newline \hspace*{1em}on Shannon diversity\textquotedbl{}) ///} & Continuation of the previous command: table title.\\
153 & \texttt{\hspace*{4.0em}title(\textquotedbl{}Callaway-{}Sant\textquotesingle{}Anna event study\textquotedbl{}, size(medium))\newline \hspace*{1em}subtitle(\textquotedbl{}\textasciigrave{}sp\textquotesingle{}; 95\% CI\textquotedbl{}, size(small)) ///} & Continuation of the previous command: table title.\\
154 & \texttt{\hspace*{4.0em}xlabel(-{}5(1)5) legend(off)} & Continuation of the graph command: axis titles, reference lines, labels and legend.\\
155 & \texttt{\hspace*{2.0em}graph export \textquotedbl{}\$T3/figures/3\_3\_event\_study\_cs\_\textasciigrave{}i\textquotesingle{}.pdf\textquotedbl{},\newline \hspace*{1em}replace} & Saves the current graph as a PDF.\\
156 & \texttt{\}} & Closes the loop / block.\\}
\subsection{3\_4\_imputation\_bjs.do}
Borusyak--Jaravel--Spiess imputation estimator on monthly data, main and paper-style samples.

\codetable{1 & \texttt{*=========================================================\newline \hspace*{1em}=================} & Divider line (comment).\\
2 & \texttt{* 3\_4\_imputation\_bjs.do} & Comment (not executed): explains the code that follows.\\
3 & \texttt{* PURPOSE : Borusyak, Jaravel \& Spiess (2024) imputation\newline \hspace*{1em}estimator on the} & Comment (not executed): explains the code that follows.\\
4 & \texttt{*           MONTHLY data: county and year-{}month fixed\newline \hspace*{1em}effects are estimated} & Comment (not executed): explains the code that follows.\\
5 & \texttt{*           on untreated observations only, untreated\newline \hspace*{1em}outcomes are imputed} & Comment (not executed): explains the code that follows.\\
6 & \texttt{*           for treated observations, and effects are\newline \hspace*{1em}averaged by event year.} & Comment (not executed): explains the code that follows.\\
7 & \texttt{* WHY     : Uses the authors\textquotesingle{} own monthly unit and\newline \hspace*{1em}fixed-{}effect structure,} & Comment (not executed): explains the code that follows.\\
8 & \texttt{*           handles the unbalanced panel, and avoids\newline \hspace*{1em}forbidden comparisons.} & Comment (not executed): explains the code that follows.\\
9 & \texttt{* SPECS   : main sample; paper-{}style sample} & Comment (not executed): explains the code that follows.\\
10 & \texttt{* OUTPUT  : tables/3\_4\_bjs.csv,\newline \hspace*{1em}figures/3\_4\_event\_study\_bjs\_*.pdf} & Comment (not executed): explains the code that follows.\\
11 & \texttt{*=========================================================\newline \hspace*{1em}=================} & Divider line (comment).\\
12 & \texttt{tempname B} & Handle for the BJS results file.\\
13 & \texttt{postfile \textasciigrave{}B\textquotesingle{} str40 spec int event\_time double(b se) using\newline \hspace*{1em}\textquotedbl{}\$T3/tables/bjs.dta\textquotedbl{}, replace} & Opens the BJS results file.\\
15 & \texttt{foreach s in main paper \{} & Runs the main and the paper-style sample.\\
16 & \texttt{\hspace*{2.0em}use \textquotedbl{}\$T3/data/staggered\_monthly.dta\textquotedbl{}, clear} & Loads the monthly staggered dataset.\\
17 & \texttt{\hspace*{2.0em}if \textquotedbl{}\textasciigrave{}s\textquotesingle{}\textquotedbl{} == \textquotedbl{}main\textquotedbl{} \{} & Main sample:\\
18 & \texttt{\hspace*{4.0em}keep if !missing(g\_main)} & Main sample only (drops the 724 missing-PSI counties).\\
19 & \texttt{\hspace*{4.0em}gen Ei = cond(g\_main == 0, ., g\_main)} & Ei = adoption year; missing = never treated (did\_imputation convention).\\
20 & \texttt{\hspace*{4.0em}local lab \textquotedbl{}Main sample\textquotedbl{}} & Label.\\
21 & \texttt{\hspace*{2.0em}\}} & Closes the loop / block.\\
22 & \texttt{\hspace*{2.0em}else \{} & Otherwise ...\\
23 & \texttt{\hspace*{4.0em}gen Ei = cond(g\_paper == 0, ., g\_paper)} & Paper-style: missing-PSI counties never treated.\\
24 & \texttt{\hspace*{4.0em}local lab \textquotedbl{}Paper-{}style sample (missing PSI = never)\textquotedbl{}} & Label.\\
25 & \texttt{\hspace*{2.0em}\}} & Closes the loop / block.\\
26 & \texttt{\hspace*{2.0em}di as res \_n \textquotedbl{}================ BJS: \textasciigrave{}lab\textquotesingle{}\newline \hspace*{1em}================\textquotedbl{}} & Prints the sample.\\
27 & \texttt{\hspace*{2.0em}* Imputation needs every county and every year-{}month to\newline \hspace*{1em}have at least one} & Comment (not executed): explains the code that follows.\\
28 & \texttt{\hspace*{2.0em}* UNTREATED observation (their fixed effects are estimated\newline \hspace*{1em}on untreated data} & Comment (not executed): explains the code that follows.\\
29 & \texttt{\hspace*{2.0em}* only). Drop the few that never do, iterating until\newline \hspace*{1em}stable.} & Comment (not executed): explains the code that follows.\\
30 & \texttt{\hspace*{2.0em}gen byte untreated = missing(Ei) | YEAR < Ei} & 1 for untreated observations (never treated, or before adoption).\\
31 & \texttt{\hspace*{2.0em}local dropped = 1} & Loop control.\\
32 & \texttt{\hspace*{2.0em}while \textasciigrave{}dropped\textquotesingle{} > 0 \{} & Repeat until nothing more is dropped:\\
33 & \texttt{\hspace*{4.0em}local n0 = \_N} & rows before ...\\
34 & \texttt{\hspace*{4.0em}bysort county: egen u\_c = max(untreated)} & ... does each county have at least one untreated observation?\\
35 & \texttt{\hspace*{4.0em}bysort year:   egen u\_t = max(untreated)} & ... does each year-month have one?\\
36 & \texttt{\hspace*{4.0em}keep if u\_c == 1 \& u\_t == 1} & Keeps only those (their fixed effects must be estimated from untreated data).\\
37 & \texttt{\hspace*{4.0em}drop u\_c u\_t} & Removes helpers.\\
38 & \texttt{\hspace*{4.0em}local dropped = \textasciigrave{}n0\textquotesingle{} -{} \_N} & Rows dropped in this round.\\
39 & \texttt{\hspace*{2.0em}\}} & Closes the loop / block.\\
41 & \texttt{\hspace*{2.0em}* overall ATT (all post-{}adoption observations, equally\newline \hspace*{1em}weighted)} & Comment (not executed): explains the code that follows.\\
42 & \texttt{\hspace*{2.0em}capture noisily did\_imputation ShannonBD county YEAR Ei,\newline \hspace*{1em}fe(county year) cluster(county) autosample} & BJS imputation: estimates county and year-month FE from untreated rows, imputes untreated Shannon for treated rows, averages the gaps; autosample drops rows that cannot be imputed.\\
43 & \texttt{\hspace*{2.0em}if \_rc \{} & If that fails ...\\
44 & \texttt{\hspace*{4.0em}di as err \textquotedbl{}year-{}month FE not imputable; falling back to\newline \hspace*{1em}county + year FE\textquotedbl{}} & ... prints why ...\\
45 & \texttt{\hspace*{4.0em}local fe \textquotedbl{}county YEAR\textquotedbl{}} & ... and switches to county + calendar-year FE.\\
46 & \texttt{\hspace*{4.0em}did\_imputation ShannonBD county YEAR Ei, fe(\textasciigrave{}fe\textquotesingle{})\newline \hspace*{1em}cluster(county) autosample} & Overall effect with the fallback FE.\\
47 & \texttt{\hspace*{2.0em}\}} & Closes the loop / block.\\
48 & \texttt{\hspace*{2.0em}else local fe \textquotedbl{}county year\textquotedbl{}} & Otherwise keeps year-month FE.\\
49 & \texttt{\hspace*{2.0em}post \textasciigrave{}B\textquotesingle{} (\textquotedbl{}\textasciigrave{}lab\textquotesingle{}\textquotedbl{}) (99) (\_b[tau]) (\_se[tau])} & Writes the overall effect (tau).\\
51 & \texttt{\hspace*{2.0em}* event study: horizons 0-{}5, pre-{}trend coefficients for 5\newline \hspace*{1em}years} & Comment (not executed): explains the code that follows.\\
52 & \texttt{\hspace*{2.0em}did\_imputation ShannonBD county YEAR Ei, fe(\textasciigrave{}fe\textquotesingle{})\newline \hspace*{1em}cluster(county) ///} & Event-study version ...\\
53 & \texttt{\hspace*{4.0em}horizons(0/5) pretrends(5) autosample} & ... effects 0-5 years after adoption and 5 pre-trend coefficients.\\
54 & \texttt{\hspace*{2.0em}di as res \textquotedbl{}Fixed effects used: \textasciigrave{}fe\textquotesingle{}\textquotedbl{}} & Logs which FE were used.\\
55 & \texttt{\hspace*{2.0em}estimates store bjs\_\textasciigrave{}s\textquotesingle{}} & Saves the results.\\
56 & \texttt{\hspace*{2.0em}forvalues k = 0/5 \{} & For k = 0 ... 5 ...\\
57 & \texttt{\hspace*{4.0em}post \textasciigrave{}B\textquotesingle{} (\textquotedbl{}\textasciigrave{}lab\textquotesingle{}\textquotedbl{}) (\textasciigrave{}k\textquotesingle{}) (\_b[tau\textasciigrave{}k\textquotesingle{}]) (\_se[tau\textasciigrave{}k\textquotesingle{}])} & Writes each post-period effect.\\
58 & \texttt{\hspace*{2.0em}\}} & Closes the loop / block.\\
59 & \texttt{\hspace*{2.0em}forvalues k = 1/5 \{} & For the 5 pre-periods ...\\
60 & \texttt{\hspace*{4.0em}post \textasciigrave{}B\textquotesingle{} (\textquotedbl{}\textasciigrave{}lab\textquotesingle{}\textquotedbl{}) (-{}\textasciigrave{}k\textquotesingle{}) (\_b[pre\textasciigrave{}k\textquotesingle{}]) (\_se[pre\textasciigrave{}k\textquotesingle{}])} & ... writes each pre-trend coefficient.\\
61 & \texttt{\hspace*{2.0em}\}} & Closes the loop / block.\\
62 & \texttt{\hspace*{2.0em}test pre1 pre2 pre3 pre4 pre5} & Joint pre-trend test.\\
63 & \texttt{\hspace*{2.0em}di as res \textquotedbl{}Joint pre-{}trend test p = \textquotedbl{} \%5.3f r(p)} & Prints p (0.239 main; 0.137 paper-style).\\
64 & \texttt{\}} & Closes the loop / block.\\
65 & \texttt{postclose \textasciigrave{}B\textquotesingle{}} & Closes the file.\\
67 & \texttt{use \textquotedbl{}\$T3/tables/bjs.dta\textquotedbl{}, clear} & Loads BJS results.\\
68 & \texttt{gen lo = b -{} 1.96*se} & Lower 95\% confidence bound.\\
69 & \texttt{gen hi = b + 1.96*se} & Upper 95\% confidence bound.\\
70 & \texttt{list, noobs sep(0)} & Lists the rows/variables on screen.\\
71 & \texttt{export delimited using \textquotedbl{}\$T3/tables/3\_4\_bjs.csv\textquotedbl{}, replace} & Exports the current data to a CSV file.\\
72 & \texttt{local i = 0} & Counter for figure file names.\\
73 & \texttt{foreach sp in \textquotedbl{}Main sample\textquotedbl{} \textquotedbl{}Paper-{}style sample (missing\newline \hspace*{1em}PSI = never)\textquotedbl{} \{} & One graph per sample.\\
74 & \texttt{\hspace*{2.0em}local ++i} & Counter + 1.\\
75 & \texttt{\hspace*{2.0em}twoway (rcap hi lo event\_time if spec==\textquotedbl{}\textasciigrave{}sp\textquotesingle{}\textquotedbl{} \& event\_time\newline \hspace*{1em}!= 99, lcolor(dkgreen)) ///} & Confidence bars ...\\
76 & \texttt{\hspace*{4.0em}(scatter b event\_time if spec==\textquotedbl{}\textasciigrave{}sp\textquotesingle{}\textquotedbl{} \& event\_time != 99,\newline \hspace*{1em}mcolor(dkgreen)), ///} & ... and estimates.\\
77 & \texttt{\hspace*{4.0em}yline(0, lpattern(dash)) xline(-{}0.5, lpattern(dot)) ///} & Continuation of the graph command: axis titles, reference lines, labels and legend.\\
78 & \texttt{\hspace*{4.0em}xtitle(\textquotedbl{}Years relative to first solar policy\textquotedbl{})\newline \hspace*{1em}ytitle(\textquotedbl{}Effect on Shannon diversity\textquotedbl{}) ///} & Continuation of the previous command: table title.\\
79 & \texttt{\hspace*{4.0em}title(\textquotedbl{}Borusyak-{}Jaravel-{}Spiess imputation event study\textquotedbl{},\newline \hspace*{1em}size(medium)) ///} & Continuation of the previous command: table title.\\
80 & \texttt{\hspace*{4.0em}subtitle(\textquotedbl{}\textasciigrave{}sp\textquotesingle{}; 95\% CI\textquotedbl{}, size(small)) xlabel(-{}5(1)5)\newline \hspace*{1em}legend(off)} & Continuation of the previous command: table title.\\
81 & \texttt{\hspace*{2.0em}graph export \textquotedbl{}\$T3/figures/3\_4\_event\_study\_bjs\_\textasciigrave{}i\textquotesingle{}.pdf\textquotedbl{},\newline \hspace*{1em}replace} & Saves the current graph as a PDF.\\
82 & \texttt{\}} & Closes the loop / block.\\}
\subsection{3\_5\_sun\_abraham.do}
Sun--Abraham interaction-weighted event study with never-treated counties as controls.

\codetable{1 & \texttt{*=========================================================\newline \hspace*{1em}=================} & Divider line (comment).\\
2 & \texttt{* 3\_5\_sun\_abraham.do} & Comment (not executed): explains the code that follows.\\
3 & \texttt{* PURPOSE : Sun \& Abraham (2021) interaction-{}weighted\newline \hspace*{1em}event study on the} & Comment (not executed): explains the code that follows.\\
4 & \texttt{*           monthly data, never-{}treated counties as the\newline \hspace*{1em}control cohort.} & Comment (not executed): explains the code that follows.\\
5 & \texttt{*           A third heterogeneity-{}robust check on the same\newline \hspace*{1em}question.} & Comment (not executed): explains the code that follows.\\
6 & \texttt{* OUTPUT  : tables/3\_5\_sun\_abraham.csv,\newline \hspace*{1em}figures/3\_5\_event\_study\_sa.pdf} & Comment (not executed): explains the code that follows.\\
7 & \texttt{*=========================================================\newline \hspace*{1em}=================} & Divider line (comment).\\
8 & \texttt{use \textquotedbl{}\$T3/data/staggered\_monthly.dta\textquotedbl{}, clear} & Loads the monthly staggered dataset.\\
9 & \texttt{keep if !missing(g\_main)} & Main sample only (drops the 724 missing-PSI counties).\\
11 & \texttt{* relative-{}time indicators for every period except the\newline \hspace*{1em}reference (-{}1)} & Comment (not executed): explains the code that follows.\\
12 & \texttt{summarize rel\_year} & Range of event time.\\
13 & \texttt{local mn = -{}r(min)} & Largest lead.\\
14 & \texttt{local mx = r(max)} & Largest lag.\\
15 & \texttt{local leads \textquotedbl{}\textquotedbl{}} & Empty list of lead dummies.\\
16 & \texttt{forvalues k = 2/\textasciigrave{}mn\textquotesingle{} \{} & For every lead k >= 2 (year -1 omitted as reference) ...\\
17 & \texttt{\hspace*{2.0em}gen byte F\textasciigrave{}k\textquotesingle{} = (rel\_year == -{}\textasciigrave{}k\textquotesingle{})} & ... dummy for k years before adoption ...\\
18 & \texttt{\hspace*{2.0em}local leads \textquotedbl{}\textasciigrave{}leads\textquotesingle{} F\textasciigrave{}k\textquotesingle{}\textquotedbl{}} & ... added to the list.\\
19 & \texttt{\}} & Closes the loop / block.\\
20 & \texttt{local lags \textquotedbl{}\textquotedbl{}} & Empty list of lags.\\
21 & \texttt{forvalues k = 0/\textasciigrave{}mx\textquotesingle{} \{} & For every lag ...\\
22 & \texttt{\hspace*{2.0em}gen byte L\textasciigrave{}k\textquotesingle{} = (rel\_year == \textasciigrave{}k\textquotesingle{})} & ... dummy for k years after adoption ...\\
23 & \texttt{\hspace*{2.0em}local lags \textquotedbl{}\textasciigrave{}lags\textquotesingle{} L\textasciigrave{}k\textquotesingle{}\textquotedbl{}} & ... added to the list.\\
24 & \texttt{\}} & Closes the loop / block.\\
25 & \texttt{gen cohort = cond(g\_main == 0, ., g\_main)} & Cohort variable; missing = never treated.\\
27 & \texttt{eventstudyinteract ShannonBD \textasciigrave{}leads\textquotesingle{} \textasciigrave{}lags\textquotesingle{},\newline \hspace*{1em}cohort(cohort) control\_cohort(never\_main) ///} & Sun-Abraham: interacts event-time dummies with cohorts, uses never-treated counties as controls, and averages cohort-specific effects with cohort-share weights ...\\
28 & \texttt{\hspace*{2.0em}absorb(county year) vce(cluster county)} & ... with county + year-month FE and county-clustered SEs.\\
30 & \texttt{matrix bb = e(b\_iw)} & Interaction-weighted estimates.\\
31 & \texttt{matrix VV = e(V\_iw)} & Their variance matrix.\\
32 & \texttt{tempname S} & Handle for the results file.\\
33 & \texttt{postfile \textasciigrave{}S\textquotesingle{} str12 estimator int event\_time double(b se)\newline \hspace*{1em}using \textquotedbl{}\$T3/tables/sa.dta\textquotedbl{}, replace} & Opens the results file.\\
34 & \texttt{forvalues k = 5(-{}1)2 \{} & For k = 5, 4, 3, 2 ...\\
35 & \texttt{\hspace*{2.0em}capture post \textasciigrave{}S\textquotesingle{} (\textquotedbl{}SA\textquotedbl{}) (-{}\textasciigrave{}k\textquotesingle{}) (bb[1,\textquotedbl{}F\textasciigrave{}k\textquotesingle{}\textquotedbl{}])\newline \hspace*{1em}(sqrt(VV[\textquotedbl{}F\textasciigrave{}k\textquotesingle{}\textquotedbl{},\textquotedbl{}F\textasciigrave{}k\textquotesingle{}\textquotedbl{}]))} & Writes each lead (if estimated).\\
36 & \texttt{\}} & Closes the loop / block.\\
37 & \texttt{post \textasciigrave{}S\textquotesingle{} (\textquotedbl{}SA\textquotedbl{}) (-{}1) (0) (0)} & Reference year.\\
38 & \texttt{forvalues k = 0/5 \{} & For k = 0 ... 5 ...\\
39 & \texttt{\hspace*{2.0em}capture post \textasciigrave{}S\textquotesingle{} (\textquotedbl{}SA\textquotedbl{}) (\textasciigrave{}k\textquotesingle{}) (bb[1,\textquotedbl{}L\textasciigrave{}k\textquotesingle{}\textquotedbl{}])\newline \hspace*{1em}(sqrt(VV[\textquotedbl{}L\textasciigrave{}k\textquotesingle{}\textquotedbl{},\textquotedbl{}L\textasciigrave{}k\textquotesingle{}\textquotedbl{}]))} & Writes each lag.\\
40 & \texttt{\}} & Closes the loop / block.\\
41 & \texttt{* overall post-{}adoption effect: simple average of event\newline \hspace*{1em}years 0-{}5} & Comment (not executed): explains the code that follows.\\
42 & \texttt{matrix w = J(1, colsof(bb), 0)} & Weight vector of zeros ...\\
43 & \texttt{local cn : colnames bb} & ... names of the estimates ...\\
44 & \texttt{local j = 0} & Counter.\\
45 & \texttt{foreach nm of local cn \{} & ... loop over them ...\\
46 & \texttt{\hspace*{2.0em}local ++j} & ... counter + 1 ...\\
47 & \texttt{\hspace*{2.0em}if inlist(\textquotedbl{}\textasciigrave{}nm\textquotesingle{}\textquotedbl{},\textquotedbl{}L0\textquotedbl{},\textquotedbl{}L1\textquotedbl{},\textquotedbl{}L2\textquotedbl{},\textquotedbl{}L3\textquotedbl{},\textquotedbl{}L4\textquotedbl{},\textquotedbl{}L5\textquotedbl{}) matrix\newline \hspace*{1em}w[1,\textasciigrave{}j\textquotesingle{}] = 1/6} & ... weight 1/6 for years 0-5.\\
48 & \texttt{\}} & Closes the loop / block.\\
49 & \texttt{matrix avg  = w*bb\textquotesingle{}} & Overall effect = average of years 0-5 (+0.024).\\
50 & \texttt{matrix vavg = w*VV*w\textquotesingle{}} & Its variance.\\
51 & \texttt{post \textasciigrave{}S\textquotesingle{} (\textquotedbl{}SA\_avg\_0\_5\textquotedbl{}) (99) (avg[1,1]) (sqrt(vavg[1,1]))} & Writes the overall effect.\\
52 & \texttt{di as res \textquotedbl{}Sun-{}Abraham average effect, years 0-{}5: \textquotedbl{} \%8.4f\newline \hspace*{1em}avg[1,1] \textquotedbl{}  (SE \textquotedbl{} \%6.4f sqrt(vavg[1,1]) \textquotedbl{})\textquotedbl{}} & Prints it.\\
53 & \texttt{postclose \textasciigrave{}S\textquotesingle{}} & Closes the file.\\
55 & \texttt{use \textquotedbl{}\$T3/tables/sa.dta\textquotedbl{}, clear} & Loads SA results.\\
56 & \texttt{gen lo = b -{} 1.96*se} & Lower 95\% confidence bound.\\
57 & \texttt{gen hi = b + 1.96*se} & Upper 95\% confidence bound.\\
58 & \texttt{list, noobs sep(0)} & Lists the rows/variables on screen.\\
59 & \texttt{export delimited using \textquotedbl{}\$T3/tables/3\_5\_sun\_abraham.csv\textquotedbl{},\newline \hspace*{1em}replace} & Exports the current data to a CSV file.\\
60 & \texttt{twoway (rcap hi lo event\_time if event\_time != 99,\newline \hspace*{1em}lcolor(dkorange)) ///} & Confidence bars ...\\
61 & \texttt{\hspace*{3.5em}(scatter b event\_time if event\_time != 99,\newline \hspace*{1em}mcolor(dkorange)), ///} & ... and estimates.\\
62 & \texttt{\hspace*{2.0em}yline(0, lpattern(dash)) xline(-{}0.5, lpattern(dot)) ///} & Continuation of the graph command: axis titles, reference lines, labels and legend.\\
63 & \texttt{\hspace*{2.0em}xtitle(\textquotedbl{}Years relative to first solar policy\textquotedbl{})\newline \hspace*{1em}ytitle(\textquotedbl{}Effect on Shannon diversity\textquotedbl{}) ///} & Continuation of the previous command: table title.\\
64 & \texttt{\hspace*{2.0em}title(\textquotedbl{}Sun-{}Abraham event study (PSI-{}data counties)\textquotedbl{},\newline \hspace*{1em}size(medium)) ///} & Continuation of the previous command: table title.\\
65 & \texttt{\hspace*{2.0em}subtitle(\textquotedbl{}Control cohort: never-{}treated; reference year =\newline \hspace*{1em}-{}1; 95\% CI\textquotedbl{}, size(small)) ///} & Continuation of the previous command: table title.\\
66 & \texttt{\hspace*{2.0em}xlabel(-{}5(1)5) legend(off)} & Continuation of the graph command: axis titles, reference lines, labels and legend.\\
67 & \texttt{graph export \textquotedbl{}\$T3/figures/3\_5\_event\_study\_sa.pdf\textquotedbl{}, replace} & Saves the current graph as a PDF.\\}
\subsection{05\_1\_effort\_test.do}
Follow-up (13 Nov): does solar policy reduce birdwatching effort (number of birdwatchers, total time)?

\codetable{1 & \texttt{*=========================================================\newline \hspace*{1em}=================} & Divider line (comment).\\
2 & \texttt{* 05\_1\_effort\_test.do} & Comment (not executed): explains the code that follows.\\
3 & \texttt{* PURPOSE : Prof. Ferraro\textquotesingle{}s question (12 Nov meeting):\newline \hspace*{1em}does solar policy} & Comment (not executed): explains the code that follows.\\
4 & \texttt{*           change birdwatching EFFORT? If PSI lowers the\newline \hspace*{1em}number of} & Comment (not executed): explains the code that follows.\\
5 & \texttt{*           birdwatchers or total birding time, the\newline \hspace*{1em}Shannon \textquotedbl{}effect\textquotedbl{} may be} & Comment (not executed): explains the code that follows.\\
6 & \texttt{*           measurement (fewer observers see fewer\newline \hspace*{1em}species), not fewer birds.} & Comment (not executed): explains the code that follows.\\
7 & \texttt{* NOTE    : One-{}way test. A decline suggests an\newline \hspace*{1em}effort/measurement channel;} & Comment (not executed): explains the code that follows.\\
8 & \texttt{*           no change does not rule it out (birders may\newline \hspace*{1em}move within a county).} & Comment (not executed): explains the code that follows.\\
9 & \texttt{* DATA    : BT = total birdwatching time in the\newline \hspace*{1em}county-{}month (raw)} & Comment (not executed): explains the code that follows.\\
10 & \texttt{*           BN = number of birdwatchers in the\newline \hspace*{1em}county-{}month (raw)} & Comment (not executed): explains the code that follows.\\
11 & \texttt{*           Duration = ln(BT/BN + 1), winsorized (the\newline \hspace*{1em}paper\textquotesingle{}s control)} & Comment (not executed): explains the code that follows.\\
12 & \texttt{*=========================================================\newline \hspace*{1em}=================} & Divider line (comment).\\
13 & \texttt{quietly do \textquotedbl{}/Users/sahiljangra/Desktop/Johns\newline \hspace*{1em}Hopkins/RA/Solar\_Birds\_Replication/00\_config.do\textquotedbl{}} & Runs another do-file: /Users/sahiljangra/Desktop/Johns Hopkins/RA/Solar\_Birds\_Replication/00\_config.do\\
14 & \texttt{use \textquotedbl{}\$SAMPLE\textquotedbl{}, clear} & Loads the analysis dataset built in step 1\_2.\\
15 & \texttt{gen lnBT = ln(BT + 1)          // log total birding time} & Log of total birdwatching time in the county-month (+1 avoids log 0).\\
16 & \texttt{gen lnBN = ln(BN)              // log number of\newline \hspace*{1em}birdwatchers (BN >= 1)} & Log of the number of birdwatchers (always at least 1).\\
17 & \texttt{label var lnBT \textquotedbl{}ln(total birdwatching time + 1)\textquotedbl{}} & Label.\\
18 & \texttt{label var lnBN \textquotedbl{}ln(number of birdwatchers)\textquotedbl{}} & Label.\\
20 & \texttt{eststo clear} & Clears previously stored estimates.\\
21 & \texttt{foreach y in lnBT lnBN Duration BN \{} & Loops over four effort outcomes.\\
22 & \texttt{\hspace*{2.0em}* paper\textquotesingle{}s fixed effects} & Comment (not executed): explains the code that follows.\\
23 & \texttt{\hspace*{2.0em}eststo \textasciigrave{}y\textquotesingle{}\_p: reghdfe \textasciigrave{}y\textquotesingle{} PI \$control\_noD, absorb(county\newline \hspace*{1em}year)       cluster(county)} & Effort on PSI with the paper's fixed effects (8 controls, no Duration).\\
24 & \texttt{\hspace*{2.0em}* urban and rural get their own monthly trends} & Comment (not executed): explains the code that follows.\\
25 & \texttt{\hspace*{2.0em}eststo \textasciigrave{}y\textquotesingle{}\_u: reghdfe \textasciigrave{}y\textquotesingle{} PI \$control\_noD, absorb(county\newline \hspace*{1em}year\#urban) cluster(county)} & Same with urban x year-month FE.\\
26 & \texttt{\hspace*{2.0em}* only counties with real PSI data} & Comment (not executed): explains the code that follows.\\
27 & \texttt{\hspace*{2.0em}eststo \textasciigrave{}y\textquotesingle{}\_r: reghdfe \textasciigrave{}y\textquotesingle{} PI \$control\_noD if\newline \hspace*{1em}in\_PI\_file==1, absorb(county year) cluster(county)} & Same using only counties with real policy data.\\
28 & \texttt{\}} & Closes the loop / block.\\
29 & \texttt{esttab lnBT\_* lnBN\_* , keep(PI) b(4) se(4) p star(* 0.10\newline \hspace*{1em}** 0.05 *** 0.01) ///} & Table of the PSI coefficient for time and number of birdwatchers ...\\
30 & \texttt{\hspace*{2.0em}mtitles(\textquotedbl{}lnBT paper\textquotedbl{} \textquotedbl{}lnBT urbYM\textquotedbl{} \textquotedbl{}lnBT realPSI\textquotedbl{} \textquotedbl{}lnBN\newline \hspace*{1em}paper\textquotedbl{} \textquotedbl{}lnBN urbYM\textquotedbl{} \textquotedbl{}lnBN realPSI\textquotedbl{})} & ... column titles.\\
31 & \texttt{esttab Duration\_* BN\_*, keep(PI) b(4) se(4) p star(* 0.10\newline \hspace*{1em}** 0.05 *** 0.01) ///} & Table for Duration and number of birdwatchers (levels) ...\\
32 & \texttt{\hspace*{2.0em}mtitles(\textquotedbl{}Dur paper\textquotedbl{} \textquotedbl{}Dur urbYM\textquotedbl{} \textquotedbl{}Dur realPSI\textquotedbl{} \textquotedbl{}BN paper\textquotedbl{}\newline \hspace*{1em}\textquotedbl{}BN urbYM\textquotedbl{} \textquotedbl{}BN realPSI\textquotedbl{})} & ... column titles.\\
34 & \texttt{* Does the Shannon effect survive holding the number of\newline \hspace*{1em}birdwatchers fixed?} & Comment (not executed): explains the code that follows.\\
35 & \texttt{* (Caution: lnBN may itself be affected by policy, i.e. a\newline \hspace*{1em}possible bad control.)} & Comment (not executed): explains the code that follows.\\
36 & \texttt{reghdfe ShannonBD PI \$control lnBN,      absorb(county\newline \hspace*{1em}year) cluster(county)} & Paper's model plus the number of birdwatchers: does the Shannon effect survive? (Python preview: +0.003, p = 0.31.)\\
37 & \texttt{reghdfe ShannonBD PI \$control lnBN lnBT, absorb(county\newline \hspace*{1em}year) cluster(county)} & Plus both number of birdwatchers and total time.\\}
\subsection{05\_2\_sensitivity\_dummies.do}
Follow-up (13 Nov): shows that absorbed fixed effects and explicit dummy variables give identical estimates and sensitivity statistics.

\codetable{1 & \texttt{*=========================================================\newline \hspace*{1em}=================} & Divider line (comment).\\
2 & \texttt{* 05\_2\_sensitivity\_dummies.do} & Comment (not executed): explains the code that follows.\\
3 & \texttt{* PURPOSE : Prof. Ferraro\textquotesingle{}s question: is Cinelli-{}Hazlett\newline \hspace*{1em}valid when the} & Comment (not executed): explains the code that follows.\\
4 & \texttt{*           fixed effects are absorbed (reghdfe) instead\newline \hspace*{1em}of entered as dummies?} & Comment (not executed): explains the code that follows.\\
5 & \texttt{* ANSWER  : Yes, by the Frisch-{}Waugh-{}Lovell theorem the\newline \hspace*{1em}coefficient, its} & Comment (not executed): explains the code that follows.\\
6 & \texttt{*           classical SE, t and the residual df are\newline \hspace*{1em}identical, and all C-{}H} & Comment (not executed): explains the code that follows.\\
7 & \texttt{*           statistics depend only on t and df. This file\newline \hspace*{1em}shows it.} & Comment (not executed): explains the code that follows.\\
8 & \texttt{* NOTE    : \textasciitilde{}2,470 regressors. Stata/BE allows at most\newline \hspace*{1em}\textasciitilde{}800 right-{}hand-{}side} & Comment (not executed): explains the code that follows.\\
9 & \texttt{*           variables, so the dummy-{}variable regression\newline \hspace*{1em}needs Stata/SE or MP.} & Comment (not executed): explains the code that follows.\\
10 & \texttt{*=========================================================\newline \hspace*{1em}=================} & Divider line (comment).\\
11 & \texttt{quietly do \textquotedbl{}/Users/sahiljangra/Desktop/Johns\newline \hspace*{1em}Hopkins/RA/Solar\_Birds\_Replication/00\_config.do\textquotedbl{}} & Runs another do-file: /Users/sahiljangra/Desktop/Johns Hopkins/RA/Solar\_Birds\_Replication/00\_config.do\\
12 & \texttt{use \textquotedbl{}\$SAMPLE\textquotedbl{}, clear} & Loads the analysis dataset built in step 1\_2.\\
13 & \texttt{bysort county: drop if \_N == 1                 // same\newline \hspace*{1em}singletons reghdfe drops} & Drops single-month counties, the same ones reghdfe drops, so both versions use identical rows.\\
15 & \texttt{* (a) absorbed fixed effects} & Comment (not executed): explains the code that follows.\\
16 & \texttt{reghdfe ShannonBD PI \$control, absorb(county year)\newline \hspace*{1em}vce(unadjusted)} & Paper's model with classical SEs (for Cinelli-Hazlett).\\
17 & \texttt{di \textquotedbl{}absorbed:  b = \textquotedbl{} \%10.7f \_b[PI] \textquotedbl{}  se = \textquotedbl{} \%9.7f \_se[PI]\newline \hspace*{1em}\textquotedbl{}  t = \textquotedbl{} \%7.4f \_b[PI]/\_se[PI] \textquotedbl{}  df = \textquotedbl{} e(df\_r)} & Prints the absorbed-FE estimate, SE, t and df.\\
19 & \texttt{* (b) explicit dummy variables (Stata/SE or MP)} & Comment (not executed): explains the code that follows.\\
20 & \texttt{capture noisily \{} & Runs the block; an error (e.g. Stata/BE variable limit) does not stop the file.\\
21 & \texttt{\hspace*{2.0em}regress ShannonBD PI \$control i.county i.year} & Same model with explicit county and year-month dummies (about 2,470 regressors; needs Stata/SE or MP).\\
22 & \texttt{\hspace*{2.0em}di \textquotedbl{}dummies:   b = \textquotedbl{} \%10.7f \_b[PI] \textquotedbl{}  se = \textquotedbl{} \%9.7f \_se[PI]\newline \hspace*{1em}\textquotedbl{}  t = \textquotedbl{} \%7.4f \_b[PI]/\_se[PI] \textquotedbl{}  df = \textquotedbl{} e(df\_r)} & Prints the dummy-variable estimate (identical).\\
23 & \texttt{\hspace*{2.0em}sensemakr ShannonBD PI \$control i.county i.year, treat(PI)\newline \hspace*{1em}benchmark(Duration) kd(1 2 3)} & Official sensemakr with explicit dummies.\\
24 & \texttt{\}} & Closes the loop / block.\\
25 & \texttt{* Expected (verified outside Stata with dummy-{}variable\newline \hspace*{1em}OLS):} & Comment (not executed): explains the code that follows.\\
26 & \texttt{*   b = -{}0.0125283, classical se = 0.0027334, t = -{}4.5834,\newline \hspace*{1em}df = 43,898} & Comment (not executed): explains the code that follows.\\
27 & \texttt{*   R2(Y\textasciitilde{}D|X) = 0.000478, RV = 0.02164, RV(alpha=.05) =\newline \hspace*{1em}0.01244} & Comment (not executed): explains the code that follows.\\}
\end{document}
